% !TEX program = xelatex
%
% Epiphany --- Core Specification
% A specification document for the Epiphany music notation platform.
% Compile with XeLaTeX.

\documentclass[11pt,letterpaper]{report}

% ---------------------------------------------------------------------------
% Packages
% ---------------------------------------------------------------------------
\usepackage{fontspec}
\usepackage{geometry}
\geometry{
  letterpaper,
  top=1.05in,
  bottom=1.05in,
  left=1.15in,
  right=1.15in,
  headheight=15pt
}

\usepackage[english]{babel}
\usepackage{microtype}
\usepackage{parskip}
\usepackage{xcolor}
\usepackage{hyperref}
\usepackage{enumitem}
\usepackage{titlesec}
\usepackage{fancyhdr}
\usepackage{booktabs}
\usepackage{array}
\usepackage{longtable}
\usepackage{listings}
\usepackage{amsmath}
\usepackage{amssymb}
\usepackage{graphicx}
\usepackage{caption}
\usepackage{subcaption}
\usepackage{tikz}
\usepackage{float}
\usepackage{tcolorbox}
\tcbuselibrary{breakable, skins}

% ---------------------------------------------------------------------------
% Color palette: refined classical-modern
%   - Deep teal as the primary structural color
%   - Antique gold as the accent and titling color
%   - Warm parchment-tinted neutrals
% ---------------------------------------------------------------------------
\definecolor{epiphanyteal}{HTML}{1A4044}      % deep teal --- primary
\definecolor{epiphanygold}{HTML}{8E6E2E}      % antique gold --- accent
\definecolor{epiphanyink}{HTML}{1F1B16}       % warm near-black for body text
\definecolor{epiphanyslate}{HTML}{6B6660}     % warm slate gray
\definecolor{epiphanycream}{HTML}{F8F4ED}     % parchment cream
\definecolor{epiphanymist}{HTML}{ECE8E0}      % slightly deeper parchment
\definecolor{epiphanycode}{HTML}{2A2520}      % warm dark for code
\definecolor{epiphanycrimson}{HTML}{7A2424}   % deep crimson for warnings

\hypersetup{
  colorlinks=true,
  linkcolor=epiphanyteal,
  citecolor=epiphanyteal,
  urlcolor=epiphanygold,
  pdftitle={Epiphany --- Core Specification},
  pdfauthor={The Epiphany Project},
  pdfsubject={Core specification for the Epiphany music notation platform},
  pdfkeywords={music notation, engraving, FOSS, Rust, microtonality, CRDT},
  bookmarksnumbered=true,
  bookmarksopen=true
}

% ---------------------------------------------------------------------------
% Typography: classical serif with refined sans for chapter titles
%   - TeX Gyre Pagella: Palatino-derived; warm, classical, highly legible
%   - TeX Gyre Heros: clean grotesque for sans elements
%   - TeX Gyre Cursor: monospace for code
% ---------------------------------------------------------------------------
\setmainfont{TeX Gyre Pagella}[
  Numbers={OldStyle, Proportional},
  Ligatures={TeX, Common}
]
\setsansfont{TeX Gyre Heros}[
  Scale=0.94,
  Ligatures={TeX, Common}
]
\setmonofont{TeX Gyre Cursor}[
  Scale=0.88,
  Ligatures={TeX}
]

% Lining figures for tables (oldstyle for body, lining for technical)
\newcommand{\tablenums}[1]{{\addfontfeatures{Numbers={Lining,Tabular}}#1}}

% A display-weight version of Pagella for the title page
\newfontfamily\titlefont{TeX Gyre Pagella}[
  Numbers={OldStyle},
  Ligatures={TeX, Common}
]

% Small-caps helpers
\newcommand{\sectionsc}[1]{{\addfontfeatures{Letters=SmallCaps}#1}}

% ---------------------------------------------------------------------------
% Section styling
% ---------------------------------------------------------------------------
\titleformat{\chapter}[display]
  {\normalfont\filright}
  {\raggedright\color{epiphanygold}\fontsize{14pt}{16pt}\selectfont
    \scshape Chapter\ \thechapter}
  {16pt}
  {\raggedright\color{epiphanyteal}\fontsize{32pt}{36pt}\selectfont\bfseries}
  [\vspace{4pt}{\color{epiphanygold}\rule{2in}{0.6pt}}]
\titlespacing*{\chapter}{0pt}{-20pt}{30pt}

\titleformat{\section}
  {\normalfont\Large\bfseries\color{epiphanyteal}}
  {\color{epiphanygold}\thesection}{1em}{}
\titleformat{\subsection}
  {\normalfont\large\bfseries\color{epiphanyteal}}
  {\color{epiphanygold}\thesubsection}{1em}{}
\titleformat{\subsubsection}
  {\normalfont\normalsize\bfseries\color{epiphanyink}}
  {\thesubsubsection}{1em}{}

% ---------------------------------------------------------------------------
% Headers and footers
% ---------------------------------------------------------------------------
\pagestyle{fancy}
\fancyhf{}
\renewcommand{\headrulewidth}{0pt}
\renewcommand{\footrulewidth}{0pt}
\fancyhead[L]{\small\scshape\color{epiphanyslate}Epiphany}
\fancyhead[R]{\small\itshape\color{epiphanyslate}\leftmark}
\fancyfoot[C]{\small\color{epiphanyslate}\thepage}

% Subtle ornamental rule under the page header
\renewcommand{\headrule}{
  \color{epiphanygold!50}\hrule width\headwidth height 0.4pt
  \vspace{1pt}
  \color{epiphanygold!30}\hrule width\headwidth height 0.2pt
}

% ---------------------------------------------------------------------------
% Code listing style
% ---------------------------------------------------------------------------
\lstdefinelanguage{Rust}{
  keywords={fn,let,mut,pub,struct,enum,impl,trait,for,in,if,else,match,return,
           use,mod,crate,self,Self,as,where,move,async,await,const,static,
           ref,type,unsafe,extern,dyn,box,break,continue,loop,while},
  keywordstyle=\color{epiphanyteal}\bfseries,
  ndkeywords={i8,i16,i32,i64,i128,u8,u16,u32,u64,u128,f32,f64,bool,char,str,
              String,Vec,Option,Result,Box,Rc,Arc,HashMap,BTreeMap,
              NonZeroU16,NonZeroU32,NonZeroU64,Duration,Timestamp},
  ndkeywordstyle=\color{epiphanygold}\bfseries,
  sensitive=true,
  comment=[l]{//},
  morecomment=[s]{/*}{*/},
  commentstyle=\color{epiphanyslate}\itshape,
  stringstyle=\color{epiphanycrimson},
  morestring=[b]",
  morestring=[b]'
}

\lstset{
  basicstyle=\ttfamily\small\color{epiphanycode},
  backgroundcolor=\color{epiphanycream},
  frame=leftline,
  rulecolor=\color{epiphanygold!60},
  framesep=8pt,
  framerule=1.5pt,
  xleftmargin=10pt,
  xrightmargin=4pt,
  breaklines=true,
  showstringspaces=false,
  numberstyle=\tiny\color{epiphanyslate},
  numbersep=10pt,
  captionpos=b,
  aboveskip=10pt,
  belowskip=10pt,
  language=Rust
}

% ---------------------------------------------------------------------------
% Custom environments
% ---------------------------------------------------------------------------

% Open question or unresolved design issue
\newtcolorbox{openquestion}[1][]{
  enhanced, breakable,
  colback=epiphanymist, colframe=epiphanycrimson,
  fonttitle=\bfseries\color{white},
  title={\scshape\hspace{2pt}Open Question},
  coltitle=white,
  colbacktitle=epiphanycrimson,
  arc=1pt, boxrule=0pt,
  leftrule=2pt,
  left=10pt, right=10pt, top=8pt, bottom=8pt,
  attach boxed title to top left={xshift=0pt, yshift=0pt},
  boxed title style={
    arc=0pt, sharp corners, boxrule=0pt,
    left=6pt, right=8pt, top=2pt, bottom=2pt,
  },
  #1
}

% Design rationale
\newtcolorbox{rationale}[1][]{
  enhanced, breakable,
  colback=epiphanymist, colframe=epiphanyteal,
  fonttitle=\bfseries\color{white},
  title={\scshape\hspace{2pt}Rationale},
  coltitle=white,
  colbacktitle=epiphanyteal,
  arc=1pt, boxrule=0pt,
  leftrule=2pt,
  left=10pt, right=10pt, top=8pt, bottom=8pt,
  attach boxed title to top left={xshift=0pt, yshift=0pt},
  boxed title style={
    arc=0pt, sharp corners, boxrule=0pt,
    left=6pt, right=8pt, top=2pt, bottom=2pt,
  },
  #1
}

% Normative requirement
%
% Numbered within chapter (this document has chapters) so that every
% requirement carries a unique, human-citable number. tcolorbox's own
% "auto counter, number within=..." keys are not used: their counter-step
% for \label purposes runs inside an internal \sbox ("phantom") whose local
% \@currentlabel assignment is discarded when that box closes, so a plain
% \label written in the box body (as every requirement's is, and per this
% pass's rules must remain) binds to the counter of whatever sectioning
% unit last really stepped -- reproducing the very bug this counter is
% meant to fix. Stepping a plain counter via a `code=` key runs directly in
% the environment's own group, so \@currentlabel sticks for the body's
% \label to see.
\newcounter{requirement}[chapter]
\renewcommand{\therequirement}{\thechapter.\arabic{requirement}}
\newtcolorbox{requirement}[1][]{
  enhanced, breakable,
  colback=white, colframe=epiphanygold,
  fonttitle=\bfseries\color{white},
  code={\refstepcounter{requirement}},
  title={\scshape\hspace{2pt}Requirement~\therequirement},
  coltitle=white,
  colbacktitle=epiphanygold,
  arc=1pt, boxrule=0pt,
  leftrule=2pt,
  left=10pt, right=10pt, top=8pt, bottom=8pt,
  attach boxed title to top left={xshift=0pt, yshift=0pt},
  boxed title style={
    arc=0pt, sharp corners, boxrule=0pt,
    left=6pt, right=8pt, top=2pt, bottom=2pt,
  },
  #1
}

% Non-goal: explicitly out of scope
\newtcolorbox{nongoal}[1][]{
  enhanced, breakable,
  colback=epiphanymist, colframe=epiphanyslate,
  fonttitle=\bfseries\color{white},
  title={\scshape\hspace{2pt}Non-Goal},
  coltitle=white,
  colbacktitle=epiphanyslate,
  arc=1pt, boxrule=0pt,
  leftrule=2pt,
  left=10pt, right=10pt, top=8pt, bottom=8pt,
  attach boxed title to top left={xshift=0pt, yshift=0pt},
  boxed title style={
    arc=0pt, sharp corners, boxrule=0pt,
    left=6pt, right=8pt, top=2pt, bottom=2pt,
  },
  #1
}

% Normative-language helpers (RFC 2119 style). MUST NOT and SHOULD NOT
% are emitted as two adjacent textbf groups separated by a regular
% control space; this prevents line breaks between the words while
% still producing a standard ASCII space character in the PDF that
% pdftotext extracts cleanly. The \nobreak prevents the line-break
% the regular space would otherwise allow.
\newcommand{\MUST}{\textbf{MUST}}
\newcommand{\MUSTNOT}{\textbf{MUST}\nobreak\ \textbf{NOT}}
\newcommand{\SHOULD}{\textbf{SHOULD}}
\newcommand{\SHOULDNOT}{\textbf{SHOULD}\nobreak\ \textbf{NOT}}
\newcommand{\MAY}{\textbf{MAY}}

% List spacing
\setlist[itemize]{topsep=2pt, itemsep=3pt, parsep=0pt}
\setlist[enumerate]{topsep=2pt, itemsep=3pt, parsep=0pt}
\setlist[description]{topsep=2pt, itemsep=5pt, parsep=0pt}

% Body text color
\AtBeginDocument{\color{epiphanyink}}

% ---------------------------------------------------------------------------
% Document
% ---------------------------------------------------------------------------
\begin{document}

% ----- Title page -----
\begin{titlepage}
  \thispagestyle{empty}
  \begin{tikzpicture}[remember picture, overlay]
    % Top ornamental rule
    \draw[epiphanygold, line width=1.2pt]
      ([yshift=-0.85in, xshift=1.15in]current page.north west) --
      ([yshift=-0.85in, xshift=-1.15in]current page.north east);
    \draw[epiphanygold, line width=0.4pt]
      ([yshift=-0.93in, xshift=1.15in]current page.north west) --
      ([yshift=-0.93in, xshift=-1.15in]current page.north east);

    % Bottom ornamental rule
    \draw[epiphanygold, line width=0.4pt]
      ([yshift=0.93in, xshift=1.15in]current page.south west) --
      ([yshift=0.93in, xshift=-1.15in]current page.south east);
    \draw[epiphanygold, line width=1.2pt]
      ([yshift=0.85in, xshift=1.15in]current page.south west) --
      ([yshift=0.85in, xshift=-1.15in]current page.south east);
  \end{tikzpicture}

  \centering
  \vspace*{0.7in}

  % Small superscript framing
  {\color{epiphanyslate}\scshape\fontsize{11pt}{13pt}\selectfont
    A music notation platform\par}
  \vspace{0.3in}

  % The main title --- Epiphany
  {\titlefont\fontsize{72pt}{76pt}\selectfont
    \color{epiphanyteal}
    Epiphany\par}

  \vspace{0.12in}

  % Subtle ornamental separator
  \begin{tikzpicture}
    \draw[epiphanygold, line width=0.5pt] (-1.3,0) -- (-0.3,0);
    \node at (0,0) {\color{epiphanygold}\large$\diamond$};
    \draw[epiphanygold, line width=0.5pt] (0.3,0) -- (1.3,0);
  \end{tikzpicture}

  \vspace{0.2in}

  {\color{epiphanyink}\fontsize{18pt}{22pt}\selectfont
    \itshape Core Specification\par}

  \vspace{0.9in}

  % Tagline
  \begin{minipage}{0.78\textwidth}
    \centering
    \color{epiphanyink}
    \fontsize{12pt}{17pt}\selectfont
    A foundational specification for an open, Rust-native music
    notation platform: pitch and time primitives, the score graph,
    the layout intermediate representation, the file format, and
    the constraint-solver interface.
  \end{minipage}

  \vspace{0.7in}

  % Status banner
  \begin{tikzpicture}
    \node[
      draw=epiphanygold,
      line width=0.5pt,
      fill=epiphanycream,
      inner sep=10pt,
      minimum width=3.5in,
    ] (status) {
      \color{epiphanyteal}
      \scshape\fontsize{10pt}{12pt}\selectfont
      Working Draft \quad$\cdot$\quad \today
    };
  \end{tikzpicture}

  \vspace{0.5in}

  \begin{minipage}{0.72\textwidth}
    \centering
    \itshape\color{epiphanyslate}
    \footnotesize
    This specification defines the foundational layer of the
    Epiphany platform: the data model, engraving engine, and file
    format on which all higher-level features are built. The user
    interface, audio engine, plugin runtime, and platform
    integrations are layered above the core and are addressed in
    separate specifications.
  \end{minipage}

  \vfill
\end{titlepage}

% ----- Front matter -----
\pagenumbering{roman}

\chapter*{Status of This Document}
\addcontentsline{toc}{chapter}{Status of This Document}

This draft is \textbf{structurally stabilized}. Substantive
review passes through the core type model, graph invariants,
concurrent reduction semantics, file-format atomicity, solver
conformance, and the determinism contract have been completed.
The architecture is no longer under active redesign; further
revisions are expected to be additive (companion specifications,
normative algorithm choices for the open canonical algorithms,
extension registry contents) rather than structural.

The remaining blockers to full conformance are companion
specifications and the canonical algorithms identified in
Appendix~\ref{app:determinism}
Section~\ref{sec:det:open}:

\begin{itemize}
  \item Several companion specifications referenced throughout
    the document have not yet been delivered:
    Appendix~\ref{app:deferred} catalogs them and distinguishes
    them from accidentally missing content.
  \item Two algorithms affect canonical state and have not
    yet received one of the dispositions in
    Appendix~\ref{app:determinism}
    Section~\ref{sec:det:open}: the tempo curve
    integration algorithm and the
    \texttt{wallclock\_to\_musical} root-finding algorithm.
    Cross-implementation byte equality of canonical score state
    \MUSTNOT{} be claimed for outputs derived from these
    algorithms until they are normatively specified or
    profile-declared (Appendix~\ref{app:determinism}
    Section~\ref{sec:det:layers}). The spelling pre-pass and the
    notational-decomposition algorithm received the
    profile-declared-by-versioned-identifier disposition in
    Pass~12
    (Requirements~\ref{req:pitch:spelling-algorithm}
    and~\ref{req:time:decomposition-algorithm}).
  \item Extension registry catalogs are referenced by typed
    identifier but not enumerated here; they are delivered as
    separate, versioned registry documents.
\end{itemize}

What conformance can be established now:

\begin{itemize}
  \item The core data model can be structurally implemented and
    tested against the invariants in
    Chapter~\ref{ch:graph} Section~\ref{sec:graph:invariants}.
  \item The operation framework
    (Chapter~\ref{ch:semops}) can be implemented end-to-end
    against the canonical reduction algorithm.
  \item The file format
    (Chapter~\ref{ch:format}) can be implemented atomically and
    crash-recoverably against the fixed prelude and superblock
    protocol.
  \item Solver implementations can be developed against the
    interface contract (Chapter~\ref{ch:solver}); reference-suite
    conformance becomes checkable when the Reference Suite
    companion is delivered.
\end{itemize}

The normative keywords \MUST, \MUSTNOT, \SHOULD, \SHOULDNOT, and \MAY{} are
to be interpreted as described in RFC 2119 when they appear in all
capital letters.

\chapter*{Conventions}
\addcontentsline{toc}{chapter}{Conventions}

\begin{description}
  \item[Code samples] Rust syntax is used for type signatures, struct
    definitions, and algorithmic illustrations. Code samples are
    illustrative unless explicitly labeled normative.
  \item[Normative requirements] Statements containing \MUST, \MUSTNOT,
    \SHOULD, \SHOULDNOT, or \MAY{} are normative. Conforming
    implementations must respect them.
  \item[Non-normative content] Discussion, rationale, examples, and
    diagrams are non-normative unless explicitly stated.
  \item[Forward references] The specification builds bottom-up. Later
    chapters depend on earlier ones; forward references are minimized but
    occasionally unavoidable and are marked when they occur.
\end{description}

\tableofcontents

\clearpage

\chapter*{Reading Guide}
\addcontentsline{toc}{chapter}{Reading Guide}

This document is long. The following tables summarize the most
load-bearing concepts and where they are defined. They are
non-normative; the chapters they reference are authoritative.

\section*{Score Graph Identity Hierarchy}

The score graph is organized as a containment hierarchy of
identified objects. Each level introduces a distinct concept of
identity.

\begin{longtable}{p{2.8cm} p{2.6cm} p{8.2cm}}
  \toprule
  \textbf{Level} & \textbf{Identifier} & \textbf{Role} \\
  \midrule
  \endhead
  Score & --- & Document root. Carries metadata, instruments,
    global staves, staff groups, tuning context, tempo map,
    canvas. \\
  Canvas & --- & Spatial root. Partitioned into regions. \\
  Region & \texttt{RegionId} & Spatial-temporal container.
    Declares time model (metric/proportional/aleatoric) and
    content model (staff-based/free-graphic/hybrid). \\
  Staff & \texttt{StaffId} & Global, abstract staff identity
    persisting across regions. Declared once at score level. \\
  Staff instance & \texttt{StaffInstanceId} & Region-local
    manifestation of a staff. One \texttt{Staff} may have
    multiple instances across regions. \\
  Voice & \texttt{VoiceId} & Polyphonic line within a staff
    instance. Carries \texttt{VoiceOrigin}: user-declared,
    imported, or system-promoted. \\
  Event & \texttt{EventId} & Rhythmic atom: pitched, unpitched,
    rest, indeterminate, trajectory, graphic, cue. Lives in the
    score's event arena. \\
  Identified pitch & \texttt{PitchId} & Stable identity for a
    pitch within an event. Enables spelling attachments,
    respelling, stable tie pairing through chord reordering. \\
  \bottomrule
\end{longtable}

\section*{Canonical versus Non-Canonical}

Many file-format and runtime objects exist in two forms.
Canonical objects define the document; non-canonical objects
accelerate access but never affect canonical state.

\begin{longtable}{p{4.4cm} p{4.0cm} p{5.2cm}}
  \toprule
  \textbf{Object} & \textbf{Status} & \textbf{Notes} \\
  \midrule
  \endhead
  Operation envelope set & Canonical &
    The replicated grow-only CRDT defining the document
    (Chapter~\ref{ch:semops}). \\
  Operation envelope block & Canonical (storage) &
    Soft 1\,MiB blocks on disk
    (Chapter~\ref{ch:format}). \\
  Operation index & Non-canonical &
    Acceleration of operation-id lookup. Rebuildable. \\
  Canonical-base snapshot & Canonical &
    The manifest's \texttt{canonical\_base}: at most one active,
    with declared causal frontier and reduction algorithm
    version. Required for pruned documents. \\
  Acceleration snapshot & Non-canonical &
    Listed in \texttt{acceleration\_snapshots}. Cache only.
    Discardable. \\
  Manifest & Canonical (index) &
    Table of roots; itself a chunk; referenced from the active
    superblock. \\
  Materialized score graph & Derivative &
    Deterministic reduction of the operation set under the
    active reduction algorithm. \\
  Layout cache, glyph cache, render artifacts &
    Non-canonical & Always discardable
    (Section~\ref{sec:semops:caches}). \\
  Conflict registry & Canonical &
    Stable, addressable, user-visible records of operations
    that could not apply cleanly. \\
  \bottomrule
\end{longtable}

\section*{Operation Lifecycle}

Every operation traverses four phases between user intent and
its visible effect.

\begin{longtable}{p{2.2cm} p{12.5cm}}
  \toprule
  \textbf{Phase} & \textbf{Description} \\
  \midrule
  \endhead
  Prepare & Local UI command validates enough to construct a
    well-formed envelope. Advisory only; failure here means the
    UI rejects the user's action. Does not bind canonical
    semantics. \\
  Commit & The envelope is added to the replicated operation
    set. After commit, the operation is part of canonical state
    and \MUSTNOT{} be discarded outside the pruning protocol. \\
  Reduce & During materialization, the operation is processed
    by the canonical reduction algorithm against working state.
    Deterministic: every replica with the same operation set
    produces the same output. \\
  Report & The reduction produces a deterministic
    \texttt{OperationEffect}: \texttt{Applied},
    \texttt{AppliedWithRepair}, \texttt{Conflicted},
    \texttt{TombstonedTarget}, or \texttt{NoOp}. Effects are
    visible graph facts. \\
  \bottomrule
\end{longtable}

\section*{Determinism Layers}

The specification distinguishes five layers of determinism.
Conflating them is the most common source of impossible
requirements.

\begin{longtable}{p{3.8cm} p{10.0cm}}
  \toprule
  \textbf{Layer} & \textbf{Contract} \\
  \midrule
  \endhead
  Canonical score determinism &
    Identical operation sets produce identical materialized
    score states across all conforming implementations, subject
    to all canonical algorithms having received the dispositions
    in Appendix~\ref{app:determinism}. \\
  Canonical serialization &
    Same canonical state $\to$ same chunk hashes, same text
    projection, same canonical iteration order. \\
  Layout determinism &
    Within one solver implementation at a fixed version:
    byte-equal output for identical input. Across
    implementations: reference-suite thresholds, not byte equal. \\
  Conformance determinism &
    Reference-suite-based. Cross-implementation byte equality
    is not required for layout. \\
  Non-canonical caches &
    May vary freely across runs, platforms, implementations.
    \MUSTNOT{} affect canonical state. \\
  \bottomrule
\end{longtable}

\section*{Solver Tiers}

Solver implementations declare a conformance tier. Tiers are
orthogonal to file-format profiles
(Chapter~\ref{ch:format}).

\begin{longtable}{p{3.0cm} p{11.0cm}}
  \toprule
  \textbf{Tier} & \textbf{Obligations} \\
  \midrule
  \endhead
  Minimal &
    Hard-constraint validity, within-implementation determinism,
    well-formed \texttt{SolveReport}, CMN + metric regions,
    Minimal reference-suite subset. Aesthetic quality may be
    mediocre. \\
  Standard &
    Minimal $+$ Standard-tier thresholds on the full reference
    suite, Measure-local invalidation under incremental solving,
    all Standard constraint families. Professional engraving
    quality for common-practice notation. \\
  Advanced &
    Standard $+$ proportional and aleatoric regions, microtonal
    accidentals, extension-contributed constraints, System-local
    invalidation under incremental solving. \\
  \bottomrule
\end{longtable}

\section*{Chapter Map}

\begin{longtable}{p{1.6cm} p{12.4cm}}
  \toprule
  \textbf{Ch.} & \textbf{Content} \\
  \midrule
  \endhead
  \ref{ch:intro} & Scope, design principles, glossary of key
    terms used throughout. \\
  \ref{ch:pitch} & Pitch as two intrinsic layers (scale
    position and acoustic), spelling subsystem with stacked
    accidentals, analytical layers. \\
  \ref{ch:time} & Rational time, time anchors with
    \texttt{AnchorOffset}, three region time models, tempo map
    with explicit segment endpoints, sounding duration vs.
    notational decomposition. \\
  \ref{ch:tuning} & Pitch spaces (including CMN and JI), tuning
    systems, score tuning context, accidental registries, glyph
    references. \\
  \ref{ch:graph} & Score graph, event arena, cross-cutting
    structures, regions, staff identity (Staff vs.
    StaffInstance), polymeter via per-staff metric grids, voice
    origin, graph invariants. \\
  \ref{ch:semops} & Operation framework, canonical reduction,
    DVV causal contexts, operation effects, tombstones, conflict
    records, transactions, re-anchoring, undo, LWW discipline. \\
  \ref{ch:layout-ir} & The four-stage layout intermediate
    representation; provenance and incremental invalidation. \\
  \ref{ch:format} & Fixed prelude, two superblock slots, atomic
    commit, content-addressed BLAKE3 chunks,
    operation-envelope blocks, snapshots as canonical bases,
    edit barriers, text projection. \\
  \ref{ch:solver} & Solver interface, SolveReport, deterministic
    budgets, NormalizedMetric, three conformance tiers,
    reference-suite conformance, six invalidation scopes. \\
  \ref{ch:perf} & Performance requirements: frame budgets, cold
    open, memory, audio-thread discipline, measurement. \\
  \ref{ch:extension} & Consolidated catalog of extension points
    by chapter of origin. \\
  \ref{app:deferred} & Intentionally deferred types: companion
    specifications, externally provided components, open
    canonical algorithms. \\
  \ref{app:glossary} & Glossary of defined terms with
    chapter cross-references. \\
  \ref{app:refs} & Bibliography and references. \\
  \ref{app:determinism} & Determinism contract: the legal code
    for reproducibility. \\
  \ref{app:bytes} & Canonical byte-layout reference: every ratified
    discriminant, derivation preimage, and primitive encoding in one
    place, for the Binary Format companion to import. \\
  \ref{app:history} & Revision history. \\
  \bottomrule
\end{longtable}

% ----- Main matter -----
\clearpage
\pagenumbering{arabic}

% ===========================================================================
\chapter{Introduction and Scope}
\label{ch:intro}

\section{Purpose}

This document specifies the core layer of \textit{Epiphany}, a music
notation platform. The core comprises the abstract musical data
model, the engraving engine that transforms that data into a typeset
layout, and the on-disk file format that serializes both. Everything
else built atop the platform---editing UI, audio engine, plugin
runtime, collaboration, export pipelines---depends on this core and
must conform to its interfaces.

The name \textit{Epiphany} evokes both the original sense of the
word---a manifestation made visible---and the experience of music
itself, which strikes the listener and the creator alike. Music
notation is, in a literal sense, the act of making the inaudible
visible: a body of work this central to that task deserves a name
of corresponding weight.

\section{Scope}

\textbf{In scope:}

\begin{itemize}
  \item Pitch, time, and duration primitives.
  \item The score graph: the in-memory representation of musical content.
  \item The layout intermediate representation produced by the engraving
    engine.
  \item The file format: the canonical on-disk projection of the score
    graph.
  \item The constraint-solver interface consumed by the engraving engine.
  \item Extension points for non-Common-Music-Notation systems, including
    microtonal pitch spaces and graphic-notation regions.
\end{itemize}

\textbf{Out of scope for this document:}

\begin{nongoal}
  The following are explicitly out of scope for this core specification.
  Each is addressed by a separate specification layered atop the core.
  \begin{itemize}
    \item User interface, input handling, and editing operations.
    \item Audio synthesis, playback, and plugin (VST3/AU/CLAP/LV2) hosting.
    \item The plugin runtime (WebAssembly capability-based extensions).
    \item Real-time collaboration transport and federation protocols.
    \item Import and export pipelines for foreign formats (MusicXML, MEI,
      LilyPond, MIDI, etc.), although the core defines the data
      sufficient to make lossless interchange feasible.
    \item Machine-learning components (transcription, OMR, engraving
      assistance).
    \item Platform-specific integrations (file pickers, sharing, accessibility
      APIs).
  \end{itemize}
\end{nongoal}

\section{Design Goals}

\begin{description}
  \item[Correctness over convenience.] The data model must represent
    music accurately, even when this is harder than a convenient
    approximation. MIDI-style pitch numbers, for example, are inadequate
    and rejected; see Chapter~\ref{ch:pitch}.
  \item[Extensibility without bolt-ons.] Microtonality, graphic notation,
    historical notations, and future notational systems must be
    expressible within the core data model, not as escape hatches.
  \item[Determinism.] Given identical input, the engraving engine must
    produce identical output. The on-disk format must serialize
    deterministically.
  \item[Incrementality.] The data model and layout engine must support
    fine-grained incremental edits without full recomputation.
  \item[Reducibility.] All editing operations on the score graph
    must have well-defined reduction rules under canonical
    deterministic reduction (Chapter~\ref{ch:semops}), enabling
    collaboration and version control as natural consequences of
    the data model rather than bolt-on features.
  \item[Performance.] Layout and rendering of a 100-page orchestral score
    must remain interactive on commodity hardware. See
    Chapter~\ref{ch:perf} for budgets.
  \item[Privacy.] The core must operate entirely offline. No core
    functionality may depend on network access.
\end{description}

\section{Document Structure}

The remainder of this document is organized bottom-up:

\begin{enumerate}
  \item Chapter~\ref{ch:pitch} defines pitch: its layered structure, the
    spelling subsystem, scale positions, and acoustic realization.
  \item Chapter~\ref{ch:time} defines time and duration primitives.
  \item Chapter~\ref{ch:tuning} defines tuning systems, pitch spaces, and
    accidental registries.
  \item Chapter~\ref{ch:graph} defines the score graph: how primitives
    compose into staves, voices, regions, and the overall score
    structure, including the canvas/region model that accommodates
    non-CMN notations.
  \item Chapter~\ref{ch:semops} defines the semantic operations on the
    score graph: the canonical set of mutations and their reduction
    semantics under concurrent edits.
  \item Chapter~\ref{ch:layout-ir} defines the layout intermediate
    representation: the contract between the score graph (input) and the
    engraving engine (consumer).
  \item Chapter~\ref{ch:format} defines the on-disk file format, including
    serialization, chunking, the operation-envelope set and its
    deterministic reduction, and schema evolution.
  \item Chapter~\ref{ch:solver} defines the constraint-solver interface
    consumed by the engraving engine.
  \item Chapter~\ref{ch:perf} states performance requirements and
    measurement methodology.
  \item Chapter~\ref{ch:extension} describes the extension points by
    which the core supports notation grammars beyond Common Music
    Notation.
\end{enumerate}

% ===========================================================================
\chapter{Pitch}
\label{ch:pitch}

This chapter specifies the pitch primitive: the data type representing
a single sounding (or to-be-sounded) note's identity in the score graph.
Pitch is the most consequential primitive in the core because nearly every
higher-level operation---transposition, key analysis, accidental rendering,
playback frequency, interchange, voice-leading analysis---is defined in
terms of it. The definition is therefore deliberately verbose.

\section{Design Principles}
\label{sec:pitch:principles}

\begin{description}
  \item[Separation of identity layers.] Pitch comprises three
    independent semantic layers: scale position (analytical identity),
    acoustic realization (sounding frequency), and spelling (notational
    appearance). Each is a first-class concern and \MUST{} be
    representable independently. Conflating them is a design error.
  \item[Spelling is attachment, not field.] A pitch carries its scale
    position and acoustic identity intrinsically. Spelling is attached
    externally through a separate indexed subsystem (see
    Section~\ref{sec:pitch:spelling}). A pitch may have zero, one, or
    multiple spelling attachments.
  \item[Grammar neutrality.] The pitch type \MUST{} support Common Music
    Notation as a fast path while remaining capable of representing
    pitches in arbitrary user-defined notation grammars. CMN is a
    privileged default, not a hardcoded ceiling.
  \item[Microtonality is not an extension.] Microtonal pitches are
    representable using the same primitive as 12-tone equal-tempered
    pitches. No separate type, no escape hatch.
  \item[Unpitched events are not pitches.] Unpitched percussion and
    similar events \MUSTNOT{} be represented using the pitch type.
    See Section~\ref{sec:pitch:not-pitch}.
\end{description}

\section{The \texttt{Pitch} Type}
\label{sec:pitch:type}

A pitch is a record of two intrinsic layers:

\begin{lstlisting}[language=Rust]
/// A pitch's intrinsic identity. Spellings are attached externally;
/// see Section~\ref{sec:pitch:spelling}.
pub struct Pitch {
    /// The analytical identity within a pitch space. Determines
    /// transposition, key analysis, and Roman-numeral behavior.
    pub scale_position: ScalePosition,

    /// The acoustic identity: what frequency this pitch produces when
    /// sounded, expressed relative to a tuning system.
    pub acoustic: AcousticPitch,
}
\end{lstlisting}

The two layers are independent. A given scale position can be acoustically
realized in many ways depending on the active tuning system; a given
acoustic realization can correspond to many scale positions in different
analytical contexts. The pitch carries both; neither is derivable from
the other in general.

\begin{rationale}
  Storing both layers, rather than deriving one from the other, costs
  a small number of bytes per pitch and saves a large amount of complexity
  downstream. Derivation requires the renderer, the analyzer, the
  transcriber, and the importer to each agree on a derivation procedure;
  storage makes the pitch self-describing.
\end{rationale}

\begin{lstlisting}[language=Rust]
/// A closed range of pitches, used for advisory declarations: an
/// instrument's playable compass and an indeterminate event's pitch hint.
/// Both bounds are full pitches.
pub struct PitchRange {
    pub lowest: Pitch,
    pub highest: Pitch,
}
\end{lstlisting}

A \texttt{PitchRange} is \emph{advisory}. Whether a given pitch lies within
it is decidable only when the pitch and both bounds share a pitch space that
defines an order (the common CMN case); cross-space or partially-ordered
comparisons are treated as ``not out of range'' --- a sound-but-incomplete
check, as with region time overlap. A range is well-formed when
\texttt{lowest} does not sort above \texttt{highest} in a common ordered
space. A solver \MAY{} flag an out-of-range pitch but \MUSTNOT{} treat the
range as a hard constraint. The type is defined with schema major~1, closing
the gap in which \texttt{IndeterminacyHints} and \texttt{Instrument} named a
\texttt{PitchRange} that no section defined.

\section{Scale Position}
\label{sec:pitch:scale-position}

A scale position locates the pitch within a \emph{pitch space}. A pitch
space is the analytical universe in which scale-degree and interval
relationships are defined.

\begin{lstlisting}[language=Rust]
pub struct ScalePosition {
    /// The pitch space this position is defined within. References an
    /// identifier in the built-in pitch-space catalog
    /// (`sec:tuning:builtin`); a score selects from this catalog and
    /// does not define its own entries.
    pub space: PitchSpaceId,

    /// The position within that space. Interpretation is space-defined.
    pub position: PitchSpacePosition,
}
\end{lstlisting}

The \texttt{PitchSpacePosition} type is a tagged union supporting the
common cases with a fast representation, plus a registry-id escape hatch
for arbitrary grammars:

\begin{lstlisting}[language=Rust]
pub enum PitchSpacePosition {
    /// Common Music Notation: a diatonic nominal plus chromatic
    /// alteration plus octave. The fast path for tonal Western music.
    Cmn {
        nominal: CmnNominal,       // C, D, E, F, G, A, B
        alteration: i8,            // chromatic steps of the enclosing space
        octave: i8,                // scientific pitch notation
    },

    /// N-tone integer position. Used for 12-tone serial music, EDO
    /// systems, and any space where pitches are indexed as integers.
    Integer {
        space_size: u16,           // e.g., 12, 19, 31, 53
        index: i32,                // absolute index; octave is implicit
    },

    /// Just-intonation lattice position. Used for JI pitch spaces.
    /// Each component is the integer power of the corresponding prime
    /// in the space's prime basis. The pitch space declares the basis
    /// ordering; by convention prime 2 is first, so the leading
    /// component carries octave information.
    JiVector {
        components: Vec<i32>,
    },

    /// A registered position whose interpretation is defined by the
    /// pitch space's grammar plugin. Used for maqamat, gamelan,
    /// historical modes, and other non-lattice, non-CMN grammars.
    Registered(PositionRegistryId),
}
\end{lstlisting}

\begin{requirement}
  \label{req:pitch:alteration-unit}
  For a \texttt{PitchSpacePosition::Cmn} enclosed by a
  \texttt{ScalePosition}, \texttt{alteration} \MUST{} be measured in steps of
  that pitch space's chromatic layer. Its zero point for each nominal \MUST{}
  be the space's \texttt{nominal\_to\_chromatic} mapping. A
  \texttt{PitchSpaceModification::CmnChromatic} value \MUST{} use that same
  unit in every pitch space whose accidental registry contains it. Thus one
  step is a semitone in \texttt{cmn-12} and a quarter-tone in
  \texttt{cmn-24}; a flat is respectively $-1$ and $-2$.

  When a bare \texttt{PitchSpacePosition::Cmn} is embedded as
  \texttt{ScoreTuningContext.reference.position}, the score's
  \texttt{default\_pitch\_space} \MUST{} supply the chromatic layer and mapping
  for this rule. Any other context embedding a bare CMN position \MUST{}
  identify the pitch space that denominates it.
\end{requirement}

\begin{requirement}
  \label{req:pitch:ji-vector-basis}
  For \texttt{PitchSpacePosition::JiVector}:
  \begin{itemize}
    \item \texttt{components.len()} \MUST{} equal the number of primes
      in the enclosing pitch space's declared prime basis.
    \item The basis ordering is established by the pitch space; the
      built-in JI pitch spaces (\texttt{ji-5limit},
      \texttt{ji-7limit}, \texttt{ji-11limit}) order primes
      ascending starting with 2.
    \item The vector $[a_0, a_1, a_2, \ldots]$ denotes the ratio
      $p_0^{a_0} \cdot p_1^{a_1} \cdot p_2^{a_2} \cdots$ where
      $p_i$ is the $i$-th prime in the basis. With the conventional
      basis $\{2, 3, 5, \ldots\}$, the first component holds
      octave information.
    \item Pitch spaces \MAY{} declare octave-reduced equivalence,
      in which case the first component is normalized to a canonical
      range. Without such a declaration, full register is preserved.
  \end{itemize}
\end{requirement}

\begin{requirement}
  \label{req:pitch:default-pitch-space}
  Every score \MUST{} select at least one pitch space from the
  built-in catalog (Section~\ref{sec:tuning:builtin}). The default
  pitch space for CMN scores is \texttt{cmn-12}, which is
  diatonic-over-chromatic in structure (seven diatonic nominals
  A--G over twelve chromatic positions), anchored to the score's
  reference pitch as resolved through the score tuning context
  (Chapter~\ref{ch:tuning}). Scores \MAY{} select additional pitch
  spaces from that catalog; pitches in different spaces cannot be
  directly compared and operations between them \MUST{} go through an
  explicit space-conversion mechanism.
\end{requirement}

\begin{rationale}
  Score-local pitch-space definition --- a score introducing a pitch
  space of its own, beyond selecting from the built-in catalog --- is
  deferred to a later schema major. The Chapter~\ref{ch:tuning} type
  surface (\texttt{PitchSpace}, \texttt{PositionStructure},
  \texttt{IntervalAlgebra}, and their kin) has never had a consumer;
  committing a registry for it to the wire now would freeze that
  surface permanently under \texttt{req:binfmt:frozen-layout}, before
  any implementation has exercised it. Selecting from the built-in
  catalog is the only reading the data model has ever supported: neither
  \texttt{Score} nor \texttt{ScoreTuningContext} carries a pitch-space
  or tuning-system registry, and none is added here.
\end{rationale}

\subsection{Octave Convention}
\label{sec:pitch:octave}

For the \texttt{Cmn} variant, the octave field uses Scientific Pitch
Notation: middle C is C4. This convention is normative throughout the
specification.

\begin{requirement}
  \label{req:pitch:scientific-pitch-octaves}
  The \texttt{octave} field of \texttt{PitchSpacePosition::Cmn} \MUST{}
  follow Scientific Pitch Notation. Middle C is C4. The lowest A on a
  standard piano is A0. C-1 is one octave below the lowest C on a
  standard piano.
\end{requirement}

\begin{rationale}
  Scientific Pitch Notation is the dominant modern convention in
  English-language music theory and notation software. MIDI octave
  numbering, where middle C is sometimes C3 and sometimes C5 depending
  on vendor, is rejected as a source of foot-guns. Helmholtz notation is
  expressible at the display layer; the data model commits to Scientific.
\end{rationale}

\subsection{The \texttt{CmnNominal} Type}
\label{sec:pitch:nominal}

\begin{lstlisting}[language=Rust]
#[repr(u8)]
pub enum CmnNominal { C = 0, D = 1, E = 2, F = 3, G = 4, A = 5, B = 6 }
\end{lstlisting}

The discriminants are normative: they define the diatonic step ordering
used by transposition algorithms.

\subsection{Transposition and the \texttt{Interval} Type}
\label{sec:pitch:interval}

Transposition acts on \emph{scale position}, not on acoustic realization.
It moves a pitch's analytical identity through the pitch space and leaves
the tuning context that sounds it untouched. A transposition therefore
requires no tuning catalog (Chapter~\ref{ch:tuning}) --- tuning decides what
frequency a scale position sounds at, not how to add an interval to one. It
requires an interval:

\begin{lstlisting}[language=Rust]
/// A signed interval: the transposition primitive. Defined here and
/// referenced by Instrument.transposition (Chapter 5), which is the
/// written-versus-sounding interval of a transposing instrument
/// (a B-flat clarinet is -1 diatonic, -2 chromatic).
pub struct TranspositionInterval {
    /// Signed diatonic steps: how far the nominal moves along C..B.
    /// A perfect fifth up is +4 (C -> G); an octave up is +7.
    pub diatonic_steps: i32,

    /// Signed steps in the enclosing pitch space's chromatic layer.
    /// A perfect fifth in cmn-12 is +7 and an octave is +12.
    /// The corresponding values are space-relative.
    pub chromatic_steps: i32,
}
\end{lstlisting}

Both components are load-bearing. The diatonic component fixes the
\emph{spelling} --- which nominal, hence which staff line the notehead
occupies --- and the chromatic component fixes the \emph{sound}. Neither
determines the other: an augmented second $(+1, +3)$ and a minor third
$(+2, +3)$ sound alike and are written differently, while a diminished
sixth $(+5, +7)$ and a perfect fifth $(+4, +7)$ are written a step apart
and sound the same. A scalar chromatic-step count cannot express the
difference, and an implementation carrying only one component is not
transposing but altering.

\begin{requirement}
  \label{req:pitch:transposition}
  Let $n$ be the \texttt{nominal}'s normative discriminant
  (Section~\ref{sec:pitch:nominal}). Let the enclosing pitch space have a
  \texttt{DiatonicOverChromatic} structure with $C$ chromatic positions per
  octave and nominal mapping $m$. A
  \texttt{PitchSpacePosition::Cmn} position's \emph{absolute chromatic
  coordinate} is
  \[
    s \;=\; m(\texttt{nominal}) \;+\; \texttt{alteration}
        \;+\; C \cdot \texttt{octave}.
  \]
  Transposing that position by \texttt{TranspositionInterval \{ d, c \}},
  where $c$ is measured in steps of the same chromatic layer, \MUST{} produce
  \begin{align*}
    \texttt{nominal}' &= \textsc{CmnNominal}\big((n + d) \bmod 7\big), \\
    \texttt{octave}'  &= \texttt{octave} + \lfloor (n + d) / 7 \rfloor, \\
    \texttt{alteration}' &= (s + c) -
      \big(m(\texttt{nominal}') + C \cdot \texttt{octave}'\big),
  \end{align*}
  where $\bmod$ and $\lfloor \cdot \rfloor$ are Euclidean (the remainder is
  non-negative). The diatonic component alone selects the nominal and the
  octave; the alteration absorbs exactly the residue.

  A transposition \MUST{} be refused --- atomically, changing nothing --- when
  any target satisfies any of:
  \begin{itemize}
    \item its \texttt{scale\_position.position} is not \texttt{Cmn}, so the
      \texttt{Interval} above has no defined action on it;
    \item its enclosing pitch space cannot be resolved to a compatible
      \texttt{DiatonicOverChromatic} structure and nominal mapping;
    \item its \texttt{acoustic.realization} is
      \texttt{AcousticRealization::AbsoluteHz}, which overrides the tuning
      system, so moving the scale position would move the notehead without
      moving the sound;
    \item the computed $\texttt{alteration}'$ or $\texttt{octave}'$ does not
      fit its field.
  \end{itemize}
  An implementation \MUSTNOT{} saturate, clamp, guess a pitch-space
  structure, or partially apply a transposition in any of these cases.
\end{requirement}

\begin{requirement}
  \label{req:pitch:space-capability-refusal}
  An implementation \MUST{} fail closed when it cannot establish the
  pitch-space structure required to interpret a position. In particular,
  transposition \MUST{} report a refusal when the enclosing CMN chromatic
  layer or nominal mapping is unavailable, and a computation claiming a
  12-chromatic pitch class or absolute semitone \MUST{} report
  unavailability unless it can establish a twelve-position chromatic layer.
  The \texttt{Cmn} position discriminant alone \MUSTNOT{} be treated as proof
  of either capability.
\end{requirement}

\begin{rationale}
  The refusal cases are the four ways a transposition can silently lie. A
  non-CMN position has no nominal to move, so a chromatic shift would have to
  invent one. An unresolved pitch space supplies neither the octave size nor
  the nominal mapping, so arithmetic would have to guess both.
  An \texttt{AbsoluteHz} pitch declares its own frequency: it can be respelled
  or it can be resounded, but a transposition that claims to do both does one
  and pretends to the other. And saturation is the worst of the four, because
  it is invisible --- a transposition that clamps has produced a pitch nobody
  asked for, reported success, and destroyed the information needed to notice.

  Refusing atomically, rather than skipping the offending target, follows from
  the same reasoning: a chord transposed except for one note is not a
  transposed chord, it is a different chord. Contrast the treatment of
  \emph{tombstoned} targets, which are skipped rather than refused --- a
  deleted pitch is not an untransposable pitch, it is a pitch the operation
  has nothing to say about.
\end{rationale}

\section{Acoustic Realization}
\label{sec:pitch:acoustic}

The acoustic layer specifies the sounding frequency. It does not carry
a frequency directly; instead it references a tuning system that, given
the pitch's scale position, resolves to a frequency.

\begin{lstlisting}[language=Rust]
pub struct AcousticPitch {
    /// The tuning system governing this pitch's frequency.
    /// May be inherited; see Section~\ref{sec:pitch:tuning-inherit}.
    pub tuning: TuningReference,

    /// How the tuning system resolves to a frequency.
    pub realization: AcousticRealization,
}

pub enum AcousticRealization {
    /// Resolve through the tuning system using the scale position alone.
    /// This is the default case for ordinary CMN.
    Implicit,

    /// An explicit offset in cents from what the tuning system would
    /// otherwise produce. For inflections, expressive retuning, or
    /// just-intonation adjustments within a tempered piece.
    CentsOffset(f64),

    /// An explicit absolute frequency in Hertz, overriding the tuning
    /// system. For spectral music or fixed-frequency electronic parts.
    AbsoluteHz(f64),
}
\end{lstlisting}

\subsection{Cents as the Offset Unit}

Cents (one cent equals one hundredth of an equal-tempered semitone) are
the normative offset unit. \texttt{f64} provides far more precision than
the ear can resolve in any context.

\subsection{Tuning Reference and Inheritance}
\label{sec:pitch:tuning-inherit}

A \texttt{TuningReference} is either an explicit tuning-system identifier
or an inheritance marker:

\begin{lstlisting}[language=Rust]
pub enum TuningReference {
    /// Inherit the tuning system from the enclosing scope.
    Inherit,

    /// Explicitly named tuning system.
    Explicit(TuningSystemId),
}
\end{lstlisting}

Resolution proceeds by walking outward: pitch, then voice, then staff,
then region, then score. The score's tuning system \MUST{} be explicit;
\texttt{Inherit} at the score level is a malformed document. Tuning
systems and their resolution rules are defined in detail in
Chapter~\ref{ch:tuning}.

\subsection{Reference Pitch}
\label{sec:pitch:reference}

A pitch's \texttt{TuningReference} resolves, through the hierarchical
walk, to a tuning system. The \emph{reference pitch} (the position
and frequency that anchor that tuning system's structure to absolute
Hertz) is resolved separately through the same hierarchical walk, and
is a property of the score's tuning context, not of the tuning system
itself.

The pitch type does not own or imply a reference frequency. Two
scores using the same tuning system (12-TET) but different references
(A4 = 440\,Hz versus A4 = 415\,Hz) share their tuning structure
entirely; only their reference pitch differs.

The full reference-pitch model, including the \texttt{ReferencePitch}
type and the score tuning context that owns it, is specified in
Chapter~\ref{ch:tuning} (Section~\ref{sec:tuning:reference}).

\section{Spelling}
\label{sec:pitch:spelling}

Spelling determines what the performer sees on the page for a given
pitch: the staff position, the accidental glyph (if any), and any
ancillary notational marks (parenthesized accidentals, courtesy
accidentals, microtonal accidentals).

Spelling \MUSTNOT{} be a field on the pitch type. Spellings are stored
externally, indexed by pitch identifier, and managed by the spelling
subsystem defined in this section.

\begin{rationale}
  Spelling as a field on pitch causes destructive overwrites on
  respelling, prevents partial spelling (where only some pitches have
  user-chosen spellings and the rest are inferred), prevents multiple
  spellings per pitch (analytical layers), and prevents provenance
  tracking. Externalizing spelling resolves all four.
\end{rationale}

\subsection{The Spelling Attachment}

\begin{lstlisting}[language=Rust]
pub struct SpellingAttachment {
    /// What this attachment applies to: a single pitch or a scope.
    pub scope: SpellingScope,

    /// The directive: explicit spelling or a rule for inference.
    pub directive: SpellingDirective,

    /// Provenance of this attachment.
    pub source: SpellingSource,

    /// Tie-break priority among attachments. Higher wins. See
    /// Section~\ref{sec:pitch:precedence} for the configurable
    /// precedence model.
    pub priority: i32,

    /// Optional analytical layer. None means the engraved layer;
    /// Some(id) means an alternative analytical layer that does not
    /// affect the engraved score.
    pub layer: Option<AnalysisLayerId>,
}

pub enum SpellingScope {
    /// Applies to one specific pitch.
    Pitch(PitchId),

    /// Applies to all pitches matching the selector within a time range.
    Range {
        start: TimeAnchor,
        end: TimeAnchor,
        voices: VoiceSelector,
    },
}

pub enum SpellingDirective {
    /// An explicit spelling for a single pitch. Only valid with
    /// SpellingScope::Pitch.
    Explicit(PitchSpelling),

    /// A rule for inferring spellings within a scope.
    Rule(SpellingRule),
}

pub struct PitchSpelling {
    /// The nominal: the staff position. For CMN, this is one of A-G.
    /// For other grammars, an identifier into the grammar's nominal
    /// registry.
    pub nominal: SpellingNominal,

    /// Stack of accidentals applied to the nominal, in stacking
    /// order: innermost (closest to the notehead) first.
    /// Empty vector means no accidental glyph is drawn, which does
    /// not necessarily mean "natural"; see
    /// Section~\ref{sec:pitch:absent-accidental}.
    pub accidentals: Vec<AccidentalId>,

    /// The octave designation.
    pub octave: i8,

    /// Optional rendering hints: parenthesized, cautionary, editorial,
    /// small-print, etc. Non-normative for sounding pitch.
    pub render_hints: SpellingRenderHints,
}

pub enum SpellingNominal {
    /// CMN nominals: the fast path.
    Cmn(CmnNominal),

    /// Integer nominals for N-tone systems.
    Integer(i32),

    /// Registered nominals for grammar-specific systems.
    Registered(NominalRegistryId),
}
\end{lstlisting}

\subsubsection{Accidental Stack Semantics}
\label{sec:pitch:accidental-stack}

\begin{requirement}
  \label{req:pitch:accidental-stack}
  The \texttt{accidentals} vector \MUST{} be interpreted as follows:

  \begin{itemize}
    \item The stored order is normative for engraving: index 0 is
      the innermost glyph (closest to the notehead), with subsequent
      indices placed further outward.
    \item For built-in additive accidental systems, the pitch-space
      modification contributed by the stack is the sum of the
      modifications of its members, and is independent of stored
      order. Reordering the stack changes engraving appearance but
      \MUSTNOT{} change sounding pitch.
    \item Custom accidental systems \MAY{} declare non-additive
      combination semantics only through their registered
      \texttt{AccidentalCombination} (Section~\ref{sec:tuning:accidentals}).
      Otherwise, accidentals are treated as commutative under their
      pitch effect.
    \item Two accidentals with the same \texttt{AccidentalId} in one
      stack constitute a malformed spelling, unless the accidental's
      registered combination semantics explicitly permits repetition.
      The conventional double-sharp and double-flat are represented
      as single accidental definitions, not as repeated singles.
    \item An empty \texttt{accidentals} vector means no glyph is
      drawn; it is distinct from a vector containing only a natural
      sign (which explicitly cancels prior alterations and is drawn).
  \end{itemize}
\end{requirement}

\subsection{Spelling Sources}

Every spelling attachment carries provenance:

\begin{lstlisting}[language=Rust]
pub enum SpellingSource {
    /// The user explicitly chose this spelling.
    UserChosen,

    /// Inferred by the spelling pre-pass from key signature and
    /// context. Lowest default precedence. The pre-pass does not
    /// store attachments (its output is a derived annotation);
    /// this variant ranks that output in the precedence order.
    Inferred,

    /// Imported from a foreign format.
    Imported { format: ForeignFormatId },

    /// Propagated from a transposition or other edit operation.
    Propagated { from: PitchId },

    /// An analytical spelling on a non-engraved layer.
    Analytical,
}
\end{lstlisting}

\begin{requirement}
  \label{req:pitch:spelling-provenance}
  Editing operations that respell a pitch \MUST{} produce attachments
  with source \texttt{UserChosen}. Editing operations that move or
  transpose a pitch \MUST{} produce attachments with source
  \texttt{Propagated}. Foreign-format importers \MUST{} produce
  attachments with source \texttt{Imported}. The spelling pre-pass
  (Section~\ref{sec:pitch:prepass}) \MUSTNOT{} mint attachments: its
  inferred spellings are derived annotations reported with
  \texttt{Inferred} provenance, promoted to a stored
  \texttt{UserChosen} attachment only by an explicit editing
  operation.
\end{requirement}

\subsection{Configurable Precedence}
\label{sec:pitch:precedence}

When multiple attachments target the same pitch on the same analysis
layer, the renderer resolves the conflict by precedence. The default
precedence is:

\begin{enumerate}
  \item \texttt{UserChosen} (highest)
  \item \texttt{Imported}
  \item \texttt{Propagated}
  \item \texttt{Inferred} (lowest)
\end{enumerate}

\begin{requirement}
  \label{req:pitch:spelling-precedence}
  Every score \MUST{} carry a \texttt{SpellingPrecedence} configuration.
  The configuration \MUST{} assign a total ordering over the variants of
  \texttt{SpellingSource}. The configuration \MAY{} differ from the
  default; for example, a workflow that treats MusicXML imports as
  authoritative may rank \texttt{Imported} above \texttt{UserChosen}.
  Conflicts are broken first by precedence, then by the \texttt{priority}
  field of the attachment (higher wins), then by canonical attachment
  order (the attachment appearing earliest in the score's canonical
  serialization wins). Attachments carry no creation timestamp: a
  timestamp tie-break would hang resolution on non-canonical state.
\end{requirement}

\subsection{The Spelling Pre-Pass}
\label{sec:pitch:prepass}

When a score is materialized (on load, after an edit, or ahead of
rendering), the spelling pre-pass computes an inferred spelling for
every spelling-eligible pitch. Its output is a \emph{canonical
derived annotation}: a deterministic function of the materialized
score graph, the active profile, and the versioned
\texttt{SpellingAlgorithmId}, recomputed on materialization and
never stored as graph state or serialized into canonical chunks.
Where a stored attachment targets the pitch, the resolved annotation
reports that attachment's spelling and provenance per the precedence
configuration (Section~\ref{sec:pitch:precedence}); the pre-pass
supplies the \texttt{Inferred}-provenance spelling that applies when
no attachment outranks it. The pre-pass \MUST{} be deterministic:
given the same materialized graph, precedence configuration, and
algorithm version, it \MUST{} produce identical annotations.

The pre-pass operates as follows:

\begin{enumerate}
  \item Collect all pitches in the score, partitioned by voice and
    measure.
  \item For each partition, in time order, resolve each pitch's spelling
    using the applicable context inputs (which of these a given
    algorithm version consumes is part of its versioned definition;
    see Requirement~\ref{req:pitch:spelling-algorithm}):
    \begin{enumerate}
      \item The active key signature for the staff.
      \item The accidental context of prior pitches in the same measure
        and voice.
      \item The melodic context (preceding interval).
      \item The harmonic context (concurrent pitches in other voices).
      \item Any active scope-level \texttt{Rule} attachments.
    \end{enumerate}
  \item Report an \texttt{Inferred}-provenance derived annotation for
    each pitch.
\end{enumerate}

Because derived annotations are not stored, a change to the algorithm
version deterministically invalidates them without any state
migration. Implementations \MAY{} cache derived annotations and
\MAY{} recompute incrementally (re-running only the affected measure
and voice, and forward through measures whose context is altered,
until the context stabilizes); a cache \MUST{} be invalidated when
the derivation key --- graph, profile, algorithm version --- changes,
and neither caching nor incrementality may be observable in the
annotations produced.

\begin{rationale}
  Deriving inferred spellings---rather than storing them as
  attachments---keeps algorithm output out of canonical state: two
  replicas at the same (graph, profile, algorithm version) agree on
  annotations byte-for-byte without ever exchanging them, an
  algorithm upgrade cannot strand stale inferred state, and
  serialization never churns on re-inference. Inferred spellings
  remain inspectable in the UI and editable by promotion: accepting
  one mints a \texttt{UserChosen} attachment through an editing
  operation, which then outranks the pre-pass.
\end{rationale}

\begin{requirement}
  \label{req:pitch:spelling-algorithm}
  \textbf{Spelling algorithm disposition (ratified Pass 12).}
  The spelling pre-pass receives the \emph{profile-declared by
  versioned identifier} disposition of Appendix~\ref{app:determinism}
  Section~\ref{sec:det:open}. The reserved identifier
  \texttt{SpellingAlgorithmId} \texttt{"default"} denotes, at
  version~1, a Temperley-style line-of-fifths preference
  algorithm: each partition's tonal context is estimated as a
  centre of gravity on the line of fifths from the accidental,
  melodic, and harmonic context, and each pitch takes the
  enharmonic spelling nearest that centre. The ascending-runs-take-
  sharps / descending-runs-take-flats convention participates
  \emph{only as a tiebreak} within the centre-of-gravity rule; an
  isolated chromatic run with no tonal context may therefore
  receive the enharmonic the convention alone would not pick, and
  a voice-leading refinement is a future algorithm version, not a
  deviation. Which of the context inputs enumerated above an
  algorithm version consumes is part of its versioned definition:
  \texttt{"default"} version~1 does \emph{not} consult declared
  key signatures (a key-aware refinement is a future version) and
  is region-time-model-independent --- pitches in aleatoric and
  proportional regions are spelled by the same rule as metric
  ones, and no region-specific spelling pass exists. A profile
  requesting any other identifier \MUST{} error;
  implementations \MUSTNOT{} silently substitute a different
  algorithm.
\end{requirement}

\begin{requirement}
  \label{req:pitch:authored-uninferred}
  \textbf{Authored annotations for inference-ineligible targets
  (ratified Pass 12).}
  A stored spelling attachment whose target the pre-pass produces
  \emph{no} inferred spelling for (an inference-ineligible pitch)
  \MUST{} still surface in the resolved derived annotations: the
  winning attachment under the precedence configuration
  (Section~\ref{sec:pitch:precedence}) is reported with its
  authored provenance, exactly as it would have been had it
  outranked an inferred spelling. An authored attachment is
  precisely how a user notates what the algorithm cannot infer;
  the derived-annotation surface \MUSTNOT{} render it invisible.
  Annotation taxonomy counts these authored-only resolutions
  distinctly from authored-over-inferred overrides. The same rule
  applies to the decomposition pre-pass
  (Section~\ref{sec:time:notrhythm}).
\end{requirement}

\subsection{Absent Accidentals}
\label{sec:pitch:absent-accidental}

A \texttt{PitchSpelling} with an empty \texttt{accidentals} vector
means no accidental glyph is drawn at this position. This is distinct
from a vector containing an explicit natural sign. The distinction
matters in two cases:

\begin{itemize}
  \item A pitch on a line affected by the key signature carries no
    accidental and sounds the key-signature pitch; drawing a natural
    would be a respelling.
  \item A pitch that follows an in-measure accidental on the same line
    and octave inherits that accidental; an empty vector means
    inheritance applies, while a vector containing a natural cancels
    the inheritance.
\end{itemize}

\begin{requirement}
  \label{req:pitch:absent-versus-natural}
  The renderer \MUST{} distinguish between an empty
  \texttt{accidentals} vector (no glyph) and a vector containing only
  a natural (explicit natural glyph) when materializing a spelling.
  Promotion from empty to natural-only, or vice versa, is a
  respelling operation.
\end{requirement}

\subsection{Analytical Layers}
\label{sec:pitch:layers}

A pitch may carry multiple spelling attachments on different analytical
layers. The engraved layer (\texttt{layer: None}) governs what the
performer sees. Other layers (\texttt{layer: Some(AnalysisLayerId)})
record alternative spellings used by analytical views, theory tools,
or pedagogical overlays.

The \texttt{AnalysisLayer} type itself is defined in
Chapter~\ref{ch:graph} (Section~\ref{sec:graph:layers}); a layer
identifier here refers to that authoritative definition.

\begin{requirement}
  \label{req:pitch:spelling-layer-fallback}
  Resolution of a pitch's spelling for a given view \MUST{} consider
  only attachments whose \texttt{layer} matches the view's active layer.
  The engraved view's active layer is \texttt{None}. Analytical views
  may set an explicit layer; if no attachment on that layer matches,
  the view \MUST{} fall back to the engraved layer.
\end{requirement}

\section{Pitch Identifiers}
\label{sec:pitch:ids}

Every pitch in the score graph carries a stable identifier:

\begin{lstlisting}[language=Rust]
pub struct PitchId(pub u128);
\end{lstlisting}

\begin{requirement}
  \label{req:pitch:pitch-id-stability}
  Pitch identifiers \MUST{} be stable across edits: a pitch's identifier
  is assigned at creation and never reassigned. Pitch identifiers
  \MUST{} be unique within a score. Pitch identifiers \MUST{} be
  collision-resistant across distributed editing scenarios; the
  recommended generation strategy is a 128-bit value combining a
  per-replica identifier and a monotonic counter, with the exact format
  defined in Chapter~\ref{ch:semops}.
\end{requirement}

\section{What Pitch Is Not}
\label{sec:pitch:not-pitch}

Several musical phenomena resemble pitches but are not represented by
the pitch type:

\begin{description}
  \item[Unpitched events.] Snare drum hits, claves, cymbal strokes, and
    other unpitched percussion are sounding events with staff positions
    but no pitched identity. They are represented by a separate
    \texttt{UnpitchedEvent} type (see Chapter~\ref{ch:graph}). They have
    no scale position and no acoustic pitch in the sense defined here.
  \item[Rests.] A rest is a duration without a sounding event. Rests are
    represented by a separate \texttt{Rest} type. They carry timing and
    voice information but no pitch.
  \item[Indeterminate-pitch events.] Events whose pitch is deliberately
    indeterminate (text instructions like ``highest possible note,''
    aleatoric pitch material, certain graphic-notation events) are
    represented by an \texttt{IndeterminatePitchEvent} type, which may
    carry hints (range, contour, density) but no concrete pitch.
  \item[Pitch trajectories.] Glissandi, portamenti, and continuous pitch
    bends are not single pitches but trajectories between or around
    pitches. They are represented as relationships between pitched
    events plus a trajectory shape, defined in Chapter~\ref{ch:graph}.
\end{description}

\begin{rationale}
  Forcing every sounding event into a pitch-shaped hole creates
  invariant-breaking edge cases throughout the codebase. Sibling types
  under a common \texttt{SoundingEvent} (or similar) abstraction keep
  each type coherent and let operations branch cleanly.
\end{rationale}

\section{Equality and Comparison}
\label{sec:pitch:eq}

Two pitches are \emph{structurally equal} if and only if their scale
positions and acoustic realizations are equal field-by-field. Structural
equality is the default \texttt{Eq} implementation.

Three additional equivalence relations are defined:

\begin{description}
  \item[Sounding equivalence.] Two pitches are sounding-equivalent if
    they resolve to the same frequency under their respective tuning
    systems. Computed, not derived from the type.
  \item[Scale-position equivalence.] Two pitches are scale-position-
    equivalent if their \texttt{ScalePosition} fields are equal, ignoring
    acoustic realization.
  \item[Enharmonic equivalence.] Two pitches are enharmonically
    equivalent if they are sounding-equivalent under 12-tone equal
    temperament, regardless of their actual tuning system. Used for
    foreign-format export and certain analytical operations.
\end{description}

\begin{requirement}
  \label{req:pitch:pitch-equivalence-functions}
  Implementations \MUST{} provide functions for each of the three
  computed equivalences. Implementations \MUSTNOT{} conflate structural
  equality with sounding or enharmonic equivalence.
\end{requirement}

\section{Forward References}

\begin{itemize}
  \item Tuning systems, pitch spaces, and accidental registries are
    defined in Chapter~\ref{ch:tuning}.
  \item The score graph, including how pitches are embedded in events,
    voices, and staves, is defined in Chapter~\ref{ch:graph}.
  \item \texttt{TimeAnchor}, \texttt{VoiceSelector}, \texttt{PitchId}
    generation, and the reduction rules governing spelling
    attachments under concurrent edits are defined in
    Chapter~\ref{ch:semops}.
\end{itemize}

% ===========================================================================
\chapter{Time and Duration}
\label{ch:time}

This chapter specifies the time and duration primitives: how positions
and durations are represented, how time signatures and tempo maps are
modelled, how notational rhythm is decomposed from sounding duration,
how tuplets and other groupings work, and how the core supports metric,
proportional, and aleatoric time models simultaneously.

\section{Design Principles}
\label{sec:time:principles}

\begin{description}
  \item[Exact rational arithmetic.] Musical positions and durations
    \MUST{} be exact rationals. Floating-point time is rejected as a
    source of accumulated error; integer-tick time is rejected as
    incompatible with arbitrary tuplet nesting. Exactness is
    non-negotiable.
  \item[Two clocks.] Musical time (measured in whole-note rationals)
    and wall-clock time (measured in fixed-point seconds) are distinct
    primitives. Conversion between them is a pure function of the
    tempo map.
  \item[Position and duration are distinct types.] A position and a
    duration share an underlying representation but are not
    interchangeable. The algebra (position $+$ duration $\rightarrow$ position;
    position $-$ position $\rightarrow$ duration; position $+$ position is
    undefined) is enforced at the type level.
  \item[Sounding duration is data; notational rhythm is derived.]
    A note's sounding duration is intrinsic; its notated decomposition
    (notehead values, ties, dots, tuplet membership) is produced by a
    deterministic pre-pass analogous to the spelling pre-pass and is
    overridable per-event by an authored attachment.
  \item[Multiple time models coexist.] A score \MUST{} support metric,
    proportional, and aleatoric regions, possibly concurrent on
    different staves. The region's declared time model determines the
    interpretation of positions within it.
  \item[Time anchors, not absolute positions.] Stored references to
    points in time \MUST{} anchor to identified objects (events,
    measures, regions) plus offsets. Absolute positions \MUSTNOT{}
    be stored where they could shift under edits.
\end{description}

\section{The Rational Time Type}
\label{sec:time:rational}

The fundamental time primitive is an exact rational number. The
specification states the contract; implementations are free to choose
representations meeting it.

\begin{requirement}
  \label{req:time:exact-rational-time}
  The rational time type \MUST{} behave as an exact rational with at
  least 32-bit numerator and 32-bit denominator range in normalized
  form. Operations \MUST{} produce exact results; on representation
  overflow, the implementation \MUST{} either widen to an exact
  unbounded representation or return an error. Silent loss of precision
  is forbidden. Implementations \MAY{} use any representation that
  meets this contract, including packed inline-or-heap-promoted designs.
\end{requirement}

\subsection{Recommended Implementation: Inline-or-Promoted}
\label{sec:time:impl}

The recommended implementation packs the common case (denominators
under $2^{32}$) into 8 bytes and promotes to a heap-allocated
arbitrary-precision rational on overflow. This is normative for
conformance only inasmuch as the observable behavior must match;
it is illustrative as a reference design.

\begin{lstlisting}[language=Rust]
/// A musical time value: position or duration. Exact rational.
pub enum RationalTime {
    /// Inline case: fits in 8 bytes.
    Small(SmallRational),

    /// Promoted case: heap-allocated arbitrary-precision rational.
    /// Used only when arithmetic overflows the inline range.
    Large(Arc<BigRational>),
}

/// The inline rational: numerator i32, denominator NonZeroU32.
/// Always normalized: gcd(|numerator|, denominator) == 1.
#[derive(Copy, Clone, PartialEq, Eq, Hash)]
pub struct SmallRational {
    numerator: i32,
    denominator: NonZeroU32,
}
\end{lstlisting}

\begin{requirement}
  \label{req:time:rational-normalization}
  All instances of the small rational \MUST{} be normalized: the
  greatest common divisor of the absolute value of the numerator and
  the denominator \MUST{} equal one, and the denominator \MUST{} be
  strictly positive. The constructor \MUST{} enforce normalization.
\end{requirement}

\subsection{Promotion and Demotion}

Arithmetic operations on two small rationals are performed in widened
intermediate types ($i64$ numerator, $u64$ denominator) and the
normalized result is examined. If the result fits in the small
representation, it is returned as \texttt{Small}; otherwise it is
returned as \texttt{Large}.

\begin{requirement}
  \label{req:time:rational-promotion}
  Arithmetic that exceeds the small representation's range \MUST{}
  silently promote to the large representation. The user of the type
  \MUST{} observe no behavioral difference between small-only and
  small-or-large arithmetic; the contract is exact arithmetic on
  exact rationals.

  Implementations \MAY{} demote a \texttt{Large} value back to
  \texttt{Small} when its normalized representation fits, but
  \MUSTNOT{} make demotion observable (e.g., by changing
  equality or hash behavior across the variant boundary).
\end{requirement}

\subsection{Equality and Ordering}

Two rational time values are equal if and only if their normalized
representations represent the same rational number. Equality between
\texttt{Small} and \texttt{Large} variants \MUST{} compare the
underlying values, not the variants.

Ordering is the standard rational ordering. Comparison between rationals
in the small representation is performed by cross-multiplication in
widened intermediate types to avoid overflow.

\subsection{Units}

The unit of musical time is the \emph{whole note}. A quarter note has
duration $\frac{1}{4}$. A triplet eighth has duration $\frac{1}{12}$.
A quintuplet sixteenth inside a duplet half (an unusual but legal
construct) has duration $\frac{1}{2} \cdot \frac{1}{2} \cdot \frac{1}{5}
= \frac{1}{20}$.

\begin{requirement}
  \label{req:time:whole-note-unit}
  The whole note \MUST{} be the unit of musical time. The duration of
  a whole note is the rational $\frac{1}{1}$. All other durations are
  expressed as rationals relative to the whole note.
\end{requirement}

\begin{rationale}
  The whole note is the only unit under which durations form a
  multiplicative algebra with $1$ as identity. Quarter-note units would
  make ``the duration of a whole note'' equal to $4$, which is
  awkward; PPQ-tick units would force a choice of resolution. The
  whole-note unit also matches the etymology of note-value names: a
  half note is half of a whole note.
\end{rationale}

\section{Position and Duration as Distinct Types}
\label{sec:time:newtypes}

The rational time representation is shared between two distinct types:
\texttt{MusicalPosition} and \texttt{MusicalDuration}. Both wrap a
\texttt{RationalTime}. Their distinction is enforced at the type level
to prevent algebraic errors.

\begin{lstlisting}[language=Rust]
/// A point in musical time, relative to the origin of a time region.
#[derive(Clone, PartialEq, Eq, Hash, PartialOrd, Ord)]
pub struct MusicalPosition(RationalTime);

/// A span of musical time.
#[derive(Clone, PartialEq, Eq, Hash, PartialOrd, Ord)]
pub struct MusicalDuration(RationalTime);

impl Add<MusicalDuration> for MusicalPosition {
    type Output = MusicalPosition;
    // ...
}

impl Add<MusicalDuration> for MusicalDuration {
    type Output = MusicalDuration;
    // ...
}

impl Sub<MusicalPosition> for MusicalPosition {
    type Output = MusicalDuration;
    // ...
}

// MusicalPosition + MusicalPosition is intentionally not implemented.
\end{lstlisting}

\begin{requirement}
  \label{req:time:position-duration-types}
  Implementations \MUST{} provide distinct types for position and
  duration. The type system \MUST{} prevent adding two positions. The
  difference of two positions \MUST{} yield a duration; the sum of a
  position and a duration \MUST{} yield a position.
\end{requirement}

\begin{rationale}
  The illustrative Rust signatures derive \texttt{Clone}, not
  \texttt{Copy}, because \texttt{RationalTime} contains an
  \texttt{Arc<BigRational>} variant that cannot satisfy
  \texttt{Copy}. Implementations \MAY{} provide \texttt{Copy}
  semantics through a compact inline-or-pointer-tagged representation
  whose \texttt{Large} case is a non-owning interior pointer rather
  than an \texttt{Arc}; such implementations satisfy the contract
  here. Conformance does not require \texttt{MusicalPosition} or
  \texttt{MusicalDuration} to be \texttt{Copy}.
\end{rationale}

\section{Wall-Clock Time}
\label{sec:time:wallclock}

Wall-clock time is the time measured by the playback engine, by SMPTE
timecode, and by external audio and video sources. It is distinct from
musical time and is represented as a fixed-point integer count of a
sub-second unit.

\begin{lstlisting}[language=Rust]
/// A point in wall-clock time, in nanoseconds from a region origin.
/// 64-bit signed: range +/- ~292 years.
#[derive(Copy, Clone, PartialEq, Eq, Hash, PartialOrd, Ord)]
pub struct WallClockTime(i64);  // nanoseconds

#[derive(Copy, Clone, PartialEq, Eq, Hash, PartialOrd, Ord)]
pub struct WallClockDuration(i64);  // nanoseconds
\end{lstlisting}

\begin{requirement}
  \label{req:time:nanosecond-wallclock}
  Wall-clock time \MUST{} be a fixed-point integer count of
  nanoseconds. Floating-point wall-clock time is forbidden in stored
  data. Implementations \MAY{} convert to floating-point for transient
  computation provided no stored or serialized value loses precision.
\end{requirement}

\begin{rationale}
  Nanosecond resolution provides over four orders of magnitude more
  precision than the human ear can resolve and aligns with the
  resolution conventions of professional audio interfaces and
  SMPTE-derived clocks. The 292-year range is sufficient for any
  conceivable musical work.
\end{rationale}

\subsection{Wall-Clock and Musical Time Coexist}

A score event may be anchored in either musical or wall-clock time.
Musical-time anchors are the default for ordinary notation; wall-clock
anchors are used for film hit points, audio reference markers, fixed
electronic parts in scores with live ensemble, and similar
sync-critical content. Both anchor types may coexist in the same
score, possibly on the same staff.

Conversion between musical and wall-clock time is performed through
the tempo map (Section~\ref{sec:time:tempomap}).

\section{Time Anchors}
\label{sec:time:anchors}

Stored references to points in time \MUST{} be expressed as
\emph{time anchors}: references to identified objects plus offsets.
Absolute positions \MUSTNOT{} be stored where they could be invalidated
by edits to the score.

\begin{lstlisting}[language=Rust]
pub enum TimeAnchor {
    /// Anchored to a specific event. Survives edits that do not
    /// delete the event.
    Event { id: EventId, offset: AnchorOffset },

    /// Anchored to a measure boundary. Survives measure reordering
    /// (the anchor follows the measure).
    Measure {
        id: MeasureId,
        position: MeasurePosition,
        offset: AnchorOffset,
    },

    /// Anchored to the start or end of a region.
    Region { id: RegionId, edge: RegionEdge, offset: AnchorOffset },

    /// Anchored to absolute wall-clock time within the score.
    /// Used for film and audio sync.
    WallClock { time: WallClockTime },
}

/// An offset applied to an anchor target. The variant kind is
/// constrained by the target's enclosing region's time model.
pub enum AnchorOffset {
    /// Offset in musical time. Valid for targets in metric regions
    /// and for aleatoric regions whose anchoring discipline
    /// admits musical coordinates.
    Musical(MusicalDuration),

    /// Offset in wall-clock time. Valid for targets in proportional
    /// regions and for aleatoric regions whose anchoring discipline
    /// admits wall-clock coordinates.
    WallClock(WallClockDuration),

    /// No offset; the anchor refers to the target's reference point
    /// exactly. Valid in any region.
    Zero,
}

pub enum MeasurePosition { Start, End }
pub enum RegionEdge { Start, End }
\end{lstlisting}

\begin{requirement}
  \label{req:time:anchor-references}
  Stored \emph{references to external time points} (cross-cutting
  endpoints, marker locations, attachment anchors, anything
  pointing at an object outside the referencing object's own
  voice and region) \MUST{} use the \texttt{TimeAnchor} type.
  Absolute musical positions \MUSTNOT{} appear as such references
  in stored data: a reference to ``measure 47, beat 3'' that does
  not survive a measure insertion in measure 12 is a broken
  reference.

  An event's own position within its owning voice and region
  \MAY{} be stored as an \texttt{EventPosition} (carrying a
  \texttt{MusicalPosition} for events in metric regions, or a
  \texttt{WallClockTime} for events in proportional regions,
  etc.). Event-local positions are part of the event's intrinsic
  data, not references to other objects. Operations that shift a
  voice or region update the contained events' positions
  directly; they do not need anchor indirection because the
  containment relationship makes the dependency explicit.

  Implementations \MAY{} compute and cache absolute positions
  from anchors for performance; such caches \MUST{} be
  invalidated when their underlying anchored objects change.
\end{requirement}

\begin{requirement}
  \label{req:time:anchor-offset-kind}
  The variant of an \texttt{AnchorOffset} \MUST{} agree with the
  time model of the anchor target's enclosing region:

  \begin{itemize}
    \item For targets in metric regions: \texttt{AnchorOffset::Musical}
      or \texttt{AnchorOffset::Zero}.
    \item For targets in proportional regions:
      \texttt{AnchorOffset::WallClock} or \texttt{AnchorOffset::Zero}.
    \item For targets in aleatoric regions: any offset variant
      consistent with the region's
      \texttt{AleatoricAnchoringDiscipline}
      (Section~\ref{sec:graph:aleatoric-discipline}).
  \end{itemize}

  Implementations \MUST{} reject anchors whose offset variant
  contradicts the target region's time model.
\end{requirement}

\subsection{Anchor Resolution}

Resolving a time anchor to an absolute position is a pure function of
the score graph: locate the anchored object, query its position within
its region, and add the offset. Resolution \MUST{} be deterministic.

\subsection{Anchor Stability Under Edits}

The chosen anchor type determines what edits the anchor survives:

\begin{itemize}
  \item An \texttt{Event} anchor survives any edit that does not delete
    the anchored event. Edits that move the event (insertions before
    it, tempo changes) move the anchor with it.
  \item A \texttt{Measure} anchor survives insertions and deletions of
    other measures, and survives reordering. It is invalidated only by
    deletion of the anchored measure.
  \item A \texttt{Region} anchor survives edits within the region. It
    is invalidated only by deletion of the region.
  \item A \texttt{WallClock} anchor is unaffected by musical edits but
    may correspond to different musical positions after tempo changes.
\end{itemize}

\begin{requirement}
  \label{req:time:orphaned-anchor-handling}
  When an anchored object is deleted, anchors targeting it \MUST{} be
  either re-anchored to a surviving object or marked as orphaned. The
  semantic operations chapter (Chapter~\ref{ch:semops}) specifies
  re-anchoring rules per operation.
\end{requirement}

\section{Time Signatures and Meter}
\label{sec:time:meter}

A time signature defines the metric organization of a passage: where
bar lines fall, where beats group, and what accent pattern is implied.
Time signatures are objects in the score graph; they are not merely
display numerator/denominator pairs.

\begin{lstlisting}[language=Rust]
pub struct TimeSignature {
    /// Display form. Determines what the performer sees.
    pub display: TimeSignatureDisplay,

    /// The total duration of one measure under this signature.
    pub measure_duration: MusicalDuration,

    /// Beat groupings within the measure. Drives beaming, accent
    /// inference, and the notational-decomposition pre-pass.
    pub beat_groups: Vec<BeatGroup>,
}

pub enum TimeSignatureDisplay {
    /// Standard: a numerator over a power-of-two denominator.
    Standard { numerator: u16, denominator: PowerOfTwo },

    /// Compound display: e.g., 3+2+3/8.
    Compound { numerators: Vec<u16>, denominator: PowerOfTwo },

    /// Irrational meter: denominator is not a power of two.
    /// e.g., 4/6, 5/12. The denominator names a non-power-of-two
    /// note value.
    Irrational { numerator: u16, denominator: NonZeroU16 },

    /// Mixed denominators: e.g., 1/4 + 1/8 + 1/16 displayed additively.
    MixedDenominators { components: Vec<(u16, NonZeroU16)> },

    /// No visible time signature (free meter or proportional).
    None,

    /// Custom symbol (cut time, common time, etc.) or
    /// grammar-specific.
    Symbolic(SymbolId),
}

pub struct BeatGroup {
    /// Duration of this beat group.
    pub duration: MusicalDuration,

    /// Subdivision hint for beaming and accent.
    pub subdivision: Option<MusicalDuration>,

    /// Accent strength relative to other beat groups in the measure.
    /// Higher values indicate stronger accents.
    pub accent: u8,
}
\end{lstlisting}

\begin{requirement}
  \label{req:time:beat-group-sum}
  The sum of the durations of a time signature's beat groups \MUST{}
  equal the measure duration. Implementations \MUST{} reject time
  signatures whose beat groups do not sum to the measure duration.
\end{requirement}

\subsection{Irrational Meters}

Time signatures whose denominator is not a power of two
(e.g., 4/6, 5/12) are first-class. Their measure duration is computed
as a rational: a 4/6 measure has duration $\frac{4}{6} = \frac{2}{3}$
whole notes. No special case in the time model; only the display layer
treats the unusual denominator.

\subsection{Polymetric and Polytemporal Scores}

A score \MAY{} have different time signatures concurrent on different
staves. Each staff carries its own meter sequence; bar lines align where
the rational positions coincide, and may misalign otherwise. The score
graph chapter (Chapter~\ref{ch:graph}) defines how concurrent meters
are represented; this chapter establishes only that the time model
supports them.

\section{Tempo and the Tempo Map}
\label{sec:time:tempomap}

The tempo map is the function mapping musical positions to wall-clock
positions. It is a piecewise definition over musical time:

\begin{lstlisting}[language=Rust]
pub struct TempoMap {
    /// Piecewise segments, in non-overlapping musical-time order.
    pub segments: Vec<TempoSegment>,

    /// Tempo before the first segment, if any. None means the first
    /// segment's start_tempo applies from time zero.
    pub initial: Option<Tempo>,
}

pub struct TempoSegment {
    /// Start of this segment, as a TimeAnchor.
    pub start: TimeAnchor,

    /// End of this segment. None means "valid until the next segment
    /// in segments order, or until the end of the score if this is
    /// the final segment." Implementations resolve None lazily.
    pub end: Option<TimeAnchor>,

    /// Tempo at the start of this segment.
    pub start_tempo: Tempo,

    /// Tempo at the end of this segment.
    /// Required if shape != Constant; ignored (and SHOULD be None)
    /// if shape == Constant. A Constant segment holds start_tempo
    /// throughout.
    pub end_tempo: Option<Tempo>,

    /// Shape of the tempo change over this segment.
    pub shape: TempoShape,
}

pub enum TempoShape {
    /// Constant tempo throughout: end_tempo MUST be None or equal to
    /// start_tempo.
    Constant,

    /// Linear interpolation from start_tempo to end_tempo.
    Linear,

    /// Exponential (continuous-rate) interpolation from start_tempo
    /// to end_tempo.
    Exponential,

    /// Arbitrary curve, parameterized by control points. The control
    /// points define intermediate tempos between start_tempo and
    /// end_tempo.
    Curve(TempoCurve),
}

pub struct Tempo {
    /// Beats per minute. The beat unit is part of the tempo
    /// specification: a quarter-note BPM and a dotted-quarter BPM
    /// at the same numeric value differ in actual rate.
    pub bpm: f64,

    /// Note value taken as one beat.
    pub beat_unit: MusicalDuration,
}
\end{lstlisting}

\begin{requirement}
  \label{req:time:tempo-segment-order}
  The tempo map's segments \MUST{} be non-overlapping when resolved
  to absolute musical positions, and \MUST{} appear in
  monotonically-increasing start order. Two adjacent segments where
  the earlier segment's resolved end equals the later segment's
  resolved start form a continuous tempo curve; the later segment's
  \texttt{start\_tempo} \SHOULD{} equal the earlier segment's resolved
  end tempo when continuity is intended.

  A segment with \texttt{end == None} extends until the next
  segment's resolved start or, if no later segment exists, until the
  end of the score. Implementations \MUST{} resolve open-ended
  segments deterministically.

  At musical positions not covered by any segment (gaps between
  segments, or before the first segment when \texttt{initial} is
  \texttt{None}), the tempo \MUST{} be the most recent applicable
  segment's terminating tempo, or the next applicable segment's
  \texttt{start\_tempo} if no earlier segment exists.
\end{requirement}

\subsection{Conversion}

Conversion between musical and wall-clock time integrates the tempo map.
For constant segments, conversion is multiplication. For linear and
exponential segments, conversion has a closed-form solution. For curve
segments, conversion is numerical.

\begin{requirement}
  \label{req:time:linear-interpolates-speed}
  For a \texttt{TempoShape::Linear} segment, the quantity that varies
  linearly across the segment is \emph{speed} --- whole notes per
  second --- as a function of fractional position within the segment,
  \emph{not} beats-per-minute and \emph{not} beat period. Speed is
  derived from a \texttt{Tempo} as
  $\text{speed} = \text{bpm} \times w(\text{beat\_unit}) / 60$, where
  $w(\cdot)$ is the beat unit's value in whole notes. Writing
  $s_0, s_1$ for the segment's start and end speeds and
  $u \in [0,1]$ for fractional position, $s(u) = s_0 + (s_1 - s_0)\,u$,
  and the wall-clock duration of the segment of musical length
  $\ell$ is the closed form
  $\int_0^1 \ell\,\mathrm{d}u / s(u) = (\ell / (s_1 - s_0))\,
  \ln(s_1 / s_0)$ (and $\ell / s_0$ when $s_1 = s_0$).

  A \texttt{TempoShape::Exponential} segment interpolates speed
  \emph{geometrically}: $s(u) = s_0\,(s_1/s_0)^{u}$, i.e. a constant
  continuous rate of change of speed.
\end{requirement}

\begin{rationale}
  ``Linear interpolation from \texttt{start\_tempo} to
  \texttt{end\_tempo}'' is ambiguous about which quantity is linear:
  bpm, period, or speed. We pin \emph{speed} because it is the only
  beat-unit-agnostic choice. A tempo map may change beat unit across
  segments (a \texttt{quarter = 120} segment abutting a
  \texttt{dotted-quarter = 80} segment); interpolating bpm would make
  the wall-clock schedule depend on the arbitrary beat unit chosen to
  notate each endpoint, whereas interpolating speed depends only on
  the sounding rate. When both endpoints share a beat unit, speed-linear
  and bpm-linear coincide, so the common case matches a conductor's
  ``linear accelerando'' intuition. Because this choice fixes the
  derived wall-clock schedule for every playback engine, it is stated
  normatively rather than left to the implementation.
\end{rationale}

\begin{requirement}
  \label{req:time:tempo-map-conversion}
  The tempo map \MUST{} expose two conversion functions:
  \texttt{musical\_to\_wallclock} and \texttt{wallclock\_to\_musical}.
  Both \MUST{} be deterministic: given identical tempo maps and inputs,
  they \MUST{} produce identical outputs across runs and platforms.
  Numerical integration and inversion, where used, \MUST{} use a
  deterministic algorithm with documented tolerance and iteration
  bounds.

  Tempo curves \MUST{} be monotonically increasing in musical time
  (i.e., tempo may vary but musical time always advances) so that
  \texttt{wallclock\_to\_musical} is single-valued. Implementations
  \MUST{} reject tempo segments whose shape produces a non-monotonic
  mapping (such curves are musically meaningless).
\end{requirement}

\begin{openquestion}
  The specific numerical-integration algorithm for
  \texttt{TempoShape::Curve} (Romberg integration, adaptive
  Gauss-Kronrod, or a normative method) and the specific
  root-finding algorithm for \texttt{wallclock\_to\_musical} are
  unresolved. Both algorithms affect canonical state and must
  therefore receive one of the dispositions in
  Appendix~\ref{app:determinism}
  Section~\ref{sec:det:open}: normatively specified in this
  document or a named companion, profile-declared by versioned
  identifier, or explicitly marked non-canonical. Until that
  disposition is made, cross-implementation byte equality of
  derived canonical state (notably tempo-driven audio
  scheduling) \MUSTNOT{} be claimed.
\end{openquestion}

\begin{rationale}
  Determinism of tempo conversion is necessary for reproducible audio
  output, click-track stability, and SMPTE sync. Non-determinism here
  would produce playback drift that compounds over long pieces.
\end{rationale}

\section{Sounding Duration and Notational Decomposition}
\label{sec:time:notrhythm}

An event has a single \emph{sounding duration} (an exact rational).
Its \emph{notational decomposition} is the sequence of notehead values,
augmentation dots, and ties used to draw that duration on the staff.

The decomposition is computed from the sounding duration by a
deterministic pre-pass as a \emph{canonical derived annotation} ---
recomputed on materialization, never stored graph state --- under the
same output model as the spelling pre-pass
(Section~\ref{sec:pitch:prepass}). The user may accept the inferred
decomposition or override it by authoring a decomposition attachment,
which takes precedence.

\subsection{The Notational Decomposition}

\begin{lstlisting}[language=Rust]
pub struct NotationalDecomposition {
    /// Sequence of notated components. Their durations sum to the
    /// event's sounding duration.
    pub components: Vec<NotatedComponent>,

    /// Provenance, analogous to SpellingSource.
    pub source: DecompositionSource,
}

pub struct NotatedComponent {
    /// Base note value: whole, half, quarter, eighth, etc.
    pub base_value: NoteValue,

    /// Number of augmentation dots (0..=4 typical).
    pub dots: u8,

    /// Tuplet membership, if any.
    pub tuplet: Option<TupletId>,

    /// Whether this component is tied to the next.
    pub tied_to_next: bool,
}

pub enum DecompositionSource {
    UserChosen,
    Inferred,
    Imported { format: ForeignFormatId },
    Propagated { from: EventId },
}
\end{lstlisting}

The structure deliberately mirrors the spelling attachment model
(Section~\ref{sec:pitch:spelling}): same sources, same source-rank
discipline, same pre-pass output model. Unlike spelling, decomposition
precedence is \emph{not configurable} (ratified Pass~12): the fixed
default source order \texttt{UserChosen} $>$ \texttt{Imported} $>$
\texttt{Propagated} $>$ \texttt{Inferred} applies, with the score's
canonical attachment order as tie-break among equal-rank authored
attachments; \texttt{DecompositionAttachment} carries no
\texttt{priority} field and the graph carries no
\texttt{DecompositionPrecedence} configuration. Making precedence
configurable would add a canonical \texttt{Score} field (a
schema-major change under the Binary Format companion's evolution
rule) for which no consumer exists; the decision may be revisited at
a future schema major if a use case appears.

\subsection{The Decomposition Pre-Pass}

When a score is materialized, the decomposition pre-pass computes
inferred decompositions, reported as derived annotations; an authored
higher-precedence attachment overrides the inferred value in the
resolved annotation. An authored attachment targeting an event the
pre-pass produces \emph{no} inferred decomposition for (ungriddable,
non-metric, or an inapplicable event kind) surfaces as the resolved
annotation on its own, per
Requirement~\ref{req:pitch:authored-uninferred} --- authoring is how
a user notates exactly what the algorithm cannot infer. The
algorithm is, in outline:

\begin{enumerate}
  \item For each event in time order within a voice and measure,
    examine the sounding duration and the current beat-group context.
  \item Decompose the duration into the longest leading note value
    that fits within the current beat group without crossing a stronger
    accent boundary.
  \item If the duration is not exhausted, produce a tie and continue
    decomposition from the next beat-group boundary.
  \item Augmentation dots are applied where they shorten the
    decomposition without crossing an accent boundary that the dot would
    obscure.
  \item Tuplet membership is recorded for events that fall within a
    tuplet group (Section~\ref{sec:time:tuplets}).
\end{enumerate}

\begin{requirement}
  \label{req:time:decomposition-determinism}
  The decomposition pre-pass \MUST{} be deterministic: given
  identical materialized graph, configuration, and algorithm
  version, it \MUST{} produce identical decompositions. Its output
  is a derived annotation, never stored graph state. Implementations
  \MAY{} cache and \MAY{} recompute incrementally (re-running only
  the affected voice and measure, with propagation only as far as
  the context is altered); neither may be observable in the
  annotations produced.
\end{requirement}

\begin{requirement}
  \label{req:time:decomposition-algorithm}
  \textbf{Decomposition algorithm disposition (ratified Pass 12).}
  The decomposition pre-pass receives the \emph{profile-declared
  by versioned identifier} disposition of
  Appendix~\ref{app:determinism} Section~\ref{sec:det:open}. The
  reserved identifier \texttt{DecompositionAlgorithmId}
  \texttt{"default"} denotes, at version~1, the integer-grid
  metric splitter, whose scope bounds are part of its normative
  definition:
  \begin{itemize}
    \item a single governing meter per region --- a mid-region
      meter change reduces cleanly into the region's grid (the
      Operation Catalog pins those semantics), but derived
      notation honours only the first governing meter until a
      wider algorithm version lands;
    \item the region origin is assumed to fall on a barline
      (anacrusis/pickup handling deferred);
    \item compound-meter beat grouping uses the dyadic default;
    \item tuplet nesting and cross-beat tuplet members are out of
      scope;
    \item at most one augmentation dot
      (\texttt{MAX\_DOTS}~$=$~1); a double-dotted value is
      written as tied components --- correct, if not the most
      compact.
  \end{itemize}
  A wider algorithm is a \emph{version bump}, which
  deterministically invalidates derived output under the
  derived-annotation model --- no state migration. Style-choice
  selection among multiple valid decompositions (e.g.,
  $\frac{3}{8}$ in $\frac{4}{4}$ as a dotted quarter versus a
  quarter tied to an eighth) is fixed by the versioned algorithm;
  user-configurable alternate rule sets, if ever wanted, are new
  registered identifiers. A profile requesting any identifier
  other than a registered one \MUST{} error; implementations
  \MUSTNOT{} silently substitute.
\end{requirement}

\section{Tuplets as Grouping Objects}
\label{sec:time:tuplets}

A tuplet is a grouping object in the score graph that records
notational intent: that a span of events forms a triplet, quintuplet,
or other irregular grouping with a visible bracket and number. Tuplets
do not modify the sounding durations of their members; sounding
durations are always concrete rationals.

\begin{lstlisting}[language=Rust]
pub struct Tuplet {
    pub id: TupletId,

    /// The actual:notated ratio. A triplet is 3:2; a quintuplet is
    /// 5:4; an irregular grouping like 7:4 has ratio 7 over 4.
    pub ratio: TupletRatio,

    /// Range of events in this tuplet, by event ID.
    pub members: Vec<EventId>,

    /// Display: bracket style, number placement, parentheses, etc.
    pub display: TupletDisplay,

    /// Nesting: if this tuplet is itself inside another, the parent.
    pub parent: Option<TupletId>,
}

pub struct TupletRatio {
    // Fields are private: a TupletRatio is constructed only through the
    // checked `new`, so a degenerate ratio is never representable (see
    // the requirement below). The fields are read via `actual()`/`notated()`.
    actual: u32,
    notated: u32,
}

impl TupletRatio {
    /// Returns `None` for a degenerate ratio (either term zero, or
    /// `actual == notated`); otherwise the well-formed ratio.
    pub fn new(actual: u32, notated: u32) -> Option<TupletRatio> { /* checked */ }
    pub fn actual(&self) -> u32 { /* ... */ }
    pub fn notated(&self) -> u32 { /* ... */ }
}
\end{lstlisting}

\begin{requirement}
  \label{req:time:tuplet-ratio-construction}
  A \texttt{TupletRatio} is \emph{degenerate} if either term is zero
  or if \texttt{actual == notated} (a ratio that expresses no
  augmentation or diminution). Degenerate ratios \MUST{} be rejected
  at construction: a conforming constructor of \texttt{TupletRatio}
  (and therefore of any \texttt{Tuplet}) \MUST{} fail rather than
  produce a value with \texttt{actual == 0}, \texttt{notated == 0}, or
  \texttt{actual == notated}. This is a construction-time rejection,
  not a runtime graph invariant (Section~\ref{sec:graph:invariants}):
  a degenerate ratio is never a representable graph state, so no
  invariant restores it after the fact.
\end{requirement}

\subsection{Tuplet Nesting}

Tuplets \MAY{} nest to arbitrary depth. The sounding duration of an
event inside nested tuplets is computed directly from its rational
duration, not from the tuplet ratios; the ratios are notational and
analytical.

\begin{rationale}
  Storing sounding durations directly (rather than as base-value times
  a stack of tuplet ratios) keeps the time model regular: every event
  has a single rational duration, full stop. Moving a tuplet, splitting
  it, or deleting it never alters the durations of its members. Tuplets
  thus become a pure annotation layer, which dramatically simplifies
  edit operations.
\end{rationale}

\subsection{Tuplet Consistency}

\begin{requirement}
  \label{req:time:tuplet-duration-consistency}
  A tuplet's notated ratio \MUST{} be consistent with the sounding
  durations of its members: the sum of member durations \MUST{} equal
  the duration that the tuplet's notated value indicates when scaled
  by the ratio. For example, a 3:2 eighth-note triplet of three members
  has total sounding duration $\frac{1}{4}$; the sum of the three
  member durations \MUST{} equal $\frac{1}{4}$.

  This is a structural invariant. Edit operations that would violate
  it \MUST{} either adjust member durations or be rejected. The
  semantic operations chapter specifies which.
\end{requirement}

\section{Region Time Models}
\label{sec:time:regions}

A region of the score declares one of three time models. The declared
model determines how positions within the region are interpreted,
displayed, and played back.

\begin{lstlisting}[language=Rust]
pub enum RegionTimeModel {
    /// Metric time: measures, beats, exact rational positions.
    /// The dominant model for Western notated music.
    Metric(MetricTimeModel),

    /// Proportional time: events placed in space, where horizontal
    /// position is wall-clock time. No measures, no beats.
    Proportional(ProportionalTimeModel),

    /// Aleatoric time: events with partial or unspecified ordering.
    Aleatoric(AleatoricTimeModel),
}
\end{lstlisting}

\subsection{Metric Time}

\begin{lstlisting}[language=Rust]
pub struct MetricTimeModel {
    /// Sequence of meter changes within this region.
    pub meters: Vec<MeterChange>,

    /// Tempo map for this region. May inherit from the score.
    pub tempo: TempoMapReference,
}
\end{lstlisting}

Metric regions have measures, beats, and time signatures as described
in Sections~\ref{sec:time:meter} and~\ref{sec:time:tempomap}. Positions
within them are \texttt{MusicalPosition} values.

\subsection{Proportional Time}

\begin{lstlisting}[language=Rust]
pub struct ProportionalTimeModel {
    /// Total wall-clock duration of the region.
    pub duration: WallClockDuration,

    /// Horizontal space per second. Drives the engraver's spacing.
    /// Stored as a layout hint; the actual spacing is decided by the
    /// constraint solver.
    pub space_per_second: SpaceUnit,

    /// Optional reference grid markings (vertical lines every N
    /// seconds, for example) for performer reference.
    pub grid: Option<ProportionalGrid>,
}
\end{lstlisting}

Events in proportional regions carry \texttt{WallClockTime} positions
directly. There are no measures, no beats, no rational durations.
Notational duration is computed from wall-clock duration by the
notational decomposition pre-pass using a proportional-mode rule set,
which generally produces unbeamed events with explicit duration
notations or graphic-duration symbols.

\begin{requirement}
  \label{req:time:proportional-region-time}
  Events in a proportional region \MUST{} use \texttt{WallClockTime}
  for positions and \texttt{WallClockDuration} for durations. The
  tempo map \MUSTNOT{} be applied to convert proportional-region
  positions; they are already in wall-clock time.
\end{requirement}

\subsection{Aleatoric Time}

\begin{lstlisting}[language=Rust]
pub struct AleatoricTimeModel {
    /// Ordering constraints among events: a directed acyclic graph
    /// where an edge a -> b means "a precedes b."
    pub ordering: EventOrderingDAG,

    /// Anchoring discipline: declares what event coordinate kinds
    /// this region permits and whether duration-bound kinds may
    /// be mixed. See Section~\ref{sec:time:aleatoric-anchoring}.
    pub anchoring: AleatoricAnchoringDiscipline,

    /// Optional per-event interval bounds. An event with start
    /// bounds [t1, t2] may begin at any time in that window.
    /// Bound kinds are constrained by the anchoring discipline.
    pub bounds: HashMap<EventId, EventBounds>,

    /// Approximate or maximum total duration, for layout.
    pub duration_hint: WallClockDuration,
}

pub enum AleatoricAnchoringDiscipline {
    /// Events in this region carry musical-time coordinates only.
    Musical,

    /// Events in this region carry wall-clock coordinates only.
    WallClock,

    /// Events may carry either kind, but each event's position and
    /// duration kinds must agree with each other. Mixed-kind
    /// duration bounds are not permitted.
    EitherPerEvent,

    /// Events may carry either kind freely, and duration bounds
    /// may mix kinds (one bound musical, the other wall-clock).
    /// Used for free-time scores synchronized to external media.
    FreelyMixed,
}

pub struct EventBounds {
    pub start: Option<TimeBounds>,
    pub end: Option<TimeBounds>,
}

pub enum TimeBounds {
    MusicalRange { min: MusicalPosition, max: MusicalPosition },
    WallClockRange { min: WallClockTime, max: WallClockTime },
    Unbounded,
}
\end{lstlisting}

Aleatoric regions support both unordered event collections (any
permutation of events is valid) and time-bracket notation (events with
specified start and end windows, à la Cage's later notation). The DAG
encoding admits both as cases.

\subsubsection{Anchoring Discipline}
\label{sec:time:aleatoric-anchoring}
\label{sec:graph:aleatoric-discipline}

\begin{requirement}
  \label{req:time:aleatoric-anchoring-discipline}
  Every aleatoric region \MUST{} declare an
  \texttt{AleatoricAnchoringDiscipline}. Events within the region
  \MUST{} carry coordinate kinds consistent with the discipline:

  \begin{itemize}
    \item \texttt{Musical}: events use \texttt{EventPosition::Musical}
      and \texttt{EventDuration::Musical} (or
      \texttt{EventDuration::Indeterminate} whose bounds are musical).
    \item \texttt{WallClock}: events use
      \texttt{EventPosition::WallClock} and
      \texttt{EventDuration::WallClock} (or
      \texttt{EventDuration::Indeterminate} whose bounds are
      wall-clock).
    \item \texttt{EitherPerEvent}: events choose per-event, but an
      event's position and duration kinds \MUST{} agree, and
      \texttt{DurationBounds} \MUST{} use a single
      \texttt{ConcreteDuration} variant.
    \item \texttt{FreelyMixed}: no agreement required;
      \texttt{DurationBounds} \MAY{} mix kinds.
  \end{itemize}

  Implementations \MUST{} reject events whose coordinate kinds
  violate the region's declared discipline.
\end{requirement}

\begin{requirement}
  \label{req:time:ordering-dag-acyclic}
  The ordering DAG \MUST{} be acyclic. Cycles \MUST{} be rejected at
  construction. Two events with no path between them in the DAG are
  unordered.
\end{requirement}

\begin{rationale}
  A DAG with optional interval bounds expresses every aleatoric form
  the project intends to support: unordered collections (empty DAG),
  partial orderings (sparse DAG), strict orderings (chain DAG), and
  time-brackets (DAG plus bounds). Higher-level constructs
  (mobile-form notation, statistical scores) can be built atop this
  primitive in later specifications without altering the core.
\end{rationale}

\subsection{Concurrent Time Models}

Different staves in the same score \MAY{} declare different time
models. A metric solo line over a proportional aleatoric texture, for
example, is expressible: each staff's region declares its own model.
Synchronization between concurrent time models is by wall-clock time
or by explicit cue points (see Chapter~\ref{ch:graph}).

\section{The Notion of ``Now''}
\label{sec:time:now}

The core does not define playback state; it does not have a notion of
``current position'' or ``transport time.'' Such state belongs to the
audio engine, which consumes the core's time model. The core defines
only the static time structure of the score.

\section{Forward References}

\begin{itemize}
  \item Event identifiers (\texttt{EventId}), measure identifiers
    (\texttt{MeasureId}), and region identifiers (\texttt{RegionId})
    are defined in Chapter~\ref{ch:graph}.
  \item The full event taxonomy (pitched, unpitched, rest, indeterminate,
    trajectory) is defined in Chapter~\ref{ch:graph}.
  \item Re-anchoring rules for time anchors when their targets are
    deleted are defined per-operation in Chapter~\ref{ch:semops}.
  \item \texttt{SpaceUnit} is defined normatively in
    Chapter~\ref{ch:graph} as of schema major~2 (it entered canonical
    state through \texttt{StaffLineConfiguration}).
    \texttt{TupletDisplay} and other purely display-oriented types are
    defined in Chapter~\ref{ch:layout-ir}.
\end{itemize}

% ===========================================================================
\chapter{Tuning Systems and Pitch Spaces}
\label{ch:tuning}

This chapter specifies the three intertwined registries that complete
the primitives layer: \emph{pitch spaces} (the analytical universes in
which scale-position arithmetic lives), \emph{accidental registries}
(the catalogs of position-modifying glyphs), and \emph{tuning systems}
(the frequency-resolution mechanisms). It also specifies the
hierarchical resolution rules that determine which tuning applies to
each pitch in a score.

\section{Design Principles}
\label{sec:tuning:principles}

\begin{description}
  \item[Pitch space and tuning system are distinct.] A pitch space
    defines what pitches exist and how they relate analytically. A
    tuning system defines what frequencies those pitches produce.
    The two are independent: a single pitch space (e.g., CMN) admits
    many tuning systems (12-TET, meantone, well-temperaments, just
    intonation, Pythagorean); a single tuning structure may be anchored
    at different reference frequencies.
  \item[Accidentals belong to pitch spaces.] An accidental defines a
    \emph{position modification} within a pitch space. The tuning system
    handles the frequency consequences of that modification. The chromatic
    semitone implied by a sharp has different cent-values in 12-TET,
    quarter-comma meantone, and Pythagorean tuning, but it is the same
    accidental denoting the same position.
  \item[Reference pitch is score-level, not tuning-level.] A score
    declares a reference pitch (position plus frequency); the tuning
    system defines structure relative to that reference. ``12-TET at
    A=440'' and ``12-TET at A=415'' share the same tuning system; the
    reference differs.
  \item[Resolution is hierarchical.] Tuning, pitch space, and accidental
    registry are resolved by walking outward from the pitch: pitch,
    voice, staff, region, score. Inheritance markers permit any level
    to defer to the enclosing scope.
  \item[Built-in defaults, extensible by users.] The core ships a
    baseline catalog of pitch spaces, tuning systems, and accidental
    registries sufficient for the common 95\% case. Scores and grammar
    plugins extend these without forking.
  \item[SMuFL primary, custom fallback.] Glyph references default to
    SMuFL codepoints. Custom glyphs are supported as fallback. SMuFL
    version is recorded per score.
\end{description}

\section{Pitch Spaces}
\label{sec:tuning:space}

A pitch space is the analytical universe in which scale-degree and
interval relationships are defined. It is the algebra; the tuning
system is the physics.

\begin{lstlisting}[language=Rust]
pub struct PitchSpace {
    pub id: PitchSpaceId,

    /// Human-readable name (e.g., "CMN 12-tone chromatic").
    pub name: String,

    /// Optional description and provenance notes.
    pub description: Option<String>,

    /// What positions exist and how they are structured.
    pub positions: PositionStructure,

    /// How intervals between positions are defined.
    pub interval_algebra: IntervalAlgebra,

    /// The accidental registry valid in this pitch space.
    pub accidental_registry: AccidentalRegistryId,

    /// The nominal registry: named position-letters
    /// (e.g., A-G for CMN).
    pub nominal_registry: NominalRegistryId,

    /// Rules governing how pitches transpose within this space.
    pub transposition: TranspositionBehavior,

    /// Default rules for the spelling pre-pass when this pitch space
    /// is active.
    pub spelling_rules: SpellingRuleSet,
}
\end{lstlisting}

\subsection{Position Structure}
\label{sec:tuning:positions}

The \texttt{PositionStructure} declares the shape of the pitch space.
Three families cover the vast majority of useful cases; a fourth is the
escape hatch for grammars defined entirely by plugin.

\begin{lstlisting}[language=Rust]
pub enum PositionStructure {
    /// Chromatic: a fixed number of equally-numbered positions per
    /// octave, with no hierarchical diatonic substructure.
    /// Includes EDOs and serial 12-tone.
    Chromatic { positions_per_octave: u16 },

    /// Diatonic over chromatic: a smaller set of named nominals
    /// (e.g., A-G for CMN) plus accidental modifications reaching
    /// the full chromatic range. CMN's structure.
    DiatonicOverChromatic {
        nominals_per_octave: u16,
        chromatic_positions_per_octave: u16,
        /// Mapping from each nominal to its position in the
        /// chromatic layer (e.g., C=0, D=2, E=4, F=5, G=7, A=9, B=11
        /// for CMN).
        nominal_to_chromatic: Vec<u16>,
    },

    /// Just intonation lattice: positions defined by exact rational
    /// ratios from a tonic, organized along prime-axis dimensions
    /// (3-limit, 5-limit, 7-limit, etc.).
    JiLattice {
        limit: u8,                     // prime limit
        generators: Vec<JiRatio>,      // one per prime dimension
    },

    /// Grammar-defined: structure is opaque to the core and resolved
    /// by a grammar plugin. Used for maqam, gamelan, raga frameworks,
    /// and other systems whose position structure is not flat or
    /// hierarchical in the above senses.
    Registered(PositionStructureRegistryId),
}
\end{lstlisting}

\begin{requirement}
  \label{req:tuning:diatonic-chromatic-mapping}
  The \texttt{DiatonicOverChromatic} variant \MUST{} have a
  \texttt{nominal\_to\_chromatic} mapping of length equal to
  \texttt{nominals\_per\_octave}, with each entry strictly less than
  \texttt{chromatic\_positions\_per\_octave}. The mapping \MUST{} be
  strictly increasing.
\end{requirement}

\subsection{Interval Algebra}

An interval algebra defines how distances between positions are
computed and how transposition is realized. For chromatic spaces the
algebra is integer arithmetic modulo the position count; for
diatonic-over-chromatic the algebra distinguishes diatonic from
chromatic intervals (a diatonic third may correspond to three or four
chromatic semitones depending on quality).

\begin{lstlisting}[language=Rust]
pub enum IntervalAlgebra {
    /// Single-axis integer arithmetic. Two positions are subtracted
    /// to yield a single integer interval.
    Chromatic,

    /// Two-axis arithmetic distinguishing diatonic and chromatic
    /// components. CMN's algebra: a major third is (3 diatonic steps,
    /// 4 chromatic steps); a diminished fourth is (3 diatonic steps,
    /// 4 chromatic steps) -- distinguishable from the major third only
    /// by spelling intent.
    DiatonicChromatic,

    /// Multi-axis JI lattice arithmetic.
    JiVector { dimensions: u8 },

    /// Grammar-defined algebra.
    Registered(IntervalAlgebraRegistryId),
}
\end{lstlisting}

\subsection{Nominal Registry}

A nominal registry catalogs the named position-letters of a pitch
space. For CMN this is the seven-element set $\{C, D, E, F, G, A, B\}$.
For 12-tone serial music the nominal registry may simply be the
integers $0$ through $11$. For maqam the nominals are the names of
maqam degrees.

\begin{lstlisting}[language=Rust]
pub struct NominalRegistry {
    pub id: NominalRegistryId,
    pub nominals: Vec<NominalDefinition>,
}

pub struct NominalDefinition {
    pub id: NominalId,

    /// Canonical name.
    pub name: String,

    /// Display glyph, typically a letter or symbol.
    pub display: GlyphReference,

    /// Default position in the pitch space's chromatic layer, if
    /// applicable.
    pub default_position: Option<u16>,
}
\end{lstlisting}

\subsection{Transposition Behavior}

Transposition behavior governs how scale positions move under interval
operations and how spellings follow.

\begin{lstlisting}[language=Rust]
pub enum TranspositionBehavior {
    /// Diatonic transposition: shift by N diatonic steps within a
    /// chosen key. Accidentals adjust to fit the destination key.
    Diatonic,

    /// Chromatic transposition: shift by N chromatic steps. Spellings
    /// follow rules of the SpellingRuleSet.
    Chromatic,

    /// Compound: support both, parameterized at the operation site.
    Compound,

    /// Grammar-defined.
    Registered(TranspositionRegistryId),
}
\end{lstlisting}

\subsection{Spelling Rule Sets}

A pitch space carries a default spelling rule set that the spelling
pre-pass (Section~\ref{sec:pitch:prepass}) consults when inferring
spellings for unspelled pitches.

\begin{lstlisting}[language=Rust]
pub struct SpellingRuleSet {
    pub id: SpellingRuleSetId,
    pub name: String,

    /// Algorithmic family this rule set belongs to. The specific
    /// algorithm is referenced by id and resolved against the
    /// implementation's registered spelling algorithms.
    pub algorithm: SpellingAlgorithmId,

    /// Algorithm-specific parameters.
    pub parameters: SpellingParameters,
}
\end{lstlisting}

\begin{openquestion}
  The default spelling algorithm is now ratified
  (Requirement~\ref{req:pitch:spelling-algorithm}:
  \texttt{"default"} = Temperley-style line-of-fifths preference,
  version~1). Still open: the catalog of \emph{additional}
  registered spelling algorithms (candidates include
  Longuet-Higgins line-of-fifths distance, Cambouropoulos
  pitch-spelling, and Meredith PS13) and their parameter schemas,
  which are normative once registered.
\end{openquestion}

\section{Accidental Registries}
\label{sec:tuning:accidentals}

An accidental registry catalogs the position-modifying glyphs valid
within a pitch space. Each accidental defines a glyph, a position
modification, and engraving metadata.

\begin{lstlisting}[language=Rust]
pub struct AccidentalRegistry {
    pub id: AccidentalRegistryId,
    pub name: String,
    pub accidentals: Vec<AccidentalDefinition>,

    /// Whether scores may extend this registry with score-local
    /// accidentals. Most registries permit this; some (e.g., a
    /// closed conformance profile) may not.
    pub extensible: bool,
}

pub struct AccidentalDefinition {
    pub id: AccidentalId,

    /// Canonical name. Used for serialization, accessibility, and
    /// foreign-format export.
    pub name: String,

    /// Glyph reference, typically SMuFL.
    pub glyph: GlyphReference,

    /// The position modification this accidental applies.
    pub modification: PitchSpaceModification,

    /// Engraving metadata.
    pub engraving: AccidentalEngraving,

    /// Combination behavior with other accidentals.
    pub combination: AccidentalCombination,
}
\end{lstlisting}

\subsection{Pitch Space Modifications}

\begin{lstlisting}[language=Rust]
pub enum PitchSpaceModification {
    /// CMN-style integer offset in chromatic steps of the enclosing
    /// pitch space (`req:pitch:alteration-unit`).
    CmnChromatic(i8),

    /// Integer step offset in an EDO position space.
    EdoSteps(i16),

    /// Modification expressed as an exact rational ratio. Used for
    /// JI accidentals such as the HEJI syntonic-comma symbols
    /// (81/80) and septimal-comma symbols (64/63).
    JiRatio { numerator: i32, denominator: NonZeroU32 },

    /// Modification expressed in cents. Used for Sagittal-style
    /// precise microtonal accidentals.
    Cents(f64),

    /// Modification defined by a grammar plugin.
    Registered(ModificationRegistryId),
}
\end{lstlisting}

\begin{requirement}
  \label{req:tuning:accidental-modification-compatibility}
  An accidental's modification \MUST{} be expressible in the interval
  algebra of every pitch space that references its registry. A
  \texttt{CmnChromatic} modification is valid only in spaces with
  \texttt{DiatonicChromatic} or compatible algebra; an \texttt{EdoSteps}
  modification is valid only in \texttt{Chromatic} or \texttt{Registered}
  spaces with appropriate compatibility. Implementations \MUST{} reject
  scores referencing an accidental in a space whose algebra does not
  admit the modification.
\end{requirement}

\subsection{Accidental Engraving Metadata}

\texttt{AccidentalEngraving} is reachable from canonical score state ---
via \texttt{AccidentalDefinition} inside
\texttt{ScoreAccidentalExtensions}, which hangs off
\texttt{ScoreTuningContext} --- so every field it carries must itself be
canonical. Its bounding box therefore cannot reuse
Chapter~\ref{ch:layout-ir}'s \texttt{BoundingBox}, which is built on
\texttt{StaffSpace} (single-precision \texttt{f32}) and belongs to the
non-canonical resolved-layout cache. This chapter defines its own box
type over \texttt{SpaceUnit} instead:

\begin{lstlisting}[language=Rust]
/// A bounding box over canonical space units, used for engraving
/// metadata that lives in canonical score state. Distinct from
/// Chapter 7's `BoundingBox` (built on `StaffSpace`, single precision,
/// for the non-canonical resolved-layout cache): this type's edges are
/// `SpaceUnit` (`CanonicalF64`), per
/// `req:determinism:canonical-floating-point`.
pub struct EngravingBoundingBox {
    pub left: SpaceUnit,
    pub right: SpaceUnit,
    pub top: SpaceUnit,
    pub bottom: SpaceUnit,
}
\end{lstlisting}

\begin{lstlisting}[language=Rust]
pub struct AccidentalEngraving {
    /// Bounding box in canonical space units, relative to the glyph's
    /// anchor point.
    pub bounding_box: EngravingBoundingBox,

    /// Where the glyph attaches to the note: typically the geometric
    /// center or a custom anchor for compound glyphs.
    pub anchor: AnchorPoint,

    /// Advance width for horizontal spacing computations.
    pub advance_width: SpaceUnit,

    /// Stacking order when multiple accidentals attach to one note.
    /// Lower values are placed closer to the notehead.
    pub stacking_order: i32,

    /// Whether this glyph should be drawn with parentheses by default
    /// (e.g., editorial or cautionary accidentals).
    pub default_parenthesized: bool,
}
\end{lstlisting}

\begin{rationale}
  Engraving metadata that lives in canonical state must carry canonical
  precision. Chapter~\ref{ch:layout-ir}'s coordinates are
  single-precision by requirement
  (Requirement~\ref{req:layoutir:staff-space-coordinates}) and reach the
  wire in the \texttt{LayoutCache}, which the Binary Format companion
  defines as an independent \emph{non-canonical} chunk --- so
  \texttt{StaffSpace} is correct there. But
  \texttt{accidental\_extensions} lives on
  \texttt{ScoreTuningContext}, which \emph{is} canonical, and
  Requirement~\ref{req:determinism:canonical-floating-point} requires
  canonical stored floats to be finite IEEE~754 binary64.
  \texttt{EngravingBoundingBox} makes \texttt{AccidentalEngraving}
  wholly core-typed, over the same \texttt{SpaceUnit} that
  \texttt{advance\_width} already uses, and removes what would
  otherwise be a backwards dependency from this chapter onto
  Chapter~\ref{ch:layout-ir}. Chapter~\ref{ch:layout-ir}'s
  \texttt{BoundingBox} is unaffected and remains the correct type for
  resolved-layout coordinates.
\end{rationale}

\subsection{Combination Behavior}

Most accidentals do not combine; a note carries one accidental at a
time. Some systems (HEJI, certain microtonal notations) permit stacking
multiple accidental glyphs on a single note to express compound
modifications.

\begin{lstlisting}[language=Rust]
pub enum AccidentalCombination {
    /// Stands alone; replaces any prior accidental on the same note.
    Solitary,

    /// May stack with members of the listed compatibility groups.
    /// Stacking order is determined by the engraving metadata.
    Stacking { compatible_groups: Vec<AccidentalGroupId> },
}
\end{lstlisting}

\subsection{Score-Local Extensions}

A score \MAY{} extend a referenced accidental registry with additional
accidental definitions, provided the registry's \texttt{extensible}
flag is true. Extensions are stored on the score and override or
augment the base registry during resolution.

\begin{lstlisting}[language=Rust]
pub struct ScoreAccidentalExtensions {
    pub base: AccidentalRegistryId,
    pub additions: Vec<AccidentalDefinition>,
    pub overrides: Vec<AccidentalDefinition>,
}
\end{lstlisting}

\section{Glyph References and SMuFL}
\label{sec:tuning:smufl}

The Standard Music Font Layout (SMuFL) is the normative source for
glyph references in the core. Every standard accidental, notehead,
articulation, dynamic, and ornament has a SMuFL codepoint. Non-standard
glyphs are supported via custom glyph references.

\begin{lstlisting}[language=Rust]
pub enum GlyphReference {
    /// SMuFL codepoint, resolved against the active SMuFL font.
    Smufl(u32),

    /// Custom glyph defined by the score or a plugin.
    Custom(CustomGlyphId),

    /// Composite glyph: multiple glyph references rendered as a
    /// single accidental. Used for compound HEJI symbols and similar.
    Composite(Vec<GlyphReference>),
}
\end{lstlisting}

\subsection{SMuFL Versioning}

\begin{requirement}
  \label{req:tuning:smufl-version-fallback}
  Every score \MUST{} declare the SMuFL version it targets. Resolving a
  SMuFL glyph reference whose codepoint is not present in the active
  font's SMuFL version \MUST{} produce a deterministic fallback (the
  reference SMuFL substitution glyph, or a clearly-marked missing-glyph
  indicator); it \MUSTNOT{} silently fail.
\end{requirement}

\begin{lstlisting}[language=Rust]
pub struct SmuflVersion {
    pub major: u16,

    /// The fractional part, normalized to hundredths.
    pub minor_centi: u16,
}
\end{lstlisting}

A SMuFL version's fractional part is stored \emph{normalized to
hundredths} in \texttt{minor\_centi}, not as the literal digits that
follow the decimal point: 1.12 $\to$ \tablenums{12}, 1.18 $\to$
\tablenums{18}, 1.20 $\to$ \tablenums{20}, 1.3 $\to$ \tablenums{30},
1.4 $\to$ \tablenums{40}. SMuFL versions are decimal fractions, and
this repository's published minor releases did not arrive in digit
order: 1.12 and 1.18 shipped before 1.3 and 1.4. A literal-digit
encoding (\texttt{minor: 3} for 1.3, \texttt{minor: 12} for 1.12) is
\textbf{non-conforming} because the derived ordering on
\texttt{(major, minor)} then sorts 1.3 and 1.4 \emph{before} 1.12,
backwards versus the real release history. The hundredths
normalization also collapses the 1.2/1.20 spelling ambiguity: both
denote the same release and both normalize to \tablenums{20}. This is
the single \texttt{SmuflVersion} type used throughout the
specification --- by \texttt{SmuflVersionRequirement} below and, in
Chapter~\ref{ch:layout-ir}, by \texttt{GlyphCatalog::smufl\_version()}
(Section~\ref{sec:layoutir:catalog}) and \texttt{GlyphCatalogIdentity}
(Section~\ref{sec:layoutir:catalog-identity}).

\begin{rationale}
  For reference, SMuFL's minor releases in their real order, each with
  its normalized \texttt{minor\_centi} value: 1.12
  (\tablenums{12}), 1.18 (\tablenums{18}), 1.20 (\tablenums{20}), 1.3
  (\tablenums{30}), 1.4 (\tablenums{40}).
\end{rationale}

\begin{lstlisting}[language=Rust]
pub struct SmuflVersionRequirement {
    /// Minimum SMuFL version required by this score.
    pub minimum: SmuflVersion,

    /// SMuFL version this score was authored against. Used for
    /// detecting whether newer glyphs are in use.
    pub authored_against: SmuflVersion,
}
\end{lstlisting}

\section{Tuning Systems}
\label{sec:tuning:system}

A tuning system maps positions in a pitch space to frequencies. Given
a reference pitch (position plus frequency) and the system's resolution
rules, every position in the space resolves to a determinate frequency.

\begin{lstlisting}[language=Rust]
pub struct TuningSystem {
    pub id: TuningSystemId,

    /// Human-readable name.
    pub name: String,

    /// The pitch space whose positions this tuning resolves.
    pub pitch_space: PitchSpaceId,

    /// How the tuning resolves positions to frequencies, given a
    /// reference pitch. The reference itself is supplied separately
    /// (see Section~\ref{sec:tuning:reference}).
    pub resolution: TuningResolution,

    /// Optional historical or provenance notes.
    pub description: Option<String>,
}
\end{lstlisting}

\subsection{Tuning Resolution}

\begin{lstlisting}[language=Rust]
pub enum TuningResolution {
    /// N-tone equal temperament: each step is the Nth root of the
    /// octave ratio. Includes 12-TET when N=12.
    EqualTemperament { divisions_per_octave: u16 },

    /// Explicit per-position ratios. Each entry is a ratio relative
    /// to the reference position.
    PerPositionRatios(Vec<PositionRatio>),

    /// Procedural definition: a registered tuning function that
    /// computes frequencies from a reference. Historical temperaments
    /// (Werckmeister, Vallotti, meantone variants) live here.
    Function {
        function: TuningFunctionId,
        parameters: TuningParameters,
    },

    /// A base resolution with per-position overrides. Used for
    /// split-accidental keyboards and similar irregular tunings.
    Overlay {
        base: Box<TuningResolution>,
        overrides: Vec<PositionOverride>,
    },

    /// Imported from a tuning file (Scala .scl, MIDI Tuning Standard,
    /// etc.). The data is captured at import; the original file is
    /// referenced for provenance but is not required at runtime.
    Imported {
        format: TuningFileFormat,
        data: Arc<ImportedTuningData>,
    },

    /// Adaptive tuning: resolution depends on the harmonic context
    /// at the moment of sounding. Used for adaptive JI and similar
    /// systems.
    Adaptive {
        function: AdaptiveTuningFunctionId,
        parameters: AdaptiveTuningParameters,
    },
}
\end{lstlisting}

\subsection{The Resolution Contract}

\begin{requirement}
  \label{req:tuning:tuning-resolution-determinism}
  Every tuning resolution \MUST{} expose a deterministic function from
  pitch-space position (plus, where applicable, harmonic context) to
  frequency in Hertz, given a reference pitch. Determinism \MUST{}
  hold across runs and platforms; numerical computations \MUSTNOT{}
  depend on floating-point rounding modes.
\end{requirement}

\subsection{Adaptive Tuning}

Adaptive tuning systems compute frequency from position \emph{plus
harmonic context}. A registered adaptive tuning function \MAY{}
resolve position per sonority --- for example, tuning a given E
differently depending on whether it sounds as the major third of a C
chord, the perfect fifth of an A chord, or a passing tone --- subject
only to the purity constraint below
(Requirement~\ref{req:tuning:adaptive-tuning-purity}). The built-in
\texttt{ji-adaptive-5limit}, at its pinned version~1
(Requirement~\ref{req:tuning:adaptive-default-version}), does not: it
consumes only the harmonic context's tonal centre.

\begin{lstlisting}[language=Rust]
pub struct HarmonicContext {
    /// Concurrently sounding pitches at the moment of resolution.
    pub concurrent: Vec<PitchId>,

    /// Recently sounding pitches, with decay weighting.
    pub recent: Vec<(PitchId, f64)>,

    /// The active key or modal center, if known.
    pub key_context: Option<KeyContext>,

    /// User-supplied context hints attached to the score.
    pub hints: Vec<ContextHint>,
}
\end{lstlisting}

The harmonic context is constructed by the audio engine (or by analysis
tooling) and passed to the tuning resolution function. Static tuning
systems ignore it; adaptive tuning systems consume it.

\begin{requirement}
  \label{req:tuning:adaptive-tuning-purity}
  Adaptive tuning resolution \MUST{} be a pure function of position
  and harmonic context. Implementations \MAY{} cache resolution
  results, but \MUST{} invalidate caches when the harmonic context
  changes.
\end{requirement}

\begin{requirement}
  \label{req:tuning:adaptive-default-version}
  \textbf{The built-in adaptive function, version 1.} The catalog
  identifier \texttt{ji-adaptive-5limit} resolves to
  \texttt{TuningResolution::Adaptive} with function identity
  \texttt{"default-v1"} --- the version is part of the
  machine-visible identifier string, not prose beside it. Version~1
  is a pure function of (position, anchor pitch class):
  \begin{itemize}
    \item the position resolves through the construction of
      Requirement~\ref{req:tuning:ji-static-construction}, transposed
      so the anchor takes the role of $1/1$;
    \item the anchor is the tonal centre supplied by the harmonic
      context (Requirement~\ref{req:tuning:adaptive-anchor-derivation}
      where that context is score-graph-derived); C (chromatic
      position~0) when no tonal centre is supplied;
    \item \texttt{concurrent}, \texttt{recent}, \texttt{hints},
      \texttt{parameters}, and mode are \emph{ignored} by version~1,
      exactly as \texttt{"default"} version~1 of
      Requirement~\ref{req:pitch:spelling-algorithm} ignores the
      context inputs it does not consult: consuming any of them is a
      new version, not a silent behaviour change under the same
      identity;
    \item every result derives from the anchor and the reference
      pitch alone. No adjustment is ever carried forward from a
      previous resolution, so the construction is comma-drift-free by
      shape, not by a correction step.
  \end{itemize}
  An unregistered or unknown \texttt{AdaptiveTuningFunctionId}
  \MUST{} be a hard error; there is no silent fallback.
  \texttt{AdaptiveTuningFunctionId} remains an extension point and
  other functions \MAY{} be registered, but the built-in
  \texttt{ji-adaptive-5limit} is bound to \texttt{"default-v1"} and
  that binding \MUSTNOT{} be overridden.
\end{requirement}

\begin{requirement}
  \label{req:tuning:adaptive-anchor-derivation}
  Where a \texttt{HarmonicContext}'s tonal centre is derived from a
  score graph, the anchor pitch class \MUST{} be computed from the
  prevailing key signature as
  \[
    \texttt{anchor\_pc} = (7 \times \texttt{fifths}) \bmod 12,
  \]
  reading \texttt{fifths} (\texttt{KeySignature}, Chapter~\ref{ch:graph})
  as the major tonic; mode is not consulted, so a signature of 0
  anchors at C whether the prevailing key is C major or A minor. The
  \emph{prevailing} key signature is the one on the staff containing
  the pitch being resolved, taken from that staff's
  \texttt{key\_sequence} (\texttt{StaffInstance::key\_sequence},
  Chapter~\ref{ch:graph}) at the latest \texttt{KeySignatureChange}
  whose anchor is at or before the pitch's onset. Both selections ---
  staff and time --- \MUST{} be made this way: without them, a
  polytonal or modulating score would resolve differently in two
  conforming implementations, which
  Requirement~\ref{req:tuning:tuning-resolution-determinism} forbids.
\end{requirement}

\section{Reference Pitch}
\label{sec:tuning:reference}

The reference pitch is the anchor that converts the tuning system's
structural ratios into absolute frequencies. It is a property of the
score, not of the tuning system.

\begin{lstlisting}[language=Rust]
pub struct ReferencePitch {
    /// The pitch-space position chosen as the reference.
    /// Conventionally A4 in CMN, but configurable.
    pub position: PitchSpacePosition,

    /// The frequency in Hertz assigned to that position.
    /// Conventional values: 440 Hz (modern standard), 415 Hz (Baroque),
    /// 430 Hz (early Classical), 442 Hz (some modern orchestras).
    pub frequency_hz: f64,
}
\end{lstlisting}

\begin{requirement}
  \label{req:tuning:reference-pitch}
  Every score \MUST{} declare a reference pitch. The reference pitch
  \MUST{} be expressible as a valid position within the score's default
  pitch space. When that position is \texttt{Cmn}, the default pitch space
  \MUST{} supply its chromatic layer and nominal mapping under
  Requirement~\ref{req:pitch:alteration-unit}. The frequency \MUST{} be
  positive and finite.
\end{requirement}

\subsection{Multiple Reference Pitches}

A score \MAY{} declare additional reference pitches at the staff or
region level, overriding the score default. This supports
mixed-instrumentation scores in which different instruments use
different concert-pitch references (e.g., a piece for period and modern
instruments with different A frequencies).

\section{Score Tuning Context and Hierarchical Resolution}
\label{sec:tuning:context}

The \emph{score tuning context} is the top-level structure declaring
the tuning environment for a score. Per-scope overrides extend it.

\begin{lstlisting}[language=Rust]
pub struct ScoreTuningContext {
    /// Default pitch space for the score.
    pub default_pitch_space: PitchSpaceId,

    /// Default tuning system. Its pitch_space must match (or be
    /// compatible with) the default_pitch_space.
    pub default_tuning_system: TuningSystemId,

    /// The reference pitch for the score.
    pub reference: ReferencePitch,

    /// Score-local accidental registry extensions.
    pub accidental_extensions: Vec<ScoreAccidentalExtensions>,

    /// SMuFL version targeting.
    pub smufl: SmuflVersionRequirement,

    /// Per-scope overrides, in the order they apply.
    pub overrides: Vec<TuningOverride>,
}

pub struct TuningOverride {
    pub scope: TuningScope,
    pub pitch_space: Option<PitchSpaceId>,
    pub tuning_system: Option<TuningSystemId>,
    pub reference: Option<ReferencePitch>,
}

pub enum TuningScope {
    Voice(VoiceId),
    Staff(StaffId),
    Region(RegionId),
    Range { start: TimeAnchor, end: TimeAnchor, voices: VoiceSelector },
}
\end{lstlisting}

\subsection{Resolution Order}

For a given pitch, resolution proceeds by walking outward through
scopes from most specific to most general until each component
(pitch space, tuning system, reference) is determined.

\begin{requirement}
  \label{req:tuning:tuning-resolution-order}
  Resolution \MUST{} proceed in the following order, halting at the
  first scope that supplies a non-inherited value for each component:

  \begin{enumerate}
    \item The pitch's own \texttt{AcousticPitch} (an explicit
      \texttt{TuningReference::Explicit} or
      \texttt{AcousticRealization::AbsoluteHz} short-circuits the
      walk).
    \item The voice containing the pitch.
    \item The staff containing the voice.
    \item Each region enclosing the pitch, from innermost to outermost.
    \item The score's default tuning context.
  \end{enumerate}

  Each of the three components (pitch space, tuning system, reference)
  is resolved independently along this chain.
\end{requirement}

\begin{rationale}
  Independent per-component resolution allows a passage to override
  only the parts that differ. A period-instrument inset in a modern
  orchestra score may override only the reference frequency (415 Hz
  instead of 440 Hz) while inheriting the pitch space and tuning
  structure from the parent context. Conversely, a microtonal passage
  in an otherwise 12-TET score may override the pitch space and
  tuning system but inherit the reference.
\end{rationale}

\subsection{Compatibility Constraints}

\begin{requirement}
  \label{req:tuning:tuning-system-compatibility}
  When a tuning override is applied, the resolved tuning system's
  declared \texttt{pitch\_space} \MUST{} either equal the resolved
  pitch space or be declared compatible with it via a registered
  compatibility mapping. Implementations \MUST{} reject configurations
  in which the resolved pitch space and the resolved tuning system's
  pitch space disagree without a compatibility mapping.
\end{requirement}

\section{Built-in Catalog}
\label{sec:tuning:builtin}

The core ships with a baseline catalog of pitch spaces, tuning systems,
and accidental registries sufficient for the common case. The catalog
is normative: a conforming implementation \MUST{} provide these
identifiers with the specified semantics.

\subsection{Built-in Pitch Spaces}

\begin{longtable}{p{3.5cm} p{10cm}}
  \toprule
  \textbf{Identifier} & \textbf{Description} \\
  \midrule
  \endhead
  \texttt{cmn-12} & Common Music Notation, 7 diatonic nominals (A--G)
  over 12 chromatic positions. Standard sharps, flats, double-sharps,
  and double-flats. The default for Western tonal and post-tonal music. \\
  \texttt{cmn-24} & CMN over 24 chromatic positions with
  \texttt{nominal\_to\_chromatic} $=[0,4,8,10,14,18,22]$ and quarter-tone
  alteration steps. A flat is $-2$ and a half-flat is $-1$. Used for
  Arabic-derived microtonal practice in CMN-compatible notation. \\
  \texttt{edo-19} & 19-tone equal division of the octave.
  Notable for its closer-to-just major thirds. \\
  \texttt{edo-22} & 22-tone equal division. Distinguishes harmonic
  intervals collapsed in 12-TET. \\
  \texttt{edo-31} & 31-tone equal division. Approximates quarter-comma
  meantone almost exactly. \\
  \texttt{edo-53} & 53-tone equal division. Excellent approximation
  of 5-limit JI. \\
  \texttt{edo-72} & 72-tone equal division. Common in research and
  contemporary microtonal practice. \\
  \texttt{ji-5limit} & 5-limit just intonation lattice with HEJI
  accidentals. Three-dimensional; prime basis $\{2, 3, 5\}$ in
  ascending order (leading component is the octave, prime 2). \\
  \texttt{ji-7limit} & 7-limit JI lattice with HEJI accidentals.
  Four-dimensional; prime basis $\{2, 3, 5, 7\}$ in ascending order. \\
  \texttt{ji-11limit} & 11-limit JI lattice with HEJI accidentals.
  Five-dimensional; prime basis $\{2, 3, 5, 7, 11\}$ in ascending
  order. \\
  \texttt{maqam-base} & Skeletal maqam framework with quarter-flat
  and quarter-sharp \texttt{CmnChromatic} accidentals, denominated by the
  space's chromatic layer under Requirement~\ref{req:pitch:alteration-unit}.
  Deep coverage of specific maqamat is the province of grammar plugins. \\
  \texttt{gamelan-slendro} & Five-tone slendro framework. \\
  \texttt{gamelan-pelog} & Seven-tone pelog framework. \\
  \bottomrule
\end{longtable}

\begin{rationale}
  The three JI spaces include prime 2 in their basis, per
  Requirement~\ref{req:pitch:ji-vector-basis}, rather than assuming octave
  equivalence. A \texttt{JiVector} is an absolute position, not a pitch
  class: without the prime-2 exponent, \texttt{ji-5limit} could not
  distinguish C4 from C5. None of these three declares octave-reduced
  equivalence, so full register is preserved, as that requirement's
  default already states.
\end{rationale}

\paragraph{Conformance note.}
At this revision the default spelling pre-pass is structurally
12-chromatic and therefore reports \texttt{spelling\_unavailable} for
\texttt{cmn-24} positions. This does not prevent storing or engraving an
explicitly authored spelling, nor does it change the space-relative
transposition algebra; automatic 24-chromatic spelling inference is deferred.

\subsection{Built-in Tuning Systems}

\begin{longtable}{p{4cm} p{9.5cm}}
  \toprule
  \textbf{Identifier} & \textbf{Description} \\
  \midrule
  \endhead
  \texttt{tet-12} & 12-tone equal temperament. The default. \\
  \texttt{tet-19} & 19-tone equal temperament. \\
  \texttt{tet-22} & 22-tone equal temperament. \\
  \texttt{tet-31} & 31-tone equal temperament. \\
  \texttt{tet-53} & 53-tone equal temperament. \\
  \texttt{tet-72} & 72-tone equal temperament. \\
  \texttt{pythagorean} & Pure-fifth (3:2) Pythagorean tuning. \\
  \texttt{meantone-1/4-comma} & Quarter-comma meantone (pure major
  thirds). \\
  \texttt{meantone-1/6-comma} & Sixth-comma meantone. \\
  \texttt{meantone-1/5-comma} & Fifth-comma meantone. \\
  \texttt{werckmeister-iii} & Werckmeister III well temperament. \\
  \texttt{werckmeister-iv} & Werckmeister IV well temperament. \\
  \texttt{vallotti} & Vallotti well temperament. \\
  \texttt{kirnberger-ii} & Kirnberger II well temperament. \\
  \texttt{kirnberger-iii} & Kirnberger III well temperament. \\
  \texttt{young-ii} & Thomas Young's second temperament. \\
  \texttt{ji-static-5limit-C} & Static 5-limit JI anchored to C tonic. \\
  \texttt{ji-static-5limit-G} & Static 5-limit JI anchored to G tonic. \\
  \texttt{ji-static-5limit-D} & Static 5-limit JI anchored to D tonic. \\
  \texttt{ji-adaptive-5limit} & Adaptive 5-limit JI; resolves based on
  harmonic context. \\
  \bottomrule
\end{longtable}

\begin{requirement}
  \label{req:tuning:builtin-tuning-catalog}
  All built-in catalog identifiers \MUST{} resolve in a conforming
  implementation. The semantics specified above are normative.
  Implementations \MAY{} provide additional built-ins; they \MUSTNOT{}
  redefine the semantics of those listed.
\end{requirement}

\subsection{Temperament Constructions}
\label{sec:tuning:temperament-constructions}

This subsection gives the normative construction for each of the ten
non-equal-tempered, non-just-intonation identifiers in the catalog
above. Per the pitch-space/tuning-system independence stated in
Section~\ref{sec:tuning:principles}, all ten are constructions over
\texttt{cmn-12}'s twelve chromatic positions per octave (C,
C$\sharp$/D$\flat$, D, \ldots, B) --- not over the open-ended
\texttt{ji-5limit} lattice pitch space, which is a separate built-in
with its own structure. Cents figures throughout are derived from the
stated ratio or fraction-of-comma, never primary.

\subsubsection{\texttt{pythagorean}}

A chain of eleven pure $3/2$ fifths (twelve notes); every other interval
is stacked fifths reduced by octaves. No comma is tempered anywhere in
the chain --- the entire Pythagorean comma ($531441/524288 \approx
23.460$ cents) is concentrated in the single interval where the chain
does not close.

The twelve-note selection is a choice of \emph{which} eleven consecutive
fifths to keep, not a consequence of ``pure fifths'' alone. The
conventional cut runs E$\flat$--B$\flat$--F--C--G--D--A--E--B--F$\sharp$--C$\sharp$--G$\sharp$
(C at the center), leaving the wolf on the diminished sixth
G$\sharp$--E$\flat$. An equally valid alternative cuts the spiral one
step the other way (chain D$\flat$--\ldots--F$\sharp$, wolf at
F$\sharp$--D$\flat$); either is a legitimate Pythagorean tuning. The
E$\flat$--G$\sharp$ cut is the construction this identifier denotes.

\begin{longtable}{p{1.8cm} p{4.5cm} p{4.5cm}}
  \toprule
  \textbf{Note} & \textbf{Ratio} & \textbf{Cents (derived)} \\
  \midrule
  \endhead
  C & $1/1$ & 0.000 \\
  C$\sharp$ & $2187/2048$ & 113.685 \\
  D & $9/8$ & 203.910 \\
  E$\flat$ & $32/27$ & 294.135 \\
  E & $81/64$ & 407.820 \\
  F & $4/3$ & 498.045 \\
  F$\sharp$ & $729/512$ & 611.730 \\
  G & $3/2$ & 701.955 \\
  G$\sharp$ & $6561/4096$ & 815.640 \\
  A & $27/16$ & 905.865 \\
  B$\flat$ & $16/9$ & 996.090 \\
  B & $243/128$ & 1109.775 \\
  \bottomrule
\end{longtable}

The closing wolf fifth G$\sharp\to$E$\flat$ (octave-reduced) is
$262144/177147 \approx 678.495$ cents.

\textbf{Closure (non-circulating, by design).} Eleven pure fifths at
701.955~c each, plus the closing wolf at 678.495~c:
$11 \times 701.955 + 678.495 = 8400.000$~c exactly (seven octaves, as
any assignment of twelve distinct pitch classes must sum to by
construction). Relative to twelve
\emph{pure} fifths ($12 \times 701.955 = 8423.460$~c), this construction
is short by exactly $8423.460 - 8400.000 = 23.460$~c --- one Pythagorean
comma, concentrated entirely on the single G$\sharp$--E$\flat$ interval
rather than distributed. That concentration is what makes the tuning
non-circulating: no pair of its twelve notes can be treated as
interchangeable fifths the way a well temperament's can.

\subsubsection{\texttt{meantone-1/4-comma}, \texttt{meantone-1/5-comma}, \texttt{meantone-1/6-comma}}

Meantone is a \emph{regular} temperament: every one of the twelve fifths
in the chain is tempered by the same fraction of the \emph{syntonic}
comma ($81/80 \approx 21.506$ cents) --- not the Pythagorean comma used
by the well temperaments below. All twelve fifths are tempered
uniformly, so there is no ``which fifths'' selection the way there is
for a well temperament; there is still a twelve-note chain-closing
choice, structurally identical to \texttt{pythagorean}'s, and the same
conventional cut places the wolf between G$\sharp$ and E$\flat$.

The tempered fifth is $(3/2) \cdot (80/81)^{1/n}$ for $1/n$-comma. For
$n=4$ this collapses to the exact form $5^{1/4}$; $n=5$ and $n=6$ are
irrational and are given only as derived cents.

\begin{longtable}{p{2.6cm} p{4.2cm} p{2.8cm} p{3.6cm}}
  \toprule
  \textbf{Variant} & \textbf{Fifth ratio} & \textbf{Cents (derived)} &
  \textbf{Narrowing vs.\ pure $3/2$} \\
  \midrule
  \endhead
  1/4-comma & $(3/2)(80/81)^{1/4} = 5^{1/4}$ exactly & 696.578 &
    5.377~c $= \tfrac{1}{4} \cdot 21.506$~c \\
  1/5-comma & $(3/2)(80/81)^{1/5}$ & 697.654 &
    4.301~c $= \tfrac{1}{5} \cdot 21.506$~c \\
  1/6-comma & $(3/2)(80/81)^{1/6}$ & 698.371 &
    3.584~c $= \tfrac{1}{6} \cdot 21.506$~c \\
  \bottomrule
\end{longtable}

Four ascending 1/4-comma fifths, lowered by two octaves, give the just
major third: $4 \times 696.578 - 2 \times 1200 = 386.31$~c, the cents of
$5/4$ --- the defining property of quarter-comma meantone.

\textbf{Closure (non-circulating, by design; all three).} Because all
twelve fifths are tempered equally, the residual wolf is the twelfth,
closing interval forced to bring the total to 8400~c (seven octaves):

\begin{longtable}{p{2.6cm} p{3.4cm} p{3.6cm} p{3.4cm}}
  \toprule
  \textbf{Variant} & \textbf{11 $\times$ tempered fifth} &
  \textbf{Wolf (12th, closing)} & \textbf{Wolf vs.\ pure $3/2$} \\
  \midrule
  \endhead
  1/4-comma & 7662.363~c & 737.637~c & $+35.682$~c (wide) \\
  1/5-comma & 7674.191~c & 725.809~c & $+23.854$~c (wide) \\
  1/6-comma & 7682.077~c & 717.923~c & $+15.968$~c (wide) \\
  \bottomrule
\end{longtable}

All three wolves land on the wide side of a pure fifth --- the opposite
sense from \texttt{pythagorean}'s narrow wolf.

\subsubsection{\texttt{werckmeister-iii}}

Fifths C--G, G--D, D--A, and B--F$\sharp$ are each narrowed by 1/4 of
the \emph{Pythagorean} comma; the remaining eight fifths are pure.
Because only four of twelve fifths are tempered and the pure fifths
absorb none of the comma, all twelve notes remain usable as a tonic:
there is no wolf.

\begin{longtable}{p{9.5cm} p{2.5cm}}
  \toprule
  \textbf{Fifths} & \textbf{Tempering} \\
  \midrule
  \endhead
  C--G, G--D, D--A, B--F$\sharp$ & narrow, 1/4 Pythagorean comma (4) \\
  A--E, E--B, F$\sharp$--C$\sharp$, C$\sharp$--G$\sharp$,
    G$\sharp$--E$\flat$, E$\flat$--B$\flat$, B$\flat$--F, F--C &
    pure (8) \\
  \bottomrule
\end{longtable}

Tempered-fifth cents: $(3/2)/(3^{12}/2^{19})^{1/4} \approx 696.090$
(narrowing 5.865~c $= \tfrac{1}{4} \cdot 23.460$~c).

\textbf{Closure.} $4 \times 5.865 = 23.460$~c exactly --- the chain
closes.

\subsubsection{\texttt{werckmeister-iv}}

Fifths C--G, D--A, E--B, F$\sharp$--C$\sharp$, and B$\flat$--F are each
narrowed by 1/3 of the Pythagorean comma; fifths G$\sharp$--D$\sharp$
and E$\flat$--B$\flat$ are each \emph{widened} by 1/3 of the Pythagorean
comma; the remaining five fifths are pure.

\begin{longtable}{p{9.5cm} p{2.5cm}}
  \toprule
  \textbf{Fifths} & \textbf{Tempering} \\
  \midrule
  \endhead
  C--G, D--A, E--B, F$\sharp$--C$\sharp$, B$\flat$--F &
    narrow, 1/3 Pythagorean comma (5) \\
  G$\sharp$--D$\sharp$, E$\flat$--B$\flat$ &
    wide, 1/3 Pythagorean comma (2) \\
  G--D, A--E, B--F$\sharp$, C$\sharp$--G$\sharp$, F--C & pure (5) \\
  \bottomrule
\end{longtable}

Narrow-fifth cents: $(3/2)/(3^{12}/2^{19})^{1/3} \approx 694.135$
(narrowing 7.820~c $= \tfrac{1}{3} \cdot 23.460$~c). Wide-fifth cents:
$(3/2) \cdot (3^{12}/2^{19})^{1/3} \approx 709.775$ (widening 7.820~c).
No wolf: all twelve notes remain usable as a tonic.

\textbf{Closure.} $5 \times 7.820 - 2 \times 7.820 = 3 \times 7.820 =
23.460$~c exactly --- the chain closes.

\subsubsection{\texttt{vallotti}}

Fifths F--C, C--G, G--D, D--A, A--E, and E--B (six consecutive) are each
narrowed by 1/6 of the Pythagorean comma; the other six fifths are pure.

\begin{longtable}{p{9.5cm} p{2.5cm}}
  \toprule
  \textbf{Fifths} & \textbf{Tempering} \\
  \midrule
  \endhead
  F--C, C--G, G--D, D--A, A--E, E--B & narrow, 1/6 Pythagorean comma (6) \\
  B--F$\sharp$, F$\sharp$--C$\sharp$, C$\sharp$--G$\sharp$,
    G$\sharp$--E$\flat$, E$\flat$--B$\flat$, B$\flat$--F & pure (6) \\
  \bottomrule
\end{longtable}

Tempered-fifth cents: $(3/2)/(3^{12}/2^{19})^{1/6} \approx 698.045$
(narrowing 3.910~c $= \tfrac{1}{6} \cdot 23.460$~c). No wolf.

\textbf{Closure.} $6 \times 3.910 = 23.460$~c exactly --- the chain
closes.

\subsubsection{\texttt{kirnberger-ii}}

Fifths D--A and A--E are each narrowed by 1/2 the syntonic comma;
fifth F$\sharp$--D$\flat$ (i.e.\ F$\sharp$--C$\sharp$, named with the
flat spelling because the chain below is built outward from
D$\flat$) is narrowed by a schisma (the difference between the
Pythagorean and syntonic commas, $32805/32768 \approx 1.9537$~c); the
remaining nine fifths are pure.

\begin{longtable}{p{9.5cm} p{2.5cm}}
  \toprule
  \textbf{Fifths} & \textbf{Tempering} \\
  \midrule
  \endhead
  D$\flat$--A$\flat$, A$\flat$--E$\flat$, E$\flat$--B$\flat$,
    B$\flat$--F, F--C, C--G, G--D, E--B, B--F$\sharp$ & pure (9) \\
  D--A, A--E & narrow, 1/2 syntonic comma (2) \\
  F$\sharp$--D$\flat$ (closing) & narrow, 1 schisma (1) \\
  \bottomrule
\end{longtable}

Cents by note, built by stacking the chain
D$\flat$--A$\flat$--E$\flat$--B$\flat$--F--C--G--D--A--E--B--F$\sharp$
from C~$= 1/1$:

\begin{longtable}{p{2.2cm} p{3.5cm}}
  \toprule
  \textbf{Note} & \textbf{Cents (derived)} \\
  \midrule
  \endhead
  C & 0.000 \\
  D$\flat$ & 90.225 \\
  D & 203.910 \\
  E$\flat$ & 294.135 \\
  E & 386.314 \\
  F & 498.045 \\
  F$\sharp$ & 590.224 \\
  G & 701.955 \\
  A$\flat$ & 792.180 \\
  A & 895.112 \\
  B$\flat$ & 996.090 \\
  B & 1088.269 \\
  \bottomrule
\end{longtable}

\textbf{Closure.} $2 \times 10.753$ (half-syntonic-comma fifths) $+ 1
\times 1.9537$ (schisma fifth) $= 21.5063 + 1.9537 = 23.4600$~c exactly
--- the chain closes.

\subsubsection{\texttt{kirnberger-iii}}

Fifths C--G, G--D, D--A, and A--E (four consecutive) are each narrowed
by 1/4 the syntonic comma; fifth F$\sharp$--D$\flat$ is narrowed by a
schisma (same position and same chain skeleton as
\texttt{kirnberger-ii}); the remaining seven fifths are pure.

\begin{longtable}{p{9.5cm} p{2.5cm}}
  \toprule
  \textbf{Fifths} & \textbf{Tempering} \\
  \midrule
  \endhead
  D$\flat$--A$\flat$, A$\flat$--E$\flat$, E$\flat$--B$\flat$,
    B$\flat$--F, F--C, E--B, B--F$\sharp$ & pure (7) \\
  C--G, G--D, D--A, A--E & narrow, 1/4 syntonic comma (4) \\
  F$\sharp$--D$\flat$ (closing) & narrow, 1 schisma (1) \\
  \bottomrule
\end{longtable}

Cents by note, same chain skeleton as \texttt{kirnberger-ii}, from
C~$= 1/1$:

\begin{longtable}{p{2.2cm} p{3.5cm}}
  \toprule
  \textbf{Note} & \textbf{Cents (derived)} \\
  \midrule
  \endhead
  C & 0.000 \\
  D$\flat$ & 90.225 \\
  D & 193.157 \\
  E$\flat$ & 294.135 \\
  E & 386.314 \\
  F & 498.045 \\
  F$\sharp$ & 590.224 \\
  G & 696.578 \\
  A$\flat$ & 792.180 \\
  A & 889.735 \\
  B$\flat$ & 996.090 \\
  B & 1088.269 \\
  \bottomrule
\end{longtable}

\textbf{Closure.} $4 \times 5.377$ (quarter-syntonic-comma fifths) $+ 1
\times 1.9537$ (schisma fifth) $= 21.5063 + 1.9537 = 23.4600$~c exactly
--- the chain closes.

\subsubsection{\texttt{young-ii}}

Thomas Young's \emph{second} temperament. Fifths C--G, G--D, D--A, A--E,
E--B, and B--F$\sharp$ (six consecutive) are each narrowed by 1/6 of the
Pythagorean comma; the remaining six fifths are pure.

\begin{longtable}{p{9.5cm} p{2.5cm}}
  \toprule
  \textbf{Fifths} & \textbf{Tempering} \\
  \midrule
  \endhead
  C--G, G--D, D--A, A--E, E--B, B--F$\sharp$ &
    narrow, 1/6 Pythagorean comma (6) \\
  F$\sharp$--C$\sharp$, C$\sharp$--G$\sharp$, G$\sharp$--E$\flat$,
    E$\flat$--B$\flat$, B$\flat$--F, F--C & pure (6) \\
  \bottomrule
\end{longtable}

This is the same six-tempered/six-pure, 1/6-Pythagorean-comma
construction as \texttt{vallotti}, rotated: \texttt{young-ii}'s
tempered run starts at C; \texttt{vallotti}'s starts at F. Tempered-fifth
cents are therefore identical to \texttt{vallotti}'s: $\approx 698.045$
(narrowing 3.910~c $= \tfrac{1}{6} \cdot 23.460$~c). No wolf.

\textbf{Closure.} $6 \times 3.910 = 23.460$~c exactly --- the chain
closes.

\subsection{The Static 5-Limit Construction}
\label{sec:tuning:ji-static}

\begin{requirement}
  \label{req:tuning:ji-static-construction}
  The static 5-limit just-intonation systems are constructed from the
  contiguous lattice block $\{3^a 5^b \mid a \in [-1,2],\ b \in
  [-1,1]\}$ --- twelve cells, octave-reduced, assigned in ascending
  order to the twelve chromatic positions of \texttt{cmn-12} starting
  from the anchor, which takes the role of $1/1$. The block is
  generated by its bounds; nothing is selected and nothing is
  discarded.
\end{requirement}

The twelve, anchor-relative (position 0 is the anchor):

\begin{longtable}{p{1.4cm} p{2.6cm} p{2.2cm} p{3.4cm}}
  \toprule
  \textbf{Step} & \textbf{Cell} & \textbf{Ratio} &
  \textbf{Cents (derived)} \\
  \midrule
  \endhead
  0 & $3^{0}5^{0}$ & $1/1$ & 0.000 \\
  1 & $3^{-1}5^{-1}$ & $16/15$ & 111.731 \\
  2 & $3^{2}5^{0}$ & $9/8$ & 203.910 \\
  3 & $3^{1}5^{-1}$ & $6/5$ & 315.641 \\
  4 & $3^{0}5^{1}$ & $5/4$ & 386.314 \\
  5 & $3^{-1}5^{0}$ & $4/3$ & 498.045 \\
  6 & $3^{2}5^{1}$ & $45/32$ & 590.224 \\
  7 & $3^{1}5^{0}$ & $3/2$ & 701.955 \\
  8 & $3^{0}5^{-1}$ & $8/5$ & 813.686 \\
  9 & $3^{-1}5^{1}$ & $5/3$ & 884.359 \\
  10 & $3^{2}5^{-1}$ & $9/5$ & 1017.596 \\
  11 & $3^{1}5^{1}$ & $15/8$ & 1088.269 \\
  \bottomrule
\end{longtable}

\texttt{ji-static-5limit-C}, \texttt{ji-static-5limit-G}, and
\texttt{ji-static-5limit-D} are this one construction at three anchors,
the named note taking the role of $1/1$. The catalog keeps exactly
these three rows; additional anchors are an implementation extension
and \MUSTNOT{} be read as a conformance obligation.

No 12-note set built this way is inversionally symmetric: $45/32$ is
present and its inverse $64/45$ is not.

\subsection{Default Score Configuration}

A newly-created score with no explicit tuning context \MUST{} default to:

\begin{itemize}
  \item Pitch space: \texttt{cmn-12}.
  \item Tuning system: \texttt{tet-12}.
  \item Reference pitch: A4 = 440\,Hz.
  \item SMuFL version: the highest SMuFL version supported by the
    implementation.
\end{itemize}

\section{Forward References}

\begin{itemize}
  \item \texttt{VoiceId}, \texttt{StaffId}, \texttt{RegionId},
    \texttt{VoiceSelector}, and the score graph's hierarchical
    organization are defined in Chapter~\ref{ch:graph}.
  \item \texttt{KeyContext} is not defined anywhere in this
    specification; completing it is out of scope for this document.
    Whatever completes it \MUST{} expose a tonal-centre chromatic
    pitch class consistent with
    Requirement~\ref{req:tuning:adaptive-anchor-derivation}.
    Chapter~\ref{ch:graph} defines \texttt{KeySignature}, not
    \texttt{KeyContext}. \texttt{ContextHint} and
    \texttt{AdaptiveTuningParameters} are likewise undefined; the
    built-in \texttt{ji-adaptive-5limit}, version~1
    (Requirement~\ref{req:tuning:adaptive-default-version}), ignores
    both, so defining them now would freeze a type surface on a
    chapter with no consumer.
  \item The specific spelling algorithms referenced by
    \texttt{SpellingAlgorithmId} remain an open question; see the
    related open question in Section~\ref{sec:pitch:prepass}.
\end{itemize}

% ===========================================================================
\chapter{The Score Graph}
\label{ch:graph}

This chapter specifies the score graph: the in-memory representation of
all musical content in a score. The graph is the canonical truth about
the music; layout, serialization, and editing operations are downstream
projections and consumers of it. This chapter is the largest in the
specification because it threads together every primitive defined in
Chapters~\ref{ch:pitch}--\ref{ch:tuning} into a coherent whole.

\section{Design Principles}
\label{sec:graph:principles}

\begin{description}
  \item[Hybrid topology.] The graph is a tree of containment overlaid
    with cross-cutting structures that hold references. The tree
    determines existence: deleting a parent deletes its descendants.
    Cross-cutting structures hold references that may become dangling
    and require re-anchoring per the rules defined in this chapter
    and in Chapter~\ref{ch:semops}.
  \item[Canvas before staff.] The spatial root of a score is the
    canvas, a 2D space partitioned into regions. Staves exist inside
    regions, not as a top-level concept. This permits regions of
    free graphic content to coexist with regions of staff-based CMN
    without forcing one paradigm to subsume the other.
  \item[Arena storage.] Events are stored in a flat arena owned by
    the score. Voices, cross-cutting structures, and analysis layers
    hold event identifiers, not events. Arena storage enables fast
    bulk iteration, simple cross-cutting reference, and clean CRDT
    semantics.
  \item[Parts are projections.] A part is a view onto the score
    defined by a part-extraction specification. Parts do not live
    in the canvas; they are separate objects that select and reshape
    canvas content for printed extraction.
  \item[Typed identifiers.] Every named object kind has its own
    identifier type. Cross-kind confusion is a compile-time error.
  \item[Identifiers are stable and deterministic.] Identifiers are
    assigned at creation, never reassigned, and generated by a scheme
    that is deterministic per replica.
\end{description}

\section{Top-Level Score Structure}
\label{sec:graph:score}

The score is the root object. All other objects in the specification
exist within or reference into a score. Conflict records are
deliberately \emph{not} part of the score graph: the
\texttt{ConflictRegistry} produced by canonical reduction is a
component of the canonical \emph{materialized state}
(Chapter~\ref{ch:semops}), recomputed with it, and persisting there
until resolved or dismissed by explicit operations.

\begin{lstlisting}[language=Rust]
pub struct Score {
    /// Bibliographic and authorship metadata.
    pub metadata: ScoreMetadata,

    /// The spatial root: regions and their staff/graphic content.
    pub canvas: Canvas,

    /// Abstract instrument definitions used by staves.
    pub instruments: Vec<Instrument>,

    /// Globally-identified staves of the score. Each region's
    /// staff instances reference these by StaffId.
    pub staves: Vec<Staff>,

    /// Optional staff groupings (grand staves, choral groups,
    /// bracketed orchestral sections). Used by engraving and parts.
    pub staff_groups: Vec<StaffGroup>,

    /// Per-part view definitions for extraction.
    pub parts: Vec<PartDefinition>,

    /// Cross-cutting structures: slurs, ties, hairpins, markers,
    /// graphic gestures, comments, and analytical annotations.
    pub cross_cutting: CrossCuttingRegistry,

    /// Tuning systems, pitch spaces, and reference pitch.
    pub tuning_context: ScoreTuningContext,

    /// The score-level tempo map. Regions may declare local tempo
    /// maps that override this.
    pub tempo_map: TempoMap,

    /// Flat arena of all events in the score.
    pub events: EventArena,

    /// Spelling attachments index (see Chapter 2).
    pub spelling_attachments: SpellingAttachmentIndex,

    /// Notational decomposition attachments index (see Chapter 3).
    pub decomposition_attachments: DecompositionAttachmentIndex,

    /// Analysis layers and view definitions.
    pub analysis_layers: Vec<AnalysisLayer>,
    pub views: Vec<ViewDefinition>,

    /// Replica identifier and identifier generation state.
    pub identity: IdentityContext,
}

pub struct StaffGroup {
    pub id: StaffGroupId,
    pub name: Option<String>,
    pub kind: StaffGroupKind,
    pub members: Vec<StaffId>,
}

pub enum StaffGroupKind {
    /// Grand staff: typically piano, organ, harp.
    GrandStaff,
    /// Bracketed group: orchestral sections.
    Bracket,
    /// Sub-bracket: divisi within a section.
    SubBracket,
    /// Choral group: SATB and similar.
    Choral,
    /// Custom group defined by a layout extension.
    Registered(StaffGroupKindRegistryId),
}
\end{lstlisting}

\subsection{Score Metadata}

\begin{lstlisting}[language=Rust]
pub struct ScoreMetadata {
    pub title: Option<String>,
    pub subtitle: Option<String>,
    pub composer: Option<String>,
    pub lyricist: Option<String>,
    pub arranger: Option<String>,
    pub copyright: Option<String>,
    pub creation_timestamp: Timestamp,
    pub modification_timestamp: Timestamp,
    pub additional: Vec<MetadataEntry>,
}

pub struct MetadataEntry {
    pub key: String,
    pub value: MetadataValue,
}

/// The value of an additional metadata entry (defined schema major 2).
/// Closed small union; growth is by appended variant.
pub enum MetadataValue {
    Text(String),
    Integer(i64),
    Flag(bool),
}

/// A calendar timestamp (defined schema major 2): nanoseconds since the
/// Unix epoch, UTC, no zone. Distinct from WallClockTime, which is
/// *performance* time within a score. The zero value is the "unset"
/// convention (genesis without a declared creation time).
pub struct Timestamp(pub i64);
\end{lstlisting}

The metadata structure is deliberately small. Extended bibliographic,
catalog, and editorial-apparatus metadata is the province of foreign-
format mappings and analytical layers, not the core.
\texttt{additional} is an \emph{ordered authored list}, not a map
(defined schema major 2): entry order is preserved verbatim in
canonical bytes, and duplicate keys are permitted --- foreign formats
legitimately carry repeated keys (multiple contributors, multiple
identifiers of one scheme). Consumers wanting map semantics take the
first entry per key.

\begin{requirement}
  \label{req:graph:metadata-timestamps}
  \textbf{Timestamps are strictly authored (ratified with schema
  major 2).} \texttt{creation\_timestamp} and
  \texttt{modification\_timestamp} are stored, \emph{authored} values:
  genesis takes the creation timestamp explicitly (the empty-document
  constructor defaults it to the epoch-zero unset convention), and a
  \texttt{SetMetadata} operation carries both verbatim as part of the
  metadata value. Implementations \MUSTNOT{} write either field
  implicitly --- an automatic modification-time update would rewrite
  canonical bytes on every edit, churning content addresses and breaking
  replica byte-agreement for identical operation histories. A client
  that wants a live modification time surfaces one from non-canonical
  state (the operation log's stamps); it does not store it here.
\end{requirement}

\section{Identifiers}
\label{sec:graph:ids}

Every named object in the graph carries a typed 128-bit identifier.

\begin{lstlisting}[language=Rust]
pub struct EventId(u128);
pub struct VoiceId(u128);
pub struct StaffId(u128);
pub struct StaffInstanceId(u128);
pub struct StaffGroupId(u128);
pub struct RegionId(u128);
pub struct InstrumentId(u128);
pub struct PartDefinitionId(u128);
pub struct MeasureId(u128);
pub struct BarlineAlignmentGroupId(u128);
pub struct SlurId(u128);
pub struct TieId(u128);
pub struct BeamId(u128);
pub struct SpannerId(u128);
pub struct MarkerId(u128);
pub struct AnalyticalAnnotationId(u128);
pub struct CommentId(u128);
pub struct GraphicObjectId(u128);
pub struct GraphicGestureId(u128);
pub struct TimeSignatureId(u128);
pub struct ConflictId(u128);
pub struct TransactionId(u128);
// ... and so on for every named object kind

/// A tagged identifier over every typed identifier kind in the
/// score graph. Used wherever an object is referenced generically
/// (cross-cutting endpoints, conflict affected_objects, repair
/// records, edit barriers). The variant tag is part of canonical
/// content: distinct variants with the same underlying u128 are
/// distinct TypedObjectIds.
pub enum TypedObjectId {
    Event(EventId),
    Pitch(PitchId),
    Voice(VoiceId),
    Staff(StaffId),
    StaffInstance(StaffInstanceId),
    StaffGroup(StaffGroupId),
    Region(RegionId),
    Instrument(InstrumentId),
    PartDefinition(PartDefinitionId),
    Measure(MeasureId),
    BarlineAlignmentGroup(BarlineAlignmentGroupId),
    Slur(SlurId),
    Tie(TieId),
    Beam(BeamId),
    Spanner(SpannerId),
    Marker(MarkerId),
    AnalyticalAnnotation(AnalyticalAnnotationId),
    Comment(CommentId),
    GraphicObject(GraphicObjectId),
    GraphicGesture(GraphicGestureId),
    TimeSignature(TimeSignatureId),
    AnalysisLayer(AnalysisLayerId),
    Tuplet(TupletId),
    RepeatStructure(RepeatStructureId),
    LyricLine(LyricLineId),
    ChordSymbol(ChordSymbolId),
    View(ViewId),
    /// Extension-defined object kind, identified by registry id
    /// plus the extension's own 128-bit identifier.
    Registered(ObjectKindRegistryId, u128),
}

impl TypedObjectId {
    /// Canonical byte form for hashing, ordering, and equality.
    /// Variant tag is encoded as a 16-bit big-endian discriminant
    /// followed by the variant payload's canonical bytes.
    pub fn canonical_bytes(&self) -> Vec<u8> { /* ... */ }
}
\end{lstlisting}

\begin{requirement}
  \label{req:graph:typed-object-id-discriminants}
  The \texttt{TypedObjectId} variant tag \MUST{} be encoded as a
  16-bit big-endian discriminant followed by the variant payload's
  canonical bytes. The discriminant of each variant is its
  declaration-order index in the following closed table; this
  assignment is normative and stable, and \MUSTNOT{} be reordered:

  \begin{center}
  \begin{tabular}{r l @{\qquad} r l}
    \toprule
    \textbf{Disc.} & \textbf{Variant} &
    \textbf{Disc.} & \textbf{Variant} \\
    \midrule
    0  & \texttt{Event}                  & 14 & \texttt{Spanner} \\
    1  & \texttt{Pitch}                   & 15 & \texttt{Marker} \\
    2  & \texttt{Voice}                   & 16 & \texttt{AnalyticalAnnotation} \\
    3  & \texttt{Staff}                   & 17 & \texttt{Comment} \\
    4  & \texttt{StaffInstance}           & 18 & \texttt{GraphicObject} \\
    5  & \texttt{StaffGroup}              & 19 & \texttt{GraphicGesture} \\
    6  & \texttt{Region}                  & 20 & \texttt{TimeSignature} \\
    7  & \texttt{Instrument}              & 21 & \texttt{AnalysisLayer} \\
    8  & \texttt{PartDefinition}          & 22 & \texttt{Tuplet} \\
    9  & \texttt{Measure}                 & 23 & \texttt{RepeatStructure} \\
    10 & \texttt{BarlineAlignmentGroup}   & 24 & \texttt{LyricLine} \\
    11 & \texttt{Slur}                    & 25 & \texttt{ChordSymbol} \\
    12 & \texttt{Tie}                     & 26 & \texttt{View} \\
    13 & \texttt{Beam}                    & 27 & \texttt{Registered} \\
    \bottomrule
  \end{tabular}
  \end{center}

  The object-kind set is closed at these 28 discriminants
  (\texttt{0}--\texttt{27}). Every named object kind in the score
  graph has a dedicated discriminant; kinds introduced by registered
  extensions \MUST{} use discriminant \texttt{27}
  (\texttt{Registered}) and are distinguished by their
  \texttt{ObjectKindRegistryId}, never by extending this table. The
  \texttt{Registered(reg, raw)} payload is encoded as the 128-bit
  \texttt{reg.canonical\_bytes()} (16 big-endian bytes) followed by
  the 16 big-endian bytes of the \texttt{raw} \texttt{u128}, for a
  total canonical form of $2 + 16 + 16 = 34$ bytes.
  \texttt{ObjectKindRegistryId} is a 128-bit value.
\end{requirement}

\subsection{Identifier Generation}

\begin{requirement}
  \label{req:graph:replica-counter-identifiers}
  Identifiers \MUST{} be 128-bit values composed of a 64-bit replica
  identifier and a 64-bit monotonic counter local to that replica.
  The replica identifier \MUST{} be unique among all replicas of the
  score; collision resistance is achieved by random initialization of
  the replica identifier at creation time with at least 64 bits of
  entropy from a cryptographically appropriate source. The counter
  \MUST{} be monotonically increasing within a replica and \MUSTNOT{}
  be reused, even if the corresponding object is deleted.
\end{requirement}

\begin{lstlisting}[language=Rust]
pub struct IdentityContext {
    /// This replica's identifier. Generated at score creation.
    pub replica_id: ReplicaId,

    /// Monotonic counter for new identifiers. Per-kind counters
    /// optional; a single counter suffices.
    pub next_counter: u64,
}

pub struct ReplicaId(u64);

impl ReplicaId {
    /// Reserved replica identifier for deterministically-derived
    /// system identifiers (system-promoted voices, content-derived
    /// conflict IDs, future deterministic synthetic identifiers).
    /// User-authored replicas MUST NOT use this value.
    pub const SYSTEM_DERIVED: ReplicaId = ReplicaId(0xffff_ffff_ffff_ffff);
}
\end{lstlisting}

\begin{rationale}
  Replica-plus-counter identifiers are deterministic given a replica
  and a creation order, which simplifies reasoning about CRDT merge
  semantics and reproducible builds of generated scores. UUIDv7 was
  considered and rejected: the timestamp embedding leaks ordering
  information that the CRDT layer would have to reconstruct anyway,
  and the random portion gives less collision resistance per bit than
  a properly initialized replica identifier.
\end{rationale}

\subsection{System-Derived Identifier Namespace}
\label{sec:graph:system-derived}

Some identifiers are not author-authored: they arise during
deterministic reduction (system-promoted voices, content-derived
conflict identifiers, attachment-deduplication identifiers). To
prevent collisions with user-authored identifiers, the spec
reserves a dedicated replica namespace.

\begin{requirement}
  \label{req:graph:system-derived-namespace}
  The replica identifier \texttt{ReplicaId::SYSTEM\_DERIVED}
  (\texttt{0xffff\_ffff\_ffff\_ffff}) is reserved for
  deterministically-derived system identifiers. User-authored
  replicas \MUSTNOT{} use this value as their \texttt{ReplicaId}.
  Implementations creating a new score \MUST{} reject this value
  when generating the local replica identifier and \MUST{}
  regenerate the random portion until a non-reserved value is
  obtained.

  Because a system-derived identifier is content-derived, the
  identified object's intrinsic content is \emph{immutable}: a
  reduction \MUST{} refuse an operation that would rewrite the
  intrinsic content of a \texttt{SYSTEM\_DERIVED}-namespace object
  in place (ratified Pass~12; the Operation Catalog pins the
  precondition and its failure reason). Rewriting the content would
  silently invalidate the id's content derivation; the sanctioned
  path is minting a replacement object.

  System-derived identifiers within the
  \texttt{ReplicaId::SYSTEM\_DERIVED} namespace \MUST{} have their
  64-bit counter portion derived deterministically by BLAKE3
  truncation:

  \begin{lstlisting}[language=Rust]
fn derive_system_counter(
    domain_tag: &[u8; 8],
    canonical_inputs: &[u8],
) -> u64 {
    let mut preimage = Vec::new();
    preimage.extend_from_slice(domain_tag);
    preimage.extend_from_slice(canonical_inputs);
    let hash = blake3(&preimage);
    u64::from_be_bytes(hash[0..8].try_into().unwrap())
}
  \end{lstlisting}

  The built-in (reserved) domain tags for system-derived identifiers
  are a closed set of three:
  \begin{itemize}
    \item \texttt{"MUSCSVCE"} for system-promoted voices
      (Section~\ref{sec:graph:promoted-voices}).
    \item \texttt{"MUSCSPCH"} for system-derived pitches (rare;
      only when promotion logic introduces a synthetic pitch;
      Section~\ref{sec:graph:system-derived-pitch}).
    \item \texttt{"MUSCSANM"} for integrity-anomaly identifiers
      (Section~\ref{sec:graph:system-collisions}). Anomalies are
      core, not an extension concern, so their tag is built in
      alongside the other two.
    \item Additional domain tags introduced by registered
      extensions \MUST{} begin with \texttt{"MUSCS"}, have length
      exactly 8 bytes, and \MUSTNOT{} collide with the three
      reserved tags above.
  \end{itemize}

  Two replicas reducing the same operation set \MUST{} derive
  identical system identifiers from identical canonical inputs.
\end{requirement}

\subsection{System-Derived Counter Collisions}
\label{sec:graph:system-collisions}

The 64-bit counter portion of a system-derived identifier is
derived by BLAKE3 truncation. Truncation to 64 bits offers
sufficient collision resistance against accidental collision
(roughly $2^{32}$ system-derived identifiers of the same kind
before a 50\% chance of any pair colliding), but is weaker than
the rest of the identifier story and is potentially vulnerable
to deliberate construction of colliding canonical inputs. A
collision means the identity layer has failed: canonical state
derived from a colliding identifier cannot be trusted.

\begin{lstlisting}[language=Rust]
/// An integrity anomaly that takes the document out of ordinary
/// canonical operation. Distinct from ConflictRecord (an
/// ordinary canonical-state fact): integrity anomalies indicate
/// that the structural assumptions underlying canonical state
/// have failed.
pub struct IntegrityAnomaly {
    pub id: IntegrityAnomalyId,
    pub kind: IntegrityAnomalyKind,
}

pub enum IntegrityAnomalyKind {
    /// A system-derived identifier counter collision.
    SystemIdentifierCollision {
        kind: ObjectKind,
        colliding_counter: u64,
        input_set_a: SerializedCanonicalInputs,
        input_set_b: SerializedCanonicalInputs,
    },
    /// An OperationSlot is Equivocated. The slot's candidates
    /// are referenced; reduction has excluded the operation.
    OperationSlotEquivocated {
        operation_id: OperationId,
    },
    /// A replica's stream has been quarantined per the HLC
    /// monotonicity rule. The AnomalousReplicaSegment is
    /// referenced.
    ReplicaStreamQuarantined {
        replica: ReplicaId,
        first_bad_counter: u64,
    },
    /// A registered anomaly defined by an extension or
    /// transport.
    Registered(IntegrityAnomalyRegistryId),
}

/// The kind of object whose system-derived identifier collided.
/// Open vocabulary, but intentionally narrower than TypedObjectId:
/// only the kinds actually minted into the system-derived namespace
/// appear here.
pub enum ObjectKind {
    Voice,   // system-promoted voices (MUSCSVCE)
    Pitch,   // content-derived synthetic pitches (MUSCSPCH)
    Registered(OperationKindRegistryId),
}
\end{lstlisting}

\begin{requirement}
  \label{req:graph:integrity-anomaly-id}
  An \texttt{IntegrityAnomalyId} \MUST{} be derived by the
  system-derived identifier function of
  Section~\ref{sec:graph:system-derived} with the reserved domain tag
  \texttt{"MUSCSANM"} over the anomaly kind's canonical bytes:
  \texttt{derive\_system\_id::<IntegrityAnomalyId>(b"MUSCSANM",
  \&kind.to\_canonical\_bytes())}. Because the identity is content-derived
  from the kind, two replicas observing the same structural failure
  derive the same anomaly identifier and therefore agree on anomaly
  identity across the network.
\end{requirement}

\begin{requirement}
  \label{req:graph:object-kind-vocab}
  \texttt{ObjectKind} (the kind field of
  \texttt{SystemIdentifierCollision}) is an \emph{open} vocabulary
  whose normative core set is exactly \texttt{Voice} and
  \texttt{Pitch}, with a \texttt{Registered} escape for
  extension-introduced system-derived kinds. It is deliberately
  narrower than \texttt{TypedObjectId}'s 28-kind table
  (Requirement~\ref{req:graph:typed-object-id-discriminants}):
  \texttt{ObjectKind} enumerates only the kinds that are actually
  minted into the \texttt{ReplicaId::SYSTEM\_DERIVED} namespace ---
  today, promoted voices and synthetic pitches --- because those are
  the only kinds for which a system-derived counter collision is
  possible. A new core variant is added here only when a new object
  kind begins being minted into the system-derived namespace.

  The discriminant byte of \texttt{ObjectKind} is its declaration-order
  index in the listing above: \texttt{Voice} $= 0$, \texttt{Pitch}
  $= 1$, \texttt{Registered} $= 2$. Because \texttt{ObjectKind} is the
  \texttt{kind} field of a \texttt{SystemIdentifierCollision} anomaly,
  this byte enters the canonical-input preimage of the resulting
  \texttt{IntegrityAnomalyId}
  (Requirement~\ref{req:graph:integrity-anomaly-id}); two replicas
  observing the same structural failure agree on its identity only if
  they agree on this byte. The assignment is therefore normative and
  stable and \MUSTNOT{} be reordered.
\end{requirement}

\begin{requirement}
  \label{req:graph:system-id-collision}
  Implementations \MUST{} perform a collision check during
  reduction whenever a new system-derived identifier is minted.
  If a newly derived 64-bit counter, within the same typed
  identifier kind and within the
  \texttt{ReplicaId::SYSTEM\_DERIVED} namespace, collides with a
  previously-derived counter from \emph{different} canonical
  inputs:

  \begin{itemize}
    \item An \texttt{IntegrityAnomaly} of kind
      \texttt{SystemIdentifierCollision} is recorded in the
      document's diagnostic integrity-anomaly register. This
      register is distinct from the canonical conflict registry:
      conflict records are ordinary canonical-state facts;
      integrity anomalies indicate the structural assumptions
      underlying canonical state have failed.
    \item Canonical reduction \MUSTNOT{} continue past the
      collision. The colliding object's identifier is held
      pending recovery: neither input set is allowed to occupy
      the collided counter in canonical state.
    \item Readers \MUST{} open such a bundle in diagnostic
      recovery mode, surfacing the collision prominently and
      prohibiting ordinary editing of the affected portion of
      canonical state.
    \item Recovery is by external action: a profile-declared
      deterministic policy, a transport-level credential
      revocation, or an explicit user reconciliation operation.
      No automatic in-spec resolution is defined.
  \end{itemize}

  Authoring implementations \SHOULD{} detect predictable
  collisions before authoring (e.g., by canonicalizing inputs
  with additional disambiguating data such as the
  causally-greatest operation visible at authoring time) and
  \SHOULDNOT{} rely on collision frequency being negligible in
  adversarial scenarios.
\end{requirement}

\begin{rationale}
  Moving system-derived collisions out of \texttt{ConflictKind}
  reflects the difference in kind between the two. A
  \texttt{ConflictRecord} is an ordinary canonical-state fact:
  the user has two valid musical edits that conflict, and the
  resolution is a normal editing action. A system-derived
  identifier collision is not a musical conflict; it is a
  structural failure of the identity layer. Treating it as a
  conflict would let canonical state grow under the collision,
  which is precisely the wrong behavior for an integrity
  failure. Treating it as a diagnostic anomaly halts canonical
  growth until external recovery resolves the structural issue.
\end{rationale}

\subsection{System-Derived Pitch Identity}
\label{sec:graph:system-derived-pitch}

When promotion logic introduces a synthetic pitch (rare), its
\texttt{PitchId} is content-addressed from the pitch's intrinsic
identity rather than allocated from a replica counter, so that all
replicas mint the same identifier for the same synthetic pitch.

\begin{requirement}
  \label{req:graph:system-derived-pitch-id}
  A system-derived \texttt{PitchId} \MUST{} be derived by the
  function of Section~\ref{sec:graph:system-derived} with domain tag
  \texttt{"MUSCSPCH"} over a fixed canonical byte form of the pitch's
  \emph{intrinsic identity}, defined as the pair (\emph{scale
  position}, \emph{acoustic realization}). All variable-width string
  components of this preimage (nominal names, accidental and tuning
  catalog references) \MUST{} be length-prefixed and \MUST{} be
  NFC-normalized at the derivation boundary, so that
  canonically-equivalent Unicode spellings derive the same identifier.

  The acoustic realization's \emph{tuning reference} is always part of
  intrinsic identity, including the \texttt{TuningReference::Inherit}
  case: \texttt{Inherit} is encoded as a distinct presence marker, not
  as the absence of a field. Consequently a synthetic pitch that
  inherits its tuning and one that names an explicit tuning are
  \emph{distinct} identities even when they would resolve to the same
  frequency in context. This makes synthetic-pitch identity stable
  under later changes to the inherited tuning environment, and aligns
  Invariant~11's uniqueness check with a fixed, context-free preimage.
\end{requirement}

\subsection{Identifier Stability}

\begin{requirement}
  \label{req:graph:identifier-non-reuse}
  An identifier, once assigned, \MUST{} never be reassigned, even
  after deletion of the corresponding object. Stale identifiers
  encountered in stored references \MUST{} be detected and reported
  by re-anchoring rules (see Chapter~\ref{ch:semops}); they
  \MUSTNOT{} silently match a different object.
\end{requirement}

\section{The Event Arena}
\label{sec:graph:arena}

Events are the rhythmic atoms of the score: pitched events, unpitched
events, rests, indeterminate events, trajectories, graphic events,
and cues. All events of all kinds in a score are stored in a single
arena.

\begin{lstlisting}[language=Rust]
pub struct EventArena {
    /// Storage backing for events. Implementation-defined; the
    /// observable contract is O(1) lookup by EventId and stable
    /// addresses for the lifetime of the event.
    storage: ArenaStorage<Event>,

    /// Auxiliary index for fast time-range queries.
    time_index: EventTimeIndex,
}
\end{lstlisting}

\begin{requirement}
  \label{req:graph:constant-time-event-lookup}
  Event lookup by \texttt{EventId} \MUST{} be $O(1)$ amortized.
  Implementations \MAY{} use any storage layout meeting this contract;
  vector-of-events with stable indices, slot-allocator schemes, and
  generational arenas are all acceptable.
\end{requirement}

\subsection{The Event Type}

\begin{lstlisting}[language=Rust]
pub enum Event {
    Pitched(PitchedEvent),
    Unpitched(UnpitchedEvent),
    Rest(Rest),
    Indeterminate(IndeterminateEvent),
    Trajectory(TrajectoryEvent),
    Graphic(GraphicEvent),
    Cue(CueEvent),
}
\end{lstlisting}

Every event variant carries:

\begin{itemize}
  \item Its own \texttt{EventId}.
  \item Its \emph{voice membership}: a \texttt{VoiceId} identifying
    the unique voice that owns it.
  \item Its \emph{position} within that voice (either a
    \texttt{MusicalPosition} or \texttt{WallClockTime}, depending on
    the enclosing region's time model).
\end{itemize}

Voice membership and position are stored on the event so that
operations addressing an event can resolve its context without
traversing parent containers.

\subsubsection{Event Position and Duration}
\label{sec:graph:event-time}

An event's temporal coordinates use union types matching its enclosing
region's time model. Position is unioned over musical and wall-clock
forms; duration is unioned over musical, wall-clock, and indeterminate
forms.

\begin{lstlisting}[language=Rust]
pub enum EventPosition {
    Musical(MusicalPosition),
    WallClock(WallClockTime),
}

pub enum EventDuration {
    Musical(MusicalDuration),
    WallClock(WallClockDuration),
    Indeterminate(DurationBounds),
}

/// Bounded interval expressing an indeterminate duration. Bounds are
/// concrete (not themselves indeterminate); this prevents recursive
/// indeterminacy in the type system.
pub struct DurationBounds {
    pub lower: Option<ConcreteDuration>,
    pub upper: Option<ConcreteDuration>,
}

pub enum ConcreteDuration {
    Musical(MusicalDuration),
    WallClock(WallClockDuration),
}
\end{lstlisting}

\begin{requirement}
  \label{req:graph:region-coordinate-kinds}
  Event position and duration variants \MUST{} agree with the time
  model of the enclosing region:

  \begin{itemize}
    \item Metric regions accept \texttt{EventPosition::Musical} and
      \texttt{EventDuration::Musical} only.
    \item Proportional regions accept \texttt{EventPosition::WallClock}
      and \texttt{EventDuration::WallClock} only.
    \item Aleatoric regions \MAY{} accept any combination consistent
      with the region's declared anchoring discipline (see
      Section~\ref{sec:graph:aleatoric-discipline}).
  \end{itemize}

  Implementations \MUST{} reject events whose coordinate variants
  contradict their enclosing region's time model.
\end{requirement}

\begin{requirement}
  \label{req:graph:duration-bound-kinds}
  If both bounds of a \texttt{DurationBounds} are present, they
  \MUST{} use the same \texttt{ConcreteDuration} variant, unless the
  enclosing aleatoric region's anchoring discipline explicitly
  permits mixed-kind bounds.
\end{requirement}

\subsubsection{Identified Pitches}
\label{sec:graph:identified-pitch}

Pitches embedded in events are wrapped in \texttt{IdentifiedPitch}
records that pair each pitch with a stable \texttt{PitchId}
(Section~\ref{sec:pitch:ids}). The identifier enables spelling
attachments, respelling operations, pitch-level reduction rules
under concurrent edits, and stable tie pairing through chord
reordering.

\begin{lstlisting}[language=Rust]
pub struct IdentifiedPitch {
    pub id: PitchId,
    pub pitch: Pitch,
}
\end{lstlisting}

\begin{requirement}
  \label{req:graph:identified-pitch-ownership}
  Every pitch embedded in an event \MUST{} be wrapped in an
  \texttt{IdentifiedPitch}. The contained \texttt{PitchId} \MUST{} be
  unique within the score's pitch-identity index. Pitch identifiers
  \MUST{} be minted atomically with their containing event: the
  operation creating an event mints \texttt{EventId} and the
  \texttt{PitchId}s of all embedded pitches in a single causal step.
\end{requirement}

\subsubsection{Pitched Events}

\begin{lstlisting}[language=Rust]
pub struct PitchedEvent {
    pub id: EventId,
    pub voice: VoiceId,
    pub position: EventPosition,
    pub duration: EventDuration,

    /// One or more identified pitches. A single pitch is a
    /// single-note event; multiple pitches are a chord.
    pub pitches: Vec<IdentifiedPitch>,

    /// Articulations attached directly to the event.
    pub articulations: Vec<ArticulationMark>,

    /// Single-event dynamic marking (e.g., sfz, fp). Spanning
    /// dynamics are cross-cutting structures, not event fields.
    pub dynamic: Option<DynamicMark>,

    /// Ornaments attached to the event. Ornament resolutions (the
    /// realized notes of a trill, mordent, etc.) are cross-cutting
    /// structures.
    pub ornaments: Vec<OrnamentMark>,

    /// Stem configuration: direction, length adjustment, hidden.
    pub stem: StemConfiguration,

    /// Whether this event is a grace note, and if so its kind.
    pub grace: Option<GraceKind>,
}

pub enum GraceKind {
    Acciaccatura,
    Appoggiatura,
    Unmeasured,
    MeasuredFraction(MusicalDuration),
}
\end{lstlisting}

\begin{requirement}
  \label{req:graph:pitched-event-nonempty}
  A \texttt{PitchedEvent} \MUST{} have at least one pitch. Empty
  pitch lists are forbidden; use \texttt{Rest} for the no-pitch case.
\end{requirement}

\subsubsection{Unpitched Events}

\begin{lstlisting}[language=Rust]
pub struct UnpitchedEvent {
    pub id: EventId,
    pub voice: VoiceId,
    pub position: EventPosition,
    pub duration: EventDuration,

    /// Staff line position for the unpitched event (e.g., snare
    /// drum at the middle line of a 5-line staff, or position 0
    /// for a single-line staff).
    pub staff_position: StaffPosition,

    /// The unpitched instrument identifier, for sound mapping.
    pub instrument_member: UnpitchedMemberId,

    pub articulations: Vec<ArticulationMark>,
    pub dynamic: Option<DynamicMark>,
    pub stem: StemConfiguration,
    pub grace: Option<GraceKind>,
}
\end{lstlisting}

\subsubsection{Rests}

\begin{lstlisting}[language=Rust]
pub struct Rest {
    pub id: EventId,
    pub voice: VoiceId,
    pub position: EventPosition,
    pub duration: EventDuration,

    /// Optional explicit vertical positioning. If None, the engraver
    /// chooses based on voice and meter context.
    pub vertical_position: Option<StaffPosition>,

    /// Whether this rest is visible. Invisible rests preserve timing
    /// but do not render.
    pub visible: bool,
}
\end{lstlisting}

\subsubsection{Indeterminate Events}

\begin{lstlisting}[language=Rust]
pub struct IndeterminateEvent {
    pub id: EventId,
    pub voice: VoiceId,
    pub position: EventPosition,
    pub duration: EventDuration,

    /// What aspect of the event is indeterminate.
    pub indeterminacy: IndeterminacyKind,

    /// Optional hints constraining the indeterminacy.
    pub hints: IndeterminacyHints,
}

pub enum IndeterminacyKind {
    /// "Any pitch in this range/contour."
    Pitch,

    /// "Any duration within these bounds." (Realized via
    /// EventDuration::Indeterminate on the event itself.)
    Duration,

    /// "Any of these alternatives."
    Choice,

    /// Combination of multiple indeterminacies.
    Compound(Vec<IndeterminacyKind>),
}

pub struct IndeterminacyHints {
    pub pitch_range: Option<PitchRange>,
    pub pitch_contour: Option<ContourHint>,
    pub density: Option<DensityHint>,
    pub textual_instruction: Option<String>,
    pub duration_bounds: Option<DurationBounds>,
    pub alternatives: Vec<EventId>,
}
\end{lstlisting}

\subsubsection{Trajectory Events}

A trajectory event represents a continuous pitch motion: a glissando,
portamento, or pitch bend.

\begin{lstlisting}[language=Rust]
pub struct TrajectoryEvent {
    pub id: EventId,
    pub voice: VoiceId,
    pub position: EventPosition,
    pub duration: EventDuration,

    /// Start and end pitch endpoints. Either may reference another
    /// event's pitch by id, or may carry an explicit identified
    /// pitch local to the trajectory.
    pub start: TrajectoryEndpoint,
    pub end: TrajectoryEndpoint,

    /// Shape of the trajectory between endpoints.
    pub shape: TrajectoryShape,

    /// Visual representation.
    pub display: TrajectoryDisplay,
}

pub enum TrajectoryEndpoint {
    /// Reference to a pitch belonging to another event.
    EventPitch(PitchId),

    /// An explicit identified pitch local to this trajectory.
    /// Such pitches participate in spelling, identity, and editing
    /// the same way embedded pitches do.
    ExplicitPitch(IdentifiedPitch),
}

pub enum TrajectoryShape {
    Linear,
    Exponential,
    Curve(Vec<TrajectoryControlPoint>),
    Stepwise(Vec<IdentifiedPitch>),
}
\end{lstlisting}

\subsubsection{Graphic Events}

A graphic event is an event whose primary content is graphic (drawn
or constructed) rather than notated through the standard staff system.
Used within hybrid and free-graphic regions for events that nevertheless
occupy a rhythmic slot in a voice.

\begin{lstlisting}[language=Rust]
pub struct GraphicEvent {
    pub id: EventId,
    pub voice: VoiceId,
    pub position: EventPosition,
    pub duration: EventDuration,

    /// Reference to the graphic object(s) realizing this event.
    pub graphics: Vec<GraphicObjectId>,

    /// Optional playback parameter bindings (see
    /// Section~\ref{sec:graph:timebinding}).
    pub playback_bindings: Vec<PlaybackBinding>,
}
\end{lstlisting}

\subsubsection{Cue Events}

A cue is a small-print rendering of music from another voice or
instrument, displayed as a guide.

\begin{lstlisting}[language=Rust]
pub struct CueEvent {
    pub id: EventId,
    pub voice: VoiceId,
    pub position: EventPosition,
    pub duration: EventDuration,

    /// The source events being cued.
    pub source: Vec<EventId>,

    /// Rendering hints: how the cue should be drawn (size, label,
    /// abbreviation of the source instrument's name, etc.).
    pub rendering: CueRendering,
}
\end{lstlisting}

\section{The Canvas}
\label{sec:graph:canvas}

The canvas is the spatial root of the score. It contains regions.

The score root and the canvas are \emph{structural givens}, not
operation products (ratified Pass~12): genesis is the creation of an
empty score together with its bundle, outside the operation set, and
neither object is ever minted, addressed, or deleted by an operation
(there is no \texttt{TypedObjectId} kind for either). This is
deliberate --- a document that could mint its own root would admit a
genesis race under concurrent editing, and every operation would need
a defined semantics against a not-yet-existing root. The decision may
be revisited only if an addressable multi-canvas model is adopted at
a future schema major.

\begin{lstlisting}[language=Rust]
pub struct Canvas {
    /// Regions in this canvas. Each region declares its time and
    /// content models and occupies a time and staff extent.
    pub regions: Vec<Region>,

    /// Canvas-level configuration: paper size, margins, default
    /// system layout, page-break behavior.
    pub layout_defaults: CanvasLayoutDefaults,
}

/// Canvas-level layout defaults. Page size and margins are in staff
/// spaces; a solver reads them as its default page geometry and MAY
/// override per solve. The type is defined with schema major 1 (its field
/// was named on the canvas from the first draft; only its type was
/// undefined).
pub struct CanvasLayoutDefaults {
    pub page_size: CanvasSize,
    pub margins: CanvasMargins,
}

/// A page size in staff spaces (1 staff space = staff height / 4).
pub struct CanvasSize {
    pub width: f64,
    pub height: f64,
}

/// Page margins in staff spaces.
pub struct CanvasMargins {
    pub top: f64,
    pub right: f64,
    pub bottom: f64,
    pub left: f64,
}
\end{lstlisting}

The default is \textbf{A4 portrait at an 8\,mm staff} (1 staff space =
2\,mm): page $105 \times 148.5$ staff spaces, $7.5$-staff-space margins,
hence a $90 \times 133.5$ content area. Each staff-space scalar is realized
canonically as a \texttt{CanonicalF64} leaf; the reference engraver maps a
\texttt{CanvasLayoutDefaults} onto its internal page geometry. This closes
the long-standing gap in which the canvas named \texttt{layout\_defaults} but
no chapter defined \texttt{CanvasLayoutDefaults}.

\subsection{Regions}

\begin{lstlisting}[language=Rust]
pub struct Region {
    pub id: RegionId,

    /// Time model: metric, proportional, or aleatoric.
    pub time_model: RegionTimeModel,

    /// Content model: staff-based, free graphic, or hybrid.
    pub content: RegionContent,

    /// Range in time that this region occupies.
    pub time_extent: TimeExtent,

    /// Range in vertical staff space that this region occupies.
    /// A region may cover all staves (full-score region) or a subset
    /// (a free-graphic gesture occupying only one staff's vertical
    /// band, for example).
    pub staff_extent: StaffExtent,

    /// Optional tempo and tuning overrides scoped to this region.
    pub local_tempo_map: Option<TempoMap>,
    pub local_tuning_overrides: Vec<TuningOverride>,

    /// Whether spanners (slurs first of all) may cross this region's
    /// boundary into an adjacent region. Default false: cross-region
    /// spanning is not permitted unless a region opts in. Added with
    /// schema major 1.
    pub permits_spanning_slurs: bool,
}
\end{lstlisting}

For a slur (or other spanner) whose endpoints lie in \emph{different}
regions, permission is governed conjunctively (ratified Pass~12):
the boundary is permeable only when \emph{both} endpoint regions set
\texttt{permits\_spanning\_slurs} --- a region that forbids spanning
is never crossed against its declaration, regardless of which side
the spanner starts on. The check is an authoring-time advisory only;
it never alters canonical reduction.

\begin{lstlisting}[language=Rust]

pub enum RegionContent {
    /// Staff-based notation: one or more staves carrying voices and
    /// events.
    StaffBased(StaffBasedContent),

    /// Free graphic content: no staves, only graphic objects placed
    /// in the region's coordinate space.
    FreeGraphic(GraphicContent),

    /// Hybrid: staves with overlaid graphic content. The graphic
    /// layer is rendered above (or below, configurably) the staff
    /// layer.
    Hybrid {
        staves: StaffBasedContent,
        overlay: GraphicContent,
        overlay_below_staves: bool,
    },
}

pub struct TimeExtent {
    pub start: TimeAnchor,
    pub end: TimeAnchor,
}

pub struct StaffExtent {
    /// Identifiers of the globally-identified staves whose content
    /// this region carries. A region carries one StaffInstance per
    /// listed StaffId. The same StaffId may appear in the
    /// staff_extent of adjacent regions; this is the mechanism by
    /// which a staff's identity persists across regions.
    pub staves: Vec<StaffId>,
}
\end{lstlisting}

\subsection{Region Overlap and Concurrency}

\begin{requirement}
  \label{req:graph:region-overlap-constraint}
  Regions \MAY{} overlap in time if their staff extents are disjoint.
  Regions \MUSTNOT{} overlap in both time and staff extent. The
  graph's construction \MUST{} reject configurations violating this
  constraint.
\end{requirement}

\begin{rationale}
  Concurrent regions on disjoint staves are essential for the hybrid
  notation feature: a metric staff and a proportional staff sounding
  simultaneously are two regions with overlapping time extents but
  disjoint staff extents. Overlapping time \emph{and} staff would mean
  two regions claim the same content, which is ambiguous.
\end{rationale}

\subsection{Staff-Based Content}

A staff-based region carries one or more staff \emph{instances}, each
of which manifests a globally-identified staff for the duration of the
region. The score's metric organization is per-staff (admitting
polymeter); barline alignment between staves is recorded as an
explicit relationship when meters coincide.

\begin{lstlisting}[language=Rust]
pub struct StaffBasedContent {
    /// Staff instances appearing in this region. Each references a
    /// globally-identified Staff and provides region-local content.
    pub staff_instances: Vec<StaffInstance>,

    /// Default metric grid for this region. Staves without their own
    /// local_metric_grid inherit this grid. None means the region
    /// has no default metric organization (acceptable for transitional
    /// regions or for regions where every staff declares its own).
    pub default_metric_grid: Option<MetricGrid>,

    /// Explicit barline alignment groups: places where measures
    /// across multiple staves are declared to coincide. Used for
    /// engraving barlines that span multiple staves under polymeter
    /// and for analytical purposes. Implementations MAY derive
    /// some alignment groups automatically from coincident measure
    /// starts; explicit groups are authoritative.
    pub barline_alignment_groups: Vec<BarlineAlignmentGroup>,

    /// System breaks declared by the user. The engraver may
    /// override these if region settings permit (see Chapter 6).
    pub user_system_breaks: Vec<TimeAnchor>,

    /// Page breaks declared by the user. Same override rule.
    pub user_page_breaks: Vec<TimeAnchor>,
}

/// A region-local manifestation of a globally-identified Staff.
pub struct StaffInstance {
    pub id: StaffInstanceId,

    /// The globally-identified staff this instance manifests.
    pub staff: StaffId,

    /// Voices appearing on this staff instance.
    pub voices: Vec<Voice>,

    /// Clef changes within this instance.
    pub clef_sequence: Vec<ClefChange>,

    /// Key signature changes within this instance.
    pub key_sequence: Vec<KeySignatureChange>,

    /// Optional local metric grid: if Some, overrides the region's
    /// default_metric_grid for this staff (polymetric case).
    /// If None, the staff inherits the region's default_metric_grid.
    pub local_metric_grid: Option<MetricGrid>,

    /// Measures belonging to this staff instance, in order. Measures
    /// are per-staff (admitting polymeter); alignment across staves
    /// is recorded in barline_alignment_groups.
    pub measures: Vec<Measure>,

    /// Optional per-instance instrument override. None means inherit
    /// from the global Staff's declared instrument.
    pub instrument_override: Option<InstrumentId>,

    /// Optional per-instance staff line configuration override.
    /// None means inherit from the global Staff's default.
    pub staff_lines_override: Option<StaffLineConfiguration>,

    /// Whether this instance is visible. Hidden instances preserve
    /// content but are not rendered (used for empty-staff suppression).
    pub visible: bool,
}

/// The clef content model (ratified schema major 2; the value types
/// landed with the visible slice and supersede the earlier ClefId
/// identifier sketch). A clef is a value, not a registry reference.
pub enum ClefShape {
    /// G clef (treble family) -- reference pitch G4.
    G,
    /// F clef (bass family) -- reference pitch F3.
    F,
    /// C clef (alto / tenor family) -- reference pitch middle C (C4).
    C,
    /// Unpitched percussion clef -- no diatonic reference.
    Percussion,
}

pub struct Clef {
    pub shape: ClefShape,
    /// The staff line (1 = bottom) the shape's reference pitch sits on.
    pub line: i8,
    /// Signed octave displacement (e.g., -1 for the vocal tenor G clef).
    pub octave_shift: i8,
}

/// A key signature on the circle of fifths, validated to the
/// conventional CMN range -7 (seven flats) ..= +7 (seven sharps).
/// Constructed only through a checked constructor; decoders reject
/// out-of-range values.
pub struct KeySignature {
    fifths: i8, // private; KeySignature::new rejects |fifths| > 7
}

/// A clef change at a point in a staff instance.
pub struct ClefChange {
    pub anchor: TimeAnchor,
    pub clef: Clef,
}

/// A key-signature change at a point in a staff instance.
pub struct KeySignatureChange {
    pub anchor: TimeAnchor,
    pub key: KeySignature,
}

/// A per-staff (or per-region-default) metric organization.
pub struct MetricGrid {
    /// Sequence of meter changes within the scope of this grid.
    pub meter_sequence: Vec<MeterChange>,
}

pub struct MeterChange {
    /// Where this meter takes effect, expressed as an anchor within
    /// the enclosing staff or region.
    pub anchor: TimeAnchor,

    /// The active time signature from this point until the next
    /// MeterChange or the end of the grid's scope.
    pub time_signature: TimeSignatureId,
}

/// A declaration that the measure boundaries of two or more staff
/// instances coincide at a particular point.
pub struct BarlineAlignmentGroup {
    pub id: BarlineAlignmentGroupId,

    /// The staff instances participating in this alignment, paired
    /// with the measures whose boundaries coincide.
    pub members: Vec<BarlineAlignmentMember>,
}

pub struct BarlineAlignmentMember {
    pub staff_instance: StaffInstanceId,
    pub measure: MeasureId,
    pub position: MeasurePosition,
}

pub struct Measure {
    pub id: MeasureId,
    pub start: TimeAnchor,
    pub time_signature: Option<TimeSignatureId>,
    pub explicit_number: Option<u32>,
    pub number_visibility: MeasureNumberVisibility,
}
\end{lstlisting}

\begin{requirement}
  \label{req:graph:staff-owned-measures}
  Measures \MUST{} belong to a single \texttt{StaffInstance}, not to
  the enclosing region. This admits polymeter: different staves in
  the same region \MAY{} carry different meter sequences and
  consequently different measure boundaries. Barline coincidences
  across staves \MUST{} be expressed via
  \texttt{BarlineAlignmentGroup} entries when explicit alignment is
  authoritative; otherwise alignment is derived at engraving time
  from coincident measure-start positions.
\end{requirement}

\section{Staves, Voices, and Instruments}
\label{sec:graph:staves}

\subsection{Instruments}

An instrument is an abstract definition of a sounding entity.

\begin{lstlisting}[language=Rust]
pub struct Instrument {
    pub id: InstrumentId,
    pub name: String,
    pub abbreviation: Option<String>,

    /// MIDI program, sound library mapping, or unpitched instrument
    /// definitions. The structure is detailed in the audio engine
    /// specification (out of scope here); the core stores the
    /// configuration opaquely.
    pub sound_config: SoundConfiguration,

    /// Default transposition (e.g., B-flat clarinet sounds a major
    /// second below written).
    pub transposition: Option<TranspositionInterval>,

    /// Default playable range. Engraving may flag out-of-range
    /// pitches; not a hard constraint.
    pub range: Option<PitchRange>,

    /// Default clef and staff line configuration for staves of this
    /// instrument. The clef is the content-bearing Clef value defined
    /// in this chapter (shape, line, octave shift) -- the value model
    /// ratified with the visible slice supersedes the earlier ClefId
    /// identifier sketch (reconciled schema major 2).
    pub default_clef: Clef,
    pub default_staff_lines: StaffLineConfiguration,

    /// For unpitched instruments: definitions of each playable
    /// member (snare, kick, ride bell, etc.) and their staff
    /// positions.
    pub unpitched_members: Vec<UnpitchedMember>,
}

/// Opaque sound configuration (defined schema major 2). The audio
/// engine specification owns the structure; the core stores the bytes
/// verbatim and never interprets them. Empty is the default.
pub struct SoundConfiguration(pub Vec<u8>);

// TranspositionInterval (defined schema major 2) is specified in
// Chapter 2, "Transposition and the Interval Type": the diatonic and
// chromatic step counts of an interval, and its action on a scale
// position. Instrument.transposition is the written-versus-sounding
// interval of a transposing instrument (a B-flat clarinet is
// -1 diatonic, -2 chromatic).
//
// Its ACTION is pinned (Chapter 2). What remains unimplemented is its
// AUTOMATIC APPLICATION at the instrument boundary: nothing in the core
// yet respells a written part into a sounding one, or resolves either to
// a frequency. Instrument.transposition is therefore still advisory --
// but for that reason, not for want of interval algebra.

/// One playable member of an unpitched instrument (defined schema
/// major 2): the identifier UnpitchedEvent.instrument_member
/// references, its display name, and its default staff position.
pub struct UnpitchedMember {
    pub member: UnpitchedMemberId,
    pub name: String,
    pub staff_position: StaffPosition,
}
\end{lstlisting}

\texttt{UnpitchedMemberId} is an \emph{instrument-scoped small value}
(a plain 32-bit integer), not a member of the 128-bit typed-identifier
family --- it is minted by the author of the instrument definition, not
by replicas, and participates in no content derivation. Member values
\emph{should} be unique within one instrument's
\texttt{unpitched\_members}; the authoring advisory layer flags a
duplicate, and resolution is by first match in list order (canonical
bytes preserve the authored order), so a duplicate is deterministic,
merely shadowed. An \texttt{UnpitchedEvent.instrument\_member} with no
matching entry resolves to \emph{no member definition}: the event still
renders (its own \texttt{staff\_position} field governs placement ---
the member's \texttt{staff\_position} is the \emph{authoring default}
copied onto new events, not a render-time override), and sound mapping
falls back to the instrument's \texttt{sound\_config}. This
no-match tolerance is load-bearing: every pre-major-2 instrument has an
empty member list while its events carry member values.
\texttt{StaffPosition}'s vertical convention (which integer is which
line, and its orientation) remains \emph{deliberately deferred} to the
Chapter~\ref{ch:layout-ir} vertical model, with the unpitched-rendering
tranche as its landing site; until then the field is preserved and
compared verbatim, and no implementation may claim conformant unpitched
\emph{rendering}.

\subsection{Staves: Identity Versus Instance}
\label{sec:graph:staff-identity}

The specification distinguishes two staff concepts:

\begin{description}
  \item[\texttt{Staff}] is the \emph{global, abstract} staff identity
    persisting across the entire score. A staff is the thing a
    performer is reading: the "Flute 1" staff, the "Piano upper
    staff." A staff is declared once in the score and is referenced
    by staff instances in any region where it appears.
  \item[\texttt{StaffInstance}] is the \emph{region-local manifestation}
    of a staff: the voices, measures, clef changes, key changes, and
    metric grid that belong to that staff for the duration of one
    region. A single \texttt{Staff} may have multiple
    \texttt{StaffInstance}s across the score; each instance belongs
    to exactly one region.
\end{description}

\begin{lstlisting}[language=Rust]
pub struct Staff {
    pub id: StaffId,

    /// Human-readable name (e.g., "Flute 1", "Violin II").
    pub name: String,

    /// Optional abbreviated name for subsequent systems.
    pub abbreviation: Option<String>,

    /// The instrument this staff renders by default. May be
    /// overridden per-instance.
    pub instrument: InstrumentId,

    /// Default clef for new instances of this staff (the
    /// content-bearing Clef value defined in this chapter; reconciled
    /// schema major 2).
    pub default_clef: Clef,

    /// Default staff line configuration.
    pub default_staff_lines: StaffLineConfiguration,

    /// Visual grouping: which staff group (if any) this staff
    /// belongs to (e.g., piano grand staff, choral group).
    pub group: Option<StaffGroupId>,
}

pub struct StaffLineConfiguration {
    pub line_count: u8,
    pub line_spacing: SpaceUnit,
    pub line_style: LineStyle,
    pub bracket: Option<StaffBracketKind>,
}

/// A dimension in staff spaces (defined schema major 2). Staff line
/// spacing is relative to the global staff space; 1.0 is a
/// normal-size staff, smaller values yield cue/ossia staves.
pub struct SpaceUnit(pub CanonicalF64);

/// A line drawing style (defined schema major 2). Shared by staff
/// lines and the slur/tie/spanner style records below.
pub enum LineStyle {
    Solid,
    Dashed,
    Dotted,
}

/// A per-staff bracket adornment (defined schema major 2), distinct
/// from StaffGroup-level bracketing.
pub enum StaffBracketKind {
    Brace,
    Bracket,
}
\end{lstlisting}

The score's top-level structure carries a list of \texttt{Staff}
objects alongside its list of \texttt{Instrument} objects; each
region's staff instances reference these by \texttt{StaffId}.

\begin{requirement}
  \label{req:graph:staff-instance-ownership}
  Every \texttt{StaffInstance.staff} \MUST{} resolve to a
  \texttt{Staff} declared at the score level. A single
  \texttt{StaffId} \MAY{} be referenced by multiple
  \texttt{StaffInstance}s in different regions: this is how a staff
  maintains continuous identity across the score's regions.

  A \texttt{StaffInstance} \MUST{} belong to exactly one region. The
  region's \texttt{staff\_instances} list \MUST{} contain the
  instance, and the instance's content is owned by that region.

  Continuous attachments (e.g., key signatures, clef changes, metric
  grids) \MAY{} appear on multiple staff instances of the same
  \texttt{Staff} when those instances are adjacent in time; the
  engraving pipeline reconstructs continuous staff appearance from
  per-instance content.
\end{requirement}

\subsection{Voices}

A voice is a polyphonic line within a staff instance. Voices hold
ordered references to events; the events themselves live in the
score's arena.

\begin{lstlisting}[language=Rust]
pub struct Voice {
    pub id: VoiceId,

    /// Ordered references to events in this voice, sorted by
    /// position. Each event has voice == this voice's id.
    pub events: Vec<EventId>,

    /// Default stem direction for this voice. Common conventions:
    /// voice 1 stems up, voice 2 stems down on the same staff.
    pub default_stem_direction: Option<StemDirection>,

    /// Whether this voice is the staff's primary voice (for purposes
    /// of rest placement, beam direction defaults, and similar
    /// engraving decisions).
    pub is_primary: bool,

    /// Provenance of this voice: user-declared, imported, or
    /// system-promoted by concurrent-edit conflict resolution.
    pub origin: VoiceOrigin,
}

pub enum VoiceOrigin {
    /// Created by an explicit user action.
    UserDeclared,

    /// Imported from a foreign format.
    Imported { format: ForeignFormatId },

    /// System-promoted to resolve a concurrent-edit collision.
    /// The original_voice is the voice into which the user had
    /// intended to insert; the promoted voice carries the event
    /// authored by losing_operation.
    SystemPromoted {
        winning_operation: OperationId,
        losing_operation: OperationId,
        original_voice: VoiceId,
    },
}
\end{lstlisting}

\begin{requirement}
  \label{req:graph:unbounded-voice-count}
  The number of voices per staff instance \MUST{} be unbounded in
  the data model. Engraving heuristics (Chapter~\ref{ch:layout-ir}
  and beyond) may flag visual issues with high voice counts; the
  graph admits any number.
\end{requirement}

\begin{requirement}
  \label{req:graph:event-voice-membership}
  Every event \MUST{} belong to exactly one voice. Voice membership
  is stored on the event (the \texttt{voice} field) and \MUST{} agree
  with the membership of the voice's \texttt{events} list:
  for any event $e$ in voice $v$'s events list, $e\texttt{.voice}$
  \MUST{} equal $v\texttt{.id}$.
\end{requirement}

\subsubsection{System-Promoted Voices}
\label{sec:graph:promoted-voices}

When concurrent operations cause an \texttt{InsertEvent} collision
that cannot be resolved by ordering alone (because both events
have positive duration and would overlap within the target voice),
the deterministically-losing event is placed in a newly allocated
voice with \texttt{origin: VoiceOrigin::SystemPromoted}. This
preserves the non-overlap invariant without rejecting user intent.

\begin{requirement}
  \label{req:graph:promoted-voice-id}
  The identifier of a system-promoted voice \MUST{} be derived
  deterministically from a fixed function of:

  \begin{enumerate}
    \item The enclosing staff instance's \texttt{StaffInstanceId}.
    \item The original target voice's \texttt{VoiceId}.
    \item The winning operation's \texttt{OperationId}.
    \item The losing operation's \texttt{OperationId}.
  \end{enumerate}

  The derivation is the system-derived identifier function of
  Section~\ref{sec:graph:system-derived} with domain tag
  \texttt{"MUSCSVCE"} over a fixed 64-byte preimage formed by
  concatenating the four identifiers in the order listed above, each
  as its 16-byte big-endian canonical form:

  \begin{lstlisting}[language=Rust]
fn derive_promoted_voice_id(
    staff_instance: StaffInstanceId,
    original_voice: VoiceId,
    winning_op: OperationId,
    losing_op: OperationId,
) -> VoiceId {
    // 4 * 16 = 64 bytes, each component big-endian.
    let mut inputs = Vec::with_capacity(64);
    inputs.extend_from_slice(&staff_instance.canonical_bytes());
    inputs.extend_from_slice(&original_voice.canonical_bytes());
    inputs.extend_from_slice(&winning_op.canonical_bytes());
    inputs.extend_from_slice(&losing_op.canonical_bytes());
    derive_system_id::<VoiceId>(b"MUSCSVCE", &inputs)
}
  \end{lstlisting}

  The enclosing staff instance is recovered from graph containment at
  reduction time; the operation does not store it redundantly.
  Determinism is required so that all replicas independently arrive at
  the same promoted voice identity. This derivation is the same one
  whose result Invariant~18 (Section~\ref{sec:graph:invariants})
  verifies; the reducer that mints the voice and the verifier that
  checks it consume one derivation.
\end{requirement}

\begin{requirement}
  \label{req:graph:promotion-generalization}
  The promotion rule generalizes from the pairwise case to any number
  of concurrent overlapping inserts into the same voice by an
  order-independent pre-pass:

  \begin{enumerate}
    \item Bucket the concurrent \texttt{InsertEvent} operations by
      their target \texttt{(staff\_instance, original\_voice)}.
    \item Within a bucket, walk the operations in ascending
      \texttt{OperationId} order, maintaining a \emph{retained set} of
      mutually non-overlapping inserts that stay in the original voice.
    \item An operation whose inserted interval does not overlap any
      member of the retained set joins the retained set. An operation
      whose interval overlaps some retained member is a \emph{loser}
      and is promoted; the lowest-\texttt{OperationId} retained member
      it overlaps is its \texttt{winning\_operation}.
  \end{enumerate}

  Because the walk is by \texttt{OperationId} and the retained set is
  built deterministically, every replica promotes the same operations
  and derives the same promoted-voice identities regardless of
  delivery order. ``Overlap'' is interval overlap of the half-open
  inserted spans \texttt{[start,\,end)}: the rule applies to
  \emph{partial} overlaps, not only to inserts that share an identical
  start position. This is a deliberate widening of the pairwise rule's
  identical-start-position language; it is the correct reading because
  any positive-duration overlap breaks the voice non-overlap invariant
  (Invariant~3), not just a shared onset.
\end{requirement}

\begin{requirement}
  \label{req:graph:promoted-voice-visibility}
  System-promoted voices are first-class graph objects: they appear
  in the staff instance's \texttt{voices} list, they participate in
  cross-cutting structures, and they are visible to users. They are
  \emph{not} hidden or transient. The layout pipeline
  (Chapter~\ref{ch:layout-ir}) \SHOULD{} provide a visual indication
  that a voice was system-promoted until the user normalizes it
  (via a later \texttt{NormalizePromotedVoice} or
  \texttt{MergeVoices} operation, specified in the operation
  catalog).

  Promoted voices \MUSTNOT{} silently become user-authored voices;
  doing so would convert a CRDT conflict resolution into musical
  authorship without the user's consent.
\end{requirement}

\begin{requirement}
  \label{req:graph:voice-event-nonoverlap}
  Within a voice, events \MUST{} be sorted by position and \MUSTNOT{}
  overlap in time. Concurrent material on the same staff \MUST{} be
  expressed via multiple voices.
\end{requirement}

\section{Cross-Cutting Structures}
\label{sec:graph:cross}

Cross-cutting structures hold references into the tree (typically
event identifiers) and represent musical phenomena that span multiple
events. The registry is partitioned by type for both type safety
and fast indexed lookup.

\begin{lstlisting}[language=Rust]
pub struct CrossCuttingRegistry {
    pub slurs: Vec<Slur>,
    pub ties: Vec<Tie>,
    pub beams: Vec<Beam>,
    pub tuplets: Vec<Tuplet>,
    pub spanners: Vec<Spanner>,
    pub markers: Vec<Marker>,
    pub repeats: Vec<RepeatStructure>,
    pub analytical: Vec<AnalyticalAnnotation>,
    pub comments: Vec<Comment>,
    pub graphic_gestures: Vec<GraphicGesture>,
    pub lyrics: Vec<LyricLine>,
    pub chord_symbols: Vec<ChordSymbol>,
}
\end{lstlisting}

\subsection{Slurs and Phrase Marks}
\label{sec:graph:slurs}

\begin{lstlisting}[language=Rust]
pub struct Slur {
    pub id: SlurId,

    /// First and last events under the slur.
    pub start_event: EventId,
    pub end_event: EventId,

    /// Slur kind: legato slur, phrase mark, articulation slur,
    /// editorial.
    pub kind: SlurKind,

    /// Optional curvature override. The engraver computes default
    /// curvature; this field overrides it.
    pub curvature_override: Option<CurvatureOverride>,

    /// Visual style: solid, dashed, dotted; line thickness curve.
    pub style: SpanStyle,
}

/// The class of a slur (defined schema major 2). The v1-to-v2
/// migration default is Legato.
pub enum SlurKind {
    /// An ordinary legato slur.
    Legato,
    /// A phrase mark (typically longer, over sub-phrases).
    Phrase,
    /// An articulation slur (e.g., over a two-note sigh figure).
    Articulation,
    /// An editorial slur (rendered distinctly, e.g., dashed).
    Editorial,
}

/// Which side of the notes a curve arcs toward (defined schema
/// major 2).
pub enum CurveDirection {
    Above,
    Below,
}

/// An authored curvature override (defined schema major 2). The
/// engraver computes default curvature; each present field overrides
/// that component of it. Consumed by the Standard engraving tier;
/// stored and preserved at every tier.
pub struct CurvatureOverride {
    pub direction: Option<CurveDirection>,
    /// Arc height at the apex.
    pub height: Option<SpaceUnit>,
}

/// The visual style of a spanning mark (defined schema major 2): one
/// shared record for Slur.style, Tie.style, and Spanner.style.
/// Defaults: solid, engraver-chosen thickness.
pub struct SpanStyle {
    pub line: LineStyle,
    /// Line thickness; None = the engraver's default.
    pub thickness: Option<SpaceUnit>,
}
\end{lstlisting}

\subsection{Ties}
\label{sec:graph:ties}

\begin{lstlisting}[language=Rust]
pub struct Tie {
    pub id: TieId,

    /// The two events being tied. The end event's pitch(es) must
    /// match the corresponding pitches in the start event.
    pub start_event: EventId,
    pub end_event: EventId,

    /// For chord ties: which pitches in the start event are tied
    /// to which in the end event, expressed by stable PitchId pairs.
    /// None means all pitches are tied by enharmonic matching in
    /// pitch-id-ascending order.
    pub pitch_pairing: Option<Vec<(PitchId, PitchId)>>,

    /// Tie class: determines the validation profile applied to this
    /// tie. The default Standard class requires same-voice immediate
    /// adjacency; other classes relax the rule for editorial,
    /// cross-voice, or notation-specific use.
    pub class: TieClass,

    /// Visual style (the shared SpanStyle record, schema major 2).
    pub style: SpanStyle,
}

pub enum TieClass {
    /// Standard CMN tie: same voice, immediate adjacency in voice
    /// order, pitches enharmonically equivalent.
    Standard,

    /// Editorial tie: connects pitches the editor considers
    /// continuous, spanning intervening rests or invisible events.
    /// Drawn with editorial style (dashed or bracketed) by default.
    Editorial,

    /// Cross-voice tie: connects pitches across voice boundaries
    /// within the same staff instance. Used for keyboard divisi,
    /// re-voicing across hands, and similar constructions.
    CrossVoice,

    /// Laissez-vibrer tie: a tie with no defined end event, drawn
    /// trailing from the start event. The end_event field is
    /// permitted to equal the start_event for this class, with
    /// pitch_pairing self-referential.
    LaissezVibrer,

    /// Registered class with custom validation behavior, defined
    /// by a notation grammar plugin.
    Registered(TieClassRegistryId),
}
\end{lstlisting}

\begin{rationale}
  Pairing by stable \texttt{PitchId} rather than by positional index
  preserves tie identity through chord-reordering edits and through
  CRDT-style concurrent modification of the chord's pitch list.
  Index-based pairing would silently retarget when pitches were
  inserted, deleted, or reordered within a chord.

  The tie class permits common editorial and contemporary-notation
  practices (editorial dashed ties spanning rests, cross-voice ties
  in keyboard scores, laissez-vibrer ties without a notated end) to
  be expressed without relaxing the structural invariants of the
  standard tie.
\end{rationale}

\begin{requirement}
  \label{req:graph:tie-class-validation}
  Tie validation is class-specific:

  \begin{itemize}
    \item \texttt{TieClass::Standard}: the start and end events
      \MUST{} be in the same voice; the start event \MUST{}
      immediately precede the end event in that voice's event
      sequence; the paired pitches \MUST{} be enharmonically
      equivalent under the active tuning system.
    \item \texttt{TieClass::Editorial}: the start and end events
      \MUST{} be in the same voice and the start \MUST{} precede the
      end in voice order, but intervening events are permitted.
      Paired pitches \MUST{} be enharmonically equivalent.
    \item \texttt{TieClass::CrossVoice}: the start and end events
      \MAY{} be in different voices on the same staff instance. The
      start's resolved position \MUST{} be less than or equal to
      the end's. Paired pitches \MUST{} be enharmonically equivalent.
    \item \texttt{TieClass::LaissezVibrer}: the end event \MAY{}
      equal the start event, in which case the tie has no notated
      destination. No adjacency requirement applies.
    \item \texttt{TieClass::Registered}: validation is defined by
      the registered class's contract; the registry \MUST{} specify
      adjacency and pitch-equivalence requirements for the class.
  \end{itemize}
\end{requirement}

\subsection{Beams}

\begin{lstlisting}[language=Rust]
pub struct Beam {
    pub id: BeamId,
    pub events: Vec<EventId>,
    pub level: u8,                  // primary beam level

    /// Sub-beam structure for inner subdivisions. Each entry
    /// describes a beam segment at a deeper level.
    pub sub_beams: Vec<SubBeam>,

    /// Beam angle and stem-length overrides.
    pub geometry_override: Option<BeamGeometryOverride>,
}

/// A beam segment at a deeper subdivision level (defined schema
/// major 2): a contiguous subset of the owning beam's events beamed
/// together at `level` (strictly deeper than the owning beam's
/// primary level).
pub struct SubBeam {
    pub level: u8,
    pub events: Vec<EventId>,
}

/// An authored beam-geometry override (defined schema major 2). Each
/// present field overrides the engraver's computed geometry. Consumed
/// by the Standard engraving tier; stored and preserved at every tier.
pub struct BeamGeometryOverride {
    /// Beam slope: staff spaces of rise per staff space of run
    /// (dimensionless, hence not a SpaceUnit).
    pub slope: Option<CanonicalF64>,
    /// Vertical displacement of the beam from its default placement
    /// (positive = up).
    pub offset: Option<SpaceUnit>,
}
\end{lstlisting}

\subsection{Spanners}

A spanner is a generic spanning mark: hairpins, octave lines, pedal
lines, trill extensions, and similar.

\begin{lstlisting}[language=Rust]
pub struct Spanner {
    pub id: SpannerId,
    pub kind: SpannerKind,
    pub start: TimeAnchor,
    pub end: TimeAnchor,

    /// Which staff or staves this spanner attaches to. Most spanners
    /// target a single staff; some (system-level dynamics, pedal
    /// lines spanning grand staff) target multiple.
    pub staves: Vec<StaffId>,

    /// Visual style (the shared SpanStyle record, schema major 2).
    pub style: SpanStyle,
}

pub enum SpannerKind {
    /// An unclassified spanning mark: renders as a plain line or
    /// bracket. First variant deliberately (schema major 2): it is
    /// the v1-to-v2 migration default -- a v1 spanner carried no
    /// kind, and Generic is the honest translation of that absence.
    Generic,
    Hairpin(HairpinDirection),
    OctaveLine(OctaveOffset),
    PedalLine(PedalKind),
    TrillExtension,
    Glissando,
    Portamento,
    TextLine(TextLineDefinition),
    Bracket(BracketKind),
}

/// Hairpin orientation (defined schema major 2).
pub enum HairpinDirection {
    Crescendo,
    Diminuendo,
}

/// Octave-line displacement in signed octaves (defined schema
/// major 2): +1 = 8va, -1 = 8vb, +2 = 15ma, -2 = 15mb. Zero is
/// representable but degenerate; the authoring advisory layer flags
/// it, reduction does not.
pub struct OctaveOffset(pub i8);

/// Pedal-line kind (defined schema major 2).
pub enum PedalKind {
    Sustain,
    Sostenuto,
    UnaCorda,
}

/// A text line's content (defined schema major 2); the dash pattern
/// comes from the spanner's style.
pub struct TextLineDefinition {
    pub text: String,
}

/// A bracket spanner's shape (defined schema major 2). Growth is by
/// appended variant.
pub enum BracketKind {
    Square,
}
\end{lstlisting}

\subsection{Markers}

A marker is a point annotation: rehearsal marks, section labels,
tempo text, fermatas (which technically attach to events but may
also be region-level), and similar.

\begin{lstlisting}[language=Rust]
pub struct Marker {
    pub id: MarkerId,
    pub anchor: TimeAnchor,
    pub kind: MarkerKind,
    pub display: MarkerDisplay,
}

pub enum MarkerKind {
    Rehearsal(RehearsalMark),
    Section(SectionLabel),
    Tempo(TempoMarking),
    Coda,
    Segno,
    DalSegno,
    DaCapo,
    Fine,
    Custom(CustomMarkerId),
}
\end{lstlisting}

\subsection{Repeat Structures}
\label{sec:graph:repeats}

\begin{lstlisting}[language=Rust]
pub struct RepeatStructure {
    pub id: RepeatStructureId,
    pub kind: RepeatKind,
    pub start: TimeAnchor,
    pub end: TimeAnchor,

    /// For voltas: which endings apply on which passes.
    pub voltas: Vec<Volta>,
}

pub enum RepeatKind {
    SimpleRepeat { count: u32 },
    DaCapo { end_target: TimeAnchor },
    DalSegno { segno: TimeAnchor, end_target: TimeAnchor },
    Volta,
}

/// One volta bracket (defined schema major 2): the passes it applies
/// on and the time span it covers. A RepeatStructure of kind Volta
/// carries one entry per ending (e.g., a first ending Volta with
/// endings = [1] and a second with endings = [2]).
pub struct Volta {
    /// The pass numbers this ending plays on (1-based), ascending.
    pub endings: Vec<u32>,
    pub start: TimeAnchor,
    pub end: TimeAnchor,
}
\end{lstlisting}

The \texttt{endings} constraints (non-empty, 1-based, strictly
ascending) are \emph{advisory}, on the \texttt{OctaveOffset} pattern:
a violating value is representable, decoders and reduction accept it,
and the authoring validation layer flags it. Rendering treats an
unlistable ending set as a plain bracket.

The v1-to-v2 migration default for a kindless v1 repeat structure is
\texttt{SimpleRepeat} with \texttt{count}~$=$~2: a v1
\texttt{RepeatStructure} \emph{meant} a repeat, and playing the span
twice is the conventional semantics of an unadorned repeat sign.
Voltas default empty.

\subsection{Analytical Annotations}

Analytical annotations record analyses of the music: Roman numerals,
form labels, motivic brackets, set-theory analyses, Schenkerian
graphs. They are first-class and may appear on dedicated analysis
layers.

\begin{lstlisting}[language=Rust]
pub struct AnalyticalAnnotation {
    pub id: AnalyticalAnnotationId,
    pub kind: AnalyticalKind,
    pub anchor: AnnotationAnchor,

    /// Optional analysis layer. None means the engraved layer;
    /// Some(id) targets a specific analysis layer.
    pub layer: Option<AnalysisLayerId>,

    pub content: AnalyticalContent,
}

pub enum AnalyticalKind {
    RomanNumeral,
    FunctionalLabel,
    SetTheoryLabel,
    SchenkerianGraph,
    FormSection,
    MotivicBracket,
    Custom(CustomAnalyticalKindId),
}

pub enum AnnotationAnchor {
    Event(EventId),
    Range { start: TimeAnchor, end: TimeAnchor },
    Region(RegionId),
}
\end{lstlisting}

\subsection{Comments}

Comments are review-mode annotations: discussion threads anchored to
points or ranges in the score.

\begin{lstlisting}[language=Rust]
pub struct Comment {
    pub id: CommentId,
    pub anchor: AnnotationAnchor,
    pub thread: Vec<CommentMessage>,
    pub resolved: bool,
}

pub struct CommentMessage {
    pub author: AuthorId,
    pub timestamp: Timestamp,
    pub body: String,
}
\end{lstlisting}

\subsection{Graphic Gestures}

A graphic gesture is a drawn or constructed graphic element that
spans events, staves, or regions and is not contained within a single
graphic region. Used for performance-direction sweeps,
multi-staff curved gestures, and similar.

\begin{lstlisting}[language=Rust]
pub struct GraphicGesture {
    pub id: GraphicGestureId,

    /// The graphic objects composing this gesture. Stored in the
    /// canvas's graphic-object storage.
    pub objects: Vec<GraphicObjectId>,

    /// Anchoring: how this gesture is positioned relative to score
    /// content.
    pub anchoring: GestureAnchoring,

    /// Optional playback parameter bindings.
    pub playback_bindings: Vec<PlaybackBinding>,
}

pub enum GestureAnchoring {
    /// Anchored to events: the gesture moves with them.
    Events(Vec<EventId>),

    /// Anchored to a time and staff range.
    Range {
        start: TimeAnchor,
        end: TimeAnchor,
        staves: Vec<StaffId>,
    },

    /// Free canvas coordinates: does not follow score edits.
    Free(CanvasCoordinates),
}
\end{lstlisting}

\section{Graphic Content}
\label{sec:graph:graphic}

Free-graphic and hybrid regions carry graphic content: vector
primitives, drawn strokes, text, and images. Graphic content is
expressed in the region's local coordinate space.

\begin{lstlisting}[language=Rust]
pub struct GraphicContent {
    pub objects: Vec<GraphicObject>,

    /// The region's local coordinate system. Region transforms
    /// map this to canvas coordinates.
    pub coordinate_system: LocalCoordinateSystem,
}
\end{lstlisting}

\subsection{Graphic Objects}

\begin{lstlisting}[language=Rust]
pub struct GraphicObject {
    pub id: GraphicObjectId,

    /// The object's geometric kind.
    pub kind: GraphicObjectKind,

    /// Transform applied within the region's coordinate space.
    pub transform: Transform2D,

    /// Layer ordering: higher values draw on top of lower.
    pub layer: i32,

    /// Visual style: stroke color, fill, line weight, opacity.
    pub style: GraphicStyle,

    /// Optional time binding. See Section~\ref{sec:graph:timebinding}.
    pub time_binding: Option<TimeBinding>,
}

pub enum GraphicObjectKind {
    Path(PathData),
    Shape(ShapePrimitive),
    Text(TextData),
    Image(ImageData),
    Group(Vec<GraphicObjectId>),
    Stroke(StrokeData),
}

pub enum ShapePrimitive {
    Line { start: Point2D, end: Point2D },
    Rectangle { origin: Point2D, size: Size2D },
    Ellipse { center: Point2D, radii: Size2D },
    Polygon(Vec<Point2D>),
    Bezier(Vec<BezierSegment>),
}

pub struct StrokeData {
    /// Pressure-sensitive stylus stroke samples: position, pressure,
    /// tilt, and time per sample.
    pub samples: Vec<StrokeSample>,
}

pub struct StrokeSample {
    pub position: Point2D,
    pub pressure: f32,        // 0.0..=1.0
    pub tilt_xy: (f32, f32),  // radians
    pub timestamp: WallClockTime,
}
\end{lstlisting}

\subsection{Coordinates}

\begin{requirement}
  \label{req:graph:region-local-coordinates}
  Graphic objects \MUST{} be stored in region-local coordinates. The
  region's transform maps local coordinates to canvas coordinates.
  This permits regions to be moved, resized, or otherwise transformed
  without rewriting the coordinates of every contained object.
\end{requirement}

\subsection{Time Bindings}
\label{sec:graph:timebinding}

A graphic object \MAY{} carry an optional time binding, which assigns
a temporal meaning to the object beyond its visual appearance. Time
bindings are the mechanism by which a drawn curve becomes a performance
instruction (a filter sweep, an amplitude envelope, a density
trajectory).

\begin{lstlisting}[language=Rust]
pub struct TimeBinding {
    /// What musical or wall-clock range this graphic spans.
    pub time_extent: TimeExtent,

    /// What parameter, if any, this graphic drives. None means the
    /// graphic is decorative; Some specifies a named parameter.
    pub parameter: Option<ParameterBinding>,

    /// How the graphic's geometry maps to parameter values.
    pub mapping: ParameterMapping,
}

pub struct ParameterBinding {
    pub target: ParameterTarget,
    pub parameter_name: String,
    pub value_range: ParameterRange,
}

pub enum ParameterTarget {
    /// A specific event's parameter (e.g., this event's velocity).
    Event(EventId),

    /// A voice-level parameter (e.g., voice expression CC).
    Voice(VoiceId),

    /// A staff-level parameter.
    Staff(StaffId),

    /// A region-level parameter.
    Region(RegionId),

    /// A score-global parameter (e.g., master gain).
    Score,
}

pub enum ParameterMapping {
    /// Vertical position in region coordinates maps linearly to
    /// parameter value.
    VerticalLinear,

    /// Vertical position with exponential mapping.
    VerticalExponential,

    /// Stroke pressure maps to parameter value.
    StrokePressure,

    /// Stroke velocity maps to parameter value.
    StrokeVelocity,

    /// Custom mapping defined by a plugin.
    Registered(MappingRegistryId),
}
\end{lstlisting}

\begin{rationale}
  Time bindings make graphic-notation playback expressible in the core
  data model without committing to specific audio engine semantics.
  The core stores what parameter a graphic drives and how its
  geometry maps to values; the audio engine spec consumes this
  binding and applies it. Decorative graphics (binding is \texttt{None})
  remain pure visual content.
\end{rationale}

\section{Parts}
\label{sec:graph:parts}

A part is a per-instrument or per-section view onto the score,
intended for printed extraction or focused performer use. Parts live
separately from the canvas; they are projections defined by a
part-extraction specification.

\begin{lstlisting}[language=Rust]
pub struct PartDefinition {
    pub id: PartDefinitionId,

    /// Human-readable name (e.g., "Flute 1", "Violins II").
    pub name: String,

    /// Which staves are included in this part, in order from top
    /// to bottom.
    pub staves: Vec<StaffId>,

    /// Part-specific layout overrides: page size, system count,
    /// cue insertions, multirest consolidation, page-turn placement.
    pub layout_overrides: PartLayoutOverrides,

    /// Per-event visibility overrides specific to this part (e.g.,
    /// hide a cue in the full score but show it in the part).
    pub visibility_overrides: Vec<VisibilityOverride>,

    /// Auto-cue insertions: source staves to draw cues from when
    /// the part has long rests.
    pub auto_cue_sources: Vec<StaffId>,
}
\end{lstlisting}

\subsection{Parts Are Projections, Not Storage}

\begin{requirement}
  \label{req:graph:part-content-projection}
  A part \MUSTNOT{} store musical content directly; it \MUST{} consist
  only of references to score content and overrides. Edits to score
  content \MUST{} propagate to all parts that reference that content.
\end{requirement}

\section{Analysis Layers and Views}
\label{sec:graph:layers}

\subsection{Analysis Layers}

Analysis layers, introduced for spelling in Section~\ref{sec:pitch:layers}
and used by analytical annotations, are first-class objects on the
score.

\begin{lstlisting}[language=Rust]
pub struct AnalysisLayer {
    pub id: AnalysisLayerId,
    pub name: String,
    pub description: Option<String>,

    /// Which views display this layer.
    pub visible_in_views: Vec<ViewId>,
}
\end{lstlisting}

\subsection{Views}

A view is a recipe for displaying a configuration of the score:
which analysis layers are active, which parts are rendered, what
overrides apply. Views are how the system supports condensed scores,
analytical views, study scores, and lead-sheet projections.

\begin{lstlisting}[language=Rust]
pub struct ViewDefinition {
    pub id: ViewId,
    pub name: String,
    pub kind: ViewKind,

    /// Which analysis layers are active in this view.
    pub active_layers: Vec<AnalysisLayerId>,

    /// View-specific overrides.
    pub overrides: ViewOverrides,
}

pub enum ViewKind {
    FullScore,
    Part(PartDefinitionId),
    CondensedScore(CondensationSpec),
    LeadSheet(LeadSheetSpec),
    Analysis(AnalysisViewSpec),
    Custom(CustomViewKindId),
}
\end{lstlisting}

\section{Graph Invariants}
\label{sec:graph:invariants}

The score graph maintains a set of structural invariants. Implementations
\MUST{} preserve these invariants across all edits.

\begin{requirement}
  \label{req:graph:score-graph-invariants}
  The following invariants hold over every well-formed score graph:

  \begin{enumerate}
    \item Every event in the arena has a valid \texttt{voice} field
      pointing to a voice that lists the event in its \texttt{events}
      vector.
    \item Every event in a voice's events list has its \texttt{voice}
      field pointing back to that voice.
    \item Events within a voice are sorted by position and do not
      overlap in time.
    \item Every event's \texttt{EventPosition} and
      \texttt{EventDuration} variants agree with the time model of
      its enclosing region (Section~\ref{sec:graph:event-time}). For
      aleatoric regions, agreement is governed by the region's
      declared \texttt{AleatoricAnchoringDiscipline}.
    \item Every voice belongs to exactly one \texttt{StaffInstance};
      every \texttt{StaffInstance} belongs to exactly one region;
      every region belongs to the canvas.
    \item Every \texttt{StaffInstance.staff} resolves to a
      \texttt{Staff} declared at the score level. A single
      \texttt{StaffId} \MAY{} be referenced by multiple
      \texttt{StaffInstance}s in different regions (this is how
      continuous staff identity is expressed across regions); a
      \texttt{StaffId} \MUSTNOT{} be referenced by two
      \texttt{StaffInstance}s in the same region.
    \item Region time and staff extents do not simultaneously overlap
      (Section~\ref{sec:graph:canvas}). A region's
      \texttt{staff\_extent} lists exactly the \texttt{StaffId}s
      manifested by the region's \texttt{staff\_instances}, with no
      duplicates.
    \item Measures belong to exactly one \texttt{StaffInstance}. A
      region's metric organization is per-staff (admitting
      polymeter); barline alignment across staves is recorded via
      \texttt{BarlineAlignmentGroup} entries when explicit alignment
      is authoritative.
    \item Every \texttt{TimeAnchor}'s \texttt{AnchorOffset} variant
      agrees with the time model of the target object's enclosing
      region (Section~\ref{sec:time:anchors}).
    \item Every cross-cutting structure's references resolve to extant
      objects in the graph, except where explicit re-anchoring rules
      permit transient dangling states during edits (see
      Chapter~\ref{ch:semops}).
    \item Every identifier in the graph is unique within its kind.
      Identifiers reserved for system-derived objects (notably
      \texttt{VoiceId}s of system-promoted voices) \MUST{} be derived
      from the deterministic function specified for their kind and
      \MUSTNOT{} collide with user-authored identifiers.
    \item Every \texttt{IdentifiedPitch} in the arena has a
      \texttt{PitchId} that is unique within the score's
      pitch-identity index.
    \item Every \texttt{SpellingScope::Pitch(PitchId)} resolves to a
      live or tombstoned \texttt{PitchId} in the pitch-identity
      index. A live target resolves to exactly one
      \texttt{IdentifiedPitch}; a tombstoned target is preserved for
      operation-log replay, undo, conflict reporting, and historical
      annotation, and \MUSTNOT{} be silently re-matched to a
      different live pitch.
    \item Every notational decomposition attachment's target
      \texttt{EventId} resolves to a live or tombstoned event;
      attachments to tombstoned events follow the same preservation
      discipline as for tombstoned pitches.
    \item For live decomposition attachments, the targeted event's
      duration matches the decomposition's component sum.
    \item Every tuplet's member event durations sum to the tuplet's
      structurally-required total
      (Section~\ref{sec:time:tuplets}). Operations that would break
      this invariant \MUST{} include compensating changes
      (Section~\ref{sec:semops:deleteevent}).
    \item Every \texttt{Tie}'s \texttt{pitch\_pairing} entries
      reference \texttt{PitchId}s that resolve to identified pitches
      belonging to the start and end events respectively. The
      class-specific adjacency, voice-membership, and
      pitch-equivalence rules of Section~\ref{sec:graph:ties} hold
      for the tie's declared \texttt{TieClass}.
    \item Every \texttt{Voice} declares an \texttt{origin} consistent
      with how it was created: \texttt{UserDeclared} for explicit
      user actions, \texttt{Imported} for foreign-format imports,
      \texttt{SystemPromoted} for voices allocated by concurrent-edit
      conflict resolution. The \texttt{VoiceId} of a
      \texttt{SystemPromoted} voice \MUST{} be the deterministic
      derivation specified in
      Section~\ref{sec:graph:promoted-voices}.
    \item Every \texttt{BarlineAlignmentGroup}'s members reference
      \texttt{StaffInstance}s within the same region and
      \texttt{Measure}s belonging to those instances. Group members
      \MUSTNOT{} reference staff instances or measures across
      region boundaries.
  \end{enumerate}

  Implementations \MUST{} reject graph configurations that violate any
  of the above invariants. Edits proposed in Chapter~\ref{ch:semops}
  must either preserve invariants directly or include compensating
  changes that restore them within the same operation.
\end{requirement}

This enumeration contains exactly \textbf{19} invariants. (Earlier
summary material, including the QUICKSTART, referred to ``18 graph
invariants''; the authoritative count is 19, matching this
enumeration and the reference implementation.) These 19 are
\emph{runtime} invariants: they hold over every well-formed graph and
are restored by compensating changes when an edit would break them.

Distinct from them are three \emph{construction-time} rejections
defined in Chapter~\ref{ch:time}, which a conforming constructor
\MUST{} enforce before a value ever enters the graph, and which no
runtime invariant restores because the offending value is never a
representable graph state:

\begin{enumerate}
  \item A \texttt{TimeSignature}'s beat groups \MUST{} sum to the
    measure duration (Section~\ref{sec:time:meter}).
  \item An \texttt{EventOrderingDAG} \MUST{} be acyclic
    (Requirement~\ref{req:time:ordering-dag-acyclic}).
  \item A \texttt{TupletRatio} \MUSTNOT{} be degenerate
    (Section~\ref{sec:time:tuplets},
    Requirement~\ref{req:time:tuplet-ratio-construction}).
\end{enumerate}

\section{Indexes and Auxiliary Structures}
\label{sec:graph:indexes}

Performance-critical operations require indexed access. The
specification states the required indexes; implementations \MAY{}
construct additional indexes provided they are kept consistent.

\begin{requirement}
  \label{req:graph:required-indexes}
  Implementations \MUST{} maintain at least the following indexes:

  \begin{description}
    \item[Event time index.] Maps (region, voice) to a sorted view
      of event positions, supporting $O(\log n)$ time-range queries.
    \item[Cross-cutting reference index.] Maps each object identifier
      to the cross-cutting structures referencing it, supporting
      $O(1)$ amortized lookup for re-anchoring.
    \item[Measure index.] Maps measure number (and measure id) to
      \texttt{MeasureId} for fast navigation.
    \item[Spelling attachment index.] Maps each \texttt{PitchId} to
      its attachments, partitioned by analysis layer.
  \end{description}

  Indexes \MUST{} be invalidated and rebuilt (or incrementally
  updated) on the corresponding edits. The semantic operations
  chapter specifies which operations invalidate which indexes.
\end{requirement}

\section{Forward References}

\begin{itemize}
  \item The complete catalog of semantic operations, their
    preconditions, postconditions, and re-anchoring rules is in
    Chapter~\ref{ch:semops}.
  \item The mapping from score graph to layout intermediate
    representation is defined in Chapter~\ref{ch:layout-ir}.
  \item The on-disk serialization of every type introduced in this
    chapter is defined in Chapter~\ref{ch:format}.
  \item Some engraving-specific types (\texttt{StemConfiguration},
    \texttt{ArticulationMark}, and similar) are introduced informally
    here and fully defined in Chapter~\ref{ch:layout-ir} where their
    interaction with the engraver is specified. The clef/key content
    model (\texttt{Clef}, \texttt{ClefShape}, \texttt{KeySignature},
    \texttt{ClefChange}, \texttt{KeySignatureChange}) and
    \texttt{LineStyle} are defined normatively \emph{in this chapter}
    as of schema major~2; the former \texttt{ClefId} identifier sketch
    is retired in favour of embedded \texttt{Clef} values.
\end{itemize}

% ===========================================================================
\chapter{Semantic Operations and Concurrent Reduction}
\label{ch:semops}

This chapter specifies the semantic operations on the score graph: the
mutations through which the data model becomes a live model. It defines
the operation framework, the canonical reduction procedure by which a
set of operations becomes a materialized score state, the object
existence model with tombstones, the conflict-record machinery, the
re-anchoring rules, transactions, and undo. The chapter ends with
representative operation specifications.

The collaboration mechanics specified here are the foundation on which
Chapter~\ref{ch:format}'s synchronization protocols build. The
materialized score is the deterministic reduction of an operation set;
the file format stores that operation set together with optional
acceleration snapshots.

\section{Design Principles}
\label{sec:semops:principles}

\begin{description}
  \item[The operation set is canonical.] A score's state is defined
    by the set of operations that have been committed to it. Any
    materialized graph is a deterministic reduction of that set;
    caches, snapshots, and partial reductions are acceleration
    structures, never the source of truth.
  \item[The replicated operation set is a CRDT; the materialized
    graph is not.] The grow-only set of operation envelopes with
    compact causal metadata is a true CRDT: replicas independently
    accumulate envelopes and converge on the same set. The
    materialized score graph is not itself claimed to be a CRDT; it
    is the deterministic reduction of that set, defined by the
    reduction algorithm specified here.
  \item[Pairwise operation merge is rejected.] Musical operations
    do not, in general, commute in intention: transpose-then-respell
    differs musically from respell-then-transpose; delete-then-tie
    differs from tie-then-delete. The framework therefore does not
    attempt to make every pair of operations algebraically commute.
    Instead, operations are reduced in a canonical order, and that
    order determines the outcome.
  \item[Conflicts, no-ops, and rejections are replicated facts.] No
    operation silently disappears because it could not apply. Every
    operation has a deterministic \texttt{OperationEffect} recorded
    as part of the materialized state; every replica reducing the
    same operation set records the same effects with the same
    reasons.
  \item[Identifier-based targeting.] Operations target objects by
    identifier, never by structural path. Identifiers are stable
    under concurrent edits; paths are not.
  \item[Tombstones preserve identity.] Deleted objects retain their
    identifiers as tombstones. References to tombstoned identifiers
    are not silently invalidated; they enter a tombstoned-target
    state unless the operation kind specifies deterministic repair.
  \item[Transactions materialize atomically.] Primitive members of a
    transaction \MUSTNOT{} be independently visible as committed
    score states. Either the whole transaction reduces successfully,
    or the transaction transitions to a conflict effect with a
    recorded conflict.
  \item[Undo is forward, not backward.] Undo is a new operation that
    computes a compensating edit against the current materialized
    state; it is not literal time travel. History remains
    append-only.
  \item[Two validation modes for advisory checking.] Operations are
    subject to strict precondition checking in \emph{authoring mode}
    (interactive edits) and lenient checking in \emph{replay mode}
    (loading historical operations or applying remote operations).
    Validation modes apply only to advisory preconditions;
    invariant preconditions are enforced unconditionally.
\end{description}

\section{The Operation Framework}
\label{sec:semops:framework}

\subsection{Operation Identity and Stamps}

Operations carry two related but distinct concepts: \emph{identity}
(who authored the operation, with what counter) and \emph{stamp}
(when the operation was committed, for ordering purposes).

\begin{lstlisting}[language=Rust]
/// Stable identity of an operation. Determined at authoring time.
pub struct OperationId {
    pub replica: ReplicaId,
    pub counter: u64,
}

/// Ordering metadata for an operation. Used to derive the canonical
/// reduction order; not part of the operation's identity.
pub struct OperationStamp {
    pub hlc: HybridLogicalClock,
    pub id: OperationId,
}

/// A hybrid logical clock combines wall-clock and logical components.
pub struct HybridLogicalClock {
    /// Physical time component, in canonical units (see Chapter 8).
    pub physical_time: WallClockTime,
    /// Logical counter, advanced when physical time does not.
    pub logical_counter: u32,
}
\end{lstlisting}

\begin{requirement}
  \label{req:semops:operation-stamp-stability}
  \texttt{OperationId} \MUST{} be stable: an operation's identity is
  fixed at the moment of authoring and \MUSTNOT{} change with
  reordering, retransmission, or merging into operation sets at
  other replicas.

  \texttt{OperationStamp} \MUST{} be set by the authoring replica
  using a hybrid logical clock that respects causal order: for any
  operation $A$ that is in operation $B$'s causal predecessor set,
  $A$'s stamp \MUST{} be strictly less than $B$'s under the
  canonical comparison (Section~\ref{sec:semops:reduction-order}).
\end{requirement}

\subsection{Causal Context via Dotted Version Vectors}

Operations carry a compact causal context, not an exhaustive
predecessor list. This is essential: exhaustive predecessor lists
scale linearly with history and become untenable.

\begin{lstlisting}[language=Rust]
/// A compact causal context: for each replica known to the authoring
/// replica, the highest contiguous counter observed, plus any
/// individual operation identifiers known but not contiguous (the
/// "dots").
pub struct CausalContext {
    /// Contiguous counter highest known per replica.
    pub vector: BTreeMap<ReplicaId, u64>,

    /// Individual operations known but not yet contiguous in the
    /// vector. Used when an authoring replica has observed an
    /// operation whose causal predecessors include operations not
    /// yet present.
    pub dots: BTreeSet<OperationId>,
}
\end{lstlisting}

\begin{requirement}
  \label{req:semops:dotted-version-vectors}
  Causal contexts \MUST{} be expressed as dotted version vectors as
  defined above. Exhaustive predecessor lists \MUSTNOT{} be used as
  the canonical causal representation; they may appear in debugging
  output or in the operation log's expanded form, but the wire and
  storage representation of causal context is the DVV.

  The per-replica counter floor is \textbf{zero-based} and
  normative: \texttt{vector[r] == n} asserts that \emph{every}
  operation \texttt{(r,\,0..=n)} is a causal predecessor --- the
  contiguous range begins at counter \texttt{0}, not \texttt{1}. A
  conforming implementation \MUSTNOT{} adopt a one-based floor: doing
  so would shift the contiguous/dot boundary and make two
  implementations disagree about which operations are pending versus
  ready. The membership check (``is operation \texttt{(r,\,c)}
  covered?'') is \texttt{c <= vector[r]} or \texttt{(r,\,c) in dots};
  an implementation \MAY{} evaluate it by walking the known operation
  ids rather than materializing the full \texttt{0..=n} range, so a
  sparse high-counter context does not force work proportional to the
  counter value.

  Interval tree clocks and other compact causal representations
  \MAY{} appear in future revisions as performance optimizations but
  are non-normative in this specification.
\end{requirement}

\begin{rationale}
  Dotted version vectors are well-understood, well-tested, and the
  right complexity tradeoff for the human-scale collaboration model
  Epiphany targets. Their normative use here closes the high-priority
  review item about exhaustive predecessor lists scaling badly with
  history.
\end{rationale}

\subsection{Operation Envelopes}

An operation is transmitted, stored, and reduced as an
\emph{envelope}: the operation's payload together with its identity,
ordering stamp, causal context, and optional transaction grouping.

\begin{lstlisting}[language=Rust]
pub struct OperationEnvelope {
    /// Stable identifier of this operation.
    pub id: OperationId,

    /// Author of this operation (may differ from replica in shared
    /// authoring sessions where multiple authors share a replica).
    pub author: AuthorId,

    /// Ordering stamp. Used for canonical reduction order; never
    /// for identity.
    pub stamp: OperationStamp,

    /// Compact causal context: what operations were visible to the
    /// authoring replica when this operation was created.
    pub causal_context: CausalContext,

    /// Optional transaction grouping. Multiple operations sharing a
    /// TransactionId are reduced atomically.
    pub transaction: Option<TransactionId>,

    /// The mutation itself.
    pub payload: OperationPayload,
}

pub enum OperationPayload {
    Primitive(OperationKind),

    /// A meta-operation that explicitly resolves a previously-recorded
    /// conflict (see Section~\ref{sec:semops:conflicts}).
    ResolveConflict(ResolveConflictPayload),

    /// A meta-operation that compensates for a previously-committed
    /// transaction; the realization of "undo" (see
    /// Section~\ref{sec:semops:undo}).
    UndoTransaction(UndoTransactionPayload),

    /// A meta-operation that resolves an equivocated operation slot
    /// by naming the chosen candidate envelope (see
    /// Section~\ref{sec:semops:equivocation}; payload schema in the
    /// Operation Catalog companion).
    ResolveEquivocation(ResolveEquivocationPayload),
}

/// The catalog of primitive operation kinds. The variants listed
/// here are normative for the core; additional kinds are introduced
/// by registered extensions via the Registered variant. The full
/// catalog including detailed payload schemas is the Operation
/// Catalog companion specification
/// (Appendix~\ref{app:deferred}).
pub enum OperationKind {
    // Event operations
    InsertEvent(InsertEventOp),
    DeleteEvent(DeleteEventOp),
    ModifyEvent(ModifyEventOp),

    // Pitch and tuning operations
    RespellPitch(RespellPitchOp),
    Transpose(TransposeOp),
    InsertIdentifiedPitch(InsertIdentifiedPitchOp),
    DeleteIdentifiedPitch(DeleteIdentifiedPitchOp),
    ModifyIdentifiedPitch(ModifyIdentifiedPitchOp),

    // Cross-cutting operations
    CreateCrossCutting(CreateCrossCuttingPayload),
    DeleteCrossCutting(DeleteCrossCuttingPayload),
    ModifyCrossCutting(ModifyCrossCuttingPayload),

    // Repeat-structure operations
    CreateRepeatStructure(CreateRepeatStructureOp),
    DeleteRepeatStructure(DeleteRepeatStructureOp),

    // Region and structure operations
    ChangeRegionTimeModel(ChangeRegionTimeModelOp),
    InsertRegion(InsertRegionOp),
    DeleteRegion(DeleteRegionOp),
    InsertStaffInstance(InsertStaffInstanceOp),
    DeleteStaffInstance(DeleteStaffInstanceOp),
    InsertStaff(InsertStaffOp),
    CreateVoice(CreateVoiceOp),
    DeleteVoice(DeleteVoiceOp),

    // Metric-model operations
    SetTimeSignature(SetTimeSignatureOp),
    SetTempoSegment(SetTempoSegmentOp),
    SetMetricGrid(SetMetricGridOp),

    // Score settings
    SetMetadata(SetMetadataOp),

    // Layout-semantic operations
    SetUserSystemBreak(SetUserSystemBreakOp),
    SetUserPageBreak(SetUserPageBreakOp),
    SetStaffLayout(SetStaffLayoutOp),

    // Transaction declaration. Carries metadata; member primitives
    // reference the transaction by id in their envelopes.
    DeclareTransaction(TransactionDescriptor),

    // Extension-defined primitive operation.
    Registered(OperationKindRegistryId, SerializedRegisteredOp),
}
\end{lstlisting}

\subsection{The Operation Set as CRDT}

\begin{requirement}
  \label{req:semops:grow-only-operation-set}
  The replicated operation set is a grow-only CRDT: replicas
  accumulate envelopes by gossip, broadcast, or any other
  delivery mechanism, and converge on the same set when they have
  observed the same envelopes. The set is set-valued, not
  multiset-valued: an envelope is a member of the set or it is
  not, with no count.

  A replica \MUSTNOT{} discard envelopes from the operation set
  except by an explicit pruning operation (Chapter~\ref{ch:format}),
  which preserves canonical state through retained base snapshots
  and the post-pruning operation set.
\end{requirement}

\subsection{OperationId Collisions and Replica Equivocation}
\label{sec:semops:equivocation}

Two received envelopes may carry the same \texttt{OperationId}.
Acceptance must not depend on arrival order: if it did, two
replicas observing the same envelopes in different orders could
diverge on which one is canonical. The operation set therefore
maps each \texttt{OperationId} to an \texttt{OperationSlot}
rather than directly to an envelope.

\begin{lstlisting}[language=Rust]
pub enum OperationSlot {
    /// Exactly one canonical envelope is known for this
    /// OperationId. The envelope reduces normally.
    Single(OperationEnvelope),

    /// Two or more distinct canonical envelopes have been
    /// observed for this OperationId. The slot is equivocated.
    /// The operation produces no canonical effect until external
    /// recovery resolves the equivocation.
    Equivocated {
        operation_id: OperationId,
        /// Every distinct canonical envelope observed for this
        /// OperationId. Identified by content hash to avoid
        /// duplicating large payloads. Sorted lexicographically by
        /// hash for deterministic enumeration.
        candidates: BTreeSet<EnvelopeHash>,
    },
}

/// BLAKE3-256 hash of a canonical envelope serialization with
/// domain tag "MUSCENVH".
pub struct EnvelopeHash(pub [u8; 32]);
\end{lstlisting}

\begin{requirement}
  \label{req:semops:operation-equivocation}
  Slot transitions are determined as follows, independently of
  arrival order:

  \begin{itemize}
    \item If no slot exists for the incoming envelope's
      \texttt{OperationId}: the incoming envelope is
      well-formedness checked
      (Section~\ref{sec:semops:acceptance}) and, if well-formed,
      a \texttt{Single} slot is created.
    \item If a \texttt{Single} slot exists and the incoming
      envelope is byte-identical (canonical serialization) to
      the slot's envelope: the duplicate is dropped silently;
      the slot is unchanged.
    \item If a \texttt{Single} slot exists and the incoming
      envelope's canonical bytes differ from the slot's
      envelope: the slot transitions to \texttt{Equivocated} with
      both envelopes' hashes in its \texttt{candidates} set. The
      previously-canonical envelope is retained in a diagnostic
      candidate store keyed by its hash; the incoming envelope is
      likewise retained. Neither envelope contributes to canonical
      reduction.
    \item If an \texttt{Equivocated} slot exists: the incoming
      envelope's hash is added to its \texttt{candidates} set if
      not already present; otherwise the duplicate is dropped.
      The envelope is retained in the candidate store. The slot
      remains \texttt{Equivocated}.
  \end{itemize}

  An equivocated slot \MUSTNOT{} contribute to canonical
  reduction. Reduction proceeds as if the operation does not
  exist: any operation that causally depends on the equivocated
  \texttt{OperationId} is reduced under the same rule as a
  missing causal predecessor (the dependent operation is
  retained but produces no effect until the equivocation is
  resolved).

  Resolution is by explicit recovery action, not by arrival
  order or content. Three resolution mechanisms are permitted:

  \begin{itemize}
    \item \emph{Transport-level credential revocation}: the
      collaboration transport revokes the equivocating replica's
      credentials and the user accepts one envelope timeline as
      canonical via an explicit reconciliation operation.
    \item \emph{Local user choice}: the user (or an authorized
      administrator) issues an explicit
      \texttt{ResolveEquivocation} operation naming the
      \texttt{OperationId} and the chosen
      \texttt{EnvelopeHash}, promoting the slot back to
      \texttt{Single}.
    \item \emph{Profile-declared deterministic policy}: a
      conformance profile \MAY{} declare a deterministic
      selection function (e.g., ``lowest canonical envelope
      hash wins''). This is a policy choice, not a default. A
      bundle reduced under a profile with a declared resolution
      policy is canonical under that profile; the same bundle
      reduced under a profile without such a policy remains
      anomalous.
  \end{itemize}

  If a bundle on disk contains an equivocated slot without a
  resolution policy or pending resolution operation, the bundle
  is anomalous. Readers \MUST{} open it, if at all, in a
  diagnostic recovery mode that surfaces the equivocation
  prominently and that prohibits ordinary editing of the
  affected portion of canonical state.
\end{requirement}

\begin{rationale}
  Replica equivocation (the same replica producing two distinct
  envelopes with the same \texttt{OperationId}) is a real
  failure mode: replica-key compromise, software bug, replay
  attack, or filesystem corruption can all produce it. The
  \texttt{OperationSlot} model makes equivocation a
  representable, order-independent state: two replicas that
  have observed the same set of envelopes converge on the same
  slot states regardless of the order of observation. Treating
  the first arrival as canonical would let a malicious
  early-arriver pin canonical state at one replica while a
  late-arriver pins different state at another replica; the
  slot model avoids that asymmetry by refusing to elevate
  either candidate without an explicit resolution action.

  Recovery is a policy choice rather than a built-in default
  because the right answer depends on context: a collaborating
  team with credential revocation will resolve differently
  from a single-user backup-restore scenario, which differs
  again from an adversarial-document forensic scenario.
\end{rationale}

\subsection{Envelope Acceptance and Malformed-Envelope Behavior}
\label{sec:semops:acceptance}

Operation envelopes arrive from local authoring and from remote
replicas. Acceptance rules determine when an envelope enters the
operation set.

\begin{requirement}
  \label{req:semops:envelope-well-formedness}
  An incoming envelope is well-formed if and only if all of the
  following hold:
  \begin{itemize}
    \item Its \texttt{OperationId} consists of a 64-bit
      \texttt{ReplicaId} that is not
      \texttt{ReplicaId::SYSTEM\_DERIVED} (operations are not
      authored by the system namespace) and a 64-bit counter.
    \item Its \texttt{stamp.id} is byte-identical to its
      top-level \texttt{id} field: an envelope's stamp \MUST{}
      address the same operation as the envelope itself. This
      eliminates an ambiguity where an authoring replica could
      otherwise stamp an envelope as belonging to a different
      operation than its identity declares.
    \item Its \texttt{stamp.hlc.physical\_time} is a finite,
      non-negative integer value, expressed in the canonical
      time unit (Chapter~\ref{ch:format}).
    \item Its causal context's DVV references only well-formed
      replica identifiers and non-negative integer counter
      values; dots are well-formed \texttt{OperationId}s.
    \item Its payload deserializes successfully against the
      profile's declared operation catalog.
  \end{itemize}

  An envelope that is not well-formed \MUST{} be rejected at
  reception. Rejection is recorded in the local diagnostic log
  but does \emph{not} enter the canonical operation set. Remote
  replicas \MAY{} retransmit; persistent rejection indicates a
  non-conforming peer.

  Acceptance \MUSTNOT{} require trust in the envelope's
  \texttt{stamp.hlc} content: a malicious or misconfigured peer
  may emit envelopes with implausible future physical times or
  zero/extreme logical counters. The canonical reduction order
  (Section~\ref{sec:semops:reduction-order}) consumes the stamp
  as ordering metadata; it does not validate plausibility.
  Implementations \MAY{} surface clock-skew warnings in
  diagnostics but \MUSTNOT{} reorder envelopes for plausibility
  reasons that would deviate from the canonical reduction order.
\end{requirement}

\subsection{Per-Replica HLC Monotonicity}
\label{sec:semops:hlc-monotonicity}

Per-replica monotonicity is a property of the \emph{set} of
accepted envelopes, not of arrival order. Out-of-order delivery
is normal in gossip systems and is not a violation.

\begin{lstlisting}[language=Rust]
/// A range of envelopes from a single replica that have been
/// identified as monotonicity-violating and excluded from
/// ordinary canonical reduction. Retained for diagnostic and
/// recovery purposes.
pub struct AnomalousReplicaSegment {
    pub replica: ReplicaId,
    /// The smallest counter at or after which the replica's
    /// stream is anomalous. All envelopes from this replica
    /// with counter >= first_bad_counter are excluded from
    /// canonical reduction.
    pub first_bad_counter: u64,
    /// The detected anomaly reason.
    pub reason: ReplicaAnomalyReason,
    /// Operation ids in the excluded segment, retained for
    /// diagnostic recovery. Sorted by counter.
    pub excluded: Vec<OperationId>,
}

pub enum ReplicaAnomalyReason {
    /// The replica produced an envelope at counter c2 with a
    /// stamp tuple strictly less than a previously-known
    /// envelope at counter c1 < c2.
    HlcMonotonicityViolation {
        violating_pair: (OperationId, OperationId),
    },
    /// A registered anomaly reason defined by an extension or
    /// transport.
    Registered(ReplicaAnomalyRegistryId),
}
\end{lstlisting}

\begin{requirement}
  \label{req:semops:per-replica-hlc-monotonicity}
  For any two envelopes accepted into the operation set that
  share the same authoring \texttt{ReplicaId}, with operation-
  counter values $c_1 < c_2$:

  \begin{quote}
    The stamp tuple
    \texttt{(stamp.hlc.physical\_time, stamp.hlc.logical\_counter,
    id.counter)} of the envelope with counter $c_1$ \MUST{} be
    lexicographically less than or equal to the corresponding
    tuple of the envelope with counter $c_2$.
  \end{quote}

  Arrival order \MUSTNOT{} be used to determine whether this
  property holds. If a replica receives counter 10 before
  counter 9, that is a property of delivery, not of monotonicity.

  If two envelopes from the same replica violate this property
  (their stamp tuples are inconsistent with their counter order),
  the document carries a \textbf{replica equivocation anomaly}.
  An \texttt{AnomalousReplicaSegment} is created (or extended)
  for the offending replica. The segment's
  \texttt{first\_bad\_counter} is the smaller of the two
  counters in the violating pair.

  Effects on canonical reduction:

  \begin{itemize}
    \item Every envelope from the offending replica with counter
      $\geq$ \texttt{first\_bad\_counter} is excluded from
      canonical reduction. Such envelopes are retained on disk
      and in memory (the CRDT property does not allow silent
      removal) but are held in the
      \texttt{AnomalousReplicaSegment}'s \texttt{excluded} list,
      not in the canonical operation set.
    \item Envelopes from the offending replica with counter
      $<$ \texttt{first\_bad\_counter} remain in the canonical
      operation set and reduce normally.
    \item Any operation in the canonical set that causally
      depends on an excluded envelope reduces under the
      missing-causal-predecessor rule: the dependent operation
      is retained, produces no canonical effect, and is held
      pending resolution.
    \item The document is marked as containing a replica-
      equivocation anomaly. Ordinary editing \MAY{} continue on
      the unaffected portion of canonical state; collaborative
      synchronization with the equivocating replica \MUST{} be
      suspended until external resolution.
  \end{itemize}

  Resolution paths mirror the \texttt{OperationSlot} equivocation
  rules (Section~\ref{sec:semops:equivocation}): transport-level
  credential revocation followed by an explicit reconciliation
  operation, an explicit local recovery action, or a
  profile-declared deterministic policy.

  Authentication of replica identities and revocation of
  compromised replicas are the responsibility of the
  collaboration transport (Appendix~\ref{app:deferred}). This
  chapter specifies only local detection, segment formation,
  exclusion from canonical reduction, and the structural shape
  of resolution.
\end{requirement}

\begin{rationale}
  Excluding the violating segment from canonical reduction
  rather than letting it reduce-but-mark-tainted is the safer
  choice: a tainted reduction still produces canonical-looking
  state derived from operations the spec has flagged as
  untrustworthy. Excluding the segment keeps canonical state
  derived only from operations that satisfy the monotonicity
  invariant, at the cost of holding pending operations that
  causally depended on the excluded segment.
\end{rationale}

\section{The Canonical Reduction}
\label{sec:semops:reduction}

The materialized score graph is the deterministic reduction of the
operation set. This section specifies how that reduction is performed.

\subsection{Operation Lifecycle}
\label{sec:semops:lifecycle}

Every operation traverses four phases:

\begin{description}
  \item[Prepare.] The local UI command validates enough to construct
    a well-formed operation envelope. \emph{Prepare is advisory
    validation only}: failure here means the UI rejects the user's
    action and never authors the operation. Prepare does not bind
    semantics; an operation that successfully prepares at one replica
    \MAY{} be reduced to a conflict effect at another.
  \item[Commit.] The operation envelope is added to the replicated
    operation set. After commit, the operation is part of canonical
    state and \MUSTNOT{} be discarded outside the pruning protocol.
  \item[Reduce.] During materialization, the operation is processed
    by the canonical reduction algorithm against the current
    materialized state. Reduction is deterministic: every replica
    with the same operation set produces the same materialized
    state.
  \item[Report.] The reduction produces a deterministic
    \texttt{OperationEffect} for each operation. Effects are visible
    graph facts: users can inspect what each operation did, and
    subsequent operations can address conflicts and no-ops directly.
\end{description}

\subsection{The Reduction Algorithm}

\begin{requirement}
  \label{req:semops:reduction-algorithm}
  The materialized score state is computed from the operation set as
  follows:

  \begin{enumerate}
    \item Sort the operation set into the canonical reduction order
      (Section~\ref{sec:semops:reduction-order}).
    \item Group consecutive operations that share a transaction id
      into transaction blocks.
    \item For each non-transaction operation, apply the operation to
      the working state, computing its \texttt{OperationEffect}.
    \item For each transaction block, apply all member operations
      atomically: either all succeed and the transaction reduces to
      \texttt{Applied} (with each member's individual effect
      recorded), or some member fails its invariant preconditions,
      in which case the whole transaction reduces to a
      \texttt{Conflicted} effect with a recorded conflict.
    \item Record every effect in the materialized state's effect log
      (which is part of canonical state, not implementation detail).
  \end{enumerate}

  Given the same operation set, conformant implementations \MUST{}
  produce score states with byte-identical structural content and
  byte-identical effect logs. Differences in implementation language,
  storage layout, threading model, or hardware \MUSTNOT{} produce
  divergent reductions.
\end{requirement}

\subsection{Canonical Reduction Order}
\label{sec:semops:reduction-order}

\begin{requirement}
  \label{req:semops:canonical-reduction-order}
  The canonical reduction order is the lexicographic ordering by:

  \begin{enumerate}
    \item Causal order: for operations $A$ and $B$, if $A$ is in
      $B$'s causal context (transitively, via DVV closure), then
      $A$ precedes $B$.
    \item For operations not ordered by causal order (i.e.,
      concurrent operations), lexicographic order by:
      \begin{enumerate}
        \item \texttt{stamp.hlc.physical\_time} (ascending)
        \item \texttt{stamp.hlc.logical\_counter} (ascending)
        \item \texttt{stamp.id.replica} (ascending, lexicographic on
          the replica identifier's canonical byte form)
        \item \texttt{stamp.id.counter} (ascending)
      \end{enumerate}
  \end{enumerate}

  Causal order strictly dominates HLC order: wall-clock skew may
  affect the ordering of concurrent operations but \MUSTNOT{} cause
  a causal predecessor to sort after a causal successor. The HLC
  authoring rule (Section~\ref{sec:semops:framework}) guarantees this
  by construction.
\end{requirement}

\subsection{Operation Effects}

\begin{lstlisting}[language=Rust]
/// The deterministic effect of an operation under canonical
/// reduction. Recorded as part of materialized state.
pub enum OperationEffect {
    /// The operation applied cleanly with no compensating changes.
    Applied,

    /// The operation applied with deterministic compensating changes
    /// (e.g., re-anchoring records, attachment migration).
    AppliedWithRepair { repairs: Vec<RepairRecord> },

    /// The operation could not apply cleanly; a conflict record was
    /// created in the conflict registry.
    Conflicted { conflict: ConflictId },

    /// The operation's target was tombstoned. The operation is
    /// preserved in the operation set but has no effect on the
    /// materialized graph beyond the recorded effect.
    TombstonedTarget { target: TypedObjectId },

    /// The operation reduces to no effect. The reason is recorded
    /// and is part of canonical state.
    NoOp { reason: NoOpReason },
}

pub enum NoOpReason {
    /// The operation's target was tombstoned by a causally-prior
    /// operation, and the target operation has no compensating
    /// repair rule.
    TargetTombstoned,

    /// The operation duplicates a causally-prior operation's effect.
    AlreadyApplied,

    /// A later operation in the canonical order subsumed this
    /// operation's effect (e.g., a later DeleteEvent on the same
    /// target supersedes an earlier RespellPitch).
    SupersededByLaterOperation { superseder: OperationId },

    /// An invariant precondition that was satisfied at authoring
    /// time fails when the operation is reduced under concurrent
    /// edits. Distinguished from Conflicted because the operation's
    /// intent is not preserved: it simply has no applicable effect.
    /// The failure reason is a typed enum, not a free-form string:
    /// canonical effects MUST NOT contain free-form text, since two
    /// implementations could otherwise produce divergent canonical
    /// state while agreeing semantically.
    PreconditionFailedUnderReduction {
        reason: PreconditionFailureReason,
    },

    /// The operation belonged to a transaction whose other members
    /// failed; the transaction was conflicted, and this operation
    /// produced no individual effect.
    TransactionConflict,
}

pub enum PreconditionFailureReason {
    /// The target object identifier did not exist in the working
    /// state at the moment of reduction.
    TargetMissing,

    /// The target object was tombstoned by a causally-prior
    /// operation. Distinguished from NoOpReason::TargetTombstoned
    /// in that the kind's rules permitted an attempt at
    /// repair-on-tombstone, but the repair itself failed precondition.
    TargetTombstoned,

    /// The operation requires a region time model (e.g., metric)
    /// different from the one currently in effect at the target
    /// position.
    WrongRegionTimeModel,

    /// A tuplet compensation declared by the operation is invalid
    /// against the current tuplet structure at the target.
    TupletCompensationInvalid,

    /// An event duration the operation specifies is invalid in
    /// the target voice or region (e.g., overlapping a tied
    /// chord-onset boundary in an incompatible way).
    EventDurationInvalid,

    /// The operation's target position falls outside the region
    /// declared by its envelope, or addresses a region that does
    /// not exist.
    PositionOutsideRegion,

    /// A pitch-space or tuning-context precondition failed: e.g.,
    /// the operation declared a pitch in a pitch space that is not
    /// active at the target region.
    PitchSpaceMismatch,

    /// A voice precondition failed: the operation targeted a
    /// voice that does not exist or has been tombstoned.
    VoiceMissing,

    /// An extension-declared precondition failed. The extension
    /// is identified; the extension's own catalog enumerates the
    /// specific precondition codes.
    ExtensionPrecondition(ExtensionPreconditionId),

    /// Registered precondition code defined by a versioned
    /// registry.
    Registered(PreconditionFailureRegistryId),
}

/// A deterministic compensating change made during reduction. Each
/// RepairRecord describes one repair the reduction performed to
/// preserve graph invariants in the presence of a tombstoning or
/// migrating operation. Recorded as part of canonical state.
pub struct RepairRecord {
    /// The kind of repair performed.
    pub kind: RepairKind,

    /// The object affected by this repair.
    pub target: TypedObjectId,
}

pub enum RepairKind {
    /// A reference was re-anchored to a surviving target per the
    /// re-anchoring rule table (Section~\ref{sec:semops:reanchor}).
    Reanchored {
        from: TypedObjectId,
        to: TypedObjectId,
        reason: ReanchorReason,
    },

    /// A spanner or beam had members removed but enough remained
    /// for the structure to survive.
    SpannerTruncated {
        removed_members: Vec<TypedObjectId>,
    },

    /// A user-content reference was orphaned (target lost; reference
    /// preserved for diagnostic purposes).
    Orphaned,

    /// A reference whose existence required its target was
    /// cascade-deleted.
    CascadeDeleted,

    /// An attachment transitioned to tombstoned-target state.
    AttachmentTombstoned,

    /// A voice was promoted to a system-derived voice identity due
    /// to concurrent insertion collision.
    VoicePromoted {
        from: VoiceId,
        to: VoiceId,
    },

    /// A tuplet was compensated according to the operation's
    /// declared TupletCompensation.
    TupletCompensated {
        compensation_kind: TupletCompensationKind,
    },

    /// A registered extension-defined repair kind.
    Registered(RepairKindRegistryId),
}
\end{lstlisting}

\begin{requirement}
  \label{req:semops:total-operation-effects}
  Every operation in the operation set \MUST{} produce exactly one
  \texttt{OperationEffect} under reduction. Effects \MUST{} be
  deterministic: every replica reducing the same operation set in
  the canonical order produces the same effect for each operation,
  including identical \texttt{NoOpReason} values where applicable.
\end{requirement}

\subsection{Cache and Snapshot Discipline}
\label{sec:semops:caches}

\begin{requirement}
  \label{req:semops:cache-transparency}
  Caches, snapshots, and partial reductions are acceleration
  structures only. They \MUSTNOT{} affect the canonical materialized
  state, which is defined solely by the operation set, the active
  conformance profile, and the normative reduction algorithm.

  Implementations \MUST{} invalidate or recompute caches whenever any
  of the following change:
  \begin{itemize}
    \item The visible operation set (additions, or pruning operations
      under the file-format protocol).
    \item The causal order derived from any envelope's causal
      context.
    \item The reduction algorithm version declared by the active
      conformance profile.
    \item The active conformance profile itself.
  \end{itemize}

  Cache hits \MUST{} return data that is bit-identical to the data a
  fresh reduction would produce. Caches that cannot guarantee this
  property \MUST{} be invalidated.
\end{requirement}

\section{Object Existence and Tombstones}
\label{sec:semops:tombstones}

\subsection{Add-Tombstone Semantics}

Object lifetimes in the materialized graph follow add-tombstone
semantics: an operation creates an object by minting a stable
identifier, and an operation deletes an object by tombstoning that
identifier. Tombstoned identifiers are retained; their resolution
yields a deletion record rather than the original object.

\begin{requirement}
  \label{req:semops:tombstone-state-machine}
  Identifiers \MUST{} have one of three states in the materialized
  graph:

  \begin{itemize}
    \item \textbf{Live.} The identifier resolves to a current object.
    \item \textbf{Tombstoned.} The identifier resolves to a deletion
      record carrying the operation that performed the deletion and
      the operation that originally minted the identifier.
    \item \textbf{Unknown.} The identifier has not been seen by this
      replica. This state is observable in collaborative scenarios
      where envelopes arrive out of order; it resolves to either Live
      or Tombstoned once causal predecessors arrive.
  \end{itemize}

  Implementations \MUSTNOT{} resurrect tombstoned identifiers
  (recreating an object with a previously-tombstoned id) except
  through an explicit inverse-of-delete operation in the operation
  set. Independent re-creation is forbidden.
\end{requirement}

\subsection{Tombstone Retention}

\begin{requirement}
  \label{req:semops:tombstone-retention}
  Tombstones \MUST{} be retained while any of the following hold:
  \begin{itemize}
    \item Pending replication: operations referencing the tombstoned
      identifier may not yet have been observed by all replicas.
    \item Pending undo: a transaction containing the deletion is
      within undo range.
    \item Active conflict resolution: a conflict record references
      the tombstoned identifier and has not been resolved.
    \item Attachment resolution: a tombstoned-target attachment
      remains in the score for historical or diagnostic purposes.
  \end{itemize}

  Tombstone garbage collection is governed by the pruning protocol
  in Chapter~\ref{ch:format}. Tombstones \MUSTNOT{} be discarded
  outside that protocol.
\end{requirement}

\subsection{Reference Resolution Across Tombstones}

\begin{requirement}
  \label{req:semops:tombstoned-reference-resolution}
  A reference to a tombstoned identifier resolves to a
  \texttt{TombstonedTarget} state rather than being silently
  invalidated. Operations targeting tombstoned identifiers reduce to
  \texttt{TombstonedTarget} or \texttt{NoOp} effects according to
  their kind's declared rules (Section~\ref{sec:semops:examples}).

  Re-anchoring rules (Section~\ref{sec:semops:reanchor}) define how
  cross-cutting structures and attachments respond to tombstoned
  references at reduction time.
\end{requirement}

\section{Conflict Records}
\label{sec:semops:conflicts}

When an operation cannot apply cleanly, the reduction produces a
\texttt{Conflicted} effect referencing a \texttt{ConflictRecord} in
the canonical conflict registry (a component of materialized state;
see below). Conflict records are first-class: stable, addressable,
visible to users, and addressable by subsequent operations.

\subsection{Conflict Record Type}

\begin{lstlisting}[language=Rust]
pub struct ConflictRecord {
    /// Content-derived identifier (Section below).
    pub id: ConflictId,

    /// The operations that participated in producing this conflict.
    /// At least two for a true conflict; one for an operation that
    /// failed precondition checking under reduction.
    pub caused_by: Vec<OperationId>,

    /// The kind of conflict, governing what resolution options
    /// exist.
    pub kind: ConflictKind,

    /// Objects affected by this conflict, for diagnostic and UI
    /// navigation purposes.
    pub affected_objects: Vec<TypedObjectId>,

    /// Current resolution state. Conflicts begin Unresolved; a
    /// later ResolveConflict operation transitions them to Resolved
    /// or Dismissed.
    pub resolution_state: ConflictResolutionState,
}

pub enum ConflictKind {
    /// Two concurrent operations attempted to write the same
    /// non-LWW field. The winning operation's effect is in the
    /// materialized state; the losing operation's effect is in the
    /// conflict record for user inspection.
    StructuralFieldCollision {
        winner: OperationId,
        loser: OperationId,
        field: FieldPath,
    },

    /// A transaction's primitive members were partially applicable
    /// under reduction; the whole transaction was rejected.
    TransactionConflict {
        transaction: TransactionId,
        failed_members: Vec<OperationId>,
    },

    /// An operation targeted a tombstoned object and could not be
    /// repaired by the kind's declared rules.
    TombstonedTarget {
        target: TypedObjectId,
        operation: OperationId,
    },

    /// Re-anchoring failed: no deterministic target could be found.
    ReanchorFailure {
        original_referent: TypedObjectId,
        referencing_object: TypedObjectId,
    },

    /// A change-region-time-model operation produced contained
    /// events whose coordinate kinds are incompatible with the new
    /// time model; the migration could not be performed
    /// deterministically.
    TimeModelMigrationFailure {
        region: RegionId,
        incompatible_events: Vec<EventId>,
    },

    /// A registered extension operation's reduction failed in a
    /// way described by the extension. Opaque to the core.
    ExtensionConflict {
        kind_id: ConflictKindRegistryId,
        details: SerializedExtensionConflict,
    },
}

pub enum ConflictResolutionState {
    Unresolved,
    Resolved { by: OperationId, action: ResolutionAction },
    Dismissed { by: OperationId },
}

// See Requirement (field-collision effect tag): for a
// StructuralFieldCollision the winner (later op) reads Conflicted;
// the earlier op retains Applied.

pub enum ResolutionAction {
    /// Accept the losing operation's effect (replacing the winner).
    AcceptLoser,
    /// Reapply the winner explicitly (no semantic change; clears the
    /// conflict).
    KeepWinner,
    /// User-authored replacement that supersedes both.
    Override { override_operation: OperationId },
    /// Re-anchor to a user-chosen target.
    Reanchor { new_target: TypedObjectId },
    /// Dismiss the conflict without changing the materialized graph:
    /// the user acknowledges it and accepts the current (winner)
    /// state. Selects the Dismissed resolution state.
    Dismiss,
    /// Custom resolution for registered conflict kinds.
    Registered(ResolutionRegistryId),
}
\end{lstlisting}

\begin{requirement}
  \label{req:semops:resolution-action-discriminants}
  The discriminant byte of \texttt{ResolutionAction} is its
  declaration-order index in the listing above: \texttt{AcceptLoser}
  $= 0$, \texttt{KeepWinner} $= 1$, \texttt{Override} $= 2$,
  \texttt{Reanchor} $= 3$, \texttt{Dismiss} $= 4$, \texttt{Registered}
  $= 5$. A \texttt{ResolutionAction} is encoded into the
  \texttt{ResolveConflict} operation payload and therefore into that
  operation's content hash, so this assignment is normative and stable
  and \MUSTNOT{} be reordered. \texttt{Dismiss} was assigned \texttt{4}
  ahead of \texttt{Registered} (\texttt{5}) when the \texttt{Dismissed}
  resolution state was made reachable by an authored operation
  (Section~\ref{sec:semops:conflict-resolution} below), fixing the wire
  form against a silent shift.
\end{requirement}

\begin{requirement}
  \label{req:semops:field-collision-effect}
  When two concurrent operations write the same non-LWW field and
  reduction records a \texttt{StructuralFieldCollision}, the
  per-operation \texttt{OperationEffect} is assigned as follows: the
  \emph{winner} --- the operation later in canonical order, whose
  choice materializes and whose processing created the conflict record
  --- reads \texttt{Conflicted\{conflict\}}; the \emph{earlier}
  (losing) operation retains the \texttt{Applied} effect it received
  when it was first reduced. The winner carries the flag because it is
  the single operation that both materializes the field and observes
  the collision, which makes the assignment independent of the order
  in which envelopes arrive at any replica.
\end{requirement}

\begin{rationale}
  The alternative --- tagging the loser \texttt{Conflicted} (or
  \texttt{NoOp\{SupersededByLaterOperation\}}) because its intent did
  not survive --- is defensible and arguably more intuitive for a UI
  that surfaces ``your edit was overridden.'' We pin
  \emph{winner-carries} instead because it is computed at the moment
  the collision is detected (when the winning operation is reduced
  against already-present state), needs no second pass to retroactively
  re-tag an earlier operation, and is manifestly order-independent. A
  UI that wishes to tell the losing author their edit was overridden
  reads the \texttt{StructuralFieldCollision} record's \texttt{loser}
  field rather than the loser operation's effect.
\end{rationale}

\subsection{The Conflict Registry}

Canonical materialized state carries a top-level conflict registry
(the score graph does not --- conflict records are a reduction
product, not authored content):

\begin{lstlisting}[language=Rust]
pub struct ConflictRegistry {
    pub records: Vec<ConflictRecord>,
}
\end{lstlisting}

\begin{requirement}
  \label{req:semops:conflict-registry-persistence}
  The conflict registry is part of canonical materialized state.
  Conflict records persist until they are resolved or dismissed by
  explicit \texttt{ResolveConflict} operations. Conflict records
  \MUSTNOT{} be silently removed; their lifetime is governed by
  the operation set, not by UI state.

  The layout pipeline (Chapter~\ref{ch:layout-ir}) \SHOULD{} provide
  visual indication of unresolved conflicts at their affected
  objects; this is the user-facing manifestation of the conflict
  registry's contents.
\end{requirement}

\subsection{ConflictId Derivation}
\label{sec:semops:conflict-id}

Conflict records are produced during deterministic reduction, not
authored by replicas. Their identifiers \MUST{} therefore be
content-derived, not minted from a local replica-plus-counter
sequence. Otherwise, two replicas reducing the same operation set
would produce conflict registries with disagreeing identifiers,
breaking canonical state.

\begin{requirement}
  \label{req:semops:conflict-id-derivation}
  Each \texttt{ConflictId} \MUST{} be derived deterministically from
  the conflict's content via BLAKE3 truncation:

  \begin{lstlisting}[language=Rust]
fn derive_conflict_id(
    kind: &ConflictKind,
    causing_operations: &[OperationId],
    affected_objects: &[TypedObjectId],
) -> ConflictId {
    let mut preimage = Vec::new();
    preimage.extend_from_slice(b"MUSCCONF");
    preimage.extend_from_slice(&kind.canonical_bytes());

    // Sort causing operations lexicographically by canonical bytes
    let mut sorted_ops = causing_operations.to_vec();
    sorted_ops.sort_by(|a, b| a.canonical_bytes().cmp(&b.canonical_bytes()));
    for op in &sorted_ops {
        preimage.extend_from_slice(&op.canonical_bytes());
    }

    // Sort affected objects lexicographically by canonical bytes
    let mut sorted_objs = affected_objects.to_vec();
    sorted_objs.sort_by(|a, b| a.canonical_bytes().cmp(&b.canonical_bytes()));
    for obj in &sorted_objs {
        preimage.extend_from_slice(&obj.canonical_bytes());
    }

    let hash = blake3(&preimage);
    ConflictId(u128::from_be_bytes(
        hash[0..16].try_into().unwrap()
    ))
}
  \end{lstlisting}

  The canonical \texttt{ConflictKind} byte representation \MUST{}
  include the kind discriminant and all kind payload (winning
  operation, losing operation, field path, transaction id,
  tombstoned target, original referent, region, extension kind
  identifier, etc.), so that distinct conflicts have distinct
  preimages by construction. Implementations \MUSTNOT{} introduce
  an ordinal or local counter into the preimage; conflicts that
  share kind, causing operations, and affected objects are by
  definition the same conflict.
\end{requirement}

\begin{rationale}
  Content-derived identifiers are the only mechanism that makes
  the conflict registry a deterministic materialized fact. Every
  replica reducing the same operation set produces the same
  conflict records with the same identifiers, addressable by
  subsequent \texttt{ResolveConflict} operations whose payloads
  carry those identifiers by reference.

  Truncation from 256 to 128 bits is safe: the birthday-collision
  probability for any practical conflict count is vanishingly
  small, and the same truncation discipline already governs other
  128-bit identifiers in the format.
\end{rationale}

\subsection{Conflict Resolution Operations}
\label{sec:semops:conflict-resolution}

\begin{lstlisting}[language=Rust]
pub struct ResolveConflictPayload {
    pub target: ConflictId,
    pub action: ResolutionAction,
}
\end{lstlisting}

\begin{requirement}
  \label{req:semops:conflict-resolution-state}
  A \texttt{ResolveConflict} operation transitions a conflict's
  state from \texttt{Unresolved} to one of two states selected by its
  \texttt{action}. The \texttt{ResolutionAction::Dismiss} action
  transitions the conflict to \texttt{Dismissed} (no semantic change
  to the materialized graph beyond the conflict record's state: the
  user acknowledges the conflict and accepts the current winner
  state). Every other \texttt{action} transitions the conflict to
  \texttt{Resolved} with that action applied. Both resolution states
  are thus reachable by an authored operation; \texttt{Dismiss} is the
  authored selector for \texttt{Dismissed}, closing the prior gap in
  which \texttt{Dismissed} was a representable state with no operation
  to reach it.

  Concurrent \texttt{ResolveConflict} operations on the same
  conflict are themselves subject to canonical reduction: the
  earlier operation under the canonical order is applied; the
  later operation reduces to \texttt{NoOp} with reason
  \texttt{AlreadyApplied} if its action is equivalent, or to
  \texttt{Conflicted} with a meta-conflict record if its action
  differs.
\end{requirement}

\section{Transactions}
\label{sec:semops:transactions}

Transactions are first-class replicated operation groups. They reduce
atomically: either every primitive member applies, or the whole
transaction conflicts.

\begin{lstlisting}[language=Rust]
pub struct TransactionDescriptor {
    pub id: TransactionId,

    /// Human-readable label, for the undo history and diagnostics.
    pub label: String,

    /// Optional categorization (used by UIs and analytics).
    pub category: Option<TransactionCategory>,
}

pub struct TransactionId(pub u128);

/// Open vocabulary: a minimal core set plus a Registered escape for
/// extension-defined categories.
pub enum TransactionCategory {
    NoteEntry,
    Structural,
    Layout,
    Import,
    Registered(OperationKindRegistryId),
}
\end{lstlisting}

\begin{requirement}
  \label{req:semops:transaction-category}
  \texttt{TransactionCategory} is an \emph{open} vocabulary. Its
  normative core set is exactly \texttt{NoteEntry}, \texttt{Structural},
  \texttt{Layout}, and \texttt{Import}; any other category \MUST{} be
  expressed as \texttt{Registered(OperationKindRegistryId)} rather than
  by extending the core set. The field is advisory (consumed by UIs and
  analytics, never by canonical reduction), so the minimal core set is
  the right tradeoff: under-specifying is cheap to extend through
  \texttt{Registered}, whereas baking speculative categories into the
  core vocabulary is not. The discriminants follow declaration order
  (\texttt{NoteEntry} $= 0$ through \texttt{Registered} $= 4$).
\end{requirement}

Each primitive operation belonging to a transaction carries the
transaction's identifier in its envelope's \texttt{transaction} field.
The transaction descriptor itself is a separate operation envelope
(\texttt{DeclareTransaction}) whose payload contains the
\texttt{TransactionDescriptor}.

\subsection{Transaction Descriptor Causal Ordering}
\label{sec:semops:transaction-causal}

The canonical reduction order is causal-first, then HLC, then
operation id (Section~\ref{sec:semops:reduction-order}). For the
descriptor to be available when its members are reduced, the
ordering between descriptor and members must be guaranteed by
causal dependency, not by HLC stamps.

\begin{requirement}
  \label{req:semops:transaction-descriptor-causality}
  Every primitive operation member of a transaction \MUST{} causally
  depend on the transaction's \texttt{DeclareTransaction} envelope:
  the descriptor's \texttt{OperationId} \MUST{} appear in the
  member's \texttt{causal\_context} (either in the per-replica DVV
  range or among its dots). Authoring implementations \MUST{}
  observe and include the descriptor in the causal context of each
  member they author.

  Delivery order need not respect this dependency: an out-of-order
  delivery in which a member arrives before its descriptor is
  normal and is handled by the canonical reduction algorithm,
  which sorts by causal order before reducing.

  At reduction time:

  \begin{itemize}
    \item If the descriptor is present in the operation set and
      causally precedes every member: the transaction is
      well-formed; it is reduced atomically per the rule below.
    \item If a primitive operation declares membership in a
      transaction whose \texttt{DeclareTransaction} envelope is
      absent from the operation set, or whose envelope is present
      but does not causally precede the member: the member is
      malformed against the transaction model. It reduces to
      \texttt{NoOp\{reason: TransactionConflict\}} and contributes
      to a synthesized conflict record of kind
      \texttt{ConflictKind::TransactionConflict} naming the
      malformed member and the missing-or-misordered descriptor.
      The member is preserved in the operation set (the CRDT
      property prohibits removal) but produces no effect on
      canonical state beyond the recorded conflict.
  \end{itemize}
\end{requirement}

\begin{rationale}
  Requiring causal dependency rather than HLC ordering matches
  how transactions actually arise at the authoring replica: a
  user begins an edit (descriptor), then commits its primitives.
  The primitives can causally observe the descriptor because they
  are authored after it on the same replica. The canonical
  reduction order then handles them correctly without any special
  transaction-extraction phase.

  This rule also makes equivocating descriptors impossible:
  a member cannot claim membership in a transaction whose
  descriptor it has not seen.
\end{rationale}

\begin{requirement}
  \label{req:semops:atomic-transaction-reduction}
  Transactions reduce atomically. During reduction:

  \begin{itemize}
    \item All primitive members of a transaction are gathered into
      a transaction block at the point of the earliest member's
      canonical position.
    \item The block is reduced as a single unit. Each member is
      evaluated against the working state produced by prior members
      within the same block.
    \item If every member's invariant preconditions are satisfied,
      the transaction reduces to \texttt{Applied}; each member's
      individual effect is recorded as part of the transaction's
      composite effect.
    \item If any member's invariant preconditions fail, the entire
      transaction reduces to \texttt{Conflicted}: a conflict record
      is created with kind
      \texttt{ConflictKind::TransactionConflict}, listing the
      failed members; no member's effect is applied; each member's
      individual effect is recorded as
      \texttt{NoOp\{reason: TransactionConflict\}}.
  \end{itemize}

  Primitive members of a transaction \MUSTNOT{} be observable as
  independently committed score states. Observers see either the
  pre-transaction state or the fully applied post-transaction
  state; they never see a partial intermediate.
\end{requirement}

\begin{rationale}
  Atomic transaction reduction is essential for the tuplet
  compensation model (Section~\ref{sec:semops:deleteevent}): the
  delete operation and its compensating insert (or rewrite, or
  cascade) are committed together as a single transaction, so the
  score is never observably broken between them.

  Local materialization may stream a transaction's primitives to the
  UI for responsive feedback, but the materialized canonical state
  visible to remote observers transitions only at transaction
  boundaries.
\end{rationale}

\section{Re-Anchoring}
\label{sec:semops:reanchor}

When an operation tombstones an object, every other object
referencing it must respond deterministically. The response is
specified as a total deterministic function.

\subsection{The Re-Anchor Function}

\begin{lstlisting}[language=Rust]
pub enum ReanchorResult {
    /// The reference is replaced with a reference to a new target.
    Reanchored {
        new_target: TypedObjectId,
        reason: ReanchorReason,
    },

    /// The reference is preserved but its target is tombstoned.
    /// The referencing object remains in the graph in a
    /// tombstoned-target state.
    TombstonedTarget,

    /// The referencing object is marked orphaned but retained.
    Orphaned,

    /// Re-anchoring cannot be performed deterministically; a
    /// conflict is recorded.
    Conflicted { conflict: ConflictId },

    /// The referencing object is cascade-deleted (tombstoned).
    CascadeDeleted,
}

pub enum ReanchorReason {
    SameVoiceNearer,
    SameStaffInstanceNearer,
    SameStaffNearer,
    SameRegionNearer,
    ExplicitFallback,
    DeclaredByExtension(ReanchorReasonRegistryId),
    /// A rank-4 (same-canvas) proximity survivor. Appended in
    /// Pass 12 (P12-C4); its wire discriminant is 6 because
    /// DeclaredByExtension already owned 5 when it was appended.
    SameCanvasNearer,
}
\end{lstlisting}

A re-anchor that selects a rank-4 (same-canvas) survivor records
\texttt{SameCanvasNearer} (ratified Pass~12); before the variant was
appended, such repairs were recorded \texttt{ExplicitFallback}, a
recording this revision supersedes.

\subsection{Total Ordering for "Nearest"}

The notion of ``nearest surviving anchor'' is defined as a total
ordering. There is no discretionary search and no implementation
freedom to choose among ``equally near'' candidates.

\begin{requirement}
  \label{req:semops:nearest-reanchor-order}
  For an object kind $K$ requiring re-anchoring to the nearest
  surviving $K$, ``nearest'' is computed as the strict lexicographic
  minimum over the surviving candidates of:

  \begin{enumerate}
    \item Containment proximity (smaller is closer): same voice
      $(0)$, same staff instance $(1)$, same staff $(2)$, same
      region $(3)$, same canvas $(4)$. Candidates further away than
      the kind's declared maximum proximity are excluded.
    \item Absolute time distance from the tombstoned referent's
      resolved position, measured under a canonical conversion
      (musical time within metric regions, wall-clock time within
      proportional regions, declared per-region for aleatoric
      regions).
    \item Direction preference: forward $(0)$ before backward $(1)$,
      unless the kind declares the opposite preference.
    \item Object identifier (lexicographic on the typed identifier's
      canonical byte form, ascending) as the final tie-break.
  \end{enumerate}

  If no candidate satisfies the kind's declared proximity bound, the
  result is \texttt{Orphaned} (for user-content kinds such as
  comments and analytical annotations) or \texttt{CascadeDeleted}
  (for kinds whose existence requires their target, such as ties
  and beams).

  Numeric proximity bounds, direction preferences, and the
  orphaned-versus-cascade choice are declared per kind in the
  re-anchoring rule table below; they are part of the kind's
  conformance contract and \MUSTNOT{} be reinterpreted by
  implementations.
\end{requirement}

\subsection{The Re-Anchoring Rule Table}

The following table specifies, for each (referencing kind, referent
kind) pair, the re-anchoring action under canonical reduction. The
table is normative.

\begin{longtable}{p{3.0cm} p{2.6cm} p{2.4cm} p{4.4cm}}
  \toprule
  \textbf{Referencing} & \textbf{Referent} & \textbf{Action} &
  \textbf{Parameters} \\
  \midrule
  \endhead

  Slur & Endpoint event &
  Re-anchor to the nearest surviving endpoint; cascade-delete only
  when no endpoint survives &
  a two-endpoint slur collapses onto its sole survivor (a
  reference-clean degenerate form); proximity-aware re-targeting
  is a deferred refinement \\

  Slur & Interior event &
  No action &
  --- \\

  Tie & Either endpoint &
  Cascade-delete &
  Tie's existence requires both endpoints. \\

  Beam & Member event &
  Truncate. Cascade-delete if fewer than two members remain &
  proximity max: same voice \\

  Tuplet & Member event &
  Operation \MUST{} provide tuplet compensation
  (Section~\ref{sec:semops:deleteevent}); otherwise the operation
  conflicts at precondition check &
  --- \\

  Spanner & Anchor &
  Re-anchor to the nearest surviving anchor; cascade-delete only
  when no anchor survives &
  as for slurs (surviving-endpoint collapse); per-spanner-kind
  proximity bounds are a deferred refinement \\

  Repeat structure & Anchor &
  Re-anchor to the nearest surviving anchor; cascade-delete only
  when no anchor survives &
  as for spanners; the rule spans \emph{every} event-referencing
  anchor site of the structure --- \texttt{start}/\texttt{end},
  the kind's jump targets (\texttt{DaCapo.end\_target},
  \texttt{DalSegno.segno}/\texttt{end\_target}), and each
  volta's span. Among multiple surviving candidates ``nearest''
  is currently the deterministic identifier-order minimum;
  proximity-aware (four-key) selection is a deferred refinement,
  as for spanners (ratified with the repeat-authoring pair,
  schema-major-2 revision) \\

  Marker & Anchor &
  Re-anchor to nearest event in same staff instance &
  proximity max: same staff instance; orphan on failure \\

  Comment & Anchor &
  Orphan &
  User content never silently deleted. \\

  Cue event & Source event &
  Cascade-delete on any source deletion &
  a multi-source cue cascades when \emph{any} source dies
  (ratified Pass~12): losing any source breaks the cue's
  quotation integrity; truncate-while-any-source-survives was
  considered and rejected \\

  Graphic gesture & Anchor event &
  Re-anchor to nearest surviving event of same staff instance;
  for Free anchoring, no action; for Range anchoring, truncate &
  proximity max: same staff instance. Truncate (ratified
  Pass~12): a dead event-anchored range endpoint moves to its
  containing region's edge --- a start endpoint to the region
  start, an end endpoint to the region end, zero offset \\

  Trajectory event & Endpoint pitch &
  Re-anchor: replace \texttt{EventPitch(PitchId)} with
  \texttt{ExplicitPitch} capturing the tombstoned pitch's last
  known value &
  --- \\

  Analytical annotation & Anchor &
  Re-anchor to time range preserving original extent; orphan if
  range cannot be reconstructed &
  orphaning is the sanctioned outcome (ratified Pass~12) for a
  wall-clock (region-relative) or indeterminate event span that
  no stored \texttt{Range} anchor form can express without
  region-origin resolution \\

  Spelling attachment & Target pitch &
  Transition to tombstoned-target state &
  Attachment is preserved for diagnostic and undo purposes. \\

  Decomposition attachment & Target event &
  Transition to tombstoned-target state &
  Attachment preserved for diagnostic and undo purposes. \\

  Voice & Containing staff instance &
  Cascade-delete &
  --- \\

  Event & Containing voice &
  Cascade-delete &
  --- \\

  Region & Canvas &
  Refuse (canvas is not deletable) &
  --- \\

  Part definition & Referenced staff &
  Remove staff from part's staff list; orphan if part becomes
  empty &
  --- \\

  \bottomrule
\end{longtable}

\begin{requirement}
  \label{req:semops:synchronous-reanchoring}
  Re-anchoring \MUST{} be performed as part of the same reduction
  step that tombstones the referent. The resulting graph state
  \MUST{} satisfy every invariant in
  Section~\ref{sec:graph:invariants}. Re-anchoring actions \MUST{}
  be recorded as \texttt{RepairRecord} entries in the triggering
  operation's effect.
\end{requirement}

\section{Undo}
\label{sec:semops:undo}

Undo is not literal time travel: it is the authoring of a new
operation that compensates for a previously-committed transaction
against the current materialized state. This preserves the append-
only operation set and is correct under concurrent edits.

\subsection{UndoTransaction Payload}

\begin{lstlisting}[language=Rust]
pub struct UndoTransactionPayload {
    /// The transaction being undone.
    pub target: TransactionId,

    /// Policy governing how to undo when the target's effects
    /// have been partially superseded by later operations.
    pub policy: UndoPolicy,
}

pub enum UndoPolicy {
    /// Compute the inverse semantic operations against the current
    /// materialized state. If any cannot apply cleanly (because the
    /// target objects have been tombstoned, fields have been
    /// further modified, etc.), the entire undo conflicts.
    StrictInverse,

    /// Undo what remains valid; record conflicts or no-ops for the
    /// rest. The undo operation succeeds, but partial.
    BestEffort,

    /// Also undo operations that are causally or semantically
    /// dependent on the target transaction. Dependency is computed
    /// from causal contexts and explicit dependency links;
    /// implementations MUST NOT infer broad musical dependence
    /// heuristically.
    Cascade,
}
\end{lstlisting}

\subsection{Undo Semantics}

\begin{requirement}
  \label{req:semops:compensating-undo}
  An \texttt{UndoTransaction} operation is committed to the operation
  set like any other operation; it does not modify history. Its
  reduction computes a compensating edit against the materialized
  state at the point of its canonical position:

  \begin{itemize}
    \item For \texttt{StrictInverse}: the inverse of each primitive
      in the target transaction is computed against the current
      working state. If every inverse applies cleanly, the undo
      transaction reduces to \texttt{Applied} with all compensating
      effects recorded. If any inverse cannot apply (target
      tombstoned, field modified, etc.), the undo transaction
      reduces to \texttt{Conflicted} with a conflict record.
    \item For \texttt{BestEffort}: inverses that apply cleanly are
      applied; inverses that cannot are recorded as \texttt{NoOp}
      effects with their reasons. The undo transaction reduces to
      \texttt{AppliedWithRepair}.
    \item For \texttt{Cascade}: the dependency set is computed first
      (causal closure plus explicit dependency links). The
      transaction's inverses, plus inverses of dependents in
      reverse order, are applied per the chosen sub-policy
      (typically \texttt{StrictInverse}).
  \end{itemize}

  Equality of the resulting graph state to the pre-target state is
  defined as \emph{content equivalence}: every live object's fields,
  every attachment, every cross-cutting structure matches. Operation
  identifiers, stamps, undo stacks, caches, and storage layout are
  \emph{not} required to match.
\end{requirement}

\begin{rationale}
  Defining undo as a forward operation eliminates the contradiction
  in earlier passes where undo had to ``restore previous state''
  while remote edits had been integrated. The compensating-edit
  approach respects the operation set's append-only nature and
  generalizes cleanly to concurrent collaboration: undoing a
  transaction whose effects have been entwined with remote edits is
  expressible as a deterministic compensating operation, possibly
  partial.
\end{rationale}

\section{LWW Discipline}
\label{sec:semops:lww}

Last-writer-wins reduction is permitted only for explicitly-marked
advisory scalar fields. Structural musical objects are governed by
their kind's operation-specific reduction rules.

\subsection{The LwwAdvisory Marker}

Fields eligible for LWW reduction are explicitly marked as
\texttt{LwwAdvisory}. No unmarked field may use LWW semantics.

\begin{lstlisting}[language=Rust]
/// Marker attribute, applied to fields whose concurrent modifications
/// are resolved by last-writer-wins under the canonical reduction
/// order.
#[attribute_only]
pub struct LwwAdvisory;

// Example field declarations:
//
//   #[lww_advisory]
//   pub title: Option<String>
//
//   #[lww_advisory]
//   pub composer: Option<String>
\end{lstlisting}

\subsection{Eligible Fields}

\begin{requirement}
  \label{req:semops:lww-advisory-fields}
  The following classes of field \MAY{} be marked
  \texttt{LwwAdvisory} and reduced by last-writer-wins:

  \begin{itemize}
    \item Score metadata: title, subtitle, composer, lyricist,
      copyright, work number, opus, dedication, free-form
      annotations.
    \item View preferences: which layers are active, which parts
      are shown, viewport, zoom level (per-replica state in
      addition; only the explicitly shared portions are LWW).
    \item Layout hint preferences: default page size, default
      margins, default staff spacing recommendations.
    \item Page-break and system-break preferences on layout-semantic
      structures (the user's stated preference, distinct from the
      engraver's decision).
    \item Optional advisory display hints on cross-cutting
      structures: line color preference, line-style preference,
      label-position preference. Endpoints, target identities,
      semantic class, and attachment ownership \MUSTNOT{} be
      LWW.
  \end{itemize}

  All other fields, including every field carrying musical
  structure, identity, ownership, or semantics, \MUST{} use
  operation-specific reduction rules defined per operation kind.
  Implementations \MUSTNOT{} default to LWW for unmarked fields.
\end{requirement}

\begin{rationale}
  LWW is acceptable for fields where the worst case of arbitrary
  deterministic resolution is benign: two authors editing the
  composer's name concurrently is a metadata collision, not a
  musical loss. LWW is catastrophic for fields where wholesale
  replacement destroys musical meaning: two authors editing a
  chord's pitch list concurrently must not produce ``one author's
  chord wins.'' The marker discipline makes the boundary explicit
  and auditable.
\end{rationale}

\section{Validation Modes}
\label{sec:semops:validation}

Operations are subject to precondition checking. Two modes
distinguish authoring contexts:

\begin{description}
  \item[Authoring mode.] Interactive edits. All preconditions are
    enforced: invariant preconditions (which preserve graph
    invariants) and advisory preconditions (which encode user-intent
    constraints, range checks, and style policy).
  \item[Replay mode.] Replaying historical operations or applying
    remote operations under reduction. Only invariant preconditions
    are enforced; advisory preconditions are skipped, since they
    represent the authoring replica's local policy at the moment of
    authoring, not invariants of the canonical state.
\end{description}

\begin{requirement}
  \label{req:semops:precondition-validation-modes}
  Every operation specification \MUST{} classify each precondition
  as invariant or advisory. Invariant preconditions \MUST{} hold in
  all modes; advisory preconditions \MUST{} hold in authoring mode
  and \MAY{} fail silently in replay mode.

  Implementations \MUSTNOT{} reclassify a precondition without a
  corresponding revision to this specification.
\end{requirement}

\section{Per-Kind Reduction Summary}
\label{sec:semops:per-kind-summary}

The following table summarizes the reduction discipline for each
representative operation kind. Detailed specifications appear in
the next section.

\begin{longtable}{p{3.4cm} p{4.0cm} p{5.0cm}}
  \toprule
  \textbf{Operation} & \textbf{Reduction Rule} &
  \textbf{Concurrency Outcomes} \\
  \midrule
  \endhead

  InsertEvent &
  Position-keyed insertion. Concurrent same-position inserts whose
  durations overlap promote the lexicographically-greater operation
  to a system-promoted voice
  (Section~\ref{sec:graph:promoted-voices}) &
  Voice promotion; non-overlap invariant preserved by allocation. \\

  DeleteEvent &
  Tombstone the event and contained pitch identifiers. Apply
  tuplet compensation from the operation. Re-anchor referencing
  structures &
  Delete-wins against concurrent field modifications on the same
  event; concurrent same-target deletes are idempotent. \\

  RespellPitch &
  Field overwrite by canonical reduction order; the later operation
  wins. Earlier respellings of the same pitch are preserved in the
  conflict record if they differ &
  Conflict record records the losing spelling for user inspection. \\

  Transpose &
  Apply in canonical reduction order. Pitch identifiers preserved.
  Operations do not commute &
  Reduction order decides; conflict record entries for non-trivial
  intervals only when explicit dependency rules require. \\

  CreateCrossCutting &
  Set union. Concurrent additions of distinct cross-cutting
  structures all succeed &
  Tombstoned endpoints cause re-anchoring or cascade-delete per
  kind. \\

  ChangeRegionTimeModel &
  Structural migration. Contained events whose coordinate kinds
  become incompatible cause the operation to conflict &
  Concurrent contained edits authored before the migration may
  conflict; explicit rebase rules per migration kind. \\

  SetUserSystemBreak &
  LwwAdvisory on per-location break preference &
  Last-writer-wins on the break state at the same anchor. \\

  \bottomrule
\end{longtable}

\section{Representative Operations}
\label{sec:semops:examples}

The remainder of the chapter specifies a representative selection of
operations, one from each major category. The full catalog is
delivered as a separate conformance specification.

These examples illustrate the contract that every operation in the
catalog \MUST{} satisfy. Each operation specification includes its
reduction rule under the canonical reduction model
(Section~\ref{sec:semops:reduction}).

\subsection{InsertEvent}
\label{sec:semops:insertevent}

\textbf{Category:} Event operations.

\begin{lstlisting}[language=Rust]
pub struct InsertEventOp {
    /// The voice the event is inserted into.
    pub voice: VoiceId,

    /// Position within the voice.
    pub position: EventPosition,

    /// The event to insert. Its EventId and any contained PitchIds
    /// are minted at operation-authoring time as a single causal
    /// step. The event's voice field must match the target voice.
    pub event: Event,
}
\end{lstlisting}

\textbf{Invariant preconditions:}
\begin{itemize}
  \item The target voice exists.
  \item The event's \texttt{voice} field equals the target voice's id.
  \item The event's \texttt{id} is not present elsewhere in the arena.
  \item Every \texttt{PitchId} contained in the event (in an
    \texttt{IdentifiedPitch}, a trajectory endpoint, or a stepwise
    trajectory shape) is unique within the score's pitch-identity
    index.
  \item The position is within the enclosing region's time extent.
  \item The event's position and duration variants agree with the
    enclosing region's time model.
  \item Inserting at the position does not produce an event overlap
    within the voice. Concurrent inserts that would otherwise produce
    overlap are resolved per the reduction rule below.
\end{itemize}

\textbf{Advisory preconditions (authoring mode only):}
\begin{itemize}
  \item For pitched events, every pitch is within the instrument's
    declared range, if any.
  \item The event's duration does not extend past a region boundary
    in a way that would require splitting.
\end{itemize}

\textbf{Effect:} The event is added to the arena; the voice's event
list is updated to include the event in sorted position. Every
contained \texttt{PitchId} is registered as live in the
pitch-identity index. Any attachments declared on the event
(spellings, decompositions) are added to their respective indexes.

\textbf{Postconditions:}
\begin{itemize}
  \item The voice's events list remains sorted and non-overlapping.
  \item The event arena's lookup index resolves the new event.
  \item The event time index includes the new event for its region
    and voice.
  \item Every contained pitch is resolvable through the
    pitch-identity index.
\end{itemize}

\textbf{Inverse:} \texttt{DeleteEvent} of the inserted event, with
captured state equal to the event itself plus any attachments added.

\textbf{Reduction rule:} Position-keyed insertion with voice promotion
on collision. Concurrent inserts at the same position whose durations
would overlap resolve as follows: the operation with the
lexicographically-greater \texttt{OperationId} is placed in a newly
allocated voice on the same staff instance, with
\texttt{VoiceOrigin::SystemPromoted} (Section~\ref{sec:graph:promoted-voices}).
The new voice's \texttt{VoiceId} is derived deterministically from the
staff instance, the original target voice, and the two operations'
identifiers, so all replicas independently allocate the same promoted
voice. The non-overlap invariant is preserved by allocation rather
than by rejection. Subsequent user normalization operations can merge
the promoted voice back into the target, rename it, or accept it as
an intentional second voice.

\subsection{DeleteEvent}
\label{sec:semops:deleteevent}

\textbf{Category:} Event operations.

\begin{lstlisting}[language=Rust]
pub struct DeleteEventOp {
    pub event: EventId,

    /// Compensation for tuplet membership, if any. Required when the
    /// target event is a member of one or more tuplets.
    pub tuplet_compensation: TupletCompensation,
}

pub enum TupletCompensation {
    /// Target is not in any tuplet; no compensation required.
    NotInTuplet,

    /// Replace the deleted event with a rest of the same duration.
    /// The rest receives a freshly-minted EventId; the deleted
    /// event's EventId and PitchIds are tombstoned.
    ReplaceWithRest { new_rest: Rest },

    /// Rewrite the enclosing tuplet(s) to remain structurally
    /// consistent. The v1 payload names the affected tuplets by
    /// id only; it does not carry rewritten values (see the
    /// precondition below).
    RewriteTuplets { tuplets: Vec<TupletId> },

    /// Cascade-delete the entire tuplet group(s) containing the
    /// target event. The listed tuplet ids are cascaded.
    CascadeDeleteTuplets { tuplets: Vec<TupletId> },
}
\end{lstlisting}

\textbf{Invariant preconditions:}
\begin{itemize}
  \item The target event exists in the arena (i.e., is live, not
    tombstoned).
  \item If the target event belongs to one or more tuplets, the
    \texttt{tuplet\_compensation} field \MUST{} be a variant other
    than \texttt{NotInTuplet}, and the chosen compensation \MUST{}
    preserve the tuplet structural consistency invariant
    (Section~\ref{sec:time:tuplets}) after the delete is applied.
  \item If the target event does not belong to any tuplet, the
    \texttt{tuplet\_compensation} field \MUST{} be
    \texttt{NotInTuplet}.
  \item For \texttt{ReplaceWithRest}: the replacement rest's
    duration \MUST{} equal the deleted event's duration; the rest's
    voice and position \MUST{} match the deleted event.
  \item For \texttt{RewriteTuplets}: the v1 payload carries only
    tuplet ids, not the rewritten ratio and member values a reducer
    would need to demonstrate that the resulting structure preserves
    the tuplet consistency invariant. Graph-aware reduction \MUST{}
    therefore refuse the variant as an ill-formed compensation (a
    precondition no-op) rather than fabricate rewritten values. A
    future payload revision carrying full rewrites (per-tuplet ratio
    and members, each of whose member-duration sums \MUST{} match
    the tuplet's structurally-required total) reopens the variant;
    until then \texttt{ReplaceWithRest} and
    \texttt{CascadeDeleteTuplets} are the applicable compensations.
  \item For \texttt{CascadeDeleteTuplets}: every tuplet listed
    \MUST{} contain the target event as a member.
\end{itemize}

\textbf{Effect:} The event is removed from the arena and from its
voice's event list. The event's \texttt{EventId} and every contained
\texttt{PitchId} are tombstoned in their respective identity
indexes: they remain resolvable to ``deleted'' markers and \MUSTNOT{}
be reused for new objects. The chosen tuplet compensation is applied
in the same delta: a replacement rest is inserted or the tuplets are
cascade-deleted as specified (a \texttt{RewriteTuplets} compensation
is refused at precondition check; see above).
Cross-cutting structures referencing the event or its pitches are
re-anchored per the rule table (Section~\ref{sec:semops:reanchor}).
Live attachments targeting the event or its pitches transition to
tombstoned-target state and are preserved for replay, undo, and
historical reference; they are not silently deleted.

\textbf{Postconditions:}
\begin{itemize}
  \item The event is no longer live in the arena.
  \item Every cross-cutting structure that referenced the event has
    been re-anchored or cascade-deleted.
  \item The event's identifier and every contained pitch identifier
    are tombstoned and resolvable only to a deletion record.
  \item The tuplet structural consistency invariant holds for every
    tuplet that was previously consistent.
\end{itemize}

\textbf{Inverse:} \texttt{InsertEvent} reconstructing the event from
captured state, re-promoting the tombstoned identifiers to live,
reversing the tuplet compensation (delete the replacement rest;
restore the original tuplet structure; un-cascade tuplet deletions),
and re-creating any cross-cutting structures that were
cascade-deleted as a consequence. Re-promotion preserves all stable
identifiers; new identifiers \MUSTNOT{} be minted by the inverse.

\textbf{Reduction rule:} Set intersection with delete-wins. Concurrent
deletion and modification of the same event resolves to deletion;
the modification is dropped. Concurrent re-anchoring on the same
referent resolves by the deterministic total ordering specified in
Section~\ref{sec:semops:reanchor}.

\begin{rationale}
  Requiring explicit tuplet compensation prevents primitive deletion
  from leaving the score in a state that violates the tuplet
  structural consistency invariant. The compensation is part of the
  same delta as the deletion, so concurrent observers see either
  the pre-delete state or a fully consistent post-delete state, never
  an intermediate broken state.
\end{rationale}

\subsection{CreateCrossCutting (Slur)}

\textbf{Category:} Cross-cutting structure operations.

\begin{lstlisting}[language=Rust]
pub struct CreateCrossCuttingOp {
    pub structure: CrossCuttingStructure,
}

pub enum CrossCuttingStructure {
    Slur(Slur),
    Tie(Tie),
    Beam(Beam),
    Spanner(Spanner),
    Marker(Marker),
    Repeat(RepeatStructure),
    Analytical(AnalyticalAnnotation),
    Comment(Comment),
    GraphicGesture(GraphicGesture),
}
\end{lstlisting}

The enum above is the data model's full cross-cutting family; the
\emph{wire} vocabulary the cross-cutting operations admit at this
revision (\texttt{CrossCuttingValue}, Binary Format companion) is the
four-kind subset \texttt{Tie}/\texttt{Slur}/\texttt{Beam}/%
\texttt{Spanner}. A \texttt{RepeatStructure} is authored by the
dedicated \texttt{CreateRepeatStructure}/%
\texttt{DeleteRepeatStructure} pair (Operation Catalog
\sectionsc{Repeat Structures}; ratified schema-major-2 revision), not
through \texttt{CreateCrossCutting}; the remaining variants await
their own authoring primitives.

\textbf{Invariant preconditions} (Slur case):
\begin{itemize}
  \item The slur's id is unique within the cross-cutting registry.
  \item The slur's start and end events both exist.
  \item The start and end events are in the same voice.
  \item The start event precedes (or equals) the end event in voice order.
\end{itemize}

\textbf{Advisory preconditions} (Slur case):
\begin{itemize}
  \item The slur does not span a region boundary, unless explicitly
    permitted by region configuration.
\end{itemize}

\textbf{Effect:} The slur is added to the cross-cutting registry's
slur list. The cross-cutting reference index is updated.

\textbf{Postconditions:}
\begin{itemize}
  \item The slur is queryable by id and by referenced events.
\end{itemize}

\textbf{Inverse:} \texttt{DeleteCrossCutting} of the created slur.

\textbf{Reduction rule:} Set union. Concurrent creation of multiple
slurs over overlapping ranges is valid: both slurs exist.

\subsection{RespellPitch}

\textbf{Category:} Pitch and tuning operations.

\begin{lstlisting}[language=Rust]
pub struct RespellPitchOp {
    pub pitch: PitchId,
    pub new_spelling: PitchSpelling,
}
\end{lstlisting}

\textbf{Invariant preconditions:}
\begin{itemize}
  \item The target pitch exists in some event in the arena (i.e.,
    is live, not tombstoned).
  \item The new spelling is valid in the active pitch space: every
    accidental in the spelling's \texttt{accidentals} stack belongs
    to the active accidental registry, the nominal exists in the
    space's nominal registry, and the octave is in range.
  \item Duplicate accidentals in the spelling's stack are absent
    unless the accidental's registered combination semantics
    explicitly permit repetition.
  \item The new spelling represents the same scale position as the
    pitch's current spelling, or is enharmonically equivalent under
    the active tuning system. (Respelling does not change sounding
    pitch.)
\end{itemize}

\textbf{Effect:} A spelling attachment with source
\texttt{UserChosen} is created (or replaces an existing
\texttt{UserChosen} attachment) targeting the pitch. The spelling
attachment index is updated. The pitch's \texttt{PitchId} is
unaffected.

\textbf{Postconditions:}
\begin{itemize}
  \item Spelling resolution for the pitch yields the new spelling,
    subject to the precedence configuration.
  \item The pitch's scale position is unchanged.
  \item The pitch's acoustic realization is unchanged.
  \item The pitch's \texttt{PitchId} is unchanged.
\end{itemize}

\textbf{Inverse:} A \texttt{RespellPitch} restoring the prior
spelling (captured), or a \texttt{RemoveAttachment} if no prior
\texttt{UserChosen} attachment existed.

\textbf{Reduction rule:} Field overwrite by canonical reduction
order. Concurrent respellings of the same live pitch resolve to the
operation later in canonical order; the materialized spelling is
that operation's choice. The losing respelling is recorded in a
\texttt{ConflictRecord} of kind
\texttt{StructuralFieldCollision} if the two respellings represent
different spellings (i.e., differ in nominal, accidental stack, or
octave). The per-operation effect assignment follows the general
field-collision rule
(Requirement~\ref{req:semops:field-collision-effect}): the later
operation --- the winner, which materializes --- reads
\texttt{Conflicted}, and the earlier operation retains
\texttt{Applied}. Identical concurrent respellings reduce
idempotently: the later operation produces
\texttt{NoOp\{reason: AlreadyApplied\}} with no conflict recorded.

\subsection{ChangeRegionTimeModel}

\textbf{Category:} Region operations.

\begin{lstlisting}[language=Rust]
pub struct ChangeRegionTimeModelOp {
    pub region: RegionId,
    pub new_time_model: RegionTimeModel,
    pub position_remapping: PositionRemapping,
}

pub enum PositionRemapping {
    /// Preserve absolute time positions where possible. Events that
    /// cannot be remapped (e.g., metric events under a proportional
    /// model with no implied tempo) are converted by the supplied
    /// conversion function.
    PreserveTime { fallback: PositionConverter },

    /// Discard event positions and reinsert by user-provided
    /// position mapping.
    Reassign(HashMap<EventId, EventPosition>),
}
\end{lstlisting}

\textbf{Invariant preconditions:}
\begin{itemize}
  \item The target region exists.
  \item The new time model is internally consistent.
  \item The position remapping covers every event in the region, or
    the remapping is \texttt{PreserveTime} with a defined fallback.
\end{itemize}

\textbf{Effect:} The region's \texttt{time\_model} is replaced. Every
event in the region has its \texttt{position} field updated according
to the remapping. The tempo map's applicability is recomputed.

\textbf{Postconditions:}
\begin{itemize}
  \item Every event in the region has a valid position in the new
    time model.
  \item The region's invariants (events sorted, non-overlapping,
    within time extent) hold.
\end{itemize}

\textbf{Inverse:} \texttt{ChangeRegionTimeModel} restoring the prior
time model, with a remapping captured from the pre-state.

\textbf{Reduction rule:} Structural migration. The operation is
applied at its canonical position. Contained events whose coordinate
kinds become incompatible with the new time model after migration
cause the operation to reduce to \texttt{Conflicted} with a
\texttt{ConflictKind::TimeModelMigrationFailure} record listing the
incompatible events; the migration does not apply. Concurrent
operations on the same region that author events in the prior time
model and reduce after the migration may themselves reduce to
\texttt{Conflicted} if their coordinate kinds disagree with the new
model. Concurrent same-target migrations resolve by canonical order;
the earlier operation applies and the later operation reduces to
\texttt{Conflicted} with a structural-field-collision conflict (the
two migrations are not generally commutable).

\begin{rationale}
  Time-model migration is a structural rebase, not a field set.
  Treating it as LWW would silently destroy contained musical
  content whose coordinate kinds depend on the prior time model.
  Conflicted reductions force the user to address the incompatibility
  explicitly.
\end{rationale}

\subsection{Transpose}

\textbf{Category:} Pitch and tuning operations. This is a compound
operation, defined as a sequence of primitives, listed here as an
illustration of compounds.

\textbf{User-facing parameters:}
\begin{itemize}
  \item Target scope (events, voice, staff, region, or selection).
  \item Interval (diatonic, chromatic, or compound).
  \item Spelling policy (preserve user spellings, respell to fit a
    target key, or apply a custom rule).
  \item Tuning policy (preserve, recompute under new tuning, or
    apply offsets).
\end{itemize}

\textbf{Primitive decomposition:} A transpose operation expands into:

\begin{enumerate}
  \item For each target pitch, a \texttt{ReplaceEvent} updating the
    inner \texttt{Pitch} values while preserving every
    \texttt{IdentifiedPitch}'s \texttt{PitchId}. Transposition does
    not mint new pitch identifiers; it rewrites the contents of
    existing identified pitches.
  \item For each pitch whose spelling is recomputed, a
    \texttt{RespellPitch} targeting the unchanged \texttt{PitchId}
    with source \texttt{Propagated}.
  \item Where the transposition crosses a clef boundary, optional
    \texttt{ChangeStaffConfiguration} operations to insert clef
    changes.
\end{enumerate}

The compound operation is recorded as a single transaction. Each
primitive within the transaction has its own re-anchoring and
reduction semantics. Identifier preservation ensures that
attachments, analytical layers, and cross-cutting references
targeting the transposed pitches continue to resolve correctly
after the operation.

\subsection{SetUserSystemBreak}

\textbf{Category:} Layout-semantic operations.

\begin{lstlisting}[language=Rust]
pub struct SetUserSystemBreakOp {
    pub anchor: TimeAnchor,
    pub region: RegionId,
    pub present: bool,
}
\end{lstlisting}

\textbf{Invariant preconditions:}
\begin{itemize}
  \item The region exists and is staff-based.
  \item The anchor resolves within the region.
\end{itemize}

\textbf{Effect:} The region's \texttt{user\_system\_breaks} list is
updated to include or exclude the anchor. Layout caches that depend
on system breaks are invalidated.

\textbf{Postconditions:}
\begin{itemize}
  \item The region's user system breaks list reflects the change.
  \item Layout output, if produced subsequently, accounts for the
    change subject to the engraver's override rules
    (Chapter~\ref{ch:layout-ir}).
\end{itemize}

\textbf{Inverse:} \texttt{SetUserSystemBreak} with the opposite
\texttt{present} value.

\textbf{Reduction rule:} LwwAdvisory on per-anchor break preference.
Two concurrent operations setting the same break to \texttt{true}
reduce idempotently. Concurrent set-true and set-false on the same
anchor resolve by canonical reduction order; the later operation's
preference is materialized. No conflict record is produced; the
break preference is an advisory field eligible for LWW reduction
(Section~\ref{sec:semops:lww}).

\section{Operation Catalog Conformance}
\label{sec:semops:catalog}

This chapter specifies the framework. The full catalog of operations
is delivered as a separate conformance specification that extends
this chapter.

\begin{requirement}
  \label{req:semops:operation-catalog-coverage}
  A conforming implementation \MUST{} provide every operation
  enumerated in the conformance catalog. Each operation \MUST{}
  satisfy the framework requirements stated in this chapter:

  \begin{itemize}
    \item Declared preconditions classified as invariant or advisory.
    \item Specified effect on the score graph.
    \item Specified postconditions, including which invariants are
      preserved.
    \item A defined inverse operation.
    \item A declared reduction rule under the canonical reduction
      model (Section~\ref{sec:semops:reduction}).
    \item Specified re-anchoring behavior for any deletions or
      moves it triggers.
  \end{itemize}
\end{requirement}

\begin{openquestion}
  The conformance catalog is presently under development as a separate
  document. It is expected to comprise approximately 60--80 primitive
  operations and an open-ended set of compound operations. The catalog
  is normative once published; this specification is non-final until
  the catalog is delivered.
\end{openquestion}

\section{Forward References}

\begin{itemize}
  \item Operation envelope storage on disk (envelope block layout,
    causal-summary metadata, the operation index, snapshot
    canonical-base semantics) is specified in
    Chapter~\ref{ch:format}.
  \item Conflict-resolution algorithms (how a
    \texttt{ResolveConflict} operation transitions a conflict's
    state, and the full conflict-record contract) are specified
    in this chapter (Section~\ref{sec:semops:conflicts}) and in
    the Operation Catalog companion
    (Appendix~\ref{app:deferred}).
  \item The interaction between layout-semantic operations (system
    breaks, page breaks, engraving overrides) and the layout pipeline
    is specified in Chapter~\ref{ch:layout-ir}.
  \item The conformance catalog of all operations is a separate
    deliverable that extends this chapter.
\end{itemize}

% ===========================================================================
\chapter{Layout Intermediate Representation}
\label{ch:layout-ir}

This chapter specifies the layout intermediate representation (Layout
IR): the data structures and pipeline that transform the score graph
into a typeset visual layout. The IR sits between the score graph and
two downstream consumers: the constraint solver
(Chapter~\ref{ch:solver}), which resolves spacing and positioning,
and the renderer (specified outside this document), which produces
final visual output (pixels, PDF, SVG).

The IR is a pipeline of distinct stages, each with its own type and
each with a well-defined contract for the next. The pipeline is the
spine of the engraving system; the engraving algorithms themselves
(beam grouping, slur curvature, accidental ordering, casting-off) are
the consumers and producers of stages and are specified outside this
document as engraving-algorithm specifications layered atop this IR.

\section{Design Principles}
\label{sec:layoutir:principles}

\begin{description}
  \item[Pipeline of stages.] The transformation from score graph to
    visual output is a pipeline of four IR stages:
    \texttt{LogicalLayoutIR} (structural projection without spatial
    information), \texttt{ConstrainedLayoutIR} (with constraint inputs
    but unresolved positions), \texttt{ResolvedLayoutIR} (with all
    positions resolved), and \texttt{RenderIR} (renderer-bound
    primitives, defined only at the interface level here). Each stage
    has its own type. Compilation between stages is unidirectional.
  \item[Shared provenance.] Every IR object at every stage carries a
    reference back to the score graph object that originated it. This
    enables selection, editing back-references, error reporting, and
    incremental layout invalidation.
  \item[Engraving decisions are explicit.] When the engraver makes a
    decision (stem direction, accidental ordering, beam consolidation),
    the decision is recorded in the IR. The decision can be inspected,
    overridden, and traced.
  \item[Overrides exist in both worlds.] User-set engraving overrides
    are authoritative in the score graph and projected into the IR
    during the logical stage. The IR may also generate transient
    overrides for downstream stages.
  \item[Spring-based spacing.] Time-to-space projection produces
    elastic spring parameters per time slot. The constraint solver
    resolves springs system-wide. This applies to both horizontal
    and vertical layout.
  \item[Staff-space units, f32 precision.] Spatial coordinates within
    the IR are expressed in staff spaces using single-precision
    floating point. Conversion to absolute units (points, millimeters)
    occurs only at the render boundary.
  \item[Glyph metrics live elsewhere.] The IR references glyphs by
    identifier and queries metrics from a font catalog. Metrics are
    not embedded in IR objects.
  \item[Incremental layout, fine-grained.] The IR supports incremental
    re-engraving via per-object dependency tracking on score graph
    objects. Edits invalidate only the layout objects whose
    dependencies changed.
  \item[Uniform region containers.] Staff-based, free-graphic, and
    hybrid regions produce structurally identical \texttt{LayoutRegion}
    containers with differing time projections and content kinds.
\end{description}

\section{Spatial Primitives}
\label{sec:layoutir:spatial}

\subsection{Units}

The IR uses two unit types:

\begin{lstlisting}[language=Rust]
/// Staff space: the fundamental unit of music engraving.
/// One staff space equals the distance between adjacent staff lines.
#[derive(Copy, Clone, PartialEq, PartialOrd, Debug)]
pub struct StaffSpace(pub f32);

/// Point: 1/72 inch. Used at the render boundary for absolute
/// dimensions (page sizes, font sizes, output coordinates).
#[derive(Copy, Clone, PartialEq, PartialOrd, Debug)]
pub struct Point(pub f32);

/// Scaling context for converting between staff spaces and points.
pub struct ScaleContext {
    /// Points per staff space. Depends on the staff size chosen for
    /// the score; typical values are 4--10 points per staff space.
    pub points_per_staff_space: f32,
}
\end{lstlisting}

\begin{requirement}
  \label{req:layoutir:staff-space-coordinates}
  Spatial coordinates within the IR (in stages
  \texttt{LogicalLayoutIR}, \texttt{ConstrainedLayoutIR}, and
  \texttt{ResolvedLayoutIR}) \MUST{} be expressed in staff spaces.
  Conversion to points occurs at the \texttt{RenderIR} boundary.

  Single-precision floating point (\texttt{f32}) \MUST{} be used for
  IR coordinates. Engraving precision requirements are bounded by
  printer resolution (typically 600--2400 DPI); single precision
  provides over 7 significant figures, comfortably exceeding any
  engraving precision requirement.
\end{requirement}

\subsection{Geometric Types}

\begin{lstlisting}[language=Rust]
#[derive(Copy, Clone, Debug)]
pub struct Point2D {
    pub x: StaffSpace,
    pub y: StaffSpace,
}

#[derive(Copy, Clone, Debug)]
pub struct Size2D {
    pub width: StaffSpace,
    pub height: StaffSpace,
}

#[derive(Copy, Clone, Debug)]
pub struct Rect {
    pub origin: Point2D,
    pub size: Size2D,
}

#[derive(Copy, Clone, Debug)]
pub struct BoundingBox {
    pub left: StaffSpace,
    pub right: StaffSpace,
    pub top: StaffSpace,
    pub bottom: StaffSpace,
}

#[derive(Copy, Clone, Debug)]
pub struct Transform2D {
    pub matrix: [[f32; 3]; 3],
}
\end{lstlisting}

\section{Provenance}
\label{sec:layoutir:provenance}

Every IR object carries a provenance record tracing it back to the
score graph object that originated it. Provenance is uniform across
all four stages; provenance records are shared structurally between
stages so that an object passing through the pipeline retains a
consistent identity.

\begin{lstlisting}[language=Rust]
pub struct Provenance {
    /// The primary source: the score graph object this IR object
    /// represents or derives from.
    pub source: TypedObjectId,

    /// For IR objects that derive from synthesizing engraving
    /// decisions rather than directly from score graph objects, the
    /// synthesis kind. Examples: an automatically inserted natural
    /// sign, a generated rest, a system break.
    pub synthesis: Option<SynthesisKind>,

    /// Additional source dependencies. An IR object may depend on
    /// multiple score graph objects (e.g., a beam depends on its
    /// member events). All dependencies are tracked for
    /// incremental layout.
    pub dependencies: Vec<TypedObjectId>,

    /// IR object's stable identifier across re-layouts where the
    /// underlying source is unchanged.
    pub stable_id: LayoutObjectId,
}

pub enum SynthesisKind {
    /// An accidental inserted to cancel a prior alteration.
    CancellationAccidental,

    /// A natural sign generated by key signature rules.
    KeySignatureNatural,

    /// A rest inserted to fill a voice gap.
    GeneratedRest,

    /// A system or page break selected by the engraver.
    EngravedBreak,

    /// A multimeasure rest combining consecutive empty measures.
    MultimeasureRest,

    /// A cautionary key or time signature at a system break.
    Cautionary,

    /// A custom synthesis declared by a layout extension.
    Registered(SynthesisRegistryId),
}

pub struct LayoutObjectId(pub u128);
\end{lstlisting}

\begin{requirement}
  \label{req:layoutir:provenance-completeness}
  Every layout object at every IR stage \MUST{} carry a non-empty
  \texttt{Provenance} record. Objects whose direct \texttt{source}
  does not correspond to a score graph object (engraver-synthesized
  objects) \MUST{} declare a \texttt{synthesis} kind. The
  \texttt{dependencies} list \MUST{} include every score graph object
  whose change should invalidate this layout object.
\end{requirement}

\begin{requirement}
  \label{req:layoutir:continuation-synthesis}
  \textbf{System-spanning continuation.} A line primitive (stroke) that spans a
  system boundary is split into per-system segments; the segments after the
  first are engraver-synthesized and \MUST{} declare a \texttt{synthesis} kind.
  The reserved registered kind
  \texttt{SynthesisKind::Registered(SYSTEM\_CONTINUATION\_SYNTHESIS)} names this
  continuation. Its \texttt{stable\_semantic\_instance\_key} (per
  Requirement~\ref{req:layoutir:object-id-derivation}) is the pair
  \texttt{(original,\,ordinal)} --- the pre-split stroke's object and the
  segment's order within the split --- which is stable for a fixed set of system
  breaks. Because \texttt{LayoutObjectId}s are non-canonical and re-derived per
  layout, this key need only be stable within a layout, not across relayouts
  that move the breaks.
\end{requirement}

\begin{requirement}
  \label{req:layoutir:object-id-derivation}
  A \texttt{LayoutObjectId} \MUST{} be stable across re-layouts whose
  underlying source is unchanged. It \MUST{} be derived by
  domain-separated BLAKE3 truncation with the reserved layout domain
  tag \texttt{"MUSCLOID"} over a key that depends on how the object is
  manifested:

  \begin{itemize}
    \item A layout object that manifests a single score-graph object
      exactly once is keyed on \texttt{source.canonical\_bytes()}.
    \item A multiply-manifested object --- one score-graph object that
      appears in more than one layout context, e.g. a staff manifested
      in two regions --- is keyed on the pair
      \texttt{(source,\,region)}, so each manifestation receives a
      distinct, stable id.
    \item A synthesized object (no direct score-graph
      \texttt{source}; \texttt{synthesis} is \texttt{Some}) is keyed
      on the triple \texttt{(source,\,synthesis\_kind,\,
      stable\_semantic\_instance\_key)}, where the
      \texttt{stable\_semantic\_instance\_key} distinguishes multiple
      synthesized objects sharing a source and synthesis kind by a
      semantically stable discriminator (not a layout-position
      ordinal, which would not survive relayout).
  \end{itemize}

  \texttt{LayoutObjectId}s are \emph{non-canonical}: they are not part
  of document state and do not enter any content hash, so the
  \texttt{"MUSCLOID"} tag lives in the layout namespace and is not one
  of the reserved \emph{canonical} system tags of
  Section~\ref{sec:graph:system-derived}. The fixed derivation is
  pinned for incremental-relayout correctness and provenance
  back-reference stability; its consumers are the solver and renderer
  (Track~A), not the interchange track. Because these ids never enter
  document state, no stored or interchanged artifact depends on this
  derivation. The v0 reference implementation realizes this derivation
  as of Pass-12 item P12-I2: \texttt{epiphany-determinism} reserves the
  built-in \texttt{"MUSCLOID"} layout tag and the \texttt{layout-ir}
  provenance derivations (single, multiply-manifested, and synthesized)
  route through it, the synthesized case no longer borrowing
  \texttt{"MUSCCONF"}. Because the ids are non-canonical, doing so
  changed layout-id values but no stored or interchanged artifact.
\end{requirement}

\section{The Stage Pipeline}
\label{sec:layoutir:pipeline}

\subsection{Pipeline Overview}

\begin{lstlisting}
ScoreGraph
   |
   |  Engraving pass:
   |  - Project structure
   |  - Make engraving decisions (stems, beams, accidentals)
   |  - Apply user overrides
   v
LogicalLayoutIR
   |
   |  Spacing pass:
   |  - Compute spring parameters per time slot
   |  - Compute glyph bounding boxes
   |  - Build collision constraints
   v
ConstrainedLayoutIR
   |
   |  Constraint solving (Chapter 9):
   |  - Resolve horizontal positions
   |  - Resolve vertical positions
   |  - Resolve page and system breaks
   v
ResolvedLayoutIR
   |
   |  Render projection (out of scope):
   |  - Convert to renderer primitives
   |  - Apply page templates
   |  - Generate draw calls
   v
RenderIR
\end{lstlisting}

\subsection{Stage Contracts}

\begin{requirement}
  \label{req:layoutir:deterministic-stage-transitions}
  Each pipeline stage \MUST{} accept its input stage as a pure
  function: identical inputs (plus identical configuration) \MUST{}
  produce identical outputs. The stages are deterministic by
  contract.

  Stage transitions \MUST{} preserve provenance: every output object
  \MUST{} carry a \texttt{Provenance} record either inherited from its
  input or newly synthesized with appropriate \texttt{synthesis} kind.
\end{requirement}

\section{LogicalLayoutIR}
\label{sec:layoutir:logical}

\texttt{LogicalLayoutIR} is the structural projection of the score
graph into layout objects, with all engraving decisions made but
spatial positions unresolved. It is the output of the engraving pass
and the input to the spacing pass.

\subsection{Top-Level Structure}

\begin{lstlisting}[language=Rust]
pub struct LogicalLayoutIR {
    /// Source score graph reference (typically a snapshot pointer
    /// or version identifier).
    pub source: ScoreVersion,

    /// Layout regions, one per score graph region.
    pub regions: Vec<LayoutRegion>,

    /// Engraving decisions made during this pass.
    pub engraving_decisions: Vec<EngravingDecision>,

    /// User overrides projected from the score graph.
    pub overrides: Vec<EngravingOverride>,

    /// Cross-region layout objects (slurs spanning regions, etc.).
    pub cross_region: Vec<CrossRegionObject>,
}
\end{lstlisting}

\subsection{Layout Regions}

\begin{lstlisting}[language=Rust]
pub struct LayoutRegion {
    pub provenance: Provenance,

    /// Coordinate system local to this region. The region's
    /// transform maps local coordinates to canvas coordinates.
    pub coordinate_system: LocalCoordinateSystem,

    /// Time axis: how time positions map to horizontal positions.
    /// A tagged union over the three built-in axis kinds plus a
    /// registered variant for extension-defined axes. The enum
    /// form is canonical for serialization, hashing, and
    /// conformance comparison; implementations MAY use a trait
    /// object internally for dynamic dispatch.
    pub time_axis: TimeAxisModel,

    /// Vertical extent: the set of staff bands this region occupies.
    pub vertical_extent: VerticalExtent,

    /// Layout objects within this region.
    pub objects: Vec<LayoutObject>,
}

/// Canonical representation of a region's time axis. Variants
/// correspond to the three built-in region time models; the
/// Registered variant carries extension-defined axes by typed
/// registry id and serialized payload.
pub enum TimeAxisModel {
    Metric(MetricTimeAxis),
    Proportional(ProportionalTimeAxis),
    Aleatoric(AleatoricTimeAxis),
    Registered(TimeAxisRegistryId, SerializedRegisteredAxis),
}
\end{lstlisting}

\subsection{The Time Axis Trait}

For dynamic dispatch and ergonomic implementation, the
\texttt{TimeAxisModel} enum's payload types share a trait:

\begin{lstlisting}[language=Rust]
pub trait TimeAxis: Send + Sync {
    /// The time model kind this axis implements.
    fn kind(&self) -> TimeAxisKind;

    /// Project a time position to a horizontal spring slot.
    /// Returns a slot identifier; spring parameters are queried
    /// separately.
    fn project(&self, time: TimePoint) -> SpringSlotId;

    /// Enumerate the spring slots in time order.
    fn slots(&self) -> &[SpringSlotId];

    /// For incremental layout: which slots are affected by a change
    /// to the given time range?
    fn affected_slots(&self, range: TimeRange) -> Vec<SpringSlotId>;
}

pub enum TimeAxisKind {
    Metric,
    Proportional,
    Aleatoric,
    Registered(TimeAxisRegistryId),
}

pub enum TimePoint {
    Musical(MusicalPosition),
    WallClock(WallClockTime),
}

pub enum TimeRange {
    Musical { start: MusicalPosition, end: MusicalPosition },
    WallClock { start: WallClockTime, end: WallClockTime },
}
\end{lstlisting}

The three built-in implementations:

\begin{lstlisting}[language=Rust]
pub struct MetricTimeAxis {
    /// Measure boundaries and their projected spring slots.
    pub measures: Vec<MeasureProjection>,

    /// Per-time-slot spring parameters, indexed by slot id.
    pub slots: Vec<SpringSlot>,
}

pub struct ProportionalTimeAxis {
    /// Wall-clock duration of the region.
    pub duration: WallClockDuration,

    /// Linear horizontal projection: space per second.
    pub space_per_second: StaffSpace,

    pub slots: Vec<SpringSlot>,
}

pub struct AleatoricTimeAxis {
    /// DAG of event ordering.
    pub ordering: EventOrderingDAG,

    /// Spring slots are organized by topological-order layers.
    pub slots: Vec<SpringSlot>,
}
\end{lstlisting}

\subsection{Layout Objects}

A \texttt{LayoutObject} is the unit of engraving. At the
\texttt{LogicalLayoutIR} stage, layout objects are composite: a
single \texttt{Note} object contains its notehead, stem, flag,
accidental, articulations, and so on. At the
\texttt{ConstrainedLayoutIR} stage these are flattened into
individual glyph objects.

\begin{lstlisting}[language=Rust]
pub enum LayoutObject {
    /// Composite note object: notehead(s), stem, flag, accidental(s),
    /// articulations, ornaments, dot(s).
    Note(NoteLayout),

    /// Composite chord object: multiple noteheads sharing stem and
    /// articulations.
    Chord(ChordLayout),

    /// Rest with its glyph kind (whole, half, quarter, ...).
    Rest(RestLayout),

    /// Beam group: multiple notes/chords beamed together.
    BeamGroup(BeamGroupLayout),

    /// Tuplet visual: bracket and number.
    TupletDisplay(TupletDisplayLayout),

    /// Slur or phrase mark: control points and style.
    Slur(SlurLayout),

    /// Tie: same as slur but connecting same-pitch events.
    Tie(TieLayout),

    /// Spanner: hairpins, octave lines, pedal lines, etc.
    Spanner(SpannerLayout),

    /// Marker: rehearsal marks, section labels, tempo marks.
    Marker(MarkerLayout),

    /// Bar line: simple, double, final, repeat, etc.
    BarLine(BarLineLayout),

    /// Clef, key signature, time signature glyphs.
    Clef(ClefLayout),
    KeySignature(KeySignatureLayout),
    TimeSignatureDisplay(TimeSignatureDisplayLayout),

    /// Staff: the staff lines themselves.
    Staff(StaffLayout),

    /// Text: lyrics, dynamics text, expressive text.
    Text(TextLayout),

    /// Free graphic: path, shape, image, stroke.
    Graphic(GraphicLayout),

    /// Multimeasure rest: consolidated empty measures.
    MultimeasureRest(MultimeasureRestLayout),

    /// Cue: small-print rendering of source material.
    Cue(CueLayout),

    /// Trajectory: glissando, portamento line.
    Trajectory(TrajectoryLayout),

    /// Generic group: implementation-defined.
    Group(GroupLayout),
}
\end{lstlisting}

\subsection{Note Layout: A Worked Example}

To illustrate the composite-object pattern at the logical stage:

\begin{lstlisting}[language=Rust]
pub struct NoteLayout {
    pub provenance: Provenance,

    /// Reference to the score graph event.
    pub event: EventId,

    /// The notehead glyph and its staff position.
    pub notehead: NoteheadLayout,

    /// The stem, if any.
    pub stem: Option<StemLayout>,

    /// Flag, if any (single-note events with eighth-or-shorter
    /// duration that are not part of a beam group).
    pub flag: Option<FlagLayout>,

    /// Accidental glyphs attached to the note, in stacking order.
    pub accidentals: Vec<AccidentalLayout>,

    /// Augmentation dots.
    pub dots: Vec<DotLayout>,

    /// Articulations attached to the note.
    pub articulations: Vec<ArticulationLayout>,

    /// Ornament glyphs.
    pub ornaments: Vec<OrnamentLayout>,

    /// Ledger lines required to display the note's staff position.
    pub ledger_lines: Vec<LedgerLineLayout>,

    /// Engraving decisions made for this note.
    pub decisions: NoteDecisions,
}

pub struct NoteDecisions {
    pub stem_direction: StemDirection,
    pub accidental_visible: bool,
    pub accidental_parenthesized: bool,
    pub notehead_shape: NoteheadShape,
    /// Voice membership influences default stem direction;
    /// overrides may change it.
    pub stem_direction_source: DecisionSource,
}

pub enum DecisionSource {
    /// Decision derived from automatic engraving rules.
    Automatic,

    /// Decision derived from a user override in the score graph.
    UserOverride(EngravingOverrideId),

    /// Decision derived from an IR-stage override.
    IrOverride,
}
\end{lstlisting}

\subsection{Engraving Decisions}

The pipeline records engraving decisions explicitly so they can be
inspected, overridden, and replayed.

\begin{lstlisting}[language=Rust]
pub struct EngravingDecision {
    pub id: EngravingDecisionId,
    pub target: LayoutObjectId,
    pub kind: EngravingDecisionKind,
    pub source: DecisionSource,
}

pub enum EngravingDecisionKind {
    StemDirection(StemDirection),
    BeamGrouping(BeamGroupingChoice),
    AccidentalOrdering(Vec<AccidentalId>),
    SlurDirection(SlurDirection),
    LedgerLineCount(u8),
    EnharmonicSpelling(PitchSpelling),
    SystemBreak,
    PageBreak,
    NoteheadShape(NoteheadShape),
    RestPosition(StaffPosition),
    // ... extensible
}
\end{lstlisting}

\section{Engraving Overrides}
\label{sec:layoutir:overrides}

Engraving overrides are user assertions that the engraver's default
decision should be replaced. Overrides are authoritative in the score
graph and projected into the IR during the logical stage.

\subsection{Override Definition}

\begin{lstlisting}[language=Rust]
pub struct EngravingOverride {
    pub id: EngravingOverrideId,

    /// What this override targets. May reference a score graph object
    /// (typical) or an IR-synthesized object (transient overrides).
    pub target: OverrideTarget,

    /// The override itself.
    pub kind: OverrideKind,

    /// Binding strength.
    pub priority: OverridePriority,

    /// Provenance: who set the override and when.
    pub origin: OverrideOrigin,
}

pub enum OverrideTarget {
    /// Targets a score graph object. Authoritative; survives reload.
    ScoreGraph(TypedObjectId),

    /// Targets an IR-synthesized object. Transient; cleared on
    /// re-engraving unless explicitly re-applied.
    IrSynthesized(LayoutObjectId),
}

pub enum OverrideKind {
    StemDirection(StemDirection),
    SlurDirection(SlurDirection),
    SlurCurvature(CurvatureOverride),
    AccidentalParenthesized(bool),
    AccidentalVisible(bool),
    NoteheadShape(NoteheadShape),
    BeamGeometry(BeamGeometryOverride),
    SystemBreak { anchor: TimeAnchor },
    PageBreak { anchor: TimeAnchor },
    HiddenObject,
    CustomPosition(Point2D),
    EnharmonicSpelling(PitchSpelling),
    LedgerLineSuppression,
    DotPlacement(DotPlacement),
    // ... extensible
}

// System- and page-break overrides address a *position*, not an
// object: the break kind carries the break's TimeAnchor, while the
// override's ScoreGraph target names the owning region. They are
// projected during the logical stage from the score graph's
// authoritative user_system_breaks / user_page_breaks lists
// (Chapter 5). Break authorship (author, timestamp) lives in the
// operation log, not the materialized break lists, so projected
// break overrides carry OverrideOrigin::Internal until the
// snapshot-undo refinement (P11-C8) surfaces authorship.

pub enum OverridePriority {
    /// The engraver MUST honor this override. Failure produces an
    /// error rather than a silent override-of-the-override.
    Hard,

    /// The engraver SHOULD honor this override. The engraver may
    /// override it only under documented exceptional conditions
    /// (system overflow, unresolvable collision); when it does, the
    /// override is recorded in the resulting decisions with source
    /// IrOverride.
    Soft,
}

pub enum OverrideOrigin {
    User { author: AuthorId, timestamp: Timestamp },
    Import { format: ForeignFormatId },
    Plugin { plugin: PluginId },
    Internal,
}
\end{lstlisting}

\subsection{Override Resolution}

\begin{requirement}
  \label{req:layoutir:override-enforcement}
  Overrides \MUST{} be applied during the engraving pass producing
  \texttt{LogicalLayoutIR}. Hard overrides \MUST{} be honored or the
  pass \MUST{} report an error. Soft overrides \MUST{} be honored
  except where doing so would produce an invalid layout; when a soft
  override is not honored, the resulting decision \MUST{} record the
  override and the reason for its non-application.
\end{requirement}

\begin{requirement}
  \label{req:layoutir:override-target-lifetime}
  Overrides whose targets reference objects in the score graph (the
  \texttt{ScoreGraph} variant) \MUST{} be preserved across
  re-engraving. Overrides targeting IR-synthesized objects (the
  \texttt{IrSynthesized} variant) are transient; they are cleared on
  re-engraving unless explicitly re-applied. Implementations \MUST{}
  not silently promote IR-synthesized overrides to score-graph
  overrides without explicit user action.
\end{requirement}

\begin{requirement}
  \label{req:layoutir:break-origin-attribution}
  \textbf{Break-override attribution.} When the engraver honours a user break
  override, the resulting decision \MUST{} carry
  \texttt{DecisionSource::UserOverride(id)} naming the override. Because a
  normalized break \emph{constraint} (\texttt{SystemBreakAt} /
  \texttt{PageBreakAt}, Section~\ref{sec:layoutir:constrained}) carries no
  override identity, this attribution is threaded through the projection: the
  \texttt{ConstrainedLayoutIR} carries a \texttt{break\_origins} sidecar,
  populated when the logical stage projects the break, mapping each user-anchored
  break slot to its override id. The normalized constraint record is
  deliberately \emph{not} widened to carry override identity --- attribution is a
  projection concern, not a constraint-solver input.
\end{requirement}

\section{ConstrainedLayoutIR}
\label{sec:layoutir:constrained}

\texttt{ConstrainedLayoutIR} is the output of the spacing pass: the
logical IR with composite objects flattened to individual glyphs,
each glyph carrying its bounding box, anchor, and constraint inputs
to the solver.

\begin{requirement}
  \label{req:layoutir:constraint-floor}
  \textbf{Minimal-tier constraint-emission floor (ratified Pass 12).}
  The spacing pass \MUST{} emit at least the following constraint
  set --- the \emph{Minimal-tier floor}, the testable acceptance
  surface for the \texttt{Minimal} conformance tier
  (Chapter~\ref{ch:solver}): no-collision chains between
  successive notehead-bearing columns; per-glyph containment
  within the owning region's frame; and one constraint per user
  break override (system and page). Richer emission sets
  (inter-band collision, kerned optical adjacency, duration-
  proportional spring preferences) belong to higher tiers, whose
  floors are defined when those tiers are ratified.
\end{requirement}

\subsection{Glyph-Level Objects}

\begin{lstlisting}[language=Rust]
pub struct GlyphObject {
    pub provenance: Provenance,

    /// The glyph to render. Reference into a font catalog; metrics
    /// are queried from the catalog rather than embedded here.
    pub glyph: GlyphReference,

    /// The spring slot this glyph belongs to (horizontal time slot).
    pub horizontal_slot: SpringSlotId,

    /// The vertical band this glyph belongs to.
    pub vertical_band: VerticalBandId,

    /// Bounding box of the glyph relative to its anchor, in staff
    /// spaces. Looked up from the font catalog at IR construction
    /// time and cached here.
    pub bounding_box: BoundingBox,

    /// Anchor point: where the glyph's reference position sits
    /// within its bounding box.
    pub anchor: Point2D,

    /// Layer ordering: higher draws on top.
    pub layer: i32,

    /// Visual style applied to the glyph.
    pub style: GlyphStyle,
}
\end{lstlisting}

\subsection{Spring Slots}

A spring slot is a time slot's spacing parameters: the constraint
solver's input for resolving horizontal positions.

\begin{lstlisting}[language=Rust]
pub struct SpringSlot {
    pub id: SpringSlotId,
    pub time: TimePoint,

    /// Minimum width: the smallest the slot may compress to. The
    /// constraint solver MUST NOT compress below this.
    pub min_width: StaffSpace,

    /// Preferred width: the natural width under unconstrained
    /// spacing. Derived from Gourlay-style proportional rules.
    pub preferred_width: StaffSpace,

    /// Maximum width, if bounded. None means unbounded above.
    pub max_width: Option<StaffSpace>,

    /// Stretch factor: how readily this slot stretches when more
    /// width is available than preferred. Higher = more elastic.
    pub stretch_factor: f32,

    /// Compress factor: how readily this slot compresses when less
    /// width is available than preferred. Higher = more compressible.
    pub compress_factor: f32,

    /// Glyphs belonging to this slot, ordered by stacking
    /// considerations (accidentals to left of notehead, etc.).
    pub members: Vec<GlyphObjectId>,
}
\end{lstlisting}

\subsection{Vertical Bands}

Vertical layout uses the same spring model. A \emph{vertical band}
is a horizontal slice of the canvas (typically a staff or an
inter-staff gap) with its own spring parameters.

\begin{lstlisting}[language=Rust]
pub struct VerticalBand {
    pub id: VerticalBandId,
    pub kind: VerticalBandKind,

    pub min_height: StaffSpace,
    pub preferred_height: StaffSpace,
    pub max_height: Option<StaffSpace>,
    pub stretch_factor: f32,
    pub compress_factor: f32,

    /// Glyphs belonging to this band.
    pub members: Vec<GlyphObjectId>,
}

pub enum VerticalBandKind {
    Staff(StaffId),
    InterStaffGap,
    InterSystemGap,
    MarginBand,
}
\end{lstlisting}

\begin{requirement}
  \label{req:layoutir:coverage-diagnostics}
  \textbf{Engraving-coverage gaps are surfaced, not guessed.} A projection
  that cannot engrave an object faithfully --- a pitch with no resolved
  spelling, a glyph the bound catalog
  (\ref{sec:layoutir:catalog}) does not carry --- \MUST{} record a
  \texttt{LayoutDiagnostic} naming the score-graph object and the kind of
  gap, and \MUST{} still place the object: a fallback notehead on the clef
  reference line, or a zero-extent traced anchor that keeps its provenance.
  It \MUSTNOT{} silently substitute a plausible shape, and it \MUSTNOT{}
  drop the object.

  Both halves matter. Dropping the object breaks the round-trip surjection
  (\ref{sec:layoutir:provenance}) --- a hit-test can no longer find it, and
  an editor can no longer select what the author wrote. Guessing a shape
  produces a score that looks engraved and is wrong, with nothing in the IR
  to say so. A diagnostic is not an error: the layout is still
  \emph{renderable}, and a conformance claim is unaffected. It is the record
  that the engraving is incomplete at a named place.
\end{requirement}

\begin{requirement}
  \label{req:layoutir:primitive-band-ownership}
  \textbf{Primitive band ownership is declared, not inferred.} Every
  primitive the projection presents to the solver in
  \texttt{ConstrainedLayoutIR} --- \texttt{GlyphObject},
  \texttt{Stroke}, and \texttt{Curve} alike --- \MUST{} declare the
  \texttt{VerticalBandId} of the band that owns it, and that band
  \MUST{} exist in \texttt{vertical\_bands}. A vertical solver \MUST{}
  take a primitive's owning staff from that declaration and
  \MUSTNOT{} infer it from the primitive's geometry. Content owned by
  no staff (a page-margin annotation, a structure spanning several
  staves) names a band whose kind is not \texttt{Staff}.

  This is a correctness requirement, not a convenience. A primitive's
  drawn position does not determine its owner: a stem shares its
  horizontal column with the staff above, and a slur's endpoints are
  \emph{lifted clear} of its own staff by construction, landing in the
  inter-staff zone where the nearest notehead routinely belongs to the
  neighbouring staff. A solver that guesses from proximity will
  attribute such a primitive to the wrong staff and, on renegotiating
  the gaps between a system's staves, tear it off the notes it belongs
  to.

  A band \MUST{} therefore exist for every staff of a region, whether
  or not any glyph landed in it, and for the region's margin ---
  otherwise a primitive that names one has nowhere to point. A band
  MAY carry no members.

  Only glyphs are \emph{members} of a band
  (\texttt{VerticalBand::members}), which is what realizes the spring
  solve; a stroke's or curve's band reference is one-way. What becomes
  of the declaration downstream is
  \ref{req:layoutir:resolved-band-ownership}.
\end{requirement}

\subsection{Constraints}

The IR carries explicit constraint inputs for the solver. Beyond
glyphs it presents the non-glyph primitives --- \texttt{Stroke} and
\texttt{Curve}, whose shapes are listed with
\ref{req:layoutir:resolved-primitives} --- since the solver positions
those too:

\begin{lstlisting}[language=Rust]
pub struct ConstrainedLayoutIR {
    pub source: ScoreVersion,

    pub regions: Vec<ConstrainedLayoutRegion>,

    /// Horizontal spring slots across all regions.
    pub horizontal_slots: Vec<SpringSlot>,

    /// Vertical bands.
    pub vertical_bands: Vec<VerticalBand>,

    /// Glyph-level layout objects.
    pub glyphs: Vec<GlyphObject>,

    /// Non-glyph line primitives (staff lines, stems, barlines,
    /// ledger lines, brackets) presented to the solver alongside
    /// the glyphs.
    pub strokes: Vec<Stroke>,

    /// Cubic-Bezier curve primitives (slurs), likewise.
    pub curves: Vec<Curve>,

    /// Additional constraints not captured by spring parameters.
    pub constraints: Vec<LayoutConstraint>,

    /// Which user break override each projected break constraint came
    /// from, so a solver can cite it in the decision it records.
    pub break_origins: Vec<BreakOrigin>,

    /// Engraving decisions, carried forward from the logical stage.
    pub engraving_decisions: Vec<EngravingDecision>,

    /// Engraving-coverage gaps, surfaced rather than hidden.
    pub diagnostics: Vec<LayoutDiagnostic>,

    /// The glyph catalog this IR's metrics were drawn from.
    pub catalog: GlyphCatalogIdentity,
}

/// A projected break constraint's user-override attribution.
pub struct BreakOrigin {
    pub slot: SpringSlotId,
    pub class: BreakClass,
    pub override_id: EngravingOverrideId,
}

/// An object the projection could not engrave faithfully.
pub struct LayoutDiagnostic {
    pub source: TypedObjectId,
    pub kind: LayoutDiagnosticKind,
}

pub enum LayoutDiagnosticKind {
    /// A pitch reached the constrained pass with no resolved (or
    /// non-CMN) spelling; its notehead is placed on the clef
    /// reference line, but its true staff position is unknown.
    MissingSpelling,

    /// A glyph the bound catalog does not carry; the object is
    /// carried as a traced anchor rather than drawn at a guessed
    /// shape.
    UnbundledGlyph(GlyphReference),
}

pub enum LayoutConstraint {
    /// Two glyphs must not overlap in 2D space.
    NoCollision { a: GlyphObjectId, b: GlyphObjectId },

    /// A glyph must align with another along an axis.
    Align {
        a: GlyphObjectId,
        b: GlyphObjectId,
        axis: Axis,
    },

    /// A glyph must be positioned within a region of space.
    PositionWithin {
        glyph: GlyphObjectId,
        region: Rect,
    },

    /// A range of slots must end at a system boundary.
    SystemBreakAt { slot: SpringSlotId, kind: BreakKind },

    /// A range of slots must end at a page boundary.
    PageBreakAt { slot: SpringSlotId, kind: BreakKind },

    /// Custom constraint declared by a layout extension.
    Registered(ConstraintRegistryId, ConstraintParameters),
}

pub enum BreakKind {
    Hard,
    Soft,
}
\end{lstlisting}

\begin{requirement}
  \label{req:layoutir:break-satisfaction}
  \textbf{Break-constraint satisfaction.} A \texttt{SystemBreakAt} at
  \texttt{slot} is \emph{satisfied} iff the final resolved layout starts a
  system at that slot; a \texttt{PageBreakAt} iff the final layout starts a page
  there. A slot that already begins a system (respectively page) --- in
  particular a region's first slot --- is trivially satisfied. Satisfaction is a
  predicate on the output \texttt{ResolvedLayoutIR}, evaluable by any consumer,
  not on the solver's internal spring state. A solver's tier claim
  (Section~\ref{sec:layoutir:resolved}) reports which \texttt{Hard} break
  constraints its layout satisfies; a \texttt{Soft} break is advisory and MAY be
  left unsatisfied.
\end{requirement}

\section{ResolvedLayoutIR}
\label{sec:layoutir:resolved}

\texttt{ResolvedLayoutIR} is the output of the constraint solver:
every glyph, stroke, and curve has a definitive position. This is the
IR consumed by the renderer.

\begin{lstlisting}[language=Rust]
pub struct ResolvedLayoutIR {
    pub source: ScoreVersion,

    /// Pages of the resolved score.
    pub pages: Vec<ResolvedPage>,

    /// Glyph-level objects with resolved positions.
    pub glyphs: Vec<ResolvedGlyph>,

    /// Non-glyph line primitives (staff lines, stems, barlines,
    /// volta brackets, ...), positioned alongside the glyphs.
    pub strokes: Vec<Stroke>,

    /// Cubic-Bezier curve primitives (slurs, ...), positioned
    /// alongside the glyphs.
    pub curves: Vec<Curve>,

    /// Engraving decisions, including any made by the constraint
    /// solver (e.g., soft system breaks placed by casting-off).
    pub engraving_decisions: Vec<EngravingDecision>,
}

pub struct ResolvedPage {
    pub provenance: Provenance,
    pub number: u32,
    pub size: Size2D,
    pub margins: Margins,
    pub systems: Vec<ResolvedSystem>,
    pub free_objects: Vec<GlyphObjectId>,
}

pub struct ResolvedSystem {
    pub provenance: Provenance,
    pub bounding_box: Rect,
    pub staves: Vec<ResolvedStaff>,
    pub measures: Vec<ResolvedMeasure>,
}

pub struct ResolvedGlyph {
    pub provenance: Provenance,
    pub glyph: GlyphReference,
    pub position: Point2D,
    pub transform: Option<Transform2D>,
    pub bounding_box: BoundingBox,
    pub style: GlyphStyle,
    pub layer: i32,
}

/// A line primitive the solver positions (a staff line, stem,
/// barline, or bracket).
pub struct Stroke {
    pub provenance: Provenance,
    pub from: Point2D,
    pub to: Point2D,
    pub thickness: StaffSpace,
    pub layer: i32,
    pub style: GlyphStyle,

    /// The vertical band this stroke belongs to: its owning staff,
    /// declared by the projection that emitted it. A vertical solver
    /// reads ownership from here rather than inferring it from
    /// geometry. Unlike a glyph, a stroke is not listed in
    /// `VerticalBand::members`; this is a one-way reference.
    pub vertical_band: VerticalBandId,
}

/// A cubic-Bezier curve primitive (a slur), its four control
/// points in `p0`..`p3` drawing order.
pub struct Curve {
    pub provenance: Provenance,
    pub p0: Point2D,
    pub p1: Point2D,
    pub p2: Point2D,
    pub p3: Point2D,
    pub thickness: StaffSpace,
    pub layer: i32,
    pub style: GlyphStyle,

    /// The line pattern the renderer strokes the path with
    /// (a slur's authored `SpanStyle.line`): solid, dashed, or dotted.
    pub line: LineStyle,

    /// The vertical band this curve belongs to (see `Stroke`). A
    /// slur's endpoints are lifted clear of its own staff, so its
    /// owner is not recoverable from its geometry.
    pub vertical_band: VerticalBandId,
}
\end{lstlisting}

\begin{requirement}
  \label{req:layoutir:resolved-primitives}
  \textbf{Non-glyph resolved primitives (ratified schema-major-2,
  Phase~F).} Beyond glyphs, \texttt{ResolvedLayoutIR} carries
  non-glyph \emph{line-stroke} primitives (staff lines, stems,
  barlines, brackets) and \emph{cubic-B\'ezier curve} primitives
  (slurs). Each \MUST{} carry provenance --- the basis of hit-testing
  (\ref{sec:layoutir:provenance}) --- and is positioned by the solver
  like a glyph: a spanning stroke or curve is re-spaced by the same
  horizontal coordinate map its endpoint columns move under. These
  primitives are \emph{non-canonical}: like layout identifiers and
  glyph positions, they are not part of document state and enter no
  content hash (\ref{req:layoutir:object-id-derivation}). A conformant
  renderer draws them alongside the glyphs; the renderer encoding
  itself is out of scope (\ref{sec:layoutir:render}).
\end{requirement}

\begin{requirement}
  \label{req:layoutir:resolved-band-ownership}
  \textbf{Band ownership through the resolved stage.} A resolved
  \texttt{Stroke} and \texttt{Curve} \MUST{} retain the
  \texttt{vertical\_band} they were given in
  \texttt{ConstrainedLayoutIR}
  (\ref{req:layoutir:primitive-band-ownership}), so a solver that
  re-spaces or relocates them --- casting off into systems, or
  renegotiating the gaps between a system's staves --- can attribute
  each one to its owning staff. A \texttt{ResolvedGlyph} carries no
  band: a glyph's ownership is consumed during the solve and its
  outcome is already baked into the glyph's resolved position, so the
  field would be dead weight in a rendering fingerprint.

  Band ownership is \emph{non-canonical layout-attribution metadata}.
  It enters no content hash and no canonical encoding
  (\ref{req:layoutir:object-id-derivation}): it draws nothing, so two
  layouts differing only in it are the same rendered layout. A
  conformant canonical encoding of \texttt{ResolvedLayoutIR} therefore
  omits \texttt{vertical\_band} even from the strokes and curves that
  carry it.
\end{requirement}

\begin{requirement}
  \label{req:layoutir:repeat-render}
  \textbf{Minimal-tier repeat and volta rendering (ratified
  schema-major-2, Phase~F).} A \texttt{RepeatStructure}
  (\ref{sec:graph:repeats}) whose kind draws barlines
  (\texttt{SimpleRepeat}, \texttt{Volta}) \MUST{} render a repeat
  barline at each of its resolved boundaries --- the precomposed
  repeat sign, either replacing a coinciding measure barline or
  standing at the boundary column when none coincides, and the
  repeat-dot pair beside a region-closing final barline (which is
  never replaced). Each \texttt{Volta} \MUST{} render a bracket
  spanning its resolved extent, with its ending numbers. A boundary
  that does not resolve to a spacing column draws \emph{no} ink: a
  traced anchor preserves the structure's provenance rather than
  placing a mark at a false position. The jump kinds
  (\texttt{DaCapo}, \texttt{DalSegno}) draw no Minimal-tier marks ---
  their segno / coda / instruction text awaits a text primitive ---
  and a repeat spanning more than one region renders in no single
  region at this tier. The precise glyph vocabulary and spacing are
  engraving-algorithm concerns (\ref{sec:layoutir:forward}).
\end{requirement}

\begin{requirement}
  \label{req:layoutir:slur-curve}
  \textbf{Slur rendering (ratified schema-major-2, Phase~F; extended
  Push~3).} A \texttt{Slur} (\ref{sec:graph:slurs}) \MUST{} render
  as a cubic-B\'ezier \texttt{Curve} arcing between its two endpoint
  event columns, honoring an authored \texttt{CurvatureOverride}
  (direction and apex height) when present and choosing them
  otherwise. A slur whose endpoint does not resolve to a column on a
  single staff of one region draws \emph{no} curve: a traced anchor
  preserves its provenance rather than a floating arc. An authored
  non-\texttt{Solid} \texttt{SpanStyle} line (dashed, dotted) is
  rendered faithfully as that pattern (its \texttt{LineStyle} rides
  the \texttt{Curve}); an implementation that instead defers the
  pattern \MUST{} surface the deferral (a layout diagnostic), never
  silently rendering solid. A slur whose span crosses a system break
  \MUST{} split into per-system sub-curves --- the first carrying the
  slur's provenance, the rest engraver-synthesized continuations
  (\ref{req:layoutir:continuation-synthesis}) --- rather than drawing
  whole in one system. The curvature-computing algorithm remains an
  engraving-specification concern (\ref{sec:layoutir:forward}).
\end{requirement}

\section{RenderIR (Interface Only)}
\label{sec:layoutir:render}

The \texttt{RenderIR} stage is the renderer's input. Its full
specification is delivered separately; this section defines only the
interface contract.

\begin{lstlisting}[language=Rust]
pub trait RenderIRProducer {
    /// Convert resolved IR to renderer-bound primitives.
    fn produce(
        &self,
        resolved: &ResolvedLayoutIR,
        scale: ScaleContext,
        config: RenderConfiguration,
    ) -> RenderIR;
}

pub struct RenderConfiguration {
    /// Target output: PDF, SVG, on-screen, print.
    pub target: RenderTarget,

    /// Color space and color management options.
    pub color: ColorConfiguration,

    /// Anti-aliasing and rasterization options.
    pub rasterization: RasterizationConfiguration,
}
\end{lstlisting}

\begin{requirement}
  \label{req:layoutir:render-provenance}
  Implementations producing \texttt{RenderIR} \MUST{} preserve
  provenance from \texttt{ResolvedLayoutIR}: every renderer primitive
  \MUST{} be traceable to its originating \texttt{ResolvedLayoutIR}
  primitive --- a \texttt{ResolvedGlyph}, \texttt{Stroke}, or
  \texttt{Curve} --- and therefore to its score graph source. This is
  the basis of hit-testing, selection, and back-reference navigation
  in the UI.

  The full RenderIR type and its production rules are specified in
  the renderer specification, outside the scope of this document.
\end{requirement}

\section{Incremental Layout and Caching}
\label{sec:layoutir:incremental}

The pipeline supports incremental layout via per-object dependency
tracking. An edit to a score graph object invalidates only the IR
objects whose provenance lists it as a dependency.

\subsection{The Layout Cache}

\begin{lstlisting}[language=Rust]
pub struct LayoutCache {
    /// Reverse index: maps score graph object ids to the layout
    /// objects depending on them.
    pub dependencies: DependencyIndex,

    /// Cached IR stages, partitioned for granular invalidation.
    pub logical: HashMap<RegionId, LogicalRegionCache>,
    pub constrained: HashMap<RegionId, ConstrainedRegionCache>,
    pub resolved: HashMap<SystemId, ResolvedSystemCache>,

    /// Cached layout-stage outputs at finer granularity for
    /// performance-critical paths.
    pub fine_cache: FineLayoutCache,
}

pub struct DependencyIndex {
    /// For each score graph object, the layout objects depending
    /// on it. Reverse-indexed to make invalidation O(k) where k is
    /// the number of dependents.
    pub forward: HashMap<TypedObjectId, Vec<LayoutObjectId>>,

    /// For each layout object, its dependencies. Used to verify and
    /// rebuild the reverse index on cache writes.
    pub reverse: HashMap<LayoutObjectId, Vec<TypedObjectId>>,
}
\end{lstlisting}

\subsection{Invalidation Rules}

\begin{requirement}
  \label{req:layoutir:dependency-invalidation}
  When a semantic operation (Chapter~\ref{ch:semops}) is applied to
  the score graph, the following invalidation \MUST{} occur:

  \begin{enumerate}
    \item Every layout object whose \texttt{Provenance} lists a
      mutated score graph object in its \texttt{source} or
      \texttt{dependencies} is invalidated.
    \item Invalidation propagates: layout objects that depend on
      invalidated layout objects (via cross-region references,
      spanning structures, etc.) are also invalidated.
    \item Invalidated objects are cleared from the cache at their
      respective stages.
    \item Invalidated systems trigger casting-off re-evaluation for
      themselves and any downstream systems whose horizontal extent
      may shift.
  \end{enumerate}

  Re-engraving runs from the first invalidated stage forward, using
  cached values for non-invalidated objects.
\end{requirement}

\subsection{Frame-Budget Requirements}

\begin{requirement}
  \label{req:layoutir:incremental-frame-budget}
  The incremental layout system \MUST{} support the frame-budget
  requirements established in Chapter~\ref{ch:perf}: an edit
  affecting a single system on a 100-page orchestral score \MUST{}
  re-layout in under 16.7\,ms on the reference hardware profile.

  Achieving this requires (a) fine-grained dependency tracking as
  specified above; (b) cached glyph metrics in the font catalog;
  (c) parallel re-engraving where the affected systems are
  independent; (d) the constraint solver supporting partial
  re-resolution of affected slots only (see Chapter~\ref{ch:solver}).
\end{requirement}

\section{Region Uniformity}
\label{sec:layoutir:region-uniformity}

Staff-based, free-graphic, and hybrid regions produce structurally
identical \texttt{LayoutRegion} containers. The differences are
expressed in:

\begin{itemize}
  \item The \texttt{time\_axis}: metric, proportional, or aleatoric.
  \item The \texttt{objects} mix: staff-based regions contain
    primarily notation objects; free-graphic regions contain primarily
    \texttt{Graphic} objects; hybrid regions contain both.
\end{itemize}

\begin{requirement}
  \label{req:layoutir:uniform-region-solving}
  All three region kinds \MUST{} use the same \texttt{LayoutRegion}
  container type. The constraint solver \MUST{} treat regions
  uniformly: it consumes spring slots and constraints without
  branching on region kind.
\end{requirement}

\begin{rationale}
  Uniform containers permit the constraint solver to be agnostic to
  region kind, simplifying its implementation. The differences
  between region kinds are encoded in the \texttt{TimeAxis}
  implementation and the object mix, not in container structure.
  This is the same insight that makes the file format and semantic
  operations layers tractable: structural uniformity at the
  containers, polymorphism in the contents.
\end{rationale}

\section{Glyph Catalog Interface}
\label{sec:layoutir:catalog}

The IR references glyphs by identifier; glyph metrics are queried
from a font catalog at IR construction. This section defines the
interface; the catalog implementation is outside the core
specification. \texttt{SmuflVersion}, below, is the same type declared
by a score's \texttt{SmuflVersionRequirement}
(Section~\ref{sec:tuning:smufl}): the glyph catalog's version and the
tuning context's declared version are one type, not two.

\begin{lstlisting}[language=Rust]
pub trait GlyphCatalog: Send + Sync {
    /// Resolve a glyph reference to its metrics.
    fn metrics(&self, glyph: &GlyphReference) -> Option<GlyphMetrics>;

    /// Resolve a glyph reference to its rendering data.
    fn render_data(&self, glyph: &GlyphReference) -> Option<GlyphRenderData>;

    /// SMuFL version supported.
    fn smufl_version(&self) -> SmuflVersion;
}

pub struct GlyphMetrics {
    pub bounding_box: BoundingBox,
    pub advance_width: StaffSpace,
    pub anchors: HashMap<AnchorName, Point2D>,
}

pub struct GlyphRenderData {
    /// Outline data: path commands suitable for vector rendering.
    pub outline: Vec<PathCommand>,

    /// Optional pre-rendered bitmap for raster targets.
    pub bitmap: Option<GlyphBitmap>,
}
\end{lstlisting}

\begin{requirement}
  \label{req:layoutir:glyph-catalog-lookups}
  Implementations \MUST{} use a glyph catalog for metric and render-data
  lookups during IR construction. Metrics \MUSTNOT{} be duplicated in
  the IR; the IR carries glyph identifiers and queries the catalog.
  Implementations \MAY{} cache metrics within IR objects as a
  performance optimization provided cached values are invalidated when
  the underlying catalog changes.
\end{requirement}

\subsection{Glyph Catalog Identity for Layout Conformance}
\label{sec:layoutir:catalog-identity}

Layout output depends on glyph metrics. For reproducible layout,
the exact glyph catalog used \MUST{} be identifiable. \texttt{smufl\_version}
below is a \texttt{SmuflVersion} (Section~\ref{sec:tuning:smufl}), ordered
by its fraction-normalized \texttt{minor\_centi} so a catalog identity's
version compares correctly against SMuFL's real release history.

\begin{lstlisting}[language=Rust]
/// A reproducibility-quality identifier for the glyph catalog
/// used to produce a layout. Required for any layout conformance
/// claim that depends on byte-equal output across runs.
pub struct GlyphCatalogIdentity {
    /// SMuFL version targeted.
    pub smufl_version: SmuflVersion,

    /// Identifier of the specific font in use (e.g., Bravura,
    /// Petaluma, Leland).
    pub font_id: FontId,

    /// Optional semantic version of the font, if the font
    /// publisher uses versioned releases.
    pub font_version: Option<SemVer>,

    /// Content hash of the metric data made available to the
    /// layout solve. The hash is over a canonical serialization
    /// of every glyph's metrics (bounding box, advance width,
    /// named anchors) for every glyph referenced by the
    /// ConstrainedLayoutIR delivered to the solver, including
    /// glyphs that are candidates but do not appear in the final
    /// ResolvedLayoutIR. Computed by BLAKE3 with the domain
    /// tag "MUSCFNTM".
    pub metrics_hash: ContentHash,
}
\end{lstlisting}

\begin{requirement}
  \label{req:layoutir:glyph-catalog-identity}
  Any layout conformance claim subject to byte-equal output
  (within-implementation determinism per
  Chapter~\ref{ch:solver}) \MUST{} declare the
  \texttt{GlyphCatalogIdentity} under which it was produced.
  Two layouts produced with different \texttt{metrics\_hash}
  values are layouts under different inputs; byte equality is
  not expected.

  Implementations that cannot identify their glyph catalog at
  this granularity \MUST{} mark their layout output as
  \emph{implementation-specific} or \emph{non-canonical};
  conformance suites \MUSTNOT{} accept such output as
  byte-equivalent to a reference solver's output.

  The \texttt{metrics\_hash} is computed over the glyphs that
  are \emph{inputs} to the layout solve (those referenced by the
  \texttt{ConstrainedLayoutIR}), not only over glyphs in the
  final layout. Two solver invocations with different available
  candidate-glyph metrics are operating on different inputs even
  if they happen to converge on the same final glyph set; this
  rule makes that difference visible in the
  \texttt{metrics\_hash}.

  The hash is computed only over glyphs the layout solve
  actually consulted, not over the entire font catalog, so that
  font additions (new glyphs unused by the solve) do not
  invalidate prior layouts.
\end{requirement}

\section{Forward References}
\label{sec:layoutir:forward}

\begin{itemize}
  \item The constraint solver consuming \texttt{ConstrainedLayoutIR}
    and producing \texttt{ResolvedLayoutIR} is specified in
    Chapter~\ref{ch:solver}.
  \item The serialization of layout cache data for cross-session
    persistence is specified in Chapter~\ref{ch:format}.
  \item The complete RenderIR specification, the engraving algorithm
    specifications (beam grouping, slur curvature, casting-off,
    accidental ordering), the font catalog implementation, and the
    renderer are delivered as separate specifications layered on top
    of this core document.
\end{itemize}

% ===========================================================================
\chapter{File Format}
\label{ch:format}

This chapter specifies the on-disk file format. Chapter~\ref{ch:semops}
established that the canonical document state is the deterministic
reduction of an operation-envelope set. This chapter specifies how
that operation set, along with acceleration structures, profile
declarations, and extension declarations, is laid out on disk in a
way that is crash-recoverable, atomically updatable, and forward-
compatible.

The byte-level encoding details (record layouts, varint conventions,
exact field bit widths) are deferred to a companion document, the
Binary Format specification. This chapter specifies \emph{what} is
stored, how the parts relate, and how the format behaves under
concurrent edits, crashes, and forward-compatibility scenarios; the
Binary Format specifies \emph{how the bits are laid out}.

The file format's name is \texttt{musc} and its file extension is
\texttt{.musc}.

\section{Design Principles}
\label{sec:format:principles}

\begin{description}
  \item[Canonical state is the operation set.] The canonical
    document is the operation-envelope set
    (Chapter~\ref{ch:semops}). After pruning, the canonical document
    is a retained base snapshot plus the operation envelopes after
    that snapshot's causal frontier. Everything else stored in the
    bundle (snapshots, layout caches, indexes, rendered artifacts)
    is an acceleration structure and \MUST{} be discardable.
  \item[Fixed prelude, atomic superblock commit.] The file begins
    with a fixed-size header followed by two fixed-size superblock
    slots. The superblocks are the only mutable on-disk objects in
    the bundle. Commits append new chunks, write a new manifest
    chunk, then flip the active superblock by writing to the
    currently-inactive slot. This makes commits atomic with respect
    to crashes.
  \item[Content-addressed immutable chunks.] Every other on-disk
    object is content-addressed by BLAKE3 of its uncompressed
    payload (with domain-separated preimage). Chunks are immutable
    once written. Duplicate content shares storage automatically.
  \item[Manifest as a table of roots.] The manifest does not carry
    user-facing data. It carries the operation-set roots, the
    retained snapshot references, blob references, profile
    declarations, extension declarations, and the text-projection
    root. User-facing metadata (title, composer, etc.) lives in
    operation envelopes or in materialized state, not in the
    physical header.
  \item[Crash recovery is by superblock selection.] Recovery is
    simply: read both superblocks, validate, and select the highest
    valid generation. Anything not reachable from the selected
    superblock is unreachable garbage and can be reclaimed by a
    later conservative GC pass.
  \item[Forward compatibility through opaque preservation and edit
    barriers.] Unknown extension chunks are preserved across reads
    and writes. Unknown required extensions either prevent open or
    force read-only mode. Unknown optional extensions are preserved
    unless an edit crosses a declared edit barrier.
  \item[Streaming reads.] The bundle is designed so that opening a
    large score touches the prelude, the active superblock, the
    manifest, and a small set of bootstrap chunks (the most recent
    canonical-base snapshot plus the envelope index). Per-region
    content loads on demand.
  \item[Deterministic text projection of semantic content.] The
    format admits a deterministic projection to a canonical
    s-expression form for diff tooling and version control. The
    text projection preserves canonical document semantics
    (operation envelopes, profile declarations, extension
    declarations, reduced state). It does \emph{not} preserve
    chunk offsets, compression choices, cache chunks, garbage
    regions, or superblock generation numbers.
  \item[No encryption, no DRM.] The format \MUSTNOT{} include
    encryption, access control, or rights-management features.
    Privacy at rest is the responsibility of the filesystem and
    operating system. The format is open and offers no mechanism by
    which user content can be locked away from its author or from
    other implementations.
\end{description}

\begin{nongoal}
  Digital rights management (DRM), content encryption,
  hardware-binding, license enforcement, and any mechanism that
  could lock user content to a particular vendor or platform are
  explicitly out of scope. The format is designed for openness,
  durability, and user ownership of data. Implementations \MUSTNOT{}
  add proprietary extensions intended to restrict access to or use
  of user content.
\end{nongoal}

\section{Durable Writes}
\label{sec:format:durable}

This chapter refers throughout to ``durable flushes'' of file
contents. The definition is platform-abstracted.

\begin{description}
  \item[Durable flush.] The platform operation that makes
    previously-written bytes durable according to the host
    filesystem contract: \texttt{fsync} on POSIX-like systems,
    \texttt{FlushFileBuffers} or \texttt{NtFlushBuffersFileEx} on
    Windows, equivalent calls on other platforms. A successful
    durable flush guarantees that bytes written before the call are
    on the persistent storage medium when the call returns.
\end{description}

\begin{requirement}
  \label{req:format:durable-flush}
  Implementations \MUST{} use the platform's durable flush
  primitive at the points required by the atomic write protocol
  (Section~\ref{sec:format:commit}). Implementations \MAY NOT{}
  rely on filesystem ordering guarantees or write-back caching
  alone; durability \MUST{} be requested explicitly at the commit
  point.
\end{requirement}

\section{The Bundle Layout}
\label{sec:format:bundle}

A \texttt{.musc} bundle is a single file whose first 576 bytes are a
fixed prelude: a 64-byte header followed by two 256-byte superblock
slots. The remainder of the file holds variable-length content:
chunk payloads, the manifest chunk, blob payloads, the text
projection (if present), and unreachable garbage from prior commits
that has not yet been collected.

\begin{lstlisting}[language=Rust]
pub struct BundleLayout {
    /// Fixed 64-byte header at offset 0.
    pub header: FixedHeader,

    /// Two fixed-size superblock slots at offsets 64 and 320.
    pub superblock_a: Superblock,  // at offset 64
    pub superblock_b: Superblock,  // at offset 320

    /// Variable-length content begins at offset 576.
    /// All chunks, the manifest, and blobs live here.
    pub body: ContentRegion,
}
\end{lstlisting}

\subsection{The Fixed Header}

The header is exactly 64 bytes at file offset zero. It identifies
the format and locates the superblock slots. The header never
changes after the file is created.

\begin{lstlisting}[language=Rust]
pub struct FixedHeader {
    /// Magic bytes: ASCII "MUSCBND\0" (8 bytes).
    pub magic: [u8; 8],

    /// Major format version. Major version changes are
    /// non-backward-compatible.
    pub format_major: u16,

    /// Minor format version. Minor version changes are
    /// backward-compatible.
    pub format_minor: u16,

    /// Length of this header in bytes (currently always 64).
    /// Reserved for future expansion via header_length growth.
    pub header_length: u32,

    /// Offset of superblock slot A. Currently always 64.
    pub superblock_a_offset: u64,

    /// Offset of superblock slot B. Currently always 320.
    pub superblock_b_offset: u64,

    /// 128-bit file UUID, set at creation of the physical bundle.
    /// Changes on Save As (the physical bundle is a new file).
    /// Distinct from document_id, which identifies the logical work.
    pub file_uuid: FileUuid,

    /// Reserved bytes at offsets 48..60. MUST be zero on write;
    /// MUST be ignored on read. Future minor versions MAY use
    /// these bytes for additional fields whose absence is
    /// interpreted as zero.
    pub reserved: [u8; 12],

    /// CRC-32C of bytes 0..60 of this header (i.e., all preceding
    /// fields plus the reserved bytes). At fixed offset 60..64.
    pub header_crc: u32,
}
\end{lstlisting}

\begin{requirement}
  \label{req:format:fixed-header-integrity}
  The header \MUST{} be exactly 64 bytes at file offset zero. Its
  contents \MUSTNOT{} change after file creation, except for the
  case where a major-version upgrade rewrites the entire file.
  Readers \MUST{} verify the magic bytes and the header CRC before
  consulting any other part of the file.

  Reserved bytes \MUST{} be zero when written by current
  implementations. Readers \MUST{} ignore the values of reserved
  bytes (future versions may use them).

  The \texttt{file\_uuid} is fixed for the lifetime of a physical
  bundle and persists across commits to that bundle. A
  filesystem-level byte copy preserves the \texttt{file\_uuid}
  (the copy is the same physical bundle duplicated). Save As,
  export-as-new-bundle, and derivative-work creation \MUST{}
  assign a fresh \texttt{file\_uuid}; the resulting bundle is a
  new physical bundle, distinct from its source.

  The \texttt{file\_uuid} identifies the physical bundle, not the
  logical work. Two bundles with the same \texttt{file\_uuid} are
  byte-level copies of the same physical bundle; two bundles
  with the same \texttt{document\_id} (defined in the manifest;
  see Section~\ref{sec:format:identity}) are versions of the
  same logical work, which may or may not share their
  \texttt{file\_uuid}.
\end{requirement}

\subsection{The Superblock Slots}

Each superblock slot is exactly 256 bytes. The bundle has two slots
(A at offset 64, B at offset 320). One is active; the other is the
write target for the next commit. Atomic commit is the act of
writing the inactive slot and durably flushing it.

\begin{lstlisting}[language=Rust]
pub struct Superblock {
    /// Magic bytes: ASCII "MUSCSUPR" (8 bytes).
    pub magic: [u8; 8],

    /// Generation counter. The active superblock is the one with
    /// the highest valid generation.
    pub generation: u64,

    /// Offset of the manifest chunk for this generation.
    pub manifest_offset: u64,

    /// Length of the manifest chunk in bytes.
    pub manifest_length: u64,

    /// BLAKE3 content hash of the manifest chunk's uncompressed
    /// payload, with domain-separated preimage
    /// (Section~\ref{sec:format:hashing}).
    pub manifest_hash: ContentHash,

    /// Schema version of the manifest at this generation.
    pub manifest_schema_version: SchemaVersion,

    /// Reduction algorithm version that produced any
    /// canonical-base snapshots in this manifest.
    pub reduction_algorithm_version: ReductionAlgorithmVersion,

    /// Profile under which this superblock is valid.
    pub profile_id: ProfileId,

    /// Commit state. Distinguishes a fully-committed superblock
    /// from one that was written but whose commit had not yet
    /// completed when the writer crashed (rare; the protocol does
    /// not normally leave such a state visible).
    pub commit_state: CommitState,

    /// Wall-clock timestamp at commit, in canonical units.
    /// Advisory only; superblock selection is by generation, not
    /// by timestamp.
    pub commit_timestamp: WallClockTime,

    /// CRC-32C of the bytes of this superblock up to (but not
    /// including) the CRC field. Lets readers reject torn writes.
    pub superblock_crc: u32,

    /// Reserved padding to 256 bytes. MUST be zero on write;
    /// MUST be ignored on read.
    pub reserved: [u8; /* computed */ 256 - HEADER_FIELDS_SIZE],
}

pub enum CommitState {
    /// Fully committed. The manifest at manifest_offset is the
    /// canonical state for this generation.
    Committed,

    /// Reserved for future use; writers MUST NOT produce
    /// non-Committed values in this version of the format.
    Reserved(u32),
}
\end{lstlisting}

\subsection{Superblock Selection}

\begin{requirement}
  \label{req:format:superblock-selection}
  On open, readers \MUST{} perform the following selection:

  \begin{enumerate}
    \item Read both superblock slots.
    \item For each slot, validate: magic bytes, CRC, that the
      referenced manifest chunk's actual BLAKE3 hash matches the
      declared \texttt{manifest\_hash}, and that
      \texttt{commit\_state} is \texttt{Committed}. A slot
      failing any of these checks is \emph{not valid} for
      ordinary selection.
    \item A slot whose \texttt{commit\_state} is not
      \texttt{Committed} is invalid for ordinary selection but
      \MAY{} be inspected in diagnostic recovery mode. Writers
      in this format version \MUSTNOT{} produce such states
      under normal commit; observing one indicates either a
      crashed writer (recovery falls back to the other slot per
      the atomic-write protocol) or a non-conforming writer.
    \item If both slots are valid and their generations differ by
      more than one, the reader \MUST{} report an integrity anomaly
      (the file may have been tampered with, accidentally merged,
      or written by a non-conforming implementation). The reader
      \MAY{} continue in read-only recovery mode using the
      highest-generation valid slot; it \MUSTNOT{} silently treat
      the file as normal.
    \item If both slots are valid and their generations differ by
      exactly one, the slot with the higher generation is active.
    \item If both slots are valid and their generations are
      \emph{equal}, the reader \MUST{} compare the slots'
      \emph{load-bearing} selection fields:
      \texttt{manifest\_hash}, \texttt{manifest\_schema\_version},
      \texttt{reduction\_algorithm\_version}, and
      \texttt{profile\_id}. The advisory fields
      \texttt{commit\_timestamp} and the physical chunk
      offset/length are \emph{excluded} from this comparison.
      \begin{itemize}
        \item If every load-bearing field is equal, the two slots
          select equivalent canonical state; the reader \MUST{}
          deterministically choose slot~A and open normally.
        \item If any load-bearing field differs, the reader \MUST{}
          report a format-layer \texttt{IntegrityAnomaly} of kind
          \texttt{DivergentSameGeneration} and open the bundle
          read-only. Two genuinely different committed states cannot
          share a generation under a conforming writer (the
          atomic-write protocol increments the generation on every
          commit), so the divergence is a structural failure, not a
          state the reader may silently resolve by picking one.
      \end{itemize}
    \item If exactly one slot is valid, that slot is active.
    \item If neither slot is valid, the file is corrupt; readers
      \MUST{} surface this as a hard error and \MUSTNOT{}
      synthesize state.
  \end{enumerate}
\end{requirement}

\begin{rationale}
  The advisory fields are excluded from the equal-generation
  comparison because they do not determine which canonical state a
  slot selects: \texttt{commit\_timestamp} is wall-clock metadata,
  and the physical offset/length merely locate the manifest chunk in
  the file. Two slots that agree on the four load-bearing fields name
  the same canonical document even if they were written at different
  times or to different offsets, so treating an advisory-only
  difference as divergence would raise false anomalies on benign
  repacks. Conversely, a difference in any load-bearing field means
  the two slots disagree about canonical state itself, which is the
  precise condition that must not be silently resolved.
\end{rationale}

\section{The Manifest}
\label{sec:format:manifest}

The manifest is a table of roots and declarations. It is itself a
chunk (content-addressed, immutable). Each commit writes a new
manifest chunk and points the inactive superblock slot at it.

\begin{lstlisting}[language=Rust]
pub struct Manifest {
    /// Identifier of the logical work. Stable across Save As
    /// operations and across copies that are intended to remain
    /// versions of the same work. May persist across forks if
    /// the user explicitly preserves it.
    pub document_id: DocumentId,

    /// Optional identifier of a shared ancestor document. Used by
    /// version-control workflows and fork-tracking to express
    /// genealogy. Absent for documents with no declared ancestry.
    pub lineage_id: Option<LineageId>,

    /// Identifier of this manifest. Each commit produces a new
    /// ManifestId; old manifests are retained in the chunk store
    /// until garbage collection removes them subject to the
    /// retention policy.
    pub manifest_id: ManifestId,

    /// Generation counter matching the superblock that references
    /// this manifest.
    pub generation: u64,

    /// Operation-envelope blocks defining the canonical document.
    /// Block ordering is irrelevant for canonical state (envelopes
    /// are a set), and the canonical manifest encoding sorts chunk
    /// references ascending by their encoded form (Appendix D), so
    /// a writer cannot order this list by stamp; stamp-ordered
    /// scanning goes through operation_block_summaries below.
    pub operation_roots: Vec<ChunkRef>,

    /// Per-block summary metadata (Section on operation-envelope
    /// blocks below), keyed by the block's ChunkId in canonical
    /// ascending order. Supplied by the operation layer and carried
    /// opaquely by the storage layer; non-canonical and optional (a
    /// block need not have an entry). Lets a reader select or skip
    /// a block by causal frontier or stamp range without decoding
    /// its envelopes.
    pub operation_block_summaries:
        BTreeMap<ChunkId, OperationBlockSummary>,

    /// Optional operation index chunk: maps OperationId to
    /// (block, offset) for fast random access. If absent,
    /// implementations rebuild on demand by scanning blocks.
    pub operation_index_root: Option<ChunkRef>,

    /// The active canonical-base snapshot, if pruning has occurred.
    /// At most one canonical base is active at a time. When
    /// present, the canonical document is the base snapshot's
    /// materialized state plus the deterministic reduction of
    /// envelopes whose OperationIds are not covered by the base's
    /// causal frontier.
    pub canonical_base: Option<SnapshotRef>,

    /// Acceleration snapshots: caches that materialize state at
    /// various causal frontiers for fast incremental loading.
    /// These MAY be discarded freely without affecting canonical
    /// state.
    pub acceleration_snapshots: Vec<SnapshotRef>,

    /// Blob references for large opaque content (audio, images,
    /// custom fonts).
    pub blob_roots: Vec<BlobRef>,

    /// Profile declarations: which conformance profiles this
    /// manifest claims to satisfy.
    pub profile_declarations: Vec<ProfileDeclaration>,

    /// Extension declarations: which extensions contributed to
    /// this manifest, with their required/optional flag and
    /// edit barriers.
    pub extension_declarations: Vec<ExtensionDeclaration>,

    /// Text projection root chunk, if a text projection is
    /// maintained in this bundle. Non-canonical accelerator.
    pub text_projection_root: Option<ChunkRef>,

    /// Optional integrity index root chunk: cross-references chunk
    /// hashes for faster integrity verification of large bundles.
    /// Non-canonical accelerator; absence is permitted.
    pub integrity_root: Option<ChunkRef>,
}
\end{lstlisting}

\subsection{Canonical and Non-Canonical Manifest Roots}
\label{sec:format:roots}

The manifest's fields partition cleanly into two groups by their
relation to canonical state.

\textbf{Canonical roots} (define the document):

\begin{itemize}
  \item \texttt{operation\_roots}: operation-envelope blocks
    containing the canonical operation set.
  \item \texttt{canonical\_base}: the active canonical base
    snapshot, if pruning has occurred.
  \item \texttt{blob\_roots}: blobs referenced by canonical
    operations or by canonical reduced state are themselves
    canonical; blobs referenced only by acceleration structures
    are non-canonical.
  \item \texttt{document\_id}, \texttt{lineage\_id},
    \texttt{profile\_declarations},
    \texttt{extension\_declarations}: structural metadata
    governing how the canonical document is interpreted.
\end{itemize}

\textbf{Non-canonical reachable roots} (accelerators only):

\begin{itemize}
  \item \texttt{operation\_index\_root}: rebuildable from
    operation envelope blocks.
  \item \texttt{acceleration\_snapshots}: caches at causal
    frontiers other than the canonical base.
  \item \texttt{text\_projection\_root}: derivable from canonical
    state.
  \item \texttt{integrity\_root}: rebuildable from chunk hashes.
\end{itemize}

\begin{requirement}
  \label{req:format:manifest-reachability}
  The manifest \MUST{} list every chunk that contributes to the
  current canonical state, transitively through its canonical
  roots (operation roots, canonical base, canonical blobs and
  extension chunks). It \MAY{} additionally list non-canonical
  reachable roots (acceleration snapshots, operation index, text
  projection, integrity index, layout caches).

  Chunks not reachable from any manifest root, canonical or
  non-canonical, are unreachable; readers \MUSTNOT{} rely on
  them and conservative garbage collection \MAY{} reclaim them
  (Section~\ref{sec:format:gc}).

  Non-canonical reachable chunks \MAY{} be rebuilt or replaced
  without altering canonical state. If a non-canonical chunk's
  hash fails verification, it \MUST{} be discarded and (if
  needed) rebuilt; failed verification of a non-canonical chunk
  is \emph{not} a corruption of the bundle.

  Failed verification of a canonical chunk (operation envelope
  block, canonical base snapshot, canonical blob, manifest)
  \emph{is} corruption and \MUST{} be surfaced as a hard error.

  User-facing metadata (title, composer, lyricist, copyright, and
  similar fields) is \emph{not} part of the manifest. Such metadata
  lives in the operation set and the materialized state. The
  manifest carries only the structural roots and declarations that
  the format needs to locate and validate content.
\end{requirement}

\subsection{File, Document, and Lineage Identity}
\label{sec:format:identity}

The format distinguishes three identifiers, each serving a
different identity question:

\begin{description}
  \item[\texttt{file\_uuid}.] Physical bundle identity. Set at
    file creation; persists for the lifetime of the physical
    file. \textbf{Changes on Save As}: a copy created by Save As
    is a new physical bundle and \MUST{} be assigned a fresh
    \texttt{file\_uuid}. Used by tools that need to identify
    ``this exact file on disk'' (caches, backup software, file
    fingerprinting).
  \item[\texttt{document\_id}.] Logical work identity. Stable
    across Save As copies that are intended to remain the same
    work (the canonical case: ``I'm saving my piece to a backup
    folder, and it's still the same piece''). \textbf{Persists
    across forks at the user's discretion}: a fork that
    represents a derivative work \SHOULD{} have a new
    \texttt{document\_id}, while a fork that is a working copy
    of the original \SHOULD{} preserve it. The format does not
    enforce a policy; the editor surfaces the choice.
  \item[\texttt{lineage\_id}.] Optional shared ancestor identity.
    Records that two documents share a common ancestor for
    version-control or genealogy purposes. Multiple documents
    \MAY{} share a \texttt{lineage\_id}; the field is set when a
    fork explicitly preserves ancestry information.
\end{description}

\begin{requirement}
  \label{req:format:file-document-lineage-identity}
  Implementations \MUST{} assign a fresh \texttt{file\_uuid} on
  any Save As operation. They \MUSTNOT{} carry the source file's
  UUID into the new file: doing so would conflate physical and
  logical identity and break tools that rely on
  \texttt{file\_uuid} for content fingerprinting.

  The choice of whether to preserve \texttt{document\_id} across
  a Save As is exposed to the user: the editor \SHOULD{} default
  to preserving \texttt{document\_id} (the common case is
  ``save a copy of the same work to a new file'') and \SHOULD{}
  offer a clear ``Save As Derivative Work'' or equivalent action
  that mints a new \texttt{document\_id}.

  The format \MUSTNOT{} treat documents with the same
  \texttt{document\_id} as automatically synchronizable: that
  decision belongs to the collaboration transport
  (Appendix~\ref{app:deferred}).
\end{requirement}

\subsection{Manifest Encoding}
\label{sec:format:manifest-encoding}

The manifest chunk is special: it is the bootstrap entry from the
superblock into the chunk graph. A cold reader needs to decode the
manifest before it has any information from the chunk store.

\begin{requirement}
  \label{req:format:uncompressed-manifest}
  In this format version, the manifest chunk \MUST{} be stored
  \emph{uncompressed}. Its \texttt{CompressionAlgorithm} on disk
  is implicitly \texttt{None}; the superblock's
  \texttt{manifest\_length} field is therefore both the on-disk
  length and the uncompressed payload length. Implementations
  \MUST{} reject as malformed any bundle whose manifest payload
  is not directly parseable as a canonical manifest chunk's
  uncompressed bytes.

  This rule eliminates a bootstrap chicken-and-egg problem: a
  reader can decode the manifest using only information present
  in the fixed header and the active superblock, without first
  consulting the manifest to learn how the manifest is compressed.

  All other chunks (operation envelope blocks, snapshots,
  extension data, layout caches, blobs, integrity indexes,
  operation indexes, text projection chunks) \MAY{} be compressed.
  Future format major versions \MAY{} permit manifest compression,
  but doing so requires explicit superblock fields
  (\texttt{manifest\_compression},
  \texttt{manifest\_uncompressed\_length}) and is a
  non-backward-compatible change.
\end{requirement}

\begin{rationale}
  The manifest is small (a few kilobytes at most for typical
  scores), so compression yields negligible space savings.
  Eliminating the bootstrap dependency is far more valuable than
  the few hundred bytes of compression headroom.
\end{rationale}

\begin{requirement}
  \label{req:format:manifest-id}
  The \texttt{manifest\_id} field \MUST{} be derived as
  \[
    \texttt{ManifestId} = \mathrm{trunc128}\!\left(
      \mathrm{BLAKE3}(\texttt{"MUSCMNIF"} \Vert \texttt{document\_id}
      \Vert \texttt{generation} \Vert \texttt{manifest\_body})\right),
  \]
  where \texttt{document\_id} is its 16-byte canonical form,
  \texttt{generation} is the 8-byte little-endian \texttt{u64}, and
  \texttt{manifest\_body} is the canonical manifest encoding with the
  \texttt{manifest\_id} field itself \emph{excluded}. The exclusion is
  normative: a conforming writer \MUST{} omit (equivalently, zero) the
  \texttt{manifest\_id} field when computing the preimage, so that the
  identifier does not reference itself. $\mathrm{trunc128}$ takes the
  leading 16 bytes of the 32-byte BLAKE3 output. Two conforming
  writers committing the same manifest body at the same generation of
  the same document therefore derive identical \texttt{ManifestId}s.
\end{requirement}

\section{Content Hashing}
\label{sec:format:hashing}

Epiphany uses BLAKE3 as its single content-hashing algorithm.

\begin{requirement}
  \label{req:format:blake3-content-hashes}
  All content hashes in the bundle \MUST{} be BLAKE3-256 outputs
  (32 bytes). Implementations \MUSTNOT{} use any other algorithm
  for content addressing; SHA-256, SHA-3, and other algorithms
  \MUSTNOT{} appear in this version of the format. Future major
  versions \MAY{} introduce additional algorithms; such changes
  require a new format major version.
\end{requirement}

\subsection{Domain-Separated Preimages}

A naive content-addressing scheme that hashes only payload bytes is
vulnerable to collisions across chunks of different semantic kinds.
Epiphany uses domain-separated preimages.

\begin{lstlisting}[language=Rust]
/// Computes the canonical hash preimage for a chunk's content.
/// The chunk's compression algorithm is metadata on the ChunkRef,
/// not part of the content identity: identical uncompressed
/// payloads produce identical hashes regardless of how they are
/// physically compressed in storage. This preserves deduplication
/// across compression choices.
fn hash_preimage(
    domain_tag: &[u8; 8],
    chunk_kind: ChunkKind,
    schema_version: SchemaVersion,
    uncompressed_length: u64,
    uncompressed_payload: &[u8],
) -> Vec<u8> {
    let mut preimage = Vec::new();
    preimage.extend_from_slice(domain_tag);
    preimage.extend_from_slice(&chunk_kind.canonical_bytes());
    preimage.extend_from_slice(&schema_version.canonical_bytes());
    preimage.extend_from_slice(&uncompressed_length.to_le_bytes());
    preimage.extend_from_slice(uncompressed_payload);
    preimage
}
\end{lstlisting}

\begin{requirement}
  \label{req:format:domain-separated-chunk-hashes}
  The content hash of a chunk \MUST{} be BLAKE3 of the canonical
  preimage above. The domain tag, chunk kind, schema version, and
  uncompressed length are part of the preimage; the compression
  algorithm is \emph{not}. Two chunks with identical uncompressed
  payloads, identical kinds, and identical schema versions \MUST{}
  produce identical hashes regardless of their compression choices.
\end{requirement}

\begin{rationale}
  Including compression algorithm in the hash would cause
  identical content stored with different compression to have
  different addresses, defeating deduplication and forcing churn
  during recompression or repack. Including domain tag, kind, and
  schema version in the preimage closes a real collision: two
  semantically different chunks with identical raw bytes are no
  longer interchangeable.

  The domain tag is a fixed 8-byte string identifying the bundle
  format and protecting against cross-format collisions if the
  same BLAKE3 hash were reused in unrelated contexts. The canonical
  domain tag for \texttt{.musc} chunks is \texttt{"MUSCCHNK"};
  manifest chunks use \texttt{"MUSCMANI"}; blob payloads use
  \texttt{"MUSCBLOB"}.
\end{rationale}

\section{Chunks}
\label{sec:format:chunks}

A chunk is an immutable, content-addressed unit of storage. Chunks
are referenced from the manifest by \texttt{ChunkRef}.

\begin{lstlisting}[language=Rust]
pub struct ChunkRef {
    /// Content identifier: BLAKE3 of the canonical preimage.
    pub id: ChunkId,

    /// Kind of chunk. The reader uses this to dispatch parsing.
    pub kind: ChunkKind,

    /// Schema version targeted by this chunk's payload encoding.
    pub schema_version: SchemaVersion,

    /// Offset of the chunk's compressed payload in the bundle file.
    pub offset: u64,

    /// Length of the on-disk compressed payload.
    pub compressed_length: u64,

    /// Length of the uncompressed payload (used to size buffers
    /// before decompression).
    pub uncompressed_length: u64,

    /// Compression algorithm. Part of the ChunkRef, not part of
    /// the chunk's identity.
    pub compression: CompressionAlgorithm,

    /// Restated content hash, redundant with the id field but
    /// included in the ChunkRef for fast verification without
    /// fetching the chunk.
    pub hash: ContentHash,
}

pub enum ChunkKind {
    /// A block of operation envelopes.
    OperationEnvelopeBlock,

    /// An operation-id-to-block index accelerator.
    OperationIndex,

    /// A materialized snapshot of canonical state (full or
    /// partial). Canonical if referenced as the manifest's
    /// canonical_base; cache (acceleration) otherwise.
    Snapshot,

    /// A large opaque blob (audio, image, font, ML model).
    Blob,

    /// Extension-defined data, preserved opaquely by core readers.
    ExtensionData,

    /// The text projection's root document.
    TextProjection,

    /// Layout cache: derivative from canonical state, always
    /// discardable.
    LayoutCache,

    /// Integrity index: cross-references chunk hashes.
    IntegrityIndex,

    /// Manifest chunk. Special: referenced from the superblock,
    /// not from another chunk.
    Manifest,
}

pub enum CompressionAlgorithm {
    /// No compression. Stored payload bytes equal uncompressed.
    None,

    /// Zstandard compression. Implementations MUST support reading.
    Zstd { level: u8 },

    /// Other algorithms reserved for future versions.
    Reserved(u8),
}

pub struct ContentHash(pub [u8; 32]);

/// Chunk identifier. A newtype around ContentHash: the chunk's
/// identifier is identical to its content hash, but the type
/// distinction is preserved to make role visible at use sites
/// (a ChunkId is what you store in a ChunkRef and what you look
/// up; a ContentHash is what you compute by hashing). Both occupy
/// the same 32 bytes.
pub struct ChunkId(pub ContentHash);
\end{lstlisting}

\begin{requirement}
  \label{req:format:chunkkind-discriminants}
  \texttt{ChunkKind::canonical\_bytes()} is a single byte: the
  variant's declaration-order discriminant. Because that byte is part
  of the chunk hash preimage (Section~\ref{sec:format:chunks},
  ``Domain-Separated Preimages''), this assignment is normative and
  stable and \MUSTNOT{} be reordered:

  \begin{center}
  \begin{tabular}{r l @{\qquad} r l}
    \toprule
    \textbf{Disc.} & \textbf{Variant} &
    \textbf{Disc.} & \textbf{Variant} \\
    \midrule
    0 & \texttt{OperationEnvelopeBlock} & 5 & \texttt{TextProjection} \\
    1 & \texttt{OperationIndex}         & 6 & \texttt{LayoutCache} \\
    2 & \texttt{Snapshot}               & 7 & \texttt{IntegrityIndex} \\
    3 & \texttt{Blob}                   & 8 & \texttt{Manifest} \\
    4 & \texttt{ExtensionData}          &   & \\
    \bottomrule
  \end{tabular}
  \end{center}

  \texttt{CompressionAlgorithm} encodes as a fixed \emph{two} bytes: a
  declaration-order discriminant byte followed by a single parameter
  byte that is always present. \texttt{None} $= 0$, and its parameter
  byte \MUST{} be zero --- a reader \MUST{} reject a non-zero one rather
  than ignore it (Section~\ref{sec:format:binary}, and the Binary Format
  companion's \sectionsc{Chunk References}). \texttt{Zstd\{level\}}
  $= 1$ with the \texttt{level} in the parameter byte;
  \texttt{Reserved(u8)} $= 2$ with the reserved value in the parameter
  byte. The parameter byte is written even for \texttt{None}, so the
  encoding is fixed-width and a reader always consumes two bytes.
  Compression is \emph{not} part of a
  chunk's content identity (it is metadata on the \texttt{ChunkRef}),
  so these discriminants are stable but do not affect chunk hashes.

  \texttt{ProfileId}'s discriminants are pinned in
  Section~\ref{sec:format:profiles}
  (Requirement~\ref{req:format:profileid-discriminants}).
\end{requirement}

\begin{requirement}
  \label{req:format:immutable-verified-chunks}
  Chunks \MUST{} be immutable. Once a chunk is written and its bytes
  durably flushed, those bytes \MUSTNOT{} be modified. New state
  is encoded as new chunks with new identifiers; old chunks become
  unreachable garbage when their referencing manifests are no
  longer reachable.

  Implementations \MUST{} verify chunk hashes on first read of each
  chunk in a session: decompress the payload, compute the canonical
  preimage's BLAKE3, and confirm the result matches the declared
  \texttt{hash}. Subsequent reads \MAY{} skip verification if the
  chunk has been validated within the current session. A chunk
  whose hash does not match \MUST{} be treated as corrupt: if it is
  a cache chunk, discard and rebuild; if it is a canonical chunk
  (operation block, retained snapshot, or blob referenced from the
  manifest), surface as hard corruption.

  Readers \MUST{} verify that decompression of a chunk's payload
  produces exactly \texttt{uncompressed\_length} bytes. Length
  mismatch is treated as corruption.
\end{requirement}

\section{Operation Envelope Blocks}
\label{sec:format:opblocks}

The canonical document is stored as operation-envelope blocks. Each
block is a chunk of kind \texttt{OperationEnvelopeBlock} whose
payload is a vector of envelopes. A block is identified by its
\texttt{ChunkId}; because chunks are content-addressed, that
identifier cannot appear inside the block payload itself. Summary
metadata for a block travels in the \emph{manifest}
(\texttt{operation\_block\_summaries}), keyed by the block's chunk
id: the summary is computed by the operation layer from the block's
envelopes, and the storage layer carries it as opaque bytes without
interpreting it.

\begin{lstlisting}[language=Rust]
pub struct OperationEnvelopeBlock {
    /// Operation envelopes in this block. Order within a block is
    /// not normative for canonical reduction (envelopes are a set)
    /// but writers SHOULD preserve authoring order within a block
    /// for diagnostic purposes.
    pub envelopes: Vec<OperationEnvelope>,
}

/// Manifest-side summary of one block, keyed by the block's
/// ChunkId. Non-canonical and optional: a reader MUST NOT rely
/// on a summary being present, and canonical state never depends
/// on one.
pub struct OperationBlockSummary {
    /// DVV summary covering the envelopes in this block. Enables
    /// readers to determine causal coverage without parsing every
    /// envelope.
    pub dvv_summary: CausalContext,

    /// Earliest stamp in the block.
    pub min_stamp: OperationStamp,

    /// Latest stamp in the block.
    pub max_stamp: OperationStamp,
}
\end{lstlisting}

\begin{requirement}
  \label{req:format:operation-block-bounds}
  Writers \SHOULD{} begin a new operation-envelope block when
  adding another envelope would cause the uncompressed block
  payload to exceed 1 \,MiB, except when an individual envelope's
  encoded size exceeds 1\,MiB, in which case the envelope occupies
  its own block.

  Readers \MUST{} accept valid blocks of any size up to a profile-
  declared maximum (default: 64\,MiB uncompressed). Profile
  declarations \MAY{} raise or lower this bound. Blocks exceeding
  the active profile's bound \MUST{} be treated as malformed.

  The \texttt{operation\_roots} field of the manifest \MUST{} list
  every operation-envelope block that contributes to the canonical
  document. The set of envelopes \emph{is} the union of all
  envelopes across all referenced blocks; block boundaries are
  storage artifacts, not semantic structure.
\end{requirement}

\subsection{The Operation Index}

The optional \texttt{operation\_index\_root} chunk maps each
\texttt{OperationId} to the \texttt{ChunkRef} of its enclosing
block plus an offset within the block. It enables $O(\log n)$
lookup of an operation by id.

\begin{requirement}
  \label{req:format:rebuildable-operation-index}
  The operation index is an acceleration structure, not canonical.
  If absent, readers rebuild it by scanning all blocks. If present
  but corrupt or stale, readers \MUST{} reject the index and
  rebuild from blocks.

  Writers \SHOULD{} rebuild or incrementally update the operation
  index at commit time when the operation set has grown
  significantly since the last commit.
\end{requirement}

\section{Snapshots and Canonical Bases}
\label{sec:format:snapshots}

A snapshot is a materialized score state at a particular causal
frontier. Snapshots have two distinct roles:

\begin{description}
  \item[Acceleration snapshot.] A cache of the materialized state at
    some point, used to skip reduction work. Listed in the manifest
    (or not), but \emph{not} declared as a canonical base.
    Discardable.
  \item[Canonical base.] A materialized state at a causal frontier
    that the manifest declares as the base after pruning. Together
    with the operation envelopes after the frontier, this
    constitutes the canonical document. Not discardable while it
    remains declared.
\end{description}

\begin{lstlisting}[language=Rust]
pub struct SnapshotRef {
    pub snapshot_id: SnapshotId,

    /// The causal frontier this snapshot materializes: the
    /// envelopes whose effects are baked into the snapshot.
    pub covers_causal_frontier: CausalContext,

    /// Reduction algorithm version under which the snapshot was
    /// produced. Snapshots produced under an earlier algorithm
    /// version cannot be used as canonical bases under a later
    /// algorithm without rebuilding.
    pub reduction_algorithm_version: ReductionAlgorithmVersion,

    /// Profile under which the snapshot was produced.
    pub profile_id: ProfileId,

    /// Root chunk of the snapshot's materialized state.
    pub root: ChunkRef,

    /// Hash of the snapshot's root chunk for fast verification.
    pub hash: ContentHash,
}
\end{lstlisting}

\subsection{The Canonical Document Identity}

\begin{requirement}
  \label{req:format:canonical-document-reduction}
  The canonical document is defined as follows:

  \begin{itemize}
    \item If \texttt{canonical\_base} is \texttt{None}, the
      canonical document is the deterministic reduction
      (Section~\ref{sec:semops:reduction}) of the union of
      envelopes in all \texttt{operation\_roots} blocks.
    \item If \texttt{canonical\_base} is \texttt{Some(base)}, the
      canonical document is the base snapshot's materialized
      state plus the deterministic reduction of all envelopes
      whose \texttt{OperationId} is \emph{not covered} by the
      base's \texttt{covers\_causal\_frontier}. An
      \texttt{OperationId} is covered by a \texttt{CausalContext}
      if and only if it appears in the dotted version vector's
      per-replica contiguous range or among its dots (i.e.,
      causal-precedes-or-equals under the DVV's induced partial
      order). HLC stamps \MUSTNOT{} be used to determine
      coverage; only the DVV frontier is canonical.
    \item A snapshot may serve as the canonical base only if its
      \texttt{reduction\_algorithm\_version} equals the active
      superblock's \texttt{reduction\_algorithm\_version} and its
      \texttt{profile\_id} matches the manifest's profile
      declarations.
  \end{itemize}

  At most one \texttt{canonical\_base} is active at any time.
  Pruning replaces the current canonical base with a new one
  whose frontier strictly dominates the previous frontier under
  the DVV partial order; the old base becomes unreachable and
  is subject to garbage collection.

  Acceleration snapshots (\texttt{acceleration\_snapshots}) \MUST{}
  be ignored for the purposes of canonical state. They exist to
  accelerate cold open and incremental reduction; if their content
  disagrees with what the canonical reduction would produce, they
  \MUST{} be rebuilt or discarded.

  Caches, indexes, layout chunks, and other materialized artifacts
  in the bundle are acceleration structures only and \MUST{} be
  ignored or rebuilt if they disagree with the canonical document.
\end{requirement}

\subsection{Pruning}

\begin{requirement}
  \label{req:format:pruning-state-preservation}
  Pruning is the act of removing operation envelopes that are
  causally covered by a new canonical base snapshot. Pruning:

  \begin{enumerate}
    \item Materializes a snapshot covering the chosen causal
      frontier (a DVV).
    \item Replaces the current \texttt{canonical\_base} field
      with a \texttt{SnapshotRef} naming the new snapshot.
    \item Removes operation-envelope blocks whose envelopes are
      \emph{entirely} covered by the new frontier. Blocks with
      mixed coverage \MUST{} be split or retained whole;
      uncovered envelopes \MUST{} remain reachable.
    \item Commits the new manifest via the atomic write protocol
      (Section~\ref{sec:format:commit}).
  \end{enumerate}

  Pruning \MUSTNOT{} alter the canonical document state: the
  reduction of (base snapshot $+$ post-frontier envelopes) under
  the active algorithm \MUST{} equal the reduction of the
  pre-pruning envelope set.

  Implementations \MUSTNOT{} prune envelopes outside the explicit
  pruning protocol. Operation envelopes are not silently garbage-
  collected based on age or other heuristics.
\end{requirement}

\section{Blobs}
\label{sec:format:blobs}

The blob store holds large opaque content: audio reference tracks,
embedded images, custom fonts, ML model checkpoints.

\begin{requirement}
  \label{req:format:blob-hash-shape}
  A \texttt{BlobId} \MUST{} be BLAKE3 over the \emph{bare} preimage
  \texttt{"MUSCBLOB"} $\Vert$ \texttt{uncompressed\_payload}: the
  8-byte domain tag immediately followed by the uncompressed payload
  bytes, with \emph{no} chunk-kind, schema-version, or length fields.
  This is deliberately \emph{unlike} the structured chunk preimage of
  Section~\ref{sec:format:chunks} (which commits to kind, schema
  version, and length): a blob is opaque payload with no schema, so
  its identity commits only to the domain tag and the bytes. As with
  chunks, the compression algorithm is metadata on the
  \texttt{BlobRef} and is \emph{not} part of the blob's identity, so
  identical payloads stored under different compression share a
  \texttt{BlobId}.
\end{requirement}

\begin{lstlisting}[language=Rust]
pub struct BlobRef {
    pub blob_id: BlobId,
    pub media_type: String,  // RFC 6838 media type
    pub offset: u64,
    pub compressed_length: u64,
    pub uncompressed_length: u64,
    pub compression: CompressionAlgorithm,
    pub hash: ContentHash,

    /// Optional declared maximum size that readers MAY use to
    /// reject blobs that exceed reader resource limits.
    pub declared_max_uncompressed_length: Option<u64>,
}
\end{lstlisting}

\begin{requirement}
  \label{req:format:lazy-bounded-blobs}
  Readers \MUST{} apply resource limits to blob loading:
  uncompressed length \MUST{} be checked against the reader's
  policy before decompression begins; \texttt{media\_type} \MUST{}
  be validated against the reader's accepted list before any
  type-specific parsing.

  Implementations \MUSTNOT{} eagerly load all blobs at open;
  blobs \MUST{} be streamed on demand. Operations requiring blob
  content \MUST{} be designed for lazy access.

  Blob content is opaque to the core. The core stores, transmits,
  and integrity-checks blobs; it does not interpret their bytes.
\end{requirement}

\section{The Atomic Write Protocol}
\label{sec:format:commit}

A commit transforms the bundle from one canonical state to another.
The protocol uses two superblock slots so that a crash at any point
during commit leaves the bundle in either the pre-commit state
(active superblock unchanged) or the post-commit state (active
superblock now points at the new manifest); no intermediate state
is visible.

\begin{requirement}
  \label{req:format:atomic-commit-sequence}
  Commits \MUST{} follow this sequence:

  \begin{enumerate}
    \item Write all new chunks (envelope blocks, possibly a new
      snapshot, possibly a new operation index, etc.) to file
      regions outside the prelude and outside currently-reachable
      chunks. Writers \MAY{} append at end-of-file or fill known-
      unreachable regions, but \MUSTNOT{} overwrite any
      currently-reachable chunk.
    \item Durably flush the file.
    \item Compute the new manifest chunk's payload and write it.
    \item Durably flush the file.
    \item Compute the new \texttt{Superblock} with generation
      $= \text{active\_generation} + 1$ and \texttt{manifest\_offset}
      pointing at the new manifest chunk.
    \item Write the new superblock to the currently-inactive
      superblock slot (a 256-byte write at a fixed offset).
    \item Durably flush the file. \emph{This is the commit point.}
      Before this flush returns successfully, the active superblock
      is the old one; after it returns successfully, the active
      superblock is the new one.
  \end{enumerate}

  After step 7, the previously-active superblock slot becomes the
  inactive slot for the next commit. Garbage collection is a
  separate, optional, deferred operation (Section~\ref{sec:format:gc}).
\end{requirement}

\subsection{Crash Recovery}

\begin{requirement}
  \label{req:format:crash-recovery}
  On open after any crash, recovery follows the superblock
  selection rule (Section~\ref{sec:format:bundle}). The bundle is
  always in a recoverable state:

  \begin{itemize}
    \item Crash between steps 1 and 7: the active superblock is
      unchanged; the new chunks and the new manifest are
      unreachable garbage; the document opens at the pre-commit
      state. A later GC reclaims the orphaned bytes.
    \item Crash during step 7: the superblock slot write may be
      torn. The CRC check rejects torn writes. Recovery falls back
      to the other slot. The document opens at the pre-commit
      state. The partially-written slot will be overwritten by the
      next successful commit.
    \item Crash after step 7: the new superblock is active. The
      document opens at the post-commit state. No special
      recovery is needed.
  \end{itemize}

  Recovery \MUSTNOT{} synthesize or guess: if neither superblock
  validates, the file is corrupt and \MUST{} be reported as such.
\end{requirement}

\section{Garbage Collection and Retention}
\label{sec:format:gc}

Garbage collection reclaims bytes occupied by chunks that are no
longer reachable from any retained manifest. Retention policy
governs which old manifests are preserved for rollback.

\subsection{Retention Policy}

\begin{lstlisting}[language=Rust]
pub struct RetentionPolicy {
    /// Maximum number of previous manifests to retain beyond the
    /// active one. 0 means "no rollback retention"; the active
    /// manifest is the only one preserved. The default declared by
    /// the Full profile is profile-defined.
    pub retain_previous_manifests: u32,

    /// Optional wall-clock duration beyond which old retained
    /// manifests may be reclaimed regardless of count. None means
    /// no time-based eviction.
    pub retain_duration: Option<WallClockDuration>,

    /// Whether to preserve manifests explicitly marked as named
    /// checkpoints (user-invoked "save a version" action) beyond
    /// the count and duration limits.
    pub retain_named_checkpoints: bool,
}
\end{lstlisting}

\begin{requirement}
  \label{req:format:manifest-retention}
  The active conformance profile \MUST{} declare a
  \texttt{RetentionPolicy}. Implementations \MUST{} honor the
  declared policy: chunks reachable from any retained manifest
  \MUSTNOT{} be garbage-collected, and retention selection
  (which manifests are retained) \MUST{} be deterministic given
  the manifest history and the policy.

  Rollback is the act of writing a new superblock that points at
  a previously-active manifest still reachable in the chunk
  store. The fixed layout has exactly two superblock slots;
  rollback does not require additional on-disk superblocks. It
  requires only that the target manifest's reachable chunks
  remain in the bundle, which the retention policy guarantees.
\end{requirement}

\subsection{Garbage Collection Rules}

\begin{requirement}
  \label{req:format:conservative-garbage-collection}
  Garbage collection is a conservative, optional, deferred
  operation. It \MUST{} preserve every byte reachable transitively
  from:

  \begin{enumerate}
    \item The active superblock's manifest.
    \item Every retained manifest selected by the active
      \texttt{RetentionPolicy}.
    \item Any pending edit's working set (chunks written but not
      yet committed).
  \end{enumerate}

  Garbage collection \MUSTNOT{} run as part of a commit's critical
  path. It \MAY{} run as a background operation or as an explicit
  user-invoked compaction. Pre-GC and post-GC bundles \MUST{}
  contain identical canonical state and identical retained-rollback
  reachability.
\end{requirement}

\section{Forward Compatibility and Edit Barriers}
\label{sec:format:fwdcompat}

The bundle must round-trip cleanly across implementations that may
not all understand every extension referenced in it.

\subsection{Extension Declarations}

\begin{lstlisting}[language=Rust]
pub struct ExtensionDeclaration {
    pub extension_id: ExtensionId,
    pub version: SemVer,

    /// If true, an implementation that does not understand this
    /// extension MUST refuse to open the bundle for editing, but
    /// MAY open read-only.
    pub required: bool,

    /// Root chunks owned by this extension. Opaque to core
    /// readers; preserved across reads and writes.
    pub preserved_chunk_roots: Vec<ChunkRef>,

    /// Object kinds this extension contributes to or depends on,
    /// for edit-barrier evaluation.
    pub affected_object_kinds: Vec<ObjectKind>,

    /// Edit barriers declared by this extension.
    pub edit_barriers: Vec<EditBarrier>,
}

pub struct EditBarrier {
    /// Scope of the barrier: whole score, region, staff instance,
    /// analysis layer, object set, pitch space, tuning context, etc.
    pub scope: BarrierScope,

    /// Object kinds protected by this barrier.
    pub affected_object_kinds: Vec<ObjectKind>,

    /// Operation kind tags prohibited by this barrier. A barrier
    /// prohibits an operation class, not one exact payload, so it
    /// references OperationKindTag rather than OperationKind.
    pub prohibited_operation_kinds: Vec<OperationKindTag>,

    /// Additional condition narrowing the barrier (e.g., applies
    /// only when a specific attachment is present).
    pub condition: BarrierCondition,
}

/// A discriminator-only view of OperationKind (Chapter ch:semops),
/// without payload. Used by edit barriers, conformance reports,
/// and diagnostic surfaces where the operation's class matters
/// but its payload does not. There is a canonical mapping from
/// OperationKind to OperationKindTag.
pub enum OperationKindTag {
    InsertEvent,
    DeleteEvent,
    ModifyEvent,
    RespellPitch,
    Transpose,
    CreateCrossCutting,
    DeleteCrossCutting,
    ModifyCrossCutting,
    ChangeRegionTimeModel,
    InsertRegion,
    DeleteRegion,
    InsertStaffInstance,
    DeleteStaffInstance,
    SetUserSystemBreak,
    SetUserPageBreak,
    DeclareTransaction,
    Registered(OperationKindRegistryId),
    // Appended after Registered (append-only vocabulary; the
    // tag's wire discriminants never reorder). Phase-2/Phase-3
    // operation-catalog growth:
    InsertIdentifiedPitch,
    DeleteIdentifiedPitch,
    ModifyIdentifiedPitch,
    CreateVoice,
    DeleteVoice,
    SetMetadata,
    SetMetricGrid,
    InsertStaff,
    SetTimeSignature,
    SetTempoSegment,
    SetStaffLayout,
    // Schema-major-2 revision (repeat authoring). The pair keeps
    // its kind names verbatim, as the cross-cutting tags do.
    CreateRepeatStructure,
    DeleteRepeatStructure,
}

pub enum BarrierScope {
    WholeScore,
    Region(RegionId),
    StaffInstance(StaffInstanceId),
    AnalysisLayer(AnalysisLayerId),
    ObjectSet(Vec<TypedObjectId>),
    PitchSpace(PitchSpaceId),
    TuningContext,
    /// Registered scope kind for grammar-specific scoping.
    Registered(BarrierScopeRegistryId),
}

/// Additional narrowing condition for an edit barrier. The barrier
/// applies only when its scope, affected object kinds, prohibited
/// operation kinds, and this condition all match the candidate edit.
pub enum BarrierCondition {
    /// The barrier applies unconditionally within its scope.
    Always,

    /// The barrier applies only while the named object exists
    /// (not tombstoned). Used when an extension's invariant
    /// depends on a specific object remaining live.
    ObjectExists(TypedObjectId),

    /// The barrier applies only when the named object carries
    /// data declared by the named extension. Used when the
    /// invariant protected by the barrier exists only for
    /// objects participating in the extension.
    ObjectHasExtensionData {
        object: TypedObjectId,
        extension: ExtensionId,
    },

    /// Conjunction: all conditions must match.
    All(Vec<BarrierCondition>),

    /// Disjunction: any condition matching is sufficient.
    Any(Vec<BarrierCondition>),

    /// Negation: the inner condition must not match.
    Not(Box<BarrierCondition>),

    /// Registered condition kind, evaluated by an extension-aware
    /// implementation. Core implementations evaluate Registered
    /// conditions as Always (conservative): they treat the barrier
    /// as active whenever an unknown extension's condition would
    /// otherwise narrow it.
    Registered(BarrierConditionRegistryId),
}
\end{lstlisting}

\subsection{Behavior Under Unknown Extensions}

\begin{requirement}
  \label{req:format:unknown-extension-preservation}
  When an implementation opens a bundle declaring an extension it
  does not understand:

  \begin{itemize}
    \item If the extension is \texttt{required = true}, the
      implementation \MUST{} either refuse to open the bundle or
      open it strictly read-only. It \MUSTNOT{} edit.
    \item If the extension is \texttt{required = false}, the
      implementation \MAY{} open the bundle for editing, but \MUST{}
      preserve the extension's \texttt{preserved\_chunk\_roots}
      across reads and writes. The chunks are treated as opaque
      bytes.
    \item Edits \MUST{} be checked against every active edit
      barrier. An edit matching a barrier's scope, affected object
      kinds, and operation kinds is prohibited unless the user
      explicitly performs an \emph{unsafe edit}.
    \item An unsafe edit is an explicit user action that
      acknowledges loss of extension data. The unsafe-edit
      operation \MUST{} tombstone the relevant extension chunks
      (so they are no longer preserved) rather than silently
      breaking extension invariants.
  \end{itemize}

  The unsafe-edit mechanism prevents unknown extensions from
  rendering a document permanently unedittable while still
  preserving extension data by default.
\end{requirement}

\begin{requirement}
  \label{req:format:barrier-matching}
  \textbf{Barrier matching for target-free and opaque operations
  (ratified Pass 12).}
  An operation with no graph target (\texttt{SetMetadata},
  \texttt{DeclareTransaction}) is matched by \emph{score-wide}
  barriers only (\texttt{BarrierScope::WholeScore}); a barrier with
  any narrower scope cannot match it, because there is no target to
  test the scope against. An opaque \texttt{Registered} operation is
  matched \emph{fully conservatively}: since its targets and effects
  are unknown to a core implementation, it matches every active
  barrier whose remaining predicates do not exclude it.
\end{requirement}

\begin{requirement}
  \label{req:format:unsafe-tombstone}
  \textbf{Unsafe-edit tombstone semantics (ratified Pass 12).}
  Crossing a barrier by an unsafe edit immediately deactivates the
  owning extension's \emph{remaining} barriers for the editing
  session --- the extension's invariants are already forfeit, so its
  other barriers no longer protect anything. The crossing \MUST{} be
  durably recorded at the next bundle commit: the pending tombstone
  set survives the session that performed the unsafe edit. A
  tombstoned extension that was declared \texttt{required = true}
  leaves the bundle openable read-only by implementations that
  depended on that extension (its invariants can no longer be
  trusted); implementations that never understood it proceed under
  the unknown-extension rules above. The manifest-side \emph{byte
  encoding} of the tombstone record is deferred to the Binary Format
  companion (the manifest schema is frozen at major~0, so the record
  must ride an append-safe channel); see the companion's open
  questions.
\end{requirement}

\section{Text Projection}
\label{sec:format:textproj}

The format admits a deterministic projection to a canonical
s-expression text form. The text projection is normative.

The form itself is supplied by the \emph{Text Projection} companion
(version~0.1.0), which this section delegates to: the syntax, the canonical
layout, the parse rules, and the round-trip conformance requirement are its.
This section states \emph{what} is projected; the companion states \emph{how}.

\begin{requirement}
  \label{req:format:textproj}
  The text projection \MUST{} preserve, deterministically and
  bidirectionally with the binary form:

  \begin{itemize}
    \item All operation envelopes (with their identities, stamps,
      causal contexts, transaction groupings, and payloads).
    \item All profile declarations.
    \item All extension declarations \emph{in full} --- identity, semantic
      version, required flag, affected object kinds, edit barriers, and
      preserved chunk roots --- with opaque payloads in a canonical text
      encoding.
    \item Every \textbf{canonical blob}: a blob referenced by a canonical
      operation or by canonical reduced state is a canonical root of this
      document (Section~\ref{sec:format:manifest}), so the projection carries
      its media type and its uncompressed payload. Blobs referenced only by
      acceleration structures are non-canonical and are not projected.
    \item The active canonical base snapshot's identity, frontier, and
      reduction-algorithm version, \emph{with its payload inline}. The core
      specification once permitted the payload to be ``referenced externally'';
      that is withdrawn, because a compacted document's base is derivable from
      nothing else and an external reference would make the projection lossy
      (Text Projection companion,
      \texttt{req:textproj:base-snapshot-inline}).
    \item All canonical reduced state --- preserved by \emph{determining} it,
      never by carrying a second literal copy. Reduced state is a deterministic
      function of the operation set and the canonical base, so a projection
      carrying both would hold two sources of truth for one fact. See the Text
      Projection companion, requirement
      \texttt{req:textproj:reduced-state-derived}.
  \end{itemize}

  The text projection \MUSTNOT{} be required to preserve:

  \begin{itemize}
    \item Chunk offsets within the file.
    \item Compression algorithm choices.
    \item Cache chunks (operation indexes, layout caches, integrity
      indexes).
    \item Garbage bytes from prior commits.
    \item Superblock generation numbers, slot assignments, or CRCs.
  \end{itemize}

  Two binary bundles whose canonical document semantics are
  identical \MUST{} project to identical text. Two text projections
  that parse to identical canonical document semantics \MUST{}
  re-serialize to binary bundles whose canonical document
  semantics are identical (though their physical layout, chunking,
  and compression \MAY{} differ).
\end{requirement}

\subsection{Use Cases}

The text projection serves three primary use cases:

\begin{description}
  \item[Version control.] Text projections diff and merge in
    line-based tools. Merge conflicts surface at the operation-
    envelope level, which is the meaningful level for collaborative
    editing.
  \item[Format inspection and debugging.] A canonical text form
    enables direct inspection of bundle contents without binary
    tooling.
  \item[Long-term archival.] Text formats survive transitions
    between binary specifications. The text projection's structure
    is stable across binary format major versions where the
    semantic content is preserved.
\end{description}

For archival to hold, the text must be self-contained. Where a document has been
compacted onto a canonical base and its prior operations pruned, that base is
derivable from nothing else in the document, so the companion carries the
snapshot payload \emph{inline} rather than by reference
(\texttt{req:textproj:base-snapshot-inline}). A reference-only projection of a
compacted document would not determine the document it claims to project.

\section{Schema Versioning}
\label{sec:format:schema}

Chunk payloads are encoded against a schema declared by each
chunk's \texttt{schema\_version} field. Schemas evolve under a
semantic versioning discipline.

\begin{lstlisting}[language=Rust]
pub struct SchemaVersion {
    pub major: u16,
    pub minor: u16,
}
\end{lstlisting}

\begin{requirement}
  \label{req:format:schema-version-compatibility}
  Schema versioning rules:

  \begin{itemize}
    \item Major version changes are non-backward-compatible.
      Readers that do not support the chunk's major schema
      version \MUST{} handle it according to the chunk's role:
      \begin{itemize}
        \item \emph{Canonical chunks} (operation envelope blocks,
          the canonical base snapshot, the manifest, and blobs
          referenced by canonical operations or reduced state)
          \MUST{} be either parseable or the bundle \MUST{} be
          opened in read-only preservation mode (or refused). A
          reader that cannot parse a canonical chunk cannot know
          the canonical document.
        \item \emph{Non-canonical chunks} (acceleration
          snapshots, operation index, layout caches, text
          projection, integrity index, extension-defined chunks
          declared as preservable opaque content)
          \MAY{} be treated as opaque preserved content; the
          reader does not need to understand them to open the
          bundle, and writes preserve them verbatim.
      \end{itemize}
    \item Unknown canonical operation payloads (operation kinds
      not in the reader's understood catalog) require the
      relevant extension support or the reader \MUST{} open the
      bundle read-only without materializing the unknown
      operations. Materializing an unknown operation \MUSTNOT{}
      proceed silently: doing so would fork canonical state
      against any reader that does understand the operation.
    \item Minor version changes are backward-compatible. A reader
      supporting major version $N$ \MUST{} accept any minor version
      $N.\textit{m}$ for $m \leq$ the reader's supported minor.
      Under the positional, frozen-layout wire form, a minor change
      \MUST{} only \emph{append discriminants to the append-safe
      vocabularies the Binary Format companion defines} (its open
      \texttt{Registered} vocabularies, plus ratified appends to
      \texttt{OperationKind}, \texttt{OperationKindTag},
      \texttt{OperationPayload}, and the closed value-layer unions), which
      leaves the bytes of every prior value untouched; it \MUSTNOT{} add or
      reorder a struct field --- adding a field, even an \texttt{Option},
      occupies a new positional slot and shifts subsequent bytes, so it is a
      \emph{major} change (Binary Format companion,
      \sectionsc{What ``Additive'' Means Here}).
    \item The manifest's \texttt{manifest\_schema\_version} declares
      the schema of the manifest itself. Each chunk's
      \texttt{schema\_version} field declares the schema of that
      chunk's payload.
    \item Schema evolution is governed by the Binary Format
      companion specification, which defines the wire encoding
      for each schema version.
  \end{itemize}
\end{requirement}

Schema major~1 is the first data-model expansion major: it adds
\texttt{Canvas.layout\_defaults}, \texttt{Instrument.range}, and
\texttt{Region.permits\_spanning\_slurs} to the graph, and unifies the
non-canonical resolved-layout length-prefix width. Its wire form and the
byte-for-byte migration from major~0 are defined in the Binary Format
companion. The canonical-base \texttt{MaterializedState} embeds none of these
values, so it is byte-identical across the bump and stays major~0.
\texttt{Canvas.layout\_defaults} and \texttt{Instrument.range} reach only the
non-canonical acceleration (full-\texttt{Score}) snapshot; but
\texttt{Region.permits\_spanning\_slurs} also reaches the \emph{canonical}
operation layer, because \texttt{CreateRegion} embeds a full \texttt{Region},
so an operation-envelope block bearing a v1 \texttt{CreateRegion} is major~1.
Consequently a major-0-only reader opens a major-1 bundle \emph{fully} only
when the bundle carries no v1 \texttt{CreateRegion} operation; otherwise it
reads the major-0 canonical base and manifest but opens read-only (it cannot
replay the v1 canonical operations), per the canonical/non-canonical rules
above. It always discards the higher-major non-canonical caches. A major-1
reader migrates a major-0 acceleration snapshot and a major-0
\texttt{CreateRegion} payload on read (default-filling the new fields).

Schema major~2 is the second data-model expansion major: it fills the
truncated bodies of the cross-cutting structures (\texttt{Slur},
\texttt{Tie}, \texttt{Beam}, \texttt{Spanner}), \texttt{RepeatStructure}
(kind and voltas), \texttt{Staff} and \texttt{StaffLineConfiguration},
\texttt{Instrument}, and \texttt{ScoreMetadata}, per this revision's
Chapter~5 definitions. The same per-payload-type rules apply: the
canonical-base \texttt{MaterializedState} embeds none of these values and
stays major~0; the filled values reach the acceleration snapshot
\emph{and}, through the eight operation payloads that embed them
(\texttt{CreateCrossCutting}, \texttt{ModifyCrossCutting},
\texttt{InsertStaff} and \texttt{InsertStaffInstance} (the Operation
Catalog's \texttt{CreateStaff} / \texttt{CreateStaffInstance}),
\texttt{SetStaffLayout}, \texttt{SetMetadata},
\texttt{CreateRepeatStructure} (the repeat-authoring revision --- it
embeds the filled \texttt{RepeatStructure} unconditionally, so it is
\emph{born at v2} and every block carrying one stamps major~2; its
\texttt{DeleteRepeatStructure} sibling carries a bare identifier, a
major-0 layout, and stamps major~0), and --- transitively,
because an embedded \texttt{Region} may carry a \texttt{StaffInstance}
whose \texttt{staff\_lines\_override} embeds the filled
\texttt{StaffLineConfiguration} --- \texttt{InsertRegion}
(\texttt{CreateRegion})), the canonical operation layer. A block is
stamped with the \emph{lowest} schema major whose layouts decode its
bytes (the companion's minimal-stamping rule --- deterministic, so
identical content hashes identically): any v2-dependent bytes make it
major~2, and a lower-major-only reader opens such a bundle read-only. The wire form
and the total default-filling v1${\to}$v2 migration are defined in the
Binary Format companion.

\section{Format Profiles}
\label{sec:format:profiles}

A format profile defines a conformance subset that implementations
\MAY{} promise to support. Profiles enable lightweight readers to
declare their capabilities precisely.

\begin{lstlisting}[language=Rust]
pub struct ProfileDeclaration {
    pub profile_id: ProfileId,
    pub version: SemVer,

    /// Constraints imposed by this profile (maximum block size,
    /// permitted compression algorithms, required extensions, etc.)
    pub constraints: ProfileConstraints,
}

pub enum ProfileId {
    /// Full profile: all features supported, no restrictions.
    Full,

    /// Read-only profile: bundles produced under this profile MAY
    /// be opened by readers that cannot edit. Editing implementations
    /// SHOULD upgrade the profile on first edit.
    ReadOnly,

    /// Lite profile: a reduced feature set for embedded or mobile
    /// readers. Restricts maximum block sizes, snapshot complexity,
    /// and permitted extensions.
    Lite,

    /// Custom profile defined by a registered profile declaration.
    Custom(ProfileRegistryId),
}
\end{lstlisting}

The constraints a profile imposes are carried in
\texttt{ProfileConstraints}, which the GC and retention rules
(Section~\ref{sec:format:gc}) reference for the mandatory
\texttt{RetentionPolicy}:

\begin{lstlisting}[language=Rust]
pub struct ProfileConstraints {
    /// Maximum uncompressed operation-envelope block size a reader
    /// must accept under this profile (default 64 MiB).
    pub max_uncompressed_block_size: u64,

    /// The RetentionPolicy this profile declares. Required: the GC
    /// and retention rules read it from the active profile.
    pub retention_policy: RetentionPolicy,
    // Additional constraint fields (permitted compression sets,
    // required-extension lists) are reserved for a later revision;
    // a reader MUST tolerate their future addition under the
    // minor-version compatibility rule.
}
\end{lstlisting}

\begin{requirement}
  \label{req:format:profile-constraints}
  Every \texttt{ProfileDeclaration} \MUST{} carry a
  \texttt{ProfileConstraints}, and every \texttt{ProfileConstraints}
  \MUST{} include a \texttt{retention\_policy} field. This satisfies
  the requirement (Section~\ref{sec:format:gc}) that the active
  conformance profile declare a \texttt{RetentionPolicy}: the policy
  lives inside the active profile's constraints, not as a free-standing
  manifest field.

  When a bundle declares more than one profile, the
  \texttt{RetentionPolicy} in force is the one declared by the
  \emph{first} profile in the manifest's \texttt{profile\_declarations}
  list (declaration order is canonical). A bundle that declares no
  profile constraints carrying a policy uses the Full profile's
  default retention policy.
\end{requirement}

\begin{requirement}
  \label{req:format:profile-support}
  Every bundle \MUST{} declare at least one profile in its manifest's
  \texttt{profile\_declarations}. An implementation \MAY{} open a
  bundle if it supports any of the declared profiles, subject to
  the required-extension rules
  (Section~\ref{sec:format:fwdcompat}).
\end{requirement}

\begin{requirement}
  \label{req:format:profileid-discriminants}
  \texttt{ProfileId} encodes as a fixed \emph{20} bytes: a
  little-endian \texttt{u32} declaration-order discriminant
  (\texttt{Full} $= 0$, \texttt{ReadOnly} $= 1$, \texttt{Lite} $= 2$,
  \texttt{Custom} $= 3$) followed by a 16-byte
  \texttt{ProfileRegistryId}. The trailing 16 bytes carry the registry
  id for \texttt{Custom} and are zero for the other three variants; the
  field is always present, so the encoding is fixed-width (the
  superblock reserves a fixed slot for it). This assignment is
  normative and stable.
  \texttt{ProfileId} participates in superblock selection
  (Section~\ref{sec:format:bundle}, the load-bearing field set), so
  its discriminants are part of that comparison.
\end{requirement}

\section{Compression}
\label{sec:format:compression}

Each chunk \MAY{} be stored compressed. Compression is per-chunk
metadata, not part of content identity.

\begin{requirement}
  \label{req:format:zstandard-support}
  Conforming implementations \MUST{} support reading chunks
  compressed with Zstandard at any level zstd defines. Writers
  \MAY{} choose to compress or not on a per-chunk basis;
  uncompressed chunks (\texttt{CompressionAlgorithm::None}) are
  always permitted.

  Implementations \MUSTNOT{} require a specific compression
  algorithm for canonical conformance. Future major format
  versions \MAY{} introduce additional algorithms via the
  \texttt{Reserved} variant; such changes require a new format
  major version.

  Decompression \MUST{} verify the output length against the
  declared \texttt{uncompressed\_length} and reject chunks whose
  decompressed size disagrees.
\end{requirement}

\section{Streaming Reads}
\label{sec:format:streaming}

Cold open of a large score \MUST{} be fast: under one second to
first interactive frame on a 100-page orchestral score, per the
performance requirements (Chapter~\ref{ch:perf}). The format's
prelude and chunk-graph design support this.

\subsection{Required Locatability}

\begin{requirement}
  \label{req:format:chunk-locatability}
  The bundle \MUST{} guarantee that:

  \begin{itemize}
    \item The fixed header is at offset zero of the file.
    \item Both superblock slots are at fixed offsets within the
      first 576 bytes.
    \item The active manifest is locatable from the active
      superblock without scanning.
    \item Operation-envelope blocks, the canonical base snapshot
      (if any), acceleration snapshots, blobs, and extension
      chunks are locatable from the manifest's roots without
      scanning.
    \item The operation index (if present) provides $O(\log n)$
      operation-id lookup; if absent, scanning a small number of
      blocks suffices.
  \end{itemize}

  Locatable means: a reader can determine the offset and length of
  the chunk in the bundle using only previously read content.
\end{requirement}

\subsection{Cold Open Procedure}

The recommended cold open sequence:

\begin{enumerate}
  \item Read the 64-byte fixed header. Verify magic and CRC.
  \item Read both 256-byte superblock slots. Verify CRC on each.
    Select the active slot per the superblock selection rule
    (Section~\ref{sec:format:bundle}).
  \item Read and verify the manifest chunk referenced by the
    active superblock. The manifest is stored uncompressed
    (Section~\ref{sec:format:manifest-encoding}), so this is a
    single read-and-parse.
  \item If the manifest's \texttt{canonical\_base} is
    \texttt{Some}, read and verify the canonical base snapshot.
    The snapshot provides the materialized state at its declared
    causal frontier.
  \item Stream the operation-envelope blocks listed in
    \texttt{operation\_roots}. Apply the canonical reduction
    (Chapter~\ref{ch:semops}) over the envelopes whose
    \texttt{OperationId} is not covered by the canonical base's
    \texttt{covers\_causal\_frontier} (if a base is present), or
    over every envelope in the operation set (if no base is
    present). The result is the materialized score graph.
  \item Invoke the layout pipeline (Chapter~\ref{ch:layout-ir})
    over the materialized graph, producing the first system's
    layout. Layout caches (if reachable from the manifest as
    accelerators) \MAY{} be consulted to skip work; cache misses
    fall back to fresh layout computation.
  \item Load remaining regions, envelope blocks, snapshot ranges,
    and layout caches lazily as the user navigates.
\end{enumerate}

For unsynchronized first-frame display, implementations \MAY{}
present the canonical base's materialized state immediately while
the post-frontier reduction completes in the background; user
edits \MUST{} be deferred until reduction completes, since the
visible state is provisional until then.

\subsection{Memory Mapping}

\begin{requirement}
  \label{req:format:safe-memory-mapping}
  Implementations \MAY{} memory-map the bundle file for efficient
  random access to chunks and blobs. The format's content-addressing
  and chunk immutability guarantees make memory mapping safe:
  chunks cannot change underneath a reader holding a pointer into
  them. Superblock slots are 256-byte aligned and never overlap
  reachable chunk content, so writers updating an inactive
  superblock cannot corrupt readers accessing the active state.
\end{requirement}

\section{Binary Format Companion}
\label{sec:format:binary}

The byte-level encoding (varint conventions, exact field bit widths,
record alignment, endianness, string encoding details, struct
layouts) is delivered as a companion specification, the Binary
Format document. The Binary Format document is normative; an
implementation cannot conform to the file format specification
without conforming to the Binary Format specification.

\subsection{Ratified Convention Baseline}
\label{sec:format:codec-baseline}

Although the full Binary Format companion is still to be written, the
canonical encoding conventions shared by the core, operations, and
bundle layers are ratified now, so those three layers do not drift
apart before the companion formalizes them. The companion
\emph{inherits} this baseline rather than re-deriving it. Every
ratified discriminant table, derivation preimage, and primitive
encoding is consolidated in
Appendix~\ref{app:bytes} (the Canonical Byte-Layout Reference); the
companion imports that appendix as its starting point rather than
recovering the layouts from the three crates.

\begin{requirement}
  \label{req:format:codec-conventions}
  The canonical encoding of every composite structure (whole
  \texttt{Score}, operation envelopes, manifest) \MUST{} follow these
  conventions:
  \begin{itemize}
    \item Integers are little-endian. A boolean is a single
      \texttt{0}/\texttt{1} byte.
    \item Counts and the length prefix of every variable-width
      \emph{leaf} (an identifier, a \texttt{RationalTime}, a string)
      are \texttt{u32} little-endian. Length-prefixing every leaf
      makes the decoder width-agnostic: it never infers a boundary
      from context.
    \item A tagged union is encoded as a single discriminant byte
      followed by the variant payload.
    \item Free-text string fields are length-prefixed UTF-8 and are
      \emph{not} NFC-folded by the codec: the in-memory value is
      preserved byte-exact, so \texttt{decode(encode(x)) == x}.
      Catalog identifiers are NFC-normalized at construction, not in
      the codec.
  \end{itemize}

  Certain fixed hashing preimages deliberately depart from the
  single-discriminant-byte convention where a wider or order-bearing
  tag is required --- notably \texttt{TypedObjectId}'s 16-bit
  big-endian discriminant
  (Requirement~\ref{req:graph:typed-object-id-discriminants}). Such
  departures are pinned at their definitions and are not overridden by
  this baseline.
\end{requirement}

\begin{requirement}
  \label{req:format:rationaltime-encoding}
  The primitive scalar layouts referenced throughout the canonical
  encodings are ratified as follows:
  \begin{itemize}
    \item \texttt{RationalTime} encodes as a sign byte
      (\texttt{0} = zero, \texttt{1} = positive, \texttt{2} =
      negative), then the \texttt{u32}-little-endian-length-prefixed
      big-endian magnitude of the numerator, then the
      \texttt{u32}-little-endian-length-prefixed big-endian magnitude
      of the denominator. The value \MUST{} always be stored reduced
      (lowest terms, positive denominator), so the encoding is
      canonical: equal rationals encode to equal bytes.
    \item Wall-clock integers (\texttt{WallClockTime},
      \texttt{WallClockDuration}) are little-endian fixed-width
      integers, matching \texttt{QuantizedCoord}.
  \end{itemize}
\end{requirement}

\begin{openquestion}
  The remainder of the Binary Format companion --- the full composite
  struct layouts, schema-version wire evolution, and varint details
  beyond the baseline above --- is under development as a separate
  specification. Its development follows this chapter and inherits the
  ratified baseline (Sections~\ref{req:format:codec-conventions}
  and~\ref{req:format:rationaltime-encoding}); the remaining byte-level
  decisions depend on the structural decisions specified here.
\end{openquestion}

\section{Forward References}

\begin{itemize}
  \item The constraint-solver interface (Chapter~\ref{ch:solver})
    operates on the in-memory \texttt{ConstrainedLayoutIR}; layout
    caches stored in the bundle are derivative and not part of the
    canonical state.
  \item The full performance requirements that this format must
    satisfy are in Chapter~\ref{ch:perf}.
  \item The Binary Format companion specifies byte-level encoding.
  \item The Text Projection companion specifies the canonical
    s-expression form.
  \item The Profile Conformance companion specifies the feature
    lists per profile.
\end{itemize}

% ===========================================================================
\chapter{Constraint-Solver Interface}
\label{ch:solver}

This chapter specifies the constraint-solver interface: the
\texttt{ConstrainedLayoutIR} the solver consumes, the
\texttt{ResolvedLayoutIR} it produces, the constraints it must
respect, the determinism contract it must satisfy, the diagnostic
information it must emit, and the conformance discipline by which
implementations are judged.

The chapter deliberately specifies the \emph{interface} and its
contracts rather than a particular algorithm. Implementations are
free to choose among approaches (Cassowary, linear programming
relaxations, custom heuristics, differentiable layout, learned
methods) and to evolve their choices over time. Conformance is
established by hard-constraint validity, deterministic behavior,
diagnostic completeness, and performance and quality on a reference
suite of test scores.

\section{Design Principles}
\label{sec:solver:principles}

\begin{description}
  \item[The interface is normative; global optimality is not.] The
    solver's input/output contract, its constraint families, its
    diagnostic surface, and the discipline by which it is evaluated
    are normative. Global aesthetic optimality (proving no better
    layout exists) is neither decidable nor practical to verify for
    arbitrary scores; it is therefore replaced by reference-suite
    thresholds.
  \item[Hard constraints are validity; quality metrics are evaluation
    axes.] Hard constraints are the conditions under which a layout
    is valid at all. Quality metrics are the axes along which valid
    layouts are evaluated and compared. The solver \MUSTNOT{}
    trade off a hard-constraint violation for a better quality
    score: a layout that violates a hard constraint is not a layout,
    regardless of its quality vector.
  \item[Determinism within an implementation; equivalence across
    implementations.] A given implementation at a fixed version
    \MUST{} produce identical output for identical input. Different
    conforming implementations \MAY{} produce different layouts;
    they conform by satisfying hard constraints, meeting diagnostic
    contracts, behaving deterministically internally, and passing
    reference-suite thresholds. Cross-implementation byte-identity
    is explicitly not required.
  \item[Pareto frontiers as design target, not conformance
    obligation.] Solvers \SHOULD{} seek layouts on or near the
    Pareto frontier of the normative quality metrics. They are
    \emph{not} required to prove they have done so universally;
    conformance is reference-suite-based.
  \item[Deterministic budgets, not wall-clock budgets.] Wall-clock
    time \MAY{} stop interactive work, but the returned layout's
    canonical form depends on deterministic counters
    (iterations, constraint evaluations, search nodes), not on
    scheduler luck.
  \item[Incremental solving as observational equivalence.]
    Incremental re-solving \MUST{} be observationally equivalent
    to full solving of the declared invalidation scope under the
    same profile and budget. Implementations are free to widen
    scopes conservatively but \MAY NOT{} return layouts that
    diverge from the canonical scoped result.
  \item[Structured diagnostics, never silent corruption.] Solver
    failures are first-class outputs with diagnostic information
    sufficient for the editor to surface meaningful feedback. The
    score graph is never invalidated by a solver failure; a solver
    error means ``we could not produce a layout,'' not ``the
    score is broken.''
\end{description}

\section{The Solver Interface}
\label{sec:solver:interface}

\begin{lstlisting}[language=Rust]
pub trait ConstraintSolver: Send + Sync {
    /// The solver's identifying tier. Used by readers/editors to
    /// decide whether this solver is sufficient for the active
    /// score's declared solver tier requirement.
    fn tier(&self) -> SolverTier;

    /// The solver's implementation version. Identical input plus
    /// identical version produces identical output (within-
    /// implementation determinism).
    fn version(&self) -> SolverVersion;

    /// Solve from scratch.
    fn solve(
        &self,
        input: &ConstrainedLayoutIR,
        config: &SolverConfig,
    ) -> SolveReport;

    /// Solve incrementally over the declared invalidation scope.
    fn solve_incremental(
        &self,
        input: &ConstrainedLayoutIR,
        prior: &SolverState,
        invalidations: &InvalidationSet,
        config: &SolverConfig,
    ) -> SolveReport;
}

pub struct SolverConfig {
    /// Conformance profile under which to solve. Selects metric
    /// thresholds, profile-specific constraints, and the active
    /// extension set.
    pub profile: SolverProfile,

    /// Deterministic budget. Wall-clock time is advisory only.
    pub budget: SolverBudget,

    /// Tie-breaking weights for selecting among layouts of
    /// equivalent quality.
    pub tie_breaking: TieBreakingWeights,
}

pub struct SolverBudget {
    /// Maximum solver iterations (algorithm-defined unit; the
    /// reference algorithm counts Cassowary pivots).
    pub max_iterations: u64,

    /// Maximum search-tree nodes explored, for algorithms that
    /// perform discrete search.
    pub max_nodes: u64,

    /// Maximum constraint-evaluation count.
    pub max_constraint_evaluations: u64,

    /// Advisory wall-clock cap, in milliseconds. The solver MAY
    /// stop early when this is exceeded, but the layout it
    /// returns under stoppage MUST be the same layout that the
    /// deterministic budget would have produced at the
    /// corresponding deterministic counter state. Wall-clock time
    /// MUST NOT change the canonical layout.
    pub advisory_wall_time_ms: Option<u64>,
}

pub struct SolverBudgetUsed {
    pub iterations: u64,
    pub nodes: u64,
    pub constraint_evaluations: u64,
    pub wall_time_ms: u64,
}
\end{lstlisting}

\begin{requirement}
  \label{req:solver:solver-purity}
  \texttt{solve} and \texttt{solve\_incremental} \MUST{} be pure
  functions of their inputs within the determinism contract
  (Section~\ref{sec:solver:determinism}). They \MUSTNOT{} consult
  external state (system clock, environment, thread scheduling,
  random sources) in ways that influence the returned layout.

  Implementations \MAY{} use multi-threaded execution internally,
  but the result \MUST{} be identical to a single-threaded
  execution under the same deterministic budget.
\end{requirement}

\section{The Solver Report}
\label{sec:solver:report}

Every solver invocation returns a \texttt{SolveReport}. There is no
separate ``error'' return; failures and partial successes appear
as report variants distinguished by \texttt{SolveStatus}.

\begin{lstlisting}[language=Rust]
pub struct SolveReport {
    pub status: SolveStatus,

    /// Whether every hard constraint is satisfied.
    pub satisfied_hard_constraints: bool,

    /// The layout produced. Always present; under failure
    /// statuses, the layout is diagnostic-only and MUST NOT be
    /// used as if it were valid.
    pub layout: ResolvedLayoutIR,

    /// Updated solver state for subsequent incremental calls.
    pub state: SolverState,

    /// Unsatisfied hard constraints, if any. Empty under
    /// SolveStatus::Solved.
    pub unsatisfied_constraints: Vec<ConstraintId>,

    /// Non-fatal warnings about the solution.
    pub warnings: Vec<SolverWarning>,

    /// Quality metric vector for the returned layout.
    pub metric_vector: QualityMetricVector,

    /// Budget consumed during this solve.
    pub budget_used: SolverBudgetUsed,
}

pub enum SolveStatus {
    /// All hard constraints satisfied, target quality reached.
    Solved,

    /// All hard constraints satisfied, but warnings were
    /// generated (e.g., a soft constraint had a large violation,
    /// or a layout decision was unusual and worth surfacing).
    SolvedWithWarnings,

    /// Deterministic budget exhausted before reaching target
    /// quality. The returned layout MUST still satisfy all hard
    /// constraints; it simply may not be optimal. If hard
    /// constraints could not be satisfied within budget, the
    /// status is Unsatisfiable, not PartialBudgetExhausted.
    PartialBudgetExhausted,

    /// Hard constraints cannot be simultaneously satisfied. The
    /// returned layout is diagnostic-only and MUST NOT be
    /// presented as a valid layout.
    Unsatisfiable,

    /// Solver bug or unexpected error. The returned layout is
    /// diagnostic-only. Editors SHOULD surface this as a defect
    /// to report.
    InternalError,
}

pub struct SolverWarning {
    pub kind: SolverWarningKind,
    pub affected_objects: Vec<TypedObjectId>,
    pub message: String,
}

pub enum SolverWarningKind {
    LargeSoftConstraintViolation { constraint: ConstraintId, magnitude: f64 },
    UnusualLayoutDecision(String),
    QualityFloorApproached { metric: QualityMetricKind },
    ExtensionWarning(ExtensionWarningId),
}
\end{lstlisting}

\begin{requirement}
  \label{req:solver:report-authority}
  The \texttt{SolveReport.layout} field is always populated, but
  its authority depends on \texttt{status}:

  \begin{itemize}
    \item \texttt{Solved} or \texttt{SolvedWithWarnings}: the
      layout is a valid layout that the editor may render.
    \item \texttt{PartialBudgetExhausted}: the layout is a valid
      layout (\texttt{satisfied\_hard\_constraints} is \texttt{true})
      and the editor may render it, while noting that quality
      is below the profile's target.
    \item \texttt{Unsatisfiable} or \texttt{InternalError}: the
      layout is diagnostic-only. Editors \MAY{} render it for the
      user to see what failed, but \MUST{} clearly indicate the
      layout is not authoritative.
  \end{itemize}

  A solver \MUSTNOT{} return \texttt{PartialBudgetExhausted} with
  \texttt{satisfied\_hard\_constraints == false}. If hard constraints
  cannot be satisfied within the budget, the correct status is
  \texttt{Unsatisfiable}, with the partial layout marked as such.
\end{requirement}

\begin{requirement}
  \label{req:solver:subconformant-report}
  \textbf{Sub-conformant report shape (ratified Pass 12).}
  A below-conformance passthrough solver (\texttt{SolverTier::Stub})
  presented with declared constraints it does not evaluate reports
  \texttt{SolvedWithWarnings} with
  \texttt{satisfied\_hard\_constraints == false} and a warning
  naming the unevaluated-constraint condition. The layout stays
  renderable; the report makes no conformance claim and satisfies
  no minimum-tier requirement. This is the one sanctioned case in
  which a renderable status carries
  \texttt{satisfied\_hard\_constraints == false}: the field reports
  \emph{evaluated} satisfaction, and a non-evaluating solver has
  nothing to claim.
\end{requirement}

\section{Constraint Families}
\label{sec:solver:families}

The solver consumes constraints in normalized form. The constraint
families enumerated here are the normative core; implementations
\MAY{} support additional families via
\texttt{LayoutConstraint::Registered}.

\subsection{Strength Levels}

\begin{lstlisting}[language=Rust]
pub enum ConstraintStrength {
    /// Hard constraint. The solver MUST satisfy this constraint,
    /// or return Unsatisfiable.
    Required,

    /// Soft constraint with an associated weight. The solver
    /// minimizes the weighted violation when optimizing.
    Preferred { weight: f64 },
}
\end{lstlisting}

\begin{requirement}
  \label{req:solver:required-constraint-strength}
  A solver \MUSTNOT{} treat a \texttt{Required} constraint as
  if it were \texttt{Preferred} for any reason, including for
  quality optimization. A layout that violates a \texttt{Required}
  constraint is not a valid layout, regardless of its quality
  vector.
\end{requirement}

\begin{requirement}
  \label{req:solver:kind-strength}
  \textbf{Strength is kind-determined (ratified Pass 12).}
  A constraint instance carries no strength field; its strength is
  determined by its kind. Break constraints take their strength
  from \texttt{BreakKind}: \texttt{Hard} $\to$ \texttt{Required},
  \texttt{Soft} $\to$ \texttt{Preferred} with weight $1.0$. Every
  other core constraint family normalizes to \texttt{Required}.
  \texttt{Registered} constraint kinds normalize conservatively to
  \texttt{Required}. A future constraint family declares its
  strength (and weight, if \texttt{Preferred}) in its normative
  definition --- the strength channel is the \emph{kind}, not the
  instance.
\end{requirement}

\subsection{Normative Constraint Families}

The constraint catalog (spring, collision, alignment, containment,
break, aesthetic, extension) is specified in
Chapter~\ref{ch:layout-ir}. Implementations \MUST{} support every
family normatively required by the active solver tier
(Section~\ref{sec:solver:tiers}).

\section{Quality Metrics}
\label{sec:solver:quality}

The solver's output is evaluated against a normative set of quality
metrics. Each metric is normalized to the closed interval
$[0.0, 1.0]$, with $0.0$ representing the best achievable value and
$1.0$ representing the worst tolerable value before a profile
threshold flags the metric as failing.

\begin{lstlisting}[language=Rust]
/// A quality metric value normalized to [0.0, 1.0].
/// Lower is always better.
pub struct NormalizedMetric(pub f64);

impl NormalizedMetric {
    /// Construct, panicking if the value is not finite or is out
    /// of range. Conforming implementations MUST construct only
    /// valid NormalizedMetric values.
    pub fn new(value: f64) -> Self {
        assert!(value.is_finite());
        assert!(0.0 <= value && value <= 1.0);
        NormalizedMetric(value)
    }
}

pub struct QualityMetricVector {
    pub collision_penalty: NormalizedMetric,
    pub spacing_distortion: NormalizedMetric,
    pub slur_shape_penalty: NormalizedMetric,
    pub beam_slope_penalty: NormalizedMetric,
    pub vertical_density_penalty: NormalizedMetric,
    pub system_break_penalty: NormalizedMetric,
    pub page_fill_efficiency: NormalizedMetric,
    pub casting_off_quality: NormalizedMetric,
    pub symbol_density_uniformity: NormalizedMetric,

    /// Additional metrics contributed by extensions. Each MUST be
    /// a valid NormalizedMetric (finite, in [0.0, 1.0], lower is
    /// better).
    pub extension_metrics: Vec<ExtensionMetric>,
}

pub struct ExtensionMetric {
    pub metric_id: ExtensionMetricId,
    pub value: NormalizedMetric,
}
\end{lstlisting}

\begin{requirement}
  \label{req:solver:normalized-metric-bounds}
  Every \texttt{NormalizedMetric} value \MUST{} satisfy:
  \begin{itemize}
    \item It is finite (not NaN, not infinity).
    \item It lies in the closed interval $[0.0, 1.0]$.
    \item Lower values denote better layouts; higher values denote
      worse layouts.
  \end{itemize}

  Per-metric normalization functions (mapping raw measurements to
  $[0.0, 1.0]$) are specified in the Quality Metric Catalog
  companion document. Implementations \MUST{} use the catalog's
  normalization; arbitrary normalization is non-conforming.

  Extension metrics \MAY{} define their own normalization function
  but \MUSTNOT{} change orientation or range: $0.0$ remains best,
  $1.0$ remains worst, and values must remain finite within range.
\end{requirement}

\subsection{Pareto Frontiers as Design Target}

Solvers should seek layouts on or near the Pareto frontier of the
normative metrics. This is the design intent. It is not the
conformance requirement.

\begin{requirement}
  \label{req:solver:pareto-efficiency}
  The normative metric set defines the axes along which layout
  quality is evaluated. Solvers \SHOULD{} produce Pareto-efficient
  layouts with respect to these metrics: layouts for which no
  metric can be improved without worsening another. Because global
  Pareto optimality is not practically verifiable for arbitrary
  scores, conformance is determined by reference-suite thresholds
  (Section~\ref{sec:solver:conformance}) rather than by proof of
  Pareto-optimality.

  When multiple Pareto-equivalent layouts exist, the configured
  \texttt{TieBreakingWeights} select among them. Tie-breaking is
  deterministic.
\end{requirement}

\begin{lstlisting}[language=Rust]
pub struct TieBreakingWeights {
    pub collision: f64,
    pub spacing: f64,
    pub slur_shape: f64,
    pub beam_slope: f64,
    pub vertical_density: f64,
    pub system_break: f64,
    pub page_fill: f64,
    pub casting_off: f64,
    pub symbol_density: f64,
}
\end{lstlisting}

\begin{requirement}
  \label{req:solver:tie-breaking-weight-defaults}
  Tie-breaking weights \MUST{} have normative defaults specified
  in the Quality Metric Catalog. Implementations \MAY{} permit
  users to customize these weights; the defaults represent the
  reference engraving aesthetic and are the basis for
  reference-suite conformance.
\end{requirement}

\section{Conformance Tiers}
\label{sec:solver:tiers}

Solver conformance is graded into tiers, each declaring a subset
of features the implementation supports. Tiers are \emph{orthogonal}
to file-format profiles (Chapter~\ref{ch:format}); a document may
declare both an authoring file profile and a minimum solver tier,
and the two axes are independent.

\begin{lstlisting}[language=Rust]
pub enum SolverTier {
    /// Not a conformance tier: an interface-only / passthrough
    /// solver that evaluates no constraints and computes no
    /// quality metrics. Ordered below every conformant tier.
    Stub,

    /// Minimal Layout Solver.
    Minimal,

    /// Standard Engraving Solver.
    Standard,

    /// Advanced / Extension-Aware Solver.
    Advanced,
}
\end{lstlisting}

The \texttt{Stub} variant exists so scaffolding, harnesses, and
not-yet-conformant pipelines can label themselves honestly: it makes
no conformance claim, and a document requiring any conformant tier
is never satisfied by a solver declaring \texttt{Stub}.

\subsection{Minimal Layout Solver}

\begin{requirement}
  \label{req:solver:minimal-tier}
  A solver claiming the Minimal tier \MUST{}:

  \begin{itemize}
    \item Satisfy every hard constraint or return
      \texttt{Unsatisfiable} (validity conformance).
    \item Be deterministic within its declared implementation
      version (Section~\ref{sec:solver:determinism}).
    \item Produce well-formed \texttt{SolveReport} values with
      accurate \texttt{satisfied\_hard\_constraints},
      \texttt{unsatisfied\_constraints}, and metric vectors
      (diagnostic conformance).
    \item Support common-music-notation regions, metric time, and
      the standard constraint families.
    \item Pass the Minimal subset of the reference suite.
  \end{itemize}

  A Minimal solver \MAY{} produce mediocre aesthetic quality;
  metric thresholds for the Minimal tier are relaxed relative to
  the Standard tier.
\end{requirement}

\subsection{Standard Engraving Solver}

\begin{requirement}
  \label{req:solver:standard-tier}
  A solver claiming the Standard tier \MUST{}:

  \begin{itemize}
    \item Satisfy every Minimal-tier obligation.
    \item Meet the Standard-tier metric thresholds on the
      reference suite for every quality metric.
    \item Support the incremental solving contract with at least
      Measure-local invalidation scope
      (Section~\ref{sec:solver:incremental}).
    \item Support every Standard-tier constraint family declared
      in the Quality Metric Catalog.
  \end{itemize}

  The Standard tier corresponds to ``professional engraving
  quality'' for common-practice notation.
\end{requirement}

\subsection{Advanced / Extension-Aware Solver}

\begin{requirement}
  \label{req:solver:advanced-tier}
  A solver claiming the Advanced tier \MUST{}:

  \begin{itemize}
    \item Satisfy every Standard-tier obligation.
    \item Support proportional and aleatoric region layout, not
      only metric regions.
    \item Support microtonal accidentals and their interactions
      with adjacent glyphs.
    \item Support extension-contributed constraint families and
      extension metric thresholds for each registered extension
      whose layout requirements are part of the Advanced
      reference suite.
    \item Support at least System-local invalidation scope
      under incremental solving.
  \end{itemize}
\end{requirement}

\subsection{Tier Declaration and Document Compatibility}

\begin{requirement}
  \label{req:solver:solver-tier-compatibility}
  Documents \MAY{} declare a minimum solver tier required for
  faithful engraving in their profile declarations
  (Chapter~\ref{ch:format}). A solver with tier lower than the
  document's declared minimum \MAY{} open the document but
  \SHOULD{} surface a clear warning that engraving quality may
  not match the document's intent.

  Solver tier is independent of file-format profile. A document
  authored under the \texttt{Lite} file profile \MAY{} declare a
  minimum solver tier of \texttt{Standard}; a document authored
  under the \texttt{Full} file profile \MAY{} declare no minimum
  tier and be displayable by a \texttt{Minimal} solver.
\end{requirement}

\section{Determinism}
\label{sec:solver:determinism}

The solver's determinism contract has two parts. \emph{Within-
implementation determinism} is a strict byte-equality obligation;
\emph{cross-implementation conformance} is reference-suite-based.

\subsection{Within-Implementation Determinism}

\begin{requirement}
  \label{req:solver:within-implementation-byte-stability}
  A solver implementation at a fixed version \MUST{} produce
  byte-for-byte identical \texttt{SolveReport} output for
  identical inputs, where ``identical inputs'' means:

  \begin{itemize}
    \item Identical \texttt{ConstrainedLayoutIR}.
    \item Identical \texttt{SolverConfig} (profile, budget,
      tie-breaking weights).
    \item Identical font and glyph metrics (referenced by version
      and content hash).
    \item Identical reduction algorithm version
      (Chapter~\ref{ch:semops}).
    \item For incremental calls: identical prior
      \texttt{SolverState} and identical invalidation set.
    \item Identical solver implementation version (i.e., the same
      binary).
  \end{itemize}

  Determinism \MUST{} hold across runs on the same machine, across
  cold and warm starts, and across hardware threads.
  Implementations \MUSTNOT{} use fast-math flags, processor-
  dependent SIMD without canonical-output guarantees, or any
  compiler options that vary floating-point results across builds
  on the same platform.

  Where parallelism is used internally, the solver \MUST{}
  serialize the merge step so that thread scheduling does not
  affect output. Deterministic budgets \MUST{} be enforced from
  shared counters that capture all parallel contributions.
\end{requirement}

\subsection{Cross-Implementation Conformance}

\begin{requirement}
  \label{req:solver:cross-implementation-conformance}
  Cross-implementation determinism (byte-equality of outputs
  across different conforming solvers) is \emph{not} required.
  Different conforming solvers \MAY{} produce different layouts
  for the same input. Conformance across implementations is
  established by:

  \begin{enumerate}
    \item All hard constraints satisfied on every reference suite
      score.
    \item The implementation is internally deterministic per the
      within-implementation rule above.
    \item \texttt{SolveReport} values are well-formed and
      diagnostically accurate.
    \item For each reference-suite score, every quality metric is
      within the tier's declared threshold.
  \end{enumerate}
\end{requirement}

\begin{rationale}
  Cross-implementation byte equality would freeze the ecosystem
  to whatever single algorithm first claimed conformance, and
  would preclude algorithmic innovation. Reference-suite
  thresholds give users a meaningful quality floor while
  permitting genuine implementation diversity. Within-
  implementation determinism remains essential for reproducible
  builds, regression testing, and predictable behavior under
  collaborative editing.
\end{rationale}

\section{Incremental Solving}
\label{sec:solver:incremental}

\subsection{Invalidation Scopes}

\begin{lstlisting}[language=Rust]
pub enum InvalidationScope {
    /// A single object's layout depends only on local context.
    ObjectLocal,
    /// Effects bounded within a single measure.
    MeasureLocal,
    /// Effects bounded within a single system.
    SystemLocal,
    /// Effects bounded within a single page.
    PageLocal,
    /// Effects bounded within a single region.
    RegionLocal,
    /// Effects may propagate across the entire score.
    WholeScore,
}

pub struct InvalidationSet {
    /// Declared scope. The solver MAY widen this scope
    /// conservatively; it MUST NOT narrow it.
    pub scope: InvalidationScope,

    /// Specific invalidated entries within the declared scope.
    pub slots: Vec<SpringSlotId>,
    pub bands: Vec<VerticalBandId>,
    pub constraints: Vec<ConstraintId>,
    pub glyphs: Vec<GlyphObjectId>,
}
\end{lstlisting}

\subsection{Observational Equivalence}

\begin{requirement}
  \label{req:solver:incremental-observational-equivalence}
  For any invalidation set with declared scope $S$,
  \texttt{solve\_incremental} \MUST{} produce a layout
  observationally equivalent to \texttt{solve} called on the
  same input restricted to $S$, under the same profile, the same
  configuration, and the same deterministic budget. ``Observational
  equivalence'' means byte-identical \texttt{ResolvedLayoutIR}
  output for objects within $S$, and unchanged output for objects
  outside $S$.

  The solver \MAY{} widen the scope conservatively: an
  \texttt{ObjectLocal} invalidation may be solved as
  \texttt{MeasureLocal} if that simplifies the implementation,
  provided the widened result is also observationally equivalent
  to full solving at the widened scope.

  The solver \MAY NOT{} narrow the scope. If
  \texttt{SystemLocal} effects might propagate (e.g., an edit at
  the system's end shifts subsequent systems' contents), the
  invalidation \MUST{} be declared at least
  \texttt{SystemLocal} or higher.
\end{requirement}

\subsection{Propagation Documentation}

\begin{requirement}
  \label{req:solver:propagation-scope-contract}
  Implementations \MUST{} document, per their published
  conformance contract, the maximum invalidation scope they can
  service for each constraint family. The documented bound is the
  implementation's contract for how far edit effects propagate
  before requiring whole-score re-solving.
\end{requirement}

\section{Conformance: The Reference Suite}
\label{sec:solver:conformance}

Solver conformance is established by performance on a normative
reference suite of test scores. Each tier has its own subset of
the suite and its own metric thresholds.

\begin{requirement}
  \label{req:solver:reference-suite-passage}
  The Reference Suite is delivered as a companion specification.
  Each suite entry consists of:

  \begin{itemize}
    \item A test score in canonical \texttt{.musc} form.
    \item Per-tier inclusion: whether the entry is required for
      Minimal, Standard, or Advanced conformance.
    \item Per-tier metric thresholds: the maximum permitted
      \texttt{NormalizedMetric} value for each quality metric.
    \item Optional fixed-expectation tests: specific layout
      properties (a particular bar's width, a specific system
      break location) that all conforming solvers \MUST{}
      reproduce. Used sparingly, only where the suite intends to
      pin an unambiguous expectation.
  \end{itemize}

  A solver claiming a given tier \MUST{} pass every entry
  required at that tier. Failure on any single entry is
  conformance failure at the claimed tier.

  Suite versions are tied to specification versions; conforming
  implementations \MUST{} declare which suite version they pass.
\end{requirement}

\begin{rationale}
  The reference suite makes conformance checkable. The earlier
  universal-Pareto formulation quantified over all possible
  layouts of every possible score, which is undecidable.
  Reference-suite thresholds replace ``no better layout exists''
  with ``this layout is within tolerance on these specific
  scores,'' which is testable, falsifiable, and operationally
  useful. The cost is that conformance is only as good as the
  suite's coverage; that is acceptable, because suite coverage is
  visible, evolvable, and can be expanded as the spec matures.
\end{rationale}

\section{Reference Algorithm}
\label{sec:solver:reference-algo}

The reference algorithm is explanatory and diagnostic. Its purpose
is to disambiguate constraint interactions where prose alone might
leave a choice open, and to generate the expected outputs published
with the reference suite. It is \emph{not} the normative algorithm
and \emph{not} the aesthetic floor.

\begin{requirement}
  \label{req:solver:algorithm-independence}
  Conforming implementations are \emph{not} required to use the
  reference algorithm or to reproduce its layouts, except where
  a reference-suite entry explicitly fixes an expected behavior
  via a fixed-expectation test. Implementations \MAY{} choose
  any algorithm satisfying the contract: hard constraints,
  determinism, diagnostics, reference-suite thresholds.
\end{requirement}

The reference algorithm is described in a non-normative companion
document. A natural starting point uses the Cassowary linear-
arithmetic constraint solver for the linear layer (positions,
alignments, containment) together with a separate numerical
optimization layer for aesthetic constraints (beam slopes, slur
curvatures) that are nonlinear in nature. Hybrid approaches and
alternative architectures (differentiable layout, constraint
programming, learned methods) are all conforming choices.

\section{Forward References}

\begin{itemize}
  \item End-to-end performance requirements, including frame budgets
    that the solver participates in, are in Chapter~\ref{ch:perf}.
    Solver performance is measured in the context of the full
    layout pipeline.
  \item The Quality Metric Catalog companion document specifies
    the exact normalization functions for each metric, the default
    tie-breaking weights, and the per-tier thresholds.
  \item The Reference Suite companion document specifies the
    test scores, per-tier inclusion, and per-tier metric
    thresholds.
  \item The reference-algorithm companion document is non-normative
    and describes one starting-point implementation.
\end{itemize}

% ===========================================================================
\chapter{Performance Requirements}
\label{ch:perf}

This chapter specifies performance contracts for the \emph{core
layer}. Targets are normative; they bound what conforming
implementations promise and what tests measure for conformance.
The chapter restates targets introduced informally in earlier
chapters and adds the measurement methodology.

\section{Core/Product Boundary}
\label{sec:perf:boundary}

Chapter~\ref{ch:intro} excluded the editor UI, the audio engine,
the renderer beyond the layout IR, and platform integrations from
the core's scope. Many performance properties of a working
notation system depend on those out-of-scope components: first
interactive frame after launch involves UI bootstrapping, edit
response involves UI event handling, scrolling and selection
involve a real renderer, audio acquisition involves a working
audio engine. The core cannot promise end-to-end timings for
behavior it does not own.

This chapter therefore distinguishes three obligation classes:

\begin{description}
  \item[Core obligations.] The core data structures, layout IR,
    solver, and file format \MUST{} admit the stated bounds. The
    core's own work (reduction, layout, file writes, snapshot
    acquisition) \MUST{} fit within the budgets allocated to it
    in the end-to-end frame budget. The core \MUST{} expose APIs
    suitable for each downstream component to do its own work
    inside its remaining budget (e.g., a lock-free immutable
    snapshot API for audio engines, an incremental invalidation
    interface for renderers, a cancellable solve API for UIs).
  \item[Product obligations.] End-to-end UI, audio, and rendering
    targets (first interactive frame after launch, end-to-end
    edit-to-pixel latency, end-to-end scroll smoothness,
    audio-thread snapshot acquisition latency under load) are the
    obligation of the product layers built atop the core. Those
    targets are specified in product-specific companion documents,
    not here.
  \item[Performance Reference Suite obligations.] The Performance
    Reference Suite (Appendix~\ref{app:deferred}) measures core
    behavior only: layout time per system, file write/read time,
    reduction throughput, solver budget consumption, snapshot
    acquisition latency. End-to-end UI traces are out of scope
    for the core's reference suite.
\end{description}

\begin{requirement}
  \label{req:perf:core-product-budget-boundary}
  Where this chapter states a performance target, the target is a
  \emph{core obligation} unless explicitly labeled product
  obligation. Core obligations \MUST{} be satisfiable by a
  conforming core implementation independent of the product
  layers built atop it; product-layer code is permitted to
  consume budget on top of the core's portion, but the core's
  own portion \MUST{} fit within the stated bound.
\end{requirement}

\section{Design Principles}
\label{sec:perf:principles}

\begin{description}
  \item[Targets are part of the contract.] Core performance
    requirements are not aspirations. A conforming core
    implementation \MUST{} meet the stated targets on the
    reference hardware profile under the stated measurement
    conditions.
  \item[Frame budgets dominate.] User experience is ultimately
    governed by frame timing. The core's portion of each frame
    budget is stated explicitly; the product layer's portion is
    specified separately.
  \item[Reference hardware profile.] All targets are stated
    against a specified hardware profile. Implementations \MAY{}
    document additional profiles (lower-end mobile, higher-end
    workstation); the reference profile is the basis of
    conformance.
  \item[Measurement methodology is normative.] Without specified
    measurement methodology, performance targets are
    unenforceable. This chapter specifies how each target is
    measured.
  \item[Tail latency matters.] Average performance is
    insufficient; targets are stated as percentiles (typically
    p99 or worst-case) to reflect what users actually experience.
\end{description}

\section{Reference Hardware Profile}
\label{sec:perf:hardware}

The reference hardware profile defines the platform against which
performance conformance is measured. Implementations \MAY{} target
additional profiles; this profile is the basis of conformance.

\begin{lstlisting}
Reference Hardware Profile (2026 baseline):

  Desktop reference:
    - Apple silicon (M-series) or x86-64 equivalent
    - 8 performance cores (or equivalent throughput)
    - 16 GB RAM
    - NVMe SSD storage
    - Discrete or integrated GPU with Metal/Vulkan support
    - Display: 1440p at 60 Hz minimum

  Tablet reference:
    - iPad Pro M-series or equivalent
    - 8 GB RAM minimum
    - SSD storage
    - Display: 120 Hz Pro Motion or equivalent
\end{lstlisting}

\begin{requirement}
  \label{req:perf:reference-hardware}
  Conforming implementations \MUST{} meet the targets in this chapter
  on the reference hardware profile. Implementations \MAY{} miss
  targets on hardware below the reference profile (e.g., older
  devices, low-end embedded targets); such gaps \MUST{} be documented
  in the implementation's published performance contract.

  The reference profile is updated periodically as commodity hardware
  evolves. The version of the reference profile against which an
  implementation conforms \MUST{} be declared.
\end{requirement}

\section{Frame Budgets}
\label{sec:perf:frame}

\subsection{Interactive Edit Latency}

\begin{requirement}
  \label{req:perf:single-system-edit-latency}
  The \emph{core's portion} of a single-system edit (operation
  envelope construction, reduction, incremental layout through
  \texttt{ResolvedLayoutIR}) \MUST{} complete within one frame
  (16.7\,ms at 60\,Hz) at the p99 percentile, measured on the
  reference hardware profile. This includes the case where the
  affected system is on a 100-page orchestral score: the core's
  incremental layout machinery (Chapter~\ref{ch:layout-ir}) and
  the constraint solver's incremental contract
  (Chapter~\ref{ch:solver}) jointly satisfy this requirement.

  On 120\,Hz displays (tablet reference), the core's portion
  \SHOULD{} complete within 8.3\,ms at p99. This is not required
  for conformance but is the target for first-class tablet
  performance.

  End-to-end edit-to-pixel latency (input handling, hit testing,
  render submission, display flip) is a product-layer obligation.
  The core's contribution is bounded by this target so that the
  product layer has the remaining frame budget to do its work.
\end{requirement}

\subsection{Multi-System Edits}

\begin{requirement}
  \label{req:perf:multi-system-edit-latency}
  Edits affecting up to ten systems (e.g., a transposition over
  a substantial passage) \MUST{} complete within four frames
  (66.7\,ms at 60\,Hz) at p99. The four-frame budget gives the
  editor headroom to display a transient ``thinking'' indication
  without violating perception of responsiveness.
\end{requirement}

\subsection{Full Re-Layout}

\begin{requirement}
  \label{req:perf:full-relayout-latency}
  Full re-layout of a 100-page orchestral score (triggered, for
  example, by a tuning system change applied score-wide) \SHOULD{}
  complete within two seconds at p99. This is a soft target; the
  hard target on full re-layout is that the operation be
  cancellable: a user beginning a different edit \MUST{} cause the
  full re-layout to abort and the new edit to proceed.
\end{requirement}

\section{Cold Open}
\label{sec:perf:open}

\begin{requirement}
  \label{req:perf:cold-open-latency}
  The \emph{core's portion} of cold open of a 100-page orchestral
  score from a \texttt{.musc} bundle \MUST{} reach
  \emph{first-system-ready} state within one second on the
  reference hardware profile.

  First-system-ready is defined as: the prelude has been read,
  the active superblock and manifest validated, the canonical
  base snapshot (if any) acquired, post-frontier envelopes
  reduced for at least the first system's coverage, and the
  first system's \texttt{ResolvedLayoutIR} produced. Rendering,
  hit-testing, and UI bootstrapping that build atop this state
  are product-layer obligations.

  Background loading of remaining systems \MAY{} continue after
  first-system-ready and \MUSTNOT{} block user interaction
  served from already-loaded state.
  not-yet-visible systems is permitted; full-score loading need not
  complete in one second.
\end{requirement}

\subsection{Cold Open Measurement}

Cold open is measured from process start (or, for warm processes
opening a new score, from the open API invocation) to the first
frame containing fully rendered content for the user's initial
viewport.

\subsection{Streaming Read Discipline}

\begin{requirement}
  \label{req:perf:streaming-cold-open}
  The cold open target \MUST{} be achieved through streaming reads
  per Chapter~\ref{ch:format}: the implementation \MUSTNOT{} load
  the full score into memory before displaying the first system.
  Implementations failing to meet the target by eager loading are
  non-conforming.
\end{requirement}

\section{Memory Budgets}
\label{sec:perf:memory}

\begin{requirement}
  \label{req:perf:core-memory-budget}
  Steady-state resident memory for a 100-page orchestral score
  \SHOULD{} remain below 500\,MB on the reference hardware profile,
  exclusive of evictable caches. Evictable caches \MAY{} push total
  resident memory higher; they \MUST{} be evicted under memory
  pressure without compromising correctness.
\end{requirement}

\subsection{Memory Categories}

For measurement, memory is partitioned into:

\begin{description}
  \item[Core data.] The score graph, attachments, indexes, and
    operation envelope set. Cannot be evicted without losing
    state.
  \item[Layout cache.] The cached IR stages. Evictable; rebuilds on
    demand but at performance cost.
  \item[Glyph and font caches.] Glyph render data and metric caches.
    Evictable; small-to-moderate size.
  \item[Audio buffers and reference tracks.] Audio data loaded for
    playback or reference. Evictable.
  \item[Implementation overhead.] Runtime, allocator metadata,
    OS-level mappings.
\end{description}

\begin{requirement}
  \label{req:perf:memory-category-reporting}
  Implementations \MUST{} expose the core-data memory consumption
  separately from evictable caches. Reporting tools \MUST{} be able
  to distinguish them. The 500\,MB target applies to core data only.
\end{requirement}

\subsection{Per-Object Memory Discipline}

\begin{requirement}
  \label{req:perf:per-object-memory-layout}
  Implementations \MUSTNOT{} embed glyph metrics in IR objects
  (Chapter~\ref{ch:layout-ir}). Implementations \MUSTNOT{} embed
  spelling-attachment payloads on pitch objects
  (Chapter~\ref{ch:pitch}); attachments are externally indexed.
  These discipline requirements bound per-object memory at known
  upper limits.
\end{requirement}

\section{Core Data Structures for Audio Thread Use}
\label{sec:perf:audio}

The audio engine is outside the core specification. The core's
obligation is to expose data structures that permit real-time
audio safety; audio-engine performance is the obligation of the
product audio layer (Section~\ref{sec:perf:boundary}).

\begin{requirement}
  \label{req:perf:audio-thread-data-access}
  The core's read-only data structures consumed by an audio
  thread (pitch resolution, tempo map lookup, event arena
  queries) \MUST{} be designed for allocation-free, lock-free
  access. Specifically:

  \begin{itemize}
    \item Pitch frequency resolution given a tuning system \MUST{}
      be performable without allocation.
    \item Tempo-map conversion (musical to wall-clock and vice
      versa) \MUST{} be performable without allocation.
    \item Time-indexed event lookup for playback \MUST{} be
      performable without allocation.
  \end{itemize}

  These are properties of the core's data layouts, not of any
  particular audio engine. A product-layer audio engine that
  consumes them is responsible for its own real-time discipline
  (no syscalls on the audio thread, no priority inversion, no
  unbounded loops); the core only promises that its data
  structures do not preclude that discipline.
\end{requirement}

\subsection{Snapshot API}

\begin{requirement}
  \label{req:perf:lock-free-snapshots}
  The core \MUST{} provide an immutable-snapshot API: the
  ability for any thread to acquire a read-only view of the
  score graph at a consistent point in time, suitable for use
  by an audio engine. Snapshot acquisition by the snapshot
  reader \MUST{} be lock-free in the sense that it never blocks
  on edit-thread progress: an edit in flight \MUSTNOT{} prevent
  the snapshot reader from obtaining a valid (possibly older)
  snapshot.

  Snapshot \emph{release} \MAY{} be deferred to a non-audio
  thread. The snapshot mechanism is implementation-defined
  (double-buffered immutable copies, persistent data structures,
  reference-counted immutable trees, hazard pointers); only the
  externally observable lock-free acquisition property is
  normative here.

  End-to-end audio-thread acquisition latency under load is a
  product-layer obligation; the core's contribution is bounded
  by the lock-free guarantee.
\end{requirement}

\section{Constraint-Solver Performance Targets}
\label{sec:perf:solver}

\begin{requirement}
  \label{req:perf:solver-latency}
  The constraint solver (Chapter~\ref{ch:solver}) \MUST{} satisfy:

  \begin{itemize}
    \item Incremental solving of a single-system edit completes
      within 10\,ms at p99 on the reference hardware profile.
      (This is a tighter target than the overall frame budget; the
      remaining 6.7\,ms accommodates the rest of the pipeline.)
    \item Cold solving of a 100-page orchestral score completes
      within two seconds at p99.
    \item Incremental solving propagation is bounded as documented
      in the implementation's published contract.
  \end{itemize}
\end{requirement}

\section{File Format Performance}
\label{sec:perf:format}

\begin{requirement}
  \label{req:perf:file-operation-latency}
  File format operations \MUST{} satisfy:

  \begin{itemize}
    \item Bundle write of a typical edit (one or more operation
      envelopes appended to the operation-envelope block stream,
      manifest rewrite, superblock flip) completes within
      50\,ms at p99 on the reference hardware profile.
    \item Bundle read of the manifest and bootstrap chunks
      (sufficient for first interactive frame) completes within
      200\,ms at p99 for a 100-page orchestral score.
    \item Operation-envelope reduction rate \MUST{} exceed
      10{,}000 envelopes per second on the reference hardware
      profile, measured during cold reduction from a fresh
      canonical base.
  \end{itemize}
\end{requirement}

\section{Measurement Methodology}
\label{sec:perf:methodology}

\begin{requirement}
  \label{req:perf:measurement-methodology}
  Performance conformance \MUST{} be measured using:

  \begin{itemize}
    \item A reference test corpus including the conformance suite
      from Chapter~\ref{ch:solver}, augmented with edit traces
      representative of typical authoring sessions.
    \item Measurements at the p99 percentile across at least 1000
      iterations per scenario.
    \item Warm runs after a cold warm-up phase, unless the target
      is explicitly cold (cold open, cold load).
    \item Release builds with full optimization enabled, but without
      fast-math or precision-compromising flags.
    \item The reference hardware profile, with no other significant
      load on the system during measurement.
  \end{itemize}

  Implementations \MUST{} publish their measurement results against
  this methodology for conformance claims.
\end{requirement}

\section{Regression Testing}
\label{sec:perf:regression}

\begin{requirement}
  \label{req:perf:performance-regression-gate}
  Conforming implementations \SHOULD{} maintain continuous
  performance regression testing against the reference corpus, with
  alerts for regressions exceeding 10\% on any target. The
  regression suite \SHOULD{} be a precondition for releases.
\end{requirement}

\section{Forward References}

\begin{itemize}
  \item The full performance reference corpus is delivered as a
    companion specification, the Performance Reference Suite.
  \item Audio engine performance targets (latency, jitter, stream
    discipline) are specified in the audio engine specification,
    outside this document.
  \item Network and synchronization performance for the collaboration
    layer is specified in the collaboration specification, outside
    this document.
\end{itemize}

% ===========================================================================
\chapter{Extension Points}
\label{ch:extension}

This chapter consolidates the extension points specified across
prior chapters into a single reference. Each extension point admits
implementation-defined or plugin-contributed content that the core
treats opaquely, subject to the contracts stated here.

The chapter is a navigation aid rather than a source of new
requirements: every extension point referenced here is defined in
its originating chapter, and the originating definition is
authoritative. This chapter pulls the points together for plugin
authors and future specification editors.

\section{Design Principles}
\label{sec:ext:principles}

\begin{description}
  \item[Extensibility without forking.] Every extension point
    permits content to be added without modifying the core
    specification or its implementations. Plugins extend; they do
    not fork.
  \item[Identifiers, not strings.] Extension points are identified
    by typed registry identifiers, not free-form strings. Registries
    are explicit, versioned, and discoverable.
  \item[Contracts, not contents.] Extension points specify what an
    extension \emph{must satisfy} (interfaces, invariants, reduction
    rules, performance bounds) rather than what an extension
    \emph{must contain}.
  \item[Capability-bound at the runtime.] Plugin runtime
    permissions (read score, write score, network, filesystem) are
    declared per extension and granted per user. This is enforced
    by the plugin runtime, which is outside this specification; the
    extension points are designed to admit such enforcement.
  \item[Forward compatibility.] Extensions referenced in a stored
    score \MUST{} be preserved across reads and writes even when the
    reading implementation does not understand them (per
    Chapter~\ref{ch:format}). Extensions thus survive distribution
    even where they cannot be exercised.
\end{description}

\section{Notation Grammar Extensions}
\label{sec:ext:grammar}

The notation grammar layer supports extension at several points:

\begin{description}
  \item[\texttt{PositionStructure::Registered}] (Chapter~\ref{ch:tuning}).
    Custom position structures for pitch spaces whose organization
    does not fit the chromatic, diatonic-over-chromatic, or JI
    lattice families. Used for maqam, gamelan, raga, and other
    non-Western grammars.
  \item[\texttt{IntervalAlgebra::Registered}] (Chapter~\ref{ch:tuning}).
    Custom interval algebras for grammar-specific arithmetic.
  \item[\texttt{TranspositionBehavior::Registered}]
    (Chapter~\ref{ch:tuning}). Custom transposition rules.
  \item[\texttt{SpellingAlgorithmId}] (Chapter~\ref{ch:pitch}).
    Registered spelling-inference algorithms for the spelling
    pre-pass.
  \item[\texttt{NominalRegistry}] (Chapter~\ref{ch:tuning}).
    Score-level or grammar-level extensions to the catalog of
    named position-letters.
  \item[\texttt{AccidentalRegistry}] (Chapter~\ref{ch:tuning}).
    Score-level extensions to the accidental catalog, subject to
    the registry's \texttt{extensible} flag.
\end{description}

\section{Tuning System Extensions}
\label{sec:ext:tuning}

\begin{description}
  \item[\texttt{TuningFunctionId}] (Chapter~\ref{ch:tuning}).
    Registered procedural tuning functions (historical
    temperaments, custom meantone variants, generative tunings).
  \item[\texttt{AdaptiveTuningFunctionId}] (Chapter~\ref{ch:tuning}).
    Registered adaptive tuning functions taking harmonic context as
    input.
  \item[\texttt{TuningFileFormat}] (Chapter~\ref{ch:tuning}).
    Recognized external tuning file formats (Scala, MIDI Tuning
    Standard, additional formats via registry).
  \item[\texttt{CompatibilityMapping}] (Chapter~\ref{ch:tuning}).
    Declared compatibility between pitch spaces and tuning systems
    whose underlying spaces differ.
\end{description}

\section{Glyph and Font Extensions}
\label{sec:ext:glyphs}

\begin{description}
  \item[\texttt{CustomGlyphId}] (Chapter~\ref{ch:tuning}).
    Custom glyphs not in SMuFL, defined per score or per plugin.
  \item[\texttt{Composite} glyph references] (Chapter~\ref{ch:tuning}).
    Composite glyphs constructed from multiple primary glyph
    references.
  \item[\texttt{ModificationRegistryId}] (Chapter~\ref{ch:tuning}).
    Registered pitch-space modifications for grammar-defined
    accidentals.
\end{description}

\section{Score Graph Extensions}
\label{sec:ext:graph}

\begin{description}
  \item[\texttt{ExtensionOp}] (Chapter~\ref{ch:semops}).
    Implementation-defined operation kinds extending the core's
    semantic operations. Extension operations \MUST{} satisfy the
    operation framework's requirements (preconditions, effects,
    inverse, reduction rule, re-anchoring).
  \item[\texttt{CustomMarkerId}] (Chapter~\ref{ch:graph}).
    Custom marker kinds beyond the core's standard catalog.
  \item[\texttt{CustomAnalyticalKindId}] (Chapter~\ref{ch:graph}).
    Custom analytical annotation kinds.
  \item[\texttt{Custom event variants}] (Chapter~\ref{ch:graph}).
    The event taxonomy is closed in the core, but the \texttt{Group}
    layout object and the \texttt{Registered} fields throughout
    permit composition of core types into extended structures.
\end{description}

\section{Layout IR Extensions}
\label{sec:ext:layout}

\begin{description}
  \item[\texttt{SynthesisRegistryId}] (Chapter~\ref{ch:layout-ir}).
    Custom synthesis kinds for engraver-generated objects.
  \item[\texttt{ConstraintRegistryId}] (Chapter~\ref{ch:layout-ir}).
    Registered layout constraint kinds extending the standard
    constraint families.
  \item[\texttt{MappingRegistryId}] (Chapter~\ref{ch:graph}).
    Custom parameter mappings for graphic-object time bindings.
  \item[\texttt{Custom view kinds}] (Chapter~\ref{ch:graph}).
    The \texttt{ViewKind::Custom} variant admits user-defined
    view kinds.
\end{description}

\section{Constraint-Solver Extensions}
\label{sec:ext:solver}

\begin{description}
  \item[\texttt{LayoutConstraint::Registered}]
    (Chapter~\ref{ch:layout-ir}). Custom constraint families
    consumed by the solver. Extension constraints \MUST{} declare
    their strength. When soft, extension constraints contribute
    to the quality vector through an extension metric subject to
    the \texttt{NormalizedMetric} discipline
    (Chapter~\ref{ch:solver}) and to the tier-specific reference-
    suite thresholds.
\end{description}

\section{File Format Extensions}
\label{sec:ext:format}

\begin{description}
  \item[\texttt{FormatProfile::Custom}] (Chapter~\ref{ch:format}).
    Profile definitions beyond Full, ReadOnly, and Lite.
  \item[\texttt{Self-describing extension envelope}]
    (Chapter~\ref{ch:format}). Plugin-contributed and
    forward-compatible content is carried in self-describing
    extension records that the core preserves opaquely.
\end{description}

\section{Extension Registry Contract}
\label{sec:ext:registry}

Every extension registry shares a common structural contract.

\begin{lstlisting}[language=Rust]
pub trait ExtensionRegistry {
    type Id;
    type Definition;

    /// Resolve an identifier to its definition.
    fn resolve(&self, id: Self::Id) -> Option<&Self::Definition>;

    /// Enumerate all registered definitions.
    fn enumerate(&self) -> Vec<&Self::Definition>;

    /// Version of the registry. Score files declare which version
    /// their references target.
    fn version(&self) -> RegistryVersion;
}
\end{lstlisting}

\begin{requirement}
  \label{req:ext:registry-versioning}
  Every extension registry \MUST{} be versioned. Stored references
  to registered definitions \MUST{} include the registry version
  they target. Loading a score with references to registry versions
  newer than the reader understands \MUST{} preserve those references
  per the forward-compatibility rules of Chapter~\ref{ch:format}.
\end{requirement}

\subsection{Identifier Stability}

\begin{requirement}
  \label{req:ext:registry-id-stability}
  Registry identifiers \MUST{} be stable across registry versions
  within a major version. Renumbering, reassignment, or repurposing
  of identifiers is forbidden within a major version. Removal of a
  definition \MUST{} reserve its identifier; the identifier \MUSTNOT{}
  be reassigned to a different definition.
\end{requirement}

\subsection{Discoverability}

Registries \MUST{} support enumeration: an implementation can list
all known definitions and their identifiers. This enables tooling
(plugin browsers, score migration assistants, conformance testers)
to operate against unknown registries by inspection.

\section{Plugin Runtime Considerations}
\label{sec:ext:runtime}

The plugin runtime is outside this specification, but extension
points are designed to admit a capability-bound runtime as described
in the feature compendium:

\begin{itemize}
  \item Extensions are isolatable: an extension's failure does not
    crash the host.
  \item Extensions are introspectable: their declared capabilities
    and inputs are inspectable before execution.
  \item Extensions are revocable: an extension can be disabled
    without invalidating the scores that reference it.
\end{itemize}

\begin{requirement}
  \label{req:ext:disabled-extension-preservation}
  Extension points \MUST{} be designed such that disabling an
  extension does not invalidate stored content referencing it. The
  content remains in the score, opaque to the host, until the
  extension is re-enabled or the references are explicitly removed
  by the user.
\end{requirement}

\section{Conformance and Extensions}
\label{sec:ext:conformance}

Extensions do not affect core conformance directly, but they
interact with it:

\begin{requirement}
  \label{req:ext:extension-conformance}
  A conforming implementation:

  \begin{itemize}
    \item \MUST{} support every extension point as a structural
      slot (i.e., scores containing extension references are
      readable and writable).
    \item \MUSTNOT{} require any specific extension to be installed
      for core conformance.
    \item \MAY{} require specific extensions for specific features
      (e.g., a particular plugin for a particular grammar).
    \item \MUST{} clearly distinguish between core conformance and
      extension-dependent features in its documentation.
  \end{itemize}
\end{requirement}

\section{Forward References}

\begin{itemize}
  \item The plugin runtime (WebAssembly-based, capability-bound)
    is specified in the Plugin Runtime specification, outside this
    document.
  \item The complete extension registry catalog is delivered as a
    companion document.
  \item Implementation-defined extension contributions are governed
    by the implementation's published extension specification.
\end{itemize}

% ===========================================================================
\appendix

\chapter{Intentionally Deferred Types and Specifications}
\label{app:deferred}

This appendix catalogs types, concepts, and specifications that
the core specification refers to but does not fully define. Each
entry is \emph{deliberately external}: it is the responsibility
of a named companion specification or a future revision of this
document. The purpose of this appendix is to distinguish
deliberately external content from accidentally missing content,
so that reviewers and implementers can quickly tell which
unresolved references represent open work and which represent
boundaries to other documents.

\section{Companion Specifications}
\label{sec:deferred:companions}

The following are companion specifications referenced throughout
the core. None has been delivered as of this revision; each is
named here with the chapter(s) that depend on it.

\begin{description}
  \item[Binary Format companion.] Byte-level encoding of all
    canonical chunks: varint conventions, exact field widths,
    record layouts, endianness rules, string encoding details.
    Referenced from Chapter~\ref{ch:format}. The on-disk
    representations specified in that chapter are structural;
    the Binary Format specifies how the bits are laid out.
  \item[Text Projection companion.] The canonical s-expression
    text projection of the bundle: grammar, deterministic
    serialization rules, base64 encoding for opaque payloads,
    handling of extension data, round-trip semantics with
    binary form. Referenced from Chapter~\ref{ch:format}
    Section~\ref{sec:format:textproj}.
  \item[Profile Conformance specification.] The complete
    catalog of file-format profiles (Full, ReadOnly, Lite, plus
    any normative custom profiles), with their constraints,
    permitted compression algorithms, maximum block sizes,
    required extensions, and conformance test discipline.
    Referenced from Chapter~\ref{ch:format}
    Section~\ref{sec:format:profiles}.
  \item[Quality Metric Catalog.] Per-metric normalization
    functions mapping raw measurements to \texttt{NormalizedMetric}
    values, default tie-breaking weights, per-tier metric
    thresholds, and the formal definition of each quality
    metric in the normative metric set. Referenced from
    Chapter~\ref{ch:solver}.
  \item[Reference Suite.] The collection of test scores against
    which solver conformance is established, with per-tier
    inclusion, per-tier metric thresholds, and any
    fixed-expectation tests. Versioned with this specification.
    Referenced from Chapter~\ref{ch:solver}
    Section~\ref{sec:solver:conformance}.
  \item[Performance Reference Suite.] The collection of edit
    traces and reference scores against which performance
    conformance is established. Referenced from
    Chapter~\ref{ch:perf}.
  \item[Operation Catalog companion.] The complete list of
    primitive operations and compound operations supported by
    the core, each with full preconditions, effects, inverse,
    and reduction rule. The representative operations in
    Chapter~\ref{ch:semops} illustrate the contract; the full
    catalog is delivered separately.
  \item[Reference Algorithm companion.] A non-normative
    description of one starting-point implementation of the
    constraint solver. Explanatory and diagnostic only;
    conforming implementations are not required to use it.
    Referenced from Chapter~\ref{ch:solver}
    Section~\ref{sec:solver:reference-algo}.
\end{description}

\section{Externally Provided Components}
\label{sec:deferred:external}

The following are real components on which the core depends but
which are owned by other specifications outside the
Epiphany core scope.

\begin{description}
  \item[SMuFL glyph catalog.] The Standard Music Font Layout
    provides the normative glyph repertoire. Referenced from
    Chapter~\ref{ch:tuning} and the layout chapters. SMuFL is
    maintained by the W3C Music Notation Community Group; the
    core spec declares its dependency on the SMuFL version
    targeted but does not redefine SMuFL.
  \item[Font metric provider.] Implementations resolve glyph
    geometric data (bounding boxes, anchor points, advance
    widths) through a font metric provider. The provider's
    interface is part of the implementation contract; specific
    font metrics are external content.
  \item[Audio engine specification.] Out of scope here per
    Chapter~\ref{ch:intro}. The core exposes data structures
    designed for allocation-free, lock-free audio-thread access
    (Chapter~\ref{ch:perf}); the audio engine's internals,
    latency obligations, and stream discipline are specified
    separately.
  \item[Plugin runtime specification.] A capability-bound,
    sandboxed runtime hosts extension code. The runtime is
    out of scope here; the extension points
    (Chapter~\ref{ch:extension}) are designed to admit such a
    runtime.
  \item[Collaboration transport.] The wire protocol by which
    operation envelopes are exchanged between replicas is
    out of scope here. Chapter~\ref{ch:semops} specifies what
    must be exchanged and how it is reduced; the transport
    layer is separate.
  \item[User interface and editor.] Editor commands, selection
    semantics, keyboard shortcuts, gesture grammar, and
    UI-level interaction patterns are out of scope. The core
    exposes a programmatic operation API; UI mapping is
    layered above.
\end{description}

\section{Open Canonical Algorithms}
\label{sec:deferred:open-algos}

The following algorithms affect canonical state but are not yet
specified normatively, profile-declared, or marked
non-canonical. Per Appendix~\ref{app:determinism}
Section~\ref{sec:det:open}, each \MUST{} receive one of those
dispositions before the document can claim full canonical
score determinism across implementations.

\begin{description}
  \item[Spelling pre-pass algorithm.] The algorithm that
    computes default spellings from key signature, harmonic
    context, and notation grammar. Referenced from
    Chapter~\ref{ch:pitch}. Several established algorithms
    exist (Longuet-Higgins, Temperley pitch-spelling preference
    rules, and others); the choice of normative algorithm or
    profile-declared algorithm set is open.
  \item[Notational decomposition algorithm.] The algorithm that
    breaks a sounding duration into notehead values, augmentation
    dots, and ties under metric context. Referenced from
    Chapter~\ref{ch:time}. Conservatoire engraving manuals
    describe the conventions verbosely; the formalization is
    open.
  \item[Tempo curve integration.] The numerical integration
    algorithm for converting musical position to wall-clock
    position under \texttt{TempoShape::Curve}. Referenced from
    Chapter~\ref{ch:time}. Possible candidates include Romberg
    integration and adaptive Gauss-Kronrod; the normative
    choice is open.
  \item[Wallclock-to-musical root-finding.] The root-finding
    algorithm for the inverse of tempo curve integration.
    Referenced from Chapter~\ref{ch:time}. Possible candidates
    include bisection with monotonicity guarantees and
    Brent's method; the normative choice is open.
\end{description}

\section{Extension Registry Catalogs}
\label{sec:deferred:registries}

The extension registries enumerated in
Chapter~\ref{ch:extension} (notation grammar extensions,
tuning system extensions, glyph and font extensions, score
graph extensions, layout IR extensions, solver extensions,
file format extensions) are referenced by typed registry
identifiers in the core. The concrete catalogs of registered
extensions (the contents of each registry) are delivered as
separate, versioned registry documents. Each registry document
specifies its identifier numbering, the version under which
each registered identifier was added, and the registered
content's contract.

\section{Support-Type Identity Semantics}
\label{sec:deferred:support-types}

Several identifier-like support types appear throughout the
specification by name and shape, but their detailed encoding is
delivered by companion specifications. Their semantic identity
rules are normative here.

\begin{longtable}{p{3.4cm} p{3.6cm} p{6.6cm}}
  \toprule
  \textbf{Type} & \textbf{Source} & \textbf{Semantic rule} \\
  \midrule
  \endhead

  \texttt{FileUuid} &
  128-bit random at bundle creation &
  Physical bundle identity. Fixed for the lifetime of the
  physical bundle; preserved by filesystem byte-copy; freshly
  minted on Save As, export-as-new-bundle, and derivative-work
  creation. \\

  \texttt{DocumentId} &
  128-bit random at logical-work creation &
  Logical work identity. Stable across Save As copies intended
  to remain versions of the same work. May persist across forks
  at user discretion. \\

  \texttt{LineageId} &
  128-bit random at fork-with-ancestry &
  Shared-ancestor identity. Optional. Multiple documents may
  share a \texttt{LineageId} declaring common ancestry for
  version-control or genealogy purposes. \\

  \texttt{ManifestId} &
  Content-derived from manifest payload &
  Unique per manifest version. Derived as
  \texttt{trunc128(BLAKE3("MUSCMNIF" || document\_id ||
  generation || manifest\_body))} with the \texttt{manifest\_id}
  field excluded from \texttt{manifest\_body}; the derivation is
  ratified normatively in
  Section~\ref{req:format:manifest-id}. \\

  \texttt{SnapshotId} &
  Content-derived from snapshot payload and frontier &
  Stable for a snapshot's root and causal frontier. Two snapshots
  with identical materialized state at identical frontiers have
  identical ids. The Binary Format companion defines the exact
  derivation. \\

  \texttt{BlobId} &
  BLAKE3 of blob content with domain tag \texttt{"MUSCBLOB"} &
  Content hash of the blob's uncompressed payload, identical to
  the blob's \texttt{ContentHash}. \\

  \texttt{AuthorId} &
  Implementation-defined; persists across replica reseatings &
  Identifies the human author of operations, distinct from
  \texttt{ReplicaId} (which identifies the device or session).
  One author may operate multiple replicas; one replica may host
  one author at a time. The collaboration transport companion
  defines authentication. \\

  \texttt{ReductionAlgorithmVersion} &
  Semantic version &
  Selects canonical reduction semantics. Two replicas with
  different \texttt{ReductionAlgorithmVersion} may produce
  different canonical states from the same operation set; the
  active superblock declares the version under which the
  bundle's canonical base was materialized
  (Section~\ref{sec:format:bundle}). \\

  \texttt{SolverProfile} &
  Registered profile identifier &
  Selects the solver's hard-constraint set, normalized-metric
  thresholds, tie-breaking weights, and active extension
  catalog (Chapter~\ref{ch:solver}). The Quality Metric Catalog
  companion defines the registered profile catalog. \\

  \texttt{ProfileId} (file format) &
  Registered profile identifier &
  Selects the file-format conformance subset (Full, ReadOnly,
  Lite, or registered Custom). The Profile Conformance companion
  defines the registered profile catalog. \\

  \texttt{ContentHash} &
  BLAKE3-256 output (32 bytes) &
  Cryptographic content hash. The canonical hash for all
  content-addressed objects in the bundle. \\

  \texttt{ChunkId} &
  Newtype around \texttt{ContentHash} &
  Same 32 bytes as the underlying \texttt{ContentHash}; the
  type distinction makes role visible at use sites. \\

  \texttt{OperationId} &
  64-bit \texttt{ReplicaId} + 64-bit counter, both
  randomly initialized; counter monotonic &
  Stable identity of an operation, set by the authoring replica
  at the moment of authoring. Never reassigned. \\

  \texttt{TransactionId} &
  128-bit, author-minted &
  Identifies a transaction declared by a
  \texttt{DeclareTransaction} envelope. Member envelopes
  reference this id in their envelope \texttt{transaction}
  field. \\

  \texttt{ConflictId} &
  Content-derived via BLAKE3 truncation with domain tag
  \texttt{"MUSCCONF"} &
  Identifier of a conflict record. Derived from canonical kind,
  sorted causing operations, sorted affected objects.
  Section~\ref{sec:semops:conflict-id}. \\

  \texttt{Typed identifiers} (EventId, PitchId, VoiceId,
  StaffId, etc.) &
  64-bit \texttt{ReplicaId} + 64-bit counter; or
  \texttt{ReplicaId::SYSTEM\_DERIVED} + BLAKE3-derived counter
  for system-derived identifiers &
  Stable identity within their typed namespace.
  Section~\ref{sec:graph:system-derived}. \\

  \bottomrule
\end{longtable}

\begin{requirement}
  \label{req:deferred:identity-semantics}
  Implementations \MUST{} respect the semantic identity rules
  declared above. They \MAY{} freely choose the bit-level
  encoding of each type (the Binary Format companion specifies
  the canonical wire encoding); they \MAY NOT{} reinterpret the
  identity semantics. In particular, two distinct works \MUSTNOT{}
  share a \texttt{DocumentId}; two byte-distinct chunks \MUSTNOT{}
  share a \texttt{ContentHash}; two distinct operations \MUSTNOT{}
  share an \texttt{OperationId}.
\end{requirement}

\section{Distinction from Open Questions}
\label{sec:deferred:vs-open}

This appendix lists content that is \emph{deliberately external}
to the core specification: it is owned by a companion document,
a future revision, or a third-party specification.

It does \emph{not} list ``open questions'' in the core's own
design space. Those are marked with the
\textsf{\textit{Open Question}} callout where they appear in
the body of the text, and they represent unresolved decisions
\emph{within} the core specification's scope. The two
categories overlap only at the four open canonical algorithms
in Section~\ref{sec:deferred:open-algos}, which are unresolved
in the core text \emph{and} whose resolution will likely take
the form of a companion specification or a profile declaration.

\chapter{Glossary}
\label{app:glossary}

This glossary defines key terms used throughout the specification.
Definitions are drawn from their introducing chapters; the chapter
reference points to the authoritative definition.

\begin{description}
  \item[Acoustic pitch] (Chapter~\ref{ch:pitch}). The frequency
    realization of a pitch under an active tuning system. One of the
    two intrinsic layers of a \texttt{Pitch}.
  \item[Adaptive tuning] (Chapter~\ref{ch:tuning}). A tuning system
    whose frequency resolution depends on harmonic context, not on
    pitch-space position alone.
  \item[Aleatoric time] (Chapter~\ref{ch:time}). A region time model
    in which events have partial or undefined ordering, expressed
    as a directed acyclic graph with optional interval bounds.
  \item[Anchor] (Chapter~\ref{ch:time}). A reference to a point in
    time, expressed relative to an identified object (event,
    measure, region, wall-clock time) plus an offset. Anchors
    survive edits that do not delete their target.
  \item[Arena (event arena)] (Chapter~\ref{ch:graph}). The flat
    storage owned by a score holding all events of all kinds, with
    $O(1)$ lookup by \texttt{EventId}.
  \item[Bundle] (Chapter~\ref{ch:format}). The on-disk container
    holding a score: header, manifest, chunk store, and blob store
    in a single \texttt{.musc} file.
  \item[Canonical reduction] (Chapter~\ref{ch:semops}). The
    deterministic procedure that materializes a score graph from
    the operation set: sort by canonical order, group transactions,
    apply each operation against working state, record effects.
  \item[Canvas] (Chapter~\ref{ch:graph}). The spatial root of a
    score: a 2D space partitioned into regions of staff-based,
    free-graphic, or hybrid content.
  \item[Causal predecessors] (Chapter~\ref{ch:semops}). The set of
    operations on which a given operation causally depends. Used by
    the CRDT layer to determine concurrency. Expressed in canonical
    form as a dotted version vector
    (\texttt{CausalContext}); see ``Dotted version vector.''
  \item[Chunk] (Chapter~\ref{ch:format}). An immutable
    content-addressed unit of storage within a bundle. Identified
    by the cryptographic hash of its uncompressed payload, with a
    domain-separated preimage.
  \item[Compound operation] (Chapter~\ref{ch:semops}). A user-facing
    operation that expands into a sequence of primitive operations
    applied as a transaction.
  \item[Conflict record] (Chapter~\ref{ch:semops}). A first-class
    graph object recording an operation that could not apply
    cleanly under canonical reduction. Stable, addressable,
    user-visible, and resolvable by subsequent
    \texttt{ResolveConflict} operations.
  \item[CRDT] (Chapters~\ref{ch:semops},~\ref{ch:format}).
    Conflict-free replicated data type. In Epiphany, the
    \emph{operation set} is a true CRDT (a grow-only set of
    operation envelopes); the materialized score graph is
    \emph{not} itself a CRDT but is the deterministic reduction
    of that set.
  \item[Cross-cutting structure] (Chapter~\ref{ch:graph}). A
    first-class object spanning multiple tree-nodes via references:
    slurs, ties, beams, spanners, markers, repeats, analytical
    annotations, comments, graphic gestures, lyrics, chord symbols.
  \item[Decomposition] (Chapter~\ref{ch:time}). The notational
    breakdown of an event's sounding duration into notehead values,
    augmentation dots, and ties. Attached externally by the
    decomposition pre-pass.
  \item[Delta] (Chapter~\ref{ch:semops}). The structural record of
    an operation's effect, sufficient to reconstruct the pre-state
    of every object the operation modifies.
  \item[Document identifier] (Chapter~\ref{ch:format}). The
    logical work identity (\texttt{DocumentId}), stable across
    Save As copies that are intended to remain versions of the
    same work. May persist across forks at the user's
    discretion. Distinct from \texttt{FileUuid}
    (physical bundle identity, changes on Save As) and
    \texttt{LineageId} (optional shared-ancestor identity).
  \item[Dotted version vector] (Chapters~\ref{ch:semops},
    \ref{ch:format}). The compact causal-context representation
    of operation history. Operational semantics (causal
    precedence, frontier coverage, monotonic growth) are
    specified in Chapter~\ref{ch:semops}; the on-disk
    serialization is specified in Chapter~\ref{ch:format}.
  \item[Edit barrier] (Chapter~\ref{ch:format}). A structured
    declaration by an extension that certain edits would
    invalidate the extension's data. Implementations not
    understanding the extension refuse matching edits unless the
    user performs an explicit unsafe edit.
  \item[Engraving decision] (Chapter~\ref{ch:layout-ir}). A choice
    made by the engraver (stem direction, beam consolidation, etc.)
    recorded explicitly in the layout IR for inspection and
    override.
  \item[Engraving override] (Chapter~\ref{ch:layout-ir}). A user
    assertion that the engraver's default decision should be
    replaced. Authoritative in the score graph; projected into the
    layout IR.
  \item[Event] (Chapter~\ref{ch:graph}). A rhythmic atom of the
    score: pitched, unpitched, rest, indeterminate, trajectory,
    graphic, or cue. Owned by exactly one voice.
  \item[Format profile] (Chapter~\ref{ch:format}). A named subset
    of the file format that conforming implementations promise to
    support. Profiles enumerated: Full, ReadOnly, Lite, Custom.
  \item[Frame budget] (Chapter~\ref{ch:perf}). The time allowed for
    a layout-and-render operation to maintain 60\,Hz (16.7\,ms) or
    120\,Hz (8.3\,ms) interactivity.
  \item[Glyph catalog identity] (Chapter~\ref{ch:layout-ir}). The
    identifier (SMuFL version, font id, optional font version,
    metrics hash) under which layout output was produced.
    Required for any byte-equal-output layout conformance claim.
  \item[Graphic content] (Chapter~\ref{ch:graph}). Free-form vector
    graphics, drawn strokes, text, and images within free-graphic
    or hybrid regions.
  \item[Hard constraint] (Chapter~\ref{ch:solver}). A constraint
    that must be satisfied; violation produces a solver error.
  \item[Hybrid logical clock] (Chapters~\ref{ch:semops},
    \ref{ch:format}). A clock combining physical wall-clock
    time with a logical counter, used as a tie-breaker in
    canonical reduction order. Operational semantics
    (per-replica monotonicity, role in canonical reduction
    ordering, malformed-envelope behavior) are specified in
    Chapter~\ref{ch:semops}; the on-disk serialization is
    specified in Chapter~\ref{ch:format}.
  \item[Identifier (typed)] (Chapter~\ref{ch:graph}). A 128-bit
    stable identifier for a named object: \texttt{EventId},
    \texttt{VoiceId}, \texttt{StaffId}, etc. Generated by
    replica-plus-counter scheme.
  \item[Incremental layout] (Chapter~\ref{ch:layout-ir}). Re-layout
    of only the IR objects whose dependencies have changed, with
    cached values reused elsewhere.
  \item[Instrument] (Chapter~\ref{ch:graph}). An abstract definition
    of a sounding entity: name, transposition, range, default
    sound, default clef, default staff count.
  \item[Layout IR] (Chapter~\ref{ch:layout-ir}). The pipeline of
    intermediate representations between score graph and renderer:
    \texttt{LogicalLayoutIR}, \texttt{ConstrainedLayoutIR},
    \texttt{ResolvedLayoutIR}, \texttt{RenderIR}.
  \item[Manifest] (Chapter~\ref{ch:format}). The mutable index of a
    bundle: a table of operation roots, optional canonical base
    snapshot, acceleration snapshots, blob roots, profile
    declarations, and extension declarations. Itself a
    content-addressed chunk; the active manifest is selected
    via the active superblock slot
    (Chapter~\ref{ch:format}).
  \item[Metric time] (Chapter~\ref{ch:time}). A region time model
    with measures, beats, and exact rational positions. The
    dominant model for Western notated music.
  \item[Notation grammar] (Chapter~\ref{ch:tuning}). The declarative
    definition of a notation system's pitch space, intervals,
    accidentals, and spelling rules. CMN is the default; others
    extend.
  \item[Operation envelope] (Chapter~\ref{ch:semops}). A
    transmitted, stored, and reducible record carrying an
    operation's payload together with its identity, ordering
    stamp, causal context, and optional transaction grouping.
    The unit of replication.
  \item[Operation kind] (Chapter~\ref{ch:semops}). The tagged
    enumeration of primitive operation variants: InsertEvent,
    DeleteEvent, RespellPitch, Transpose, ChangeRegionTimeModel,
    and so on, plus \texttt{Registered} for extension-defined
    primitives. The full catalog is delivered separately.
  \item[Operation log] (Chapter~\ref{ch:format}). The physical
    storage of the canonical operation-envelope set on disk, as
    operation-envelope blocks. The canonical document is the
    deterministic reduction of this envelope set, or, after
    pruning, of the manifest's \texttt{canonical\_base}
    snapshot plus envelopes whose \texttt{OperationId}s are not
    covered by the base's causal frontier.
  \item[Operation set] (Chapter~\ref{ch:semops}). The replicated
    grow-only CRDT of operation envelopes. The canonical
    document is its deterministic reduction.
  \item[Pareto frontier (design target)]
    (Chapter~\ref{ch:solver}). The set of layouts for which no
    quality metric can be improved without worsening another.
    Solvers \SHOULD{} seek layouts near the frontier; conformance
    is reference-suite-based, not proof of Pareto optimality.
  \item[Part] (Chapter~\ref{ch:graph}). A per-instrument or
    per-section view onto the score for printed extraction. Parts
    are projections, not storage.
  \item[Pitch space] (Chapter~\ref{ch:tuning}). The analytical
    universe of a pitch system: what positions exist, how they
    relate, what accidentals are valid. Distinct from tuning system.
  \item[Pitch spelling] (Chapter~\ref{ch:pitch}). The notational
    appearance of a pitch: nominal, accidental, octave. Attached
    externally by the spelling pre-pass.
  \item[Primitive operation] (Chapter~\ref{ch:semops}). The atomic
    unit of editing. Primitive operations compose into compound
    operations and transactions.
  \item[Provenance] (Chapter~\ref{ch:layout-ir}). The record on
    each layout IR object tracing back to its score graph source,
    enabling selection, editing back-references, and incremental
    layout invalidation.
  \item[Proportional time] (Chapter~\ref{ch:time}). A region time
    model in which horizontal position is wall-clock time, with no
    measures or beats.
  \item[Re-anchoring] (Chapter~\ref{ch:semops}). The normative
    response of a cross-cutting structure when its referent is
    deleted: cascade delete, truncate, re-anchor, mark orphaned,
    or refuse.
  \item[Reference pitch] (Chapter~\ref{ch:tuning}). The anchor
    converting tuning system structure to absolute frequencies:
    a position and a frequency in Hertz. A property of the score,
    not the tuning system.
  \item[Region] (Chapter~\ref{ch:graph}). A spatial-temporal
    container in the canvas declaring a time model (metric,
    proportional, aleatoric) and a content model (staff-based,
    free-graphic, hybrid).
  \item[Repair record] (Chapter~\ref{ch:semops}). A deterministic
    compensating change made during reduction to preserve graph
    invariants (re-anchoring, spanner truncation, attachment
    tombstoning, voice promotion, tuplet compensation). Recorded
    as part of canonical state via
    \texttt{OperationEffect::AppliedWithRepair}.
  \item[Replica] (Chapter~\ref{ch:graph},~\ref{ch:format}). An
    instance of a score on a device. Each replica has a unique
    identifier and authors operations with monotonic local
    sequence numbers. The replica identifier
    \texttt{ReplicaId::SYSTEM\_DERIVED} is reserved for
    deterministically-derived system identifiers.
  \item[Retention policy] (Chapter~\ref{ch:format}). The
    declared rule governing how many old manifests (and their
    reachable chunks) are preserved in the bundle for rollback
    purposes. Rollback is achieved through retained manifests in
    the chunk store, not by retaining multiple on-disk
    superblock generations.
  \item[Scale position] (Chapter~\ref{ch:pitch}). The analytical
    identity of a pitch within a pitch space. One of the two
    intrinsic layers of a \texttt{Pitch}.
  \item[Score graph] (Chapter~\ref{ch:graph}). The in-memory
    representation of all musical content in a score: canonical
    truth from which layout, serialization, and editing operations
    derive.
  \item[Semantic operation] (Chapter~\ref{ch:semops}). A mutation
    of the score graph: insert event, delete measure, transpose,
    respell pitch, etc.
  \item[SMuFL] (Chapter~\ref{ch:tuning}). Standard Music Font
    Layout. The normative source for glyph references in the core.
  \item[Snapshot] (Chapter~\ref{ch:format}). A materialized score
    state. As an acceleration snapshot (listed in
    \texttt{acceleration\_snapshots}), derivative and discardable.
    As a canonical-base snapshot (named in the manifest's
    \texttt{canonical\_base} field with declared causal frontier
    and reduction algorithm version), the canonical base for
    pruned documents, in which case the canonical state is the
    base snapshot plus envelopes whose \texttt{OperationId}s are
    not covered by the base's causal frontier.
  \item[Soft constraint] (Chapter~\ref{ch:solver}). A constraint
    with a weight; violation is permitted but penalized in
    quality metrics.
  \item[Sounding duration] (Chapter~\ref{ch:time}). The exact
    rational duration an event occupies in time. Intrinsic;
    notational decomposition is attached separately.
  \item[Spring slot] (Chapter~\ref{ch:layout-ir}). A time slot's
    horizontal spacing parameters (min, preferred, max widths;
    stretch, compress factors) consumed by the constraint solver.
  \item[Staff] (Chapter~\ref{ch:graph}). The global, abstract
    staff identity persisting across the entire score. A staff is
    declared once at the score level and referenced by staff
    instances in any region where it appears.
  \item[Staff instance] (Chapter~\ref{ch:graph}). The
    region-local manifestation of a staff: the voices, measures,
    clef changes, key changes, and metric grid that belong to
    that staff for one region. A single \texttt{Staff} may have
    multiple \texttt{StaffInstance}s across the score.
  \item[Tempo map] (Chapter~\ref{ch:time}). The function mapping
    musical positions to wall-clock positions, piecewise over
    musical time with constant, linear, exponential, or curve
    segments.
  \item[Time anchor] (Chapter~\ref{ch:time}). See ``Anchor.''
  \item[Time axis] (Chapter~\ref{ch:layout-ir}). The polymorphic
    layout-IR component projecting time positions to horizontal
    positions. Three implementations: metric, proportional,
    aleatoric.
  \item[Transaction] (Chapter~\ref{ch:semops}). A replicated
    operation group that reduces atomically: either every member
    applies and the transaction reduces to \texttt{Applied}, or
    the entire transaction reduces to \texttt{Conflicted} with a
    recorded conflict.
  \item[Trajectory] (Chapter~\ref{ch:graph}). A continuous
    pitch-motion event: glissando, portamento, pitch bend. A
    sibling to pitched events under the event taxonomy.
  \item[Tuning system] (Chapter~\ref{ch:tuning}). The mechanism
    mapping pitch-space positions to frequencies, given a
    reference pitch. Distinct from pitch space.
  \item[Tuplet] (Chapter~\ref{ch:time}). A grouping object
    annotating irregular rhythmic groupings. Does not modify
    sounding durations; purely notational and analytical.
  \item[Typed object identifier]
    (Chapter~\ref{ch:graph}). A tagged union
    (\texttt{TypedObjectId}) over every identifier kind in the
    score graph. Used wherever an object is referenced
    generically (cross-cutting endpoints, conflict
    \texttt{affected\_objects}, repair records, edit barriers).
  \item[Validation mode] (Chapter~\ref{ch:semops}). The strictness
    setting for operation precondition checking: strict authoring
    or lenient replay.
  \item[Vertical band] (Chapter~\ref{ch:layout-ir}). A horizontal
    slice of the canvas with spring parameters for vertical
    spacing: staff bands, inter-staff gaps, inter-system gaps,
    margin bands.
  \item[View] (Chapter~\ref{ch:graph}). A recipe for displaying a
    configuration of the score: which layers are active, which
    parts are rendered, what overrides apply.
  \item[Voice] (Chapter~\ref{ch:graph}). A polyphonic line within
    a staff. Owns a sequence of events.
  \item[Wall-clock time] (Chapter~\ref{ch:time}). Time measured by
    the playback engine or external sync source, in fixed-point
    nanoseconds. Distinct from musical time.
\end{description}

\chapter{Bibliography and References}
\label{app:refs}

\section*{Engraving Treatises}

\begin{description}
  \item[Behind Bars.] Elaine Gould. \emph{Behind Bars: The Definitive
    Guide to Music Notation.} Faber Music, 2011. The principal
    modern reference for music engraving conventions.
  \item[Music Notation in the Twentieth Century.] Kurt Stone.
    \emph{Music Notation in the Twentieth Century: A Practical
    Guidebook.} W. W. Norton, 1980. Authoritative reference for
    contemporary notation.
  \item[Essentials of Music Notation.] Tom Ross. \emph{The Art of
    Music Engraving and Processing.} Hansen Books, 1970. Classical
    treatise on engraving conventions, still influential.
  \item[Music Notation.] Gardner Read. \emph{Music Notation: A
    Manual of Modern Practice.} Crescendo Publishers, 1969.
    Reference for engraving practices and history.
\end{description}

\section*{Music Theory and Analysis}

\begin{description}
  \item[The Cognition of Basic Musical Structures.]
    David Temperley. MIT Press, 2001. Source of preference-rule
    approaches to spelling and meter analysis.
  \item[Pitch Spelling and the Line of Fifths.]
    H. Christopher Longuet-Higgins. Foundational work on
    pitch-spelling algorithms.
  \item[Tonal Pitch Space.] Fred Lerdahl. Oxford University Press,
    2001. Pitch-space theory as analytical framework.
  \item[A Generative Theory of Tonal Music.]
    Fred Lerdahl and Ray Jackendoff. MIT Press, 1983. Foundational
    work in formal music theory and the source of preference-rule
    methodology.
\end{description}

\section*{Tuning, Microtonality, and Just Intonation}

\begin{description}
  \item[Genesis of a Music.] Harry Partch. University of Wisconsin
    Press, 2nd ed. 1974. Foundational text in extended just
    intonation and microtonal practice.
  \item[Tuning, Timbre, Spectrum, Scale.] William Sethares.
    Springer, 2nd ed. 2004. Modern treatment of tuning systems
    and their psychoacoustic foundations.
  \item[Helmholtz-Ellis JI Pitch Notation.] Marc Sabat and
    Wolfgang von Schweinitz. The HEJI accidental system referenced
    by the core's built-in JI pitch spaces.
  \item[Sagittal: A Microtonal Notation System.] George D. Secor
    and David C. Keenan. The Sagittal microtonal accidental system.
\end{description}

\section*{Constraint-Based Layout}

\begin{description}
  \item[Cassowary Linear-Arithmetic Constraint Solver.]
    Greg J. Badros, Alan Borning, and Peter J. Stuckey.
    \emph{The Cassowary Linear Arithmetic Constraint Solving
    Algorithm.} ACM Transactions on Computer-Human Interaction,
    8(4):267-306, 2001. Source of the constraint-solver model
    referenced in Chapter~\ref{ch:solver}.
  \item[Music Engraving with Constraint Programming.]
    Multiple authors. Survey of constraint-based approaches to
    music engraving.
  \item[The Gourlay Algorithm.] J. S. Gourlay. \emph{A Language
    for Music Printing.} Communications of the ACM, 1986. Source
    of the proportional spacing model referenced in
    Chapter~\ref{ch:layout-ir}.
\end{description}

\section*{CRDT and Distributed Systems}

\begin{description}
  \item[A Comprehensive Study of Convergent and Commutative
    Replicated Data Types.] Marc Shapiro, Nuno Pregui\c{c}a,
    Carlos Baquero, and Marek Zawirski. INRIA Technical Report
    RR-7506, 2011. Foundational survey of CRDTs.
  \item[Dotted Version Vectors.] Nuno Pregui\c{c}a et al.
    Refinement of vector clocks used in the file format
    (Chapter~\ref{ch:format}).
  \item[Hybrid Logical Clocks.] Sandeep S. Kulkarni et al.
    Clocks combining physical and logical time, used for total
    ordering of operation envelopes
    (Chapter~\ref{ch:semops}).
\end{description}

\section*{Prior-Art Notation Formats}

\begin{description}
  \item[MusicXML.] Recordare LLC, later W3C Music Notation Community
    Group. The dominant interchange format; surveyed for the
    interchange profile work referenced in
    Chapter~\ref{ch:graph}.
  \item[MEI (Music Encoding Initiative).] An XML-based encoding
    designed for scholarly editions and musicology. Influenced the
    spec's handling of editorial apparatus and analytical layers.
  \item[MNX.] W3C Music Notation Community Group draft. The
    successor effort to MusicXML; informed several decisions on
    layered data models.
  \item[LilyPond.] David Bauer, Han-Wen Nienhuys, Jan Nieuwenhuizen.
    Open-source music engraving software using a text-based input
    format and a sophisticated rule-based engraver. Significant
    influence on the spec's engraving philosophy.
  \item[SMuFL.] W3C Music Notation Community Group. The Standard
    Music Font Layout, referenced as the normative glyph catalog
    in Chapter~\ref{ch:tuning}.
\end{description}

\section*{Additional References}

\begin{description}
  \item[The Anatomy of a Note.] Multiple papers on the data
    representation of musical notation in computer systems.
  \item[Notation in Computer Music.] Various proceedings of the
    International Computer Music Conference (ICMC) and the New
    Interfaces for Musical Expression (NIME) workshops.
\end{description}

\chapter{Determinism Contract}
\label{app:determinism}

This appendix consolidates the reproducibility obligations that
appear throughout the specification. It defines exactly which
values are admissible in canonical state, how they are serialized,
how comparisons work, and where implementation freedom begins.

The appendix is the document's legal code for reproducibility:
where the body of the specification says ``deterministic,'' this
appendix says what that means.

\section{Thesis}
\label{sec:det:thesis}

\begin{requirement}
  \label{req:determinism:platform-independent-state}
  Canonical document state \MUST{} be independent of platform,
  CPU, locale, thread scheduling, hash-map iteration order,
  floating-point environment, compression settings, and
  wall-clock timing.

  Where exact equality is required, this specification defines
  canonical representations. Where numerical approximation is
  unavoidable, this specification defines deterministic
  tolerances, quantization, and reporting rules.

  Implementations \MAY{} freely choose internal data structures,
  algorithms, and parallelism strategies. The canonical externally
  observable result \MUSTNOT{} depend on those choices.
\end{requirement}

\section{Layers of Determinism}
\label{sec:det:layers}

The specification distinguishes five layers. Each layer has its
own determinism contract; conflating them is the most common
source of impossible requirements.

\begin{description}
  \item[Canonical score determinism] (Chapter~\ref{ch:semops}).
    The reduction of an operation-envelope set to a materialized
    score graph is fully deterministic \emph{when every algorithm
    contributing to that reduction has received one of the
    dispositions in Section~\ref{sec:det:open}} (normatively
    specified, profile-declared by versioned identifier, or
    explicitly non-canonical). Under that condition, identical
    operation sets produce identical materialized score states
    across all conforming implementations, with identical object
    identifiers, tombstones, conflicts, and effects. For
    algorithms still in open-question state
    (Section~\ref{sec:det:open}), this byte-equal cross-
    implementation guarantee \MUSTNOT{} be claimed; the output
    of such algorithms is either non-canonical or governed by the
    profile's declared algorithm version.
  \item[Canonical serialization determinism]
    (Chapter~\ref{ch:format}). The bundle format's chunk hashes,
    manifest encoding, text projection, and the ordering of
    canonical maps and sets are deterministic. Two implementations
    producing the same canonical state \MUST{} produce identical
    chunk content hashes and identical text-projection bytes.
  \item[Layout determinism] (Chapter~\ref{ch:solver}). Within a
    single solver implementation at a fixed version, with fixed
    profile, fixed budget, and fixed font/glyph metrics, layout
    output is byte-identical across runs. Across different solver
    implementations, layout output \MAY{} differ aesthetically
    within reference-suite thresholds.
  \item[Conformance determinism]. Reference-suite-based:
    conforming implementations need not produce byte-identical
    output to each other; they need to fall within declared
    metric thresholds on every reference-suite entry. See
    Chapter~\ref{ch:solver} Section~\ref{sec:solver:conformance}.
  \item[Non-canonical cache determinism]. Implementations
    maintain caches (operation index, layout cache, glyph metric
    cache, render artifact cache). These \MAY{} vary across
    runs, platforms, and implementations. They \MUSTNOT{} affect
    canonical state and \MUST{} be discardable
    (Chapter~\ref{ch:semops} Section~\ref{sec:semops:caches}).
\end{description}

\begin{requirement}
  \label{req:determinism:determinism-layer-scope}
  Byte-equality obligations \MUST{} be applied only to the layer
  for which they are stated. Implementations \MUSTNOT{} interpret
  ``deterministic'' more strongly than its layer specifies, and
  \MUSTNOT{} interpret it more weakly either.
\end{requirement}

\begin{rationale}
  The conditional in canonical score determinism is essential.
  Several algorithms affecting canonical state remain open
  questions in this draft: the spelling pre-pass
  (Chapter~\ref{ch:pitch}), the notational decomposition
  algorithm (Chapter~\ref{ch:time}), and the
  \texttt{wallclock\_to\_musical} root-finding algorithm
  (Chapter~\ref{ch:time}). Until those algorithms receive
  Section~\ref{sec:det:open} dispositions, cross-implementation
  byte equality cannot be promised for outputs derived from them.

  This is not a weakening of the determinism contract; it is a
  precise statement of what the contract currently covers. Once
  the open algorithms are specified or profile-declared, the
  conditional resolves and full cross-implementation byte
  equality applies to all of canonical state.
\end{rationale}

\section{Floating-Point Values in Canonical State}
\label{sec:det:fp}

\subsection{Permitted Forms}

\begin{requirement}
  \label{req:determinism:canonical-floating-point}
  Canonical stored floating-point values \MUST{} be finite
  IEEE 754 binary64 values. NaN and infinity \MUSTNOT{} appear
  in canonical chunks: implementations \MUST{} reject them at
  serialization time and \MUST{} treat their presence on read as
  data corruption.

  The value \texttt{-0.0} (negative zero) \MUST{} be canonicalized
  to \texttt{+0.0} before storage. Two floating-point values that
  differ only in zero sign are equal under canonical comparison.

  Floating-point values appear in advisory, acoustic, tuning,
  tempo, layout, and quality-metric contexts. Where exact
  identity is needed (object identifiers, operation ordering,
  graph membership, hash identity), floating point \MUSTNOT{}
  be used.
\end{requirement}

\subsection{Serialization}

\begin{requirement}
  \label{req:determinism:floating-point-serialization}
  Canonical \texttt{f64} values are serialized as little-endian
  IEEE 754 binary64 octets after applying the
  \texttt{-0.0 $\to$ +0.0} canonicalization. The eight serialized
  bytes are the canonical representation; two values whose
  canonical bytes are equal are canonically equal.

  Hash preimages incorporating floating-point values
  (Section~\ref{sec:format:hashing}) \MUST{} use the canonical
  serialized bytes, not platform-native byte orders or
  representations.
\end{requirement}

\subsection{Equality}

\begin{requirement}
  \label{req:determinism:floating-point-equality}
  Canonical equality of two floating-point values is byte equality
  of their canonical serialized representations. The IEEE 754
  notion that \texttt{NaN $\ne$ NaN} is irrelevant in canonical
  state, because canonical state cannot contain NaN.

  For \texttt{NormalizedMetric} values
  (Chapter~\ref{ch:solver}), additional invariants apply: the
  value is finite, in $[0.0, 1.0]$, and lower is better.
  Comparisons of \texttt{NormalizedMetric} use IEEE 754 ordered
  comparison; the \texttt{-0.0 $\to$ +0.0} canonicalization
  applies before comparison.
\end{requirement}

\section{Rounding and CPU Behavior}
\label{sec:det:rounding}

\begin{requirement}
  \label{req:determinism:floating-point-execution}
  Numerical operations whose results enter canonical state
  \MUST{} use IEEE 754 round-to-nearest, ties-to-even. The
  ambient floating-point rounding mode (which on most platforms
  defaults to round-to-nearest-ties-to-even but can be changed)
  \MUSTNOT{} affect canonical results: implementations \MUST{}
  either save/restore the rounding mode or use operations
  insensitive to it.

  Implementations \MUSTNOT{} compile canonical numerical
  algorithms with \texttt{-ffast-math}, \texttt{-funsafe-math-
  optimizations}, or equivalent flags that permit
  non-IEEE-conformant transformations. Implementations \MUSTNOT{}
  reorder, reassociate, or distribute floating-point expressions
  in ways that can change the canonical result.

  Implementations \MUSTNOT{} silently fuse multiplications and
  additions into fused-multiply-add operations unless the
  algorithm explicitly requests FMA. Implementations \MUSTNOT{}
  depend on, or be sensitive to, ambient floating-point exception
  flags or host rounding mode.

  Where implementations use SIMD, GPU acceleration, or platform
  math libraries (\texttt{libm}, vendor math libraries) in
  canonical algorithms, they \MUST{} produce results identical
  to a strict-IEEE single-threaded baseline. Differences in
  transcendental function implementations (\texttt{sin},
  \texttt{cos}, \texttt{exp}, \texttt{log}) across libm versions
  are a real source of nondeterminism; implementations using
  transcendentals in canonical algorithms \MUST{} use a
  documented portable implementation (e.g., a vendored math
  library or a software fallback path).
\end{requirement}

\section{Exact and Quantized Representations}
\label{sec:det:exact}

The specification prefers exact representations over
floating-point wherever possible. The following are exact by
design:

\begin{description}
  \item[Time and duration.] \texttt{MusicalPosition} and
    \texttt{MusicalDuration} are rational
    (Chapter~\ref{ch:time}). Inline-or-promoted representation;
    arithmetic is exact.
  \item[Operation ordering.] Causal order from DVV; total order
    by HLC and identifier (Chapter~\ref{ch:semops}). All integer-
    valued.
  \item[Identifiers.] All graph identifiers are 128-bit values
    composed of replica plus counter (Chapter~\ref{ch:graph}).
  \item[Hashes.] All canonical hashes are BLAKE3-256
    (Chapter~\ref{ch:format}).
\end{description}

\subsection{Quantized Layout Coordinates}

Spatial layout coordinates require numeric computation, but
canonical output must be stable. The specification quantizes
canonical coordinates to a fixed grid.

\begin{lstlisting}[language=Rust]
/// A canonical spatial coordinate, quantized to 1/1024 of a
/// staff space. Used for canonical serialized output of
/// ResolvedLayoutIR positions. Internal solvers MAY use floating
/// point during computation; canonical serialization rounds to
/// QuantizedCoord.
pub struct QuantizedCoord {
    /// Coordinate value in units of 1/1024 staff space.
    pub units: i64,
}
\end{lstlisting}

\begin{requirement}
  \label{req:determinism:layout-coordinate-quantization}
  The canonical spatial coordinate grid is $1/1024$ staff space
  per unit. Layout coordinates emitted into canonical
  \texttt{ResolvedLayoutIR} \MUST{} be quantized to this grid
  before serialization. The quantization rule is round-to-nearest,
  ties-to-even on the integer unit value.

  Implementations \MAY{} use floating-point coordinates internally
  during layout computation. The act of quantization at
  serialization time absorbs all floating-point variation from
  the canonical output. Two implementations whose internal
  computations agree to better than $1/2048$ staff space at
  every coordinate produce identical canonical output after
  quantization.

  The grid spacing is fixed at $1/1024$ staff space for this
  format version. Future major versions \MAY{} adjust the
  spacing; doing so is a non-backward-compatible change.
\end{requirement}

\begin{rationale}
  $1/1024$ staff space is finer than visual resolution at any
  practical printing density (a staff space is typically 1.5--2\,mm,
  so $1/1024$ is roughly 1.5--2\,$\mu$m, well below print
  resolution). The grid is fine enough that quantization is
  inaudible musically and invisible visually, but coarse enough
  that floating-point noise from different solver implementations
  routinely rounds away.
\end{rationale}

\section{Tolerance Classes}
\label{sec:det:tolerances}

The specification refers to numerical tolerances in several
places (acoustic comparison, layout quality thresholds, tempo
integration, solver residual). All such tolerances belong to one
of a small set of named classes, each with declared semantics.

\begin{lstlisting}[language=Rust]
pub enum ToleranceClass {
    /// Acoustic pitch comparison, in cents.
    AcousticCents,

    /// Layout coordinate comparison, in staff spaces.
    LayoutCoordinate,

    /// Quality metric comparison; absolute on the [0.0, 1.0]
    /// NormalizedMetric scale.
    QualityMetric,

    /// Tempo integration residual: maximum permitted error in
    /// musical_to_wallclock or wallclock_to_musical conversion.
    TempoIntegration,

    /// Solver residual: maximum permitted constraint violation
    /// for a soft constraint to be considered satisfied.
    SolverResidual,
}

pub struct Tolerance {
    pub class: ToleranceClass,

    /// Absolute tolerance. The unit is implied by the class.
    pub absolute: f64,

    /// Optional relative tolerance, applied to nonzero values.
    pub relative: Option<f64>,

    /// What the tolerance governs.
    pub governance: ToleranceGovernance,
}

pub enum ToleranceGovernance {
    /// Tolerance affects canonical equality (rare; usually only
    /// for diagnostic comparison).
    Equality,

    /// Tolerance is a validation threshold (constraint considered
    /// satisfied if violation is below tolerance).
    Validation,

    /// Tolerance affects only diagnostic output; canonical state
    /// is unaffected.
    Diagnostic,
}
\end{lstlisting}

\begin{requirement}
  \label{req:determinism:named-tolerances}
  Specifications and conformance documents \MUSTNOT{} introduce
  ad-hoc epsilon constants. Every numerical tolerance that
  affects normative behavior \MUST{} be declared as a named
  \texttt{Tolerance} value belonging to a defined
  \texttt{ToleranceClass}, with an explicit unit, absolute and
  optional relative bounds, and governance category.

  Tolerances \MUSTNOT{} apply to identity: \texttt{PitchId},
  \texttt{EventId}, \texttt{OperationId}, graph membership,
  hash identity, and operation ordering are exact, never within
  tolerance.

  The Quality Metric Catalog, the Reference Suite, and the
  Performance Reference Suite enumerate the specific tolerance
  values for each profile and tier. Their values are normative
  per the companion specifications.
\end{requirement}

\section{Ordered Iteration over Sets and Maps}
\label{sec:det:ordering}

Canonical output frequently depends on iterating a collection.
Hash-map and B-tree implementation order leaks platform-specific
behavior into canonical output if not controlled.

\begin{requirement}
  \label{req:determinism:canonical-collection-order}
  Whenever canonical output, canonical serialization, or canonical
  hashing depends on iterating a set, map, or other collection,
  iteration \MUST{} occur in a specified total order. The
  following total orders are normative:

  \begin{itemize}
    \item Typed identifiers (\texttt{EventId},
      \texttt{PitchId}, \texttt{VoiceId}, etc.): lexicographic
      ascending on the identifier's canonical byte form
      (16 bytes: 8-byte replica, 8-byte counter, big-endian).
    \item Operation envelopes: causal order first (from DVV
      closure), then HLC physical time, then HLC logical
      counter, then \texttt{OperationId} as final tie-break.
    \item Chunk references: ascending by \texttt{ChunkKind}
      discriminant, then by content hash lexicographic, then by
      file offset.
    \item Conflict records: ascending by \texttt{ConflictId}.
    \item Extension declarations: ascending by
      \texttt{ExtensionId}, then by semantic version
      lexicographic.
    \item Strings (text fields): byte-lexicographic on the UTF-8
      NFC representation (Section~\ref{sec:det:text}).
    \item Rationals (\texttt{RationalTime} sequences): by
      numeric value ascending.
    \item Composite keys: lexicographic on the canonical
      representation of the key tuple's component fields, in
      declaration order.
  \end{itemize}

  Hash-map iteration order, B-tree implementation differences,
  and similar implementation artifacts \MUSTNOT{} leak into
  canonical output.
\end{requirement}

\section{Text and Unicode}
\label{sec:det:text}

Locale-dependent text behavior is a routine source of canonical
divergence. The specification fixes this at the encoding layer.

\begin{requirement}
  \label{req:determinism:unicode-canonicalization}
  Canonical text fields \MUST{} be encoded as UTF-8 with Unicode
  NFC (Normalization Form Canonical Composition) applied. Two
  texts whose NFC byte representations are equal are canonically
  equal regardless of decomposition state.

  Comparisons of canonical text fields for identity \MUST{} be
  byte comparisons of NFC-encoded UTF-8. User-facing display
  collation (sorting by locale rules, case-insensitive matching,
  diacritic folding) is non-canonical and \MUSTNOT{} affect
  canonical state.

  Locale settings \MUSTNOT{} affect:
  \begin{itemize}
    \item Parsing of canonical text projections.
    \item Serialization of canonical text fields.
    \item Sorting of canonical map keys.
    \item Decimal formatting in canonical text projections (the
      text projection uses a fixed decimal-point convention;
      locale comma-vs-period preferences are irrelevant).
    \item Date and time formatting in canonical data (canonical
      timestamps use a fixed format; locale formats are for UI
      only).
  \end{itemize}

  Two implementations operating under different locale settings
  on the same canonical document \MUST{} produce identical
  canonical output, including identical chunk hashes and
  identical text projections.
\end{requirement}

\section{Randomness and Parallelism}
\label{sec:det:rng}

\subsection{Randomness}

\begin{requirement}
  \label{req:determinism:canonical-random-seeds}
  Randomness \MUSTNOT{} affect canonical output unless the
  random seed is an explicit canonical input.

  Algorithms whose results enter canonical state and that use
  pseudorandom generation (e.g., probabilistic constraint solvers,
  randomized rounding) \MUST{} consume their seed as a declared
  algorithm parameter that becomes part of the algorithm's
  conformance configuration. Identical seed plus identical input
  yields identical output.

  Identifier minting (\texttt{ReplicaId} initialization) is the
  only conforming use of platform-level randomness in this
  specification, and even there the random output enters
  canonical state only via the resulting identifiers.
\end{requirement}

\subsection{Parallelism}

\begin{requirement}
  \label{req:determinism:parallel-baseline-equivalence}
  Parallel implementations of canonical algorithms \MUST{}
  produce results identical to a strict single-threaded baseline.
  Race timing, work-stealing order, and thread count \MUSTNOT{}
  affect canonical state.

  Where parallel reductions are used (summation over many
  contributions, parallel constraint evaluation), the merge step
  \MUST{} be deterministic: ordered reduction over a canonical
  ordering of contributions, not arbitrary tree-merge order.

  Implementations \MAY{} use parallel execution as a performance
  optimization. They \MAY NOT{} use it as a license for
  nondeterminism.
\end{requirement}

\section{Compression and File Bytes}
\label{sec:det:compression}

\begin{requirement}
  \label{req:determinism:compression-independent-identity}
  Canonical content identity is the uncompressed payload bytes
  with the domain-separated preimage
  (Chapter~\ref{ch:format}). Compression algorithm choice and
  compression level \MUSTNOT{} affect canonical identity.
  Identical uncompressed content has identical
  \texttt{ContentHash} regardless of how it was physically
  compressed.

  Canonical chunk payload encoding \MUST{} be deterministic: the
  same in-memory canonical structure \MUST{} produce identical
  uncompressed bytes across runs and across implementations,
  within the same schema version. The Binary Format companion
  specification defines the canonical encoding.

  Compression bytes (the physical bytes stored on disk after
  zstd or another algorithm is applied) are \emph{not} canonical
  unless a profile explicitly requires bit-identical bundle
  reproduction. Implementations \MAY{} recompress chunks (e.g.,
  during compaction or repack) without altering canonical state.

  Readers \MUST{} verify decompressed payload against the
  declared content hash before treating the payload as
  authoritative.
\end{requirement}

\section{Open Algorithm Hooks}
\label{sec:det:open}

Two algorithms remain open questions in the body of the
specification: tempo integration for arbitrary curves
and root-finding for \texttt{wallclock\_to\_musical}
(Chapter~\ref{ch:time}). These algorithms affect canonical
state. The spelling pre-pass (Chapter~\ref{ch:pitch}) and the
notational decomposition algorithm (Chapter~\ref{ch:time})
were resolved in Pass~12 to the profile-declared disposition
below, each with a ratified default identifier
(Requirements~\ref{req:pitch:spelling-algorithm}
and~\ref{req:time:decomposition-algorithm}).

\begin{requirement}
  \label{req:determinism:canonical-algorithm-registration}
  Any algorithm that affects canonical state \MUST{} satisfy one
  of:

  \begin{itemize}
    \item It is specified normatively in this document or a
      named companion specification. Conforming implementations
      \MUST{} implement the specified algorithm.
    \item It is profile-declared: the active conformance profile
      names a specific algorithm by versioned identifier. Two
      conforming implementations on the same profile use the
      same algorithm. Profiles may differ in their algorithm
      choices.
    \item It is explicitly marked non-canonical or advisory: its
      output appears in the score only in non-canonical layers
      (caches, diagnostics, user-facing rendering hints) and
      \MUSTNOT{} appear in canonical chunks.
  \end{itemize}

  Implementations \MUSTNOT{} introduce algorithms affecting
  canonical state outside of these three categories. ``Vendor's
  spelling preference'' or ``proprietary decomposition heuristic''
  \MUSTNOT{} silently determine canonical content.
\end{requirement}

\section{Conformance Statement}
\label{sec:det:conformance}

\begin{requirement}
  \label{req:determinism:conformance-disclosure}
  Conforming implementations \MUST{} declare, as part of their
  published conformance statement:

  \begin{itemize}
    \item Their platform and floating-point library, with the
      libm version or vendored math library reference.
    \item Their parallel execution strategy and the demonstrated
      equivalence to single-threaded baseline.
    \item Their handling of the \texttt{-0.0} canonicalization
      and NaN/infinity rejection.
    \item Their adherence to round-to-nearest-ties-to-even.
    \item Their compliance with the canonical iteration orders
      in Section~\ref{sec:det:ordering}.
    \item Their NFC normalization implementation reference.
    \item Their declared algorithms for any open-question area
      (spelling, decomposition, tempo curves), with versioned
      identifiers.
  \end{itemize}
\end{requirement}

% ===========================================================================
\chapter{Canonical Byte-Layout Reference}
\label{app:bytes}

\providecommand{\bul}{\_\discretionary{}{}{}}% breakable underscore for long identifiers
\providecommand{\cwb}{\discretionary{}{}{}}% zero-width break for long camelCase identifiers

\section{Status of this reference}

This appendix consolidates, in one place, every byte layout pinned by
the Pass 11 ratification: the discriminant tables, derivation
preimages, and primitive encodings that the core, operations, and
bundle layers already implement and lock with golden-bytes tests. Its
purpose is to give the Binary Format companion
(Section~\ref{sec:format:binary}) a single point to import, rather than
re-deriving these layouts from three separate crates.

This reference is a \emph{consolidation, not a second source of truth.}
Every entry cites the requirement that defines it normatively; where
this appendix and a cited requirement appear to disagree, the cited
requirement governs. The appendix introduces no obligation beyond what
its citations already impose. The layouts that Pass 11 deliberately
left to the interchange companions are listed in
Section~\ref{sec:bytes:deferred} so the boundary is explicit.

\section{Shared encoding conventions}
\label{sec:bytes:conventions}

All composite canonical encodings follow the convention baseline that
the Binary Format companion inherits
(Requirement~\ref{req:format:codec-conventions}): little-endian
integers; a boolean as a single \texttt{0}/\texttt{1} byte; \texttt{u32}
little-endian counts and a length prefix on every variable-width leaf;
a tagged union as a single discriminant byte followed by its variant
payload; and free-text strings as length-prefixed UTF-8 that the codec
does \emph{not} NFC-fold (catalog identifiers are NFC-normalized at
construction, not in the codec). Certain hashing preimages deliberately
depart from the single-discriminant-byte convention where a wider or
order-bearing tag is required --- notably the \texttt{TypedObjectId}
16-bit big-endian discriminant; such departures are pinned at their
definitions and noted in the tables below.

\section{System-derived identifier function}
\label{sec:bytes:system-derived}

Several identifiers below are content-addressed rather than
counter-allocated. They share one function
(Section~\ref{sec:graph:system-derived}): the 64-bit counter of a
\texttt{ReplicaId::SYSTEM\_DERIVED}
(\texttt{0xffff\_ffff\_ffff\_ffff}) identifier is
$\mathrm{BE\_u64}(\mathrm{BLAKE3}(\texttt{domain\_tag} \Vert
\texttt{canonical\_inputs})[0..8])$. The domain tag is exactly eight
bytes. The reserved built-in tags \emph{for this derivation} are the
closed set \{\texttt{MUSCSVCE}, \texttt{MUSCSPCH}, \texttt{MUSCSANM}\}
--- the three system-derived-identifier tags only; the other reserved
domain tags catalogued in Section~\ref{sec:bytes:tags} feed plain
hashing preimages, not this counter. Tags introduced by registered
extensions \MUST{} begin with \texttt{MUSCS}, be exactly eight bytes,
and not collide with the reserved three.

\section{Discriminant tables, derivations, and encodings}
\label{sec:bytes:tables}

\subsection{Object-graph identity (canonical document state)}

\begin{longtable}{>{\raggedright\arraybackslash}p{3.5cm} >{\raggedright\arraybackslash}p{6.8cm} p{2.2cm}}
  \toprule
  \textbf{Layout} & \textbf{Canonical bytes / preimage} &
  \textbf{Requirement} \\
  \midrule
  \endhead
  \texttt{TypedObjectId} tag &
  16-bit big-endian discriminant (\texttt{0}--\texttt{27}) $\Vert$
  variant payload. \texttt{Registered} $= 27$ encodes
  \texttt{reg}~(16 BE) $\Vert$ \texttt{raw}~(16 BE), 34 bytes total.
  \texttt{ObjectKindRegistryId} is 128-bit &
  \ref{req:graph:typed-object-id-discriminants} \\[2pt]

  Promoted \texttt{VoiceId} &
  System-derived under \texttt{MUSCSVCE} over a 64-byte preimage:
  \texttt{staff\_instance} $\Vert$ \texttt{original\_voice} $\Vert$
  \texttt{winning\_op} $\Vert$ \texttt{losing\_op}, each 16-byte
  big-endian &
  \S\ref{sec:graph:promoted-voices} \\[2pt]

  System-derived \texttt{PitchId} &
  System-derived under \texttt{MUSCSPCH} over the intrinsic identity
  (scale position, acoustic realization); variable-width strings
  length-prefixed and NFC-normalized at the boundary; the tuning
  reference is always present (\texttt{Inherit} as a distinct presence
  marker) &
  \ref{req:graph:system-derived-pitch-id} \\[2pt]

  \texttt{Integrity\cwb AnomalyId} &
  System-derived under \texttt{MUSCSANM} over
  \texttt{Integrity\cwb AnomalyKind::\cwb to\bul canonical\bul bytes()} &
  \ref{req:graph:integrity-anomaly-id} \\[2pt]

  \texttt{ObjectKind} discriminant &
  Single byte: \texttt{Voice} $= 0$, \texttt{Pitch} $= 1$,
  \texttt{Registered} $= 2$. Enters the \texttt{IntegrityAnomalyId}
  preimage &
  \ref{req:graph:object-kind-vocab} \\
  \bottomrule
\end{longtable}

\subsection{Bundle layer (chunk store, manifest, blobs)}

\begin{longtable}{>{\raggedright\arraybackslash}p{3.5cm} >{\raggedright\arraybackslash}p{6.8cm} p{2.2cm}}
  \toprule
  \textbf{Layout} & \textbf{Canonical bytes / preimage} &
  \textbf{Requirement} \\
  \midrule
  \endhead
  Chunk content hash &
  BLAKE3 of \texttt{domain\bul tag} $\Vert$ \texttt{ChunkKind} byte
  $\Vert$ \texttt{SchemaVersion} (4) $\Vert$
  \texttt{uncompressed\bul length} (u64 LE) $\Vert$ \texttt{payload};
  tag \texttt{MUSCCHNK} for chunks, \texttt{MUSCMANI} for the manifest.
  Compression is \emph{not} in the preimage &
  \S\ref{sec:format:hashing} \\[2pt]

  \texttt{ChunkKind} &
  Single declaration-order byte, \texttt{0}--\texttt{8} (in the chunk
  hash preimage, so stable) &
  \ref{req:format:chunkkind-discriminants} \\[2pt]

  \texttt{Compression\cwb Algorithm} &
  Fixed two bytes: discriminant $\Vert$ parameter. \texttt{None} $=$
  \texttt{[0,0]}, \texttt{Zstd\{level\}} $=$ \texttt{[1,level]},
  \texttt{Reserved(v)} $=$ \texttt{[2,v]}. Not part of chunk identity &
  \ref{req:format:chunkkind-discriminants} \\[2pt]

  \texttt{SchemaVersion} &
  \texttt{major}~(u16 LE) $\Vert$ \texttt{minor}~(u16 LE), four bytes &
  \S\ref{sec:format:hashing} \\[2pt]

  \texttt{BlobId} &
  Bare BLAKE3 of \texttt{MUSCBLOB} $\Vert$
  \texttt{uncompressed\bul payload}: the domain tag immediately
  followed by payload, with no kind, schema, or length fields &
  \ref{req:format:blob-hash-shape} \\[2pt]

  \texttt{ManifestId} &
  \texttt{trunc128} (leading 16 bytes) of BLAKE3 of \texttt{MUSCMNIF}
  $\Vert$ \texttt{document\bul id} (16) $\Vert$ \texttt{generation}
  (u64 LE) $\Vert$ \texttt{manifest\bul body}; the
  \texttt{manifest\bul id} field is excluded from the body &
  \ref{req:format:manifest-id} \\[2pt]

  \texttt{ProfileId} &
  Fixed 20 bytes: discriminant (u32 LE; \texttt{Full} $= 0$,
  \texttt{ReadOnly} $= 1$, \texttt{Lite} $= 2$, \texttt{Custom} $= 3$)
  $\Vert$ \texttt{ProfileRegistryId}~(16, zero unless \texttt{Custom}).
  Load-bearing in superblock selection &
  \ref{req:format:profileid-discriminants} \\
  \bottomrule
\end{longtable}

\subsection{Operation layer (transaction metadata)}

\begin{longtable}{>{\raggedright\arraybackslash}p{3.5cm} >{\raggedright\arraybackslash}p{6.8cm} p{2.2cm}}
  \toprule
  \textbf{Layout} & \textbf{Canonical bytes} & \textbf{Requirement} \\
  \midrule
  \endhead
  \texttt{Transaction\cwb Category} &
  Single byte: \texttt{NoteEntry} $= 0$, \texttt{Structural} $= 1$,
  \texttt{Layout} $= 2$, \texttt{Import} $= 3$, \texttt{Registered}
  $= 4$ &
  \ref{req:semops:transaction-category} \\[2pt]

  \texttt{ResolutionAction} &
  Single byte: \texttt{AcceptLoser} $= 0$, \texttt{KeepWinner} $= 1$,
  \texttt{Override} $= 2$, \texttt{Reanchor} $= 3$, \texttt{Dismiss}
  $= 4$, \texttt{Registered} $= 5$. Encoded into the operation content
  hash &
  \ref{req:semops:resolution-action-discriminants} \\
  \bottomrule
\end{longtable}

\subsection{Primitive scalar encodings}

\begin{longtable}{>{\raggedright\arraybackslash}p{3.5cm} >{\raggedright\arraybackslash}p{6.8cm} p{2.2cm}}
  \toprule
  \textbf{Layout} & \textbf{Canonical bytes} & \textbf{Requirement} \\
  \midrule
  \endhead
  \texttt{RationalTime} &
  Sign byte (\texttt{0} zero, \texttt{1} positive, \texttt{2} negative)
  $\Vert$ numerator magnitude (u32-LE length-prefixed, big-endian)
  $\Vert$ denominator magnitude (u32-LE length-prefixed, big-endian).
  Always stored reduced (lowest terms, positive denominator), so equal
  rationals encode to equal bytes &
  \ref{req:format:rationaltime-encoding} \\[2pt]

  Wall-clock integers &
  Fixed-width little-endian (matching \texttt{QuantizedCoord}) &
  \ref{req:format:rationaltime-encoding} \\
  \bottomrule
\end{longtable}

\section{Domain-tag registry}
\label{sec:bytes:tags}

Every reserved built-in eight-byte domain tag, in one place. The nine
\emph{canonical} tags ---
whose preimages produce identifiers and content hashes that are part of
the interoperable, durably persisted form every conforming
implementation must reproduce bit-identically --- are
listed first; \texttt{MUSCFNTM} and \texttt{MUSCLOID} are
\emph{non-canonical} (layout-only) and listed last, for completeness
(see Section~\ref{sec:bytes:deferred}). Registered extensions may
additionally mint \emph{system-derived} tags, which \MUST{} begin with
\texttt{MUSCS}, be exactly eight bytes, and not collide with a reserved
tag (Section~\ref{sec:bytes:system-derived}).

\begin{longtable}{>{\raggedright\arraybackslash}p{2.6cm} >{\raggedright\arraybackslash}p{10.4cm}}
  \toprule
  \textbf{Tag} & \textbf{Use} \\
  \midrule
  \endhead
  \texttt{MUSCCHNK} & Chunk content-hash preimage
    (\S\ref{sec:format:hashing}) \\
  \texttt{MUSCMANI} & Manifest chunk content-hash preimage
    (\S\ref{sec:format:hashing}) \\
  \texttt{MUSCBLOB} & Blob identity, bare preimage
    (Requirement~\ref{req:format:blob-hash-shape}) \\
  \texttt{MUSCMNIF} & Manifest-id derivation
    (Requirement~\ref{req:format:manifest-id}) \\
  \texttt{MUSCCONF} & \texttt{ConflictId} derivation
    (\S\ref{sec:semops:conflict-id}) \\
  \texttt{MUSCENVH} & Operation-envelope hash (\texttt{EnvelopeHash})
    (\S\ref{sec:semops:equivocation}) \\
  \texttt{MUSCSVCE} & System-promoted voice id
    (\S\ref{sec:graph:promoted-voices}) \\
  \texttt{MUSCSPCH} & System-derived pitch id
    (Requirement~\ref{req:graph:system-derived-pitch-id}) \\
  \texttt{MUSCSANM} & Integrity-anomaly id
    (Requirement~\ref{req:graph:integrity-anomaly-id}) \\
  \texttt{MUSCFNTM} & Font-metrics hash for layout conformance ---
    \emph{non-canonical} (\S\ref{sec:layoutir:catalog-identity}) \\
  \texttt{MUSCLOID} & Layout-object id --- \emph{non-canonical},
    wired by the reference implementation as of P12-I2
    (Requirement~\ref{req:layoutir:object-id-derivation}) \\
  \bottomrule
\end{longtable}

\section{Reference-implementation locks}
\label{sec:bytes:goldens}

Each ratified layout above is anchored by a test in the reference
implementation that fails deliberately if the layout drifts. Most are
\emph{golden-bytes} tests pinning the literal encoding; the two noted
exceptions (\texttt{BlobId} and \texttt{RationalTime}) are anchored by
round-trip and canonicalization tests, which catch a drift in decode or
normalized form but not every literal-byte change. These anchors make
the spec-to-code correspondence checkable.

\begin{longtable}{>{\raggedright\arraybackslash}p{3.7cm} >{\raggedright\arraybackslash}p{9.3cm}}
  \toprule
  \textbf{Layout} & \textbf{Anchoring test} \\
  \midrule
  \endhead
  \texttt{TypedObjectId} &
    {\footnotesize\texttt{epiphany-core/src/ids.rs}:
    \texttt{typed\bul object\bul id\bul byte\bul form\bul is\bul locked}} \\
  Promoted \texttt{VoiceId} &
    {\footnotesize\texttt{epiphany-core/src/graph.rs}:
    \texttt{promoted\bul voice\bul id\bul byte\bul form\bul is\bul locked}} \\
  System \texttt{PitchId} &
    {\footnotesize\texttt{epiphany-core/src/pitch.rs}:
    \texttt{system\bul pitch\bul id\bul byte\bul form\bul is\bul locked}} \\
  \texttt{Integrity\cwb AnomalyId} &
    {\footnotesize\texttt{epiphany-ops/src/anomaly.rs}:
    \texttt{integrity\bul anomaly\bul id\bul byte\bul form\bul is\bul locked}} \\
  \texttt{ObjectKind} &
    {\footnotesize\texttt{epiphany-ops/src/support.rs}:
    \texttt{object\bul kind\bul discriminants\bul are\bul golden}} \\
  \texttt{ChunkKind} &
    {\footnotesize\texttt{epiphany-bundle/src/chunk.rs}:
    \texttt{chunk\bul kind\bul discriminants\bul are\bul golden}} \\
  \texttt{Compression\cwb Algorithm} &
    {\footnotesize\texttt{epiphany-bundle/src/chunk.rs}:
    \texttt{compression\bul algorithm\bul encoding\bul is\bul golden}} \\
  \texttt{ProfileId} &
    {\footnotesize\texttt{epiphany-bundle/src/superblock.rs}:
    \texttt{profile\bul id\bul discriminants\bul are\bul golden}} \\
  \texttt{ManifestId} &
    {\footnotesize\texttt{epiphany-bundle/src/ids.rs}:
    \texttt{manifest\bul id\bul is\bul content\bul derived\bul and\bul deterministic}} \\
  \texttt{Transaction\cwb Category} &
    {\footnotesize\texttt{epiphany-ops/src/payload.rs}:
    \texttt{transaction\bul category\bul discriminants\bul are\bul golden}} \\
  \texttt{ResolutionAction} &
    {\footnotesize\texttt{epiphany-ops/src/conflict.rs}:
    \texttt{resolution\bul action\bul discriminants\bul are\bul golden}} \\
  \texttt{BlobId} &
    {\footnotesize\texttt{epiphany-determinism/src/hash.rs}:
    \texttt{ContentHash::of\bul blob} and its round-trip locks
    \emph{(round-trip, not a literal-byte golden)}} \\
  \texttt{RationalTime} &
    {\footnotesize\texttt{epiphany-core/src/time.rs}:
    \texttt{equal\bul rationals\bul encode\bul identically}
    \emph{(canonical-equality lock, not a literal-byte golden)}} \\
  \bottomrule
\end{longtable}

The \texttt{MUSCLOID} layout-object-id derivation
(Requirement~\ref{req:layoutir:object-id-derivation}), wired in P12-I2,
is anchored differently because it is \emph{non-canonical}: its ids
never enter document state or a content hash, so there is no durable
byte layout to golden-lock. It is instead pinned by a derivation
reference-lock test (\texttt{epiphany-layout-ir/src/provenance.rs}:
\texttt{stable\bul id\bul uses\bul the\bul ratified\bul muscloid\bul derivation}),
which fixes the tagged preimage; by the tag-spelling lock
(\texttt{epiphany-determinism/src/domain.rs}:
\texttt{exact\bul tag\bul spellings\bul match\bul spec}); and by the
renderer's provenance goldens, whose \texttt{data-prov} hex changes if
the derivation drifts.

\section{Layouts deferred to the companions}
\label{sec:bytes:deferred}

The following are intentionally \emph{not} ratified here; they remain
provisional pending the Binary Format companion and the Operation
Catalog, and are recorded so the companion authors know precisely which
layouts they own versus inherit:

\begin{itemize}
  \item \texttt{OperationKindTag} / \texttt{OperationKind} discriminants
    and the full per-operation payload encodings. The reference
    implementation assigns provisional declaration-order discriminants
    (\texttt{0}--\texttt{16}) locked only for mutual distinctness, not
    to literal values; pinning the literal wire form is the
    Operation Catalog's and Binary Format companion's job.
  \item The full composite struct layouts --- whole-\texttt{Score},
    operation envelopes, the manifest body's field order --- together
    with schema-version wire evolution and varint details beyond the
    convention baseline
    (Requirement~\ref{req:format:codec-conventions}). The companion
    inherits the baseline and formalizes these.
\end{itemize}

\chapter{Revision History}
\label{app:history}

\begin{longtable}{p{2cm} p{2.5cm} p{9cm}}
  \toprule
  \textbf{Date} & \textbf{Section} & \textbf{Change} \\
  \midrule
  \endhead
  \today & All & Initial scaffold; chapter structure and front matter
  established. \\
  \today & Ch.~\ref{ch:pitch} & Pitch chapter written: layered pitch type,
  scale position, acoustic realization, spelling subsystem with attachments,
  precedence, pre-pass, and analytical layers. Chapter structure expanded:
  pitch separated from time/tuning; semantic operations chapter added. \\
  \today & Ch.~\ref{ch:time} & Time and duration chapter written: exact
  rational time with inline-or-promoted representation, distinct position
  and duration newtypes, wall-clock time as fixed-point nanoseconds, time
  anchors, time signatures with explicit beat groups including irrational
  meters, tempo maps with deterministic conversion, notational
  decomposition pre-pass, tuplets as grouping objects, and three region
  time models (metric, proportional, aleatoric). \\
  \today & Ch.~\ref{ch:tuning} & Tuning systems and pitch spaces chapter
  written: pitch spaces as analytical universes distinct from tuning
  systems, position structures (chromatic, diatonic-over-chromatic, JI
  lattice, registered), interval algebras, nominal and accidental
  registries, SMuFL glyph references with versioning, tuning resolution
  variants including adaptive tuning, score-level reference pitch,
  hierarchical resolution rules, and the normative built-in catalog of
  pitch spaces and tuning systems. \\
  \today & Ch.~\ref{ch:graph} & Score graph chapter written: hybrid
  tree-plus-cross-cutting topology, canvas-region-staff-voice-event
  hierarchy with regions carrying time and content models, typed 128-bit
  identifiers with replica-plus-counter generation, flat event arena,
  complete event taxonomy (pitched, unpitched, rest, indeterminate,
  trajectory, graphic, cue), cross-cutting registry partitioned by type,
  graphic content with region-local coordinates and time bindings, parts
  as projections separate from the canvas, analysis layers and views,
  graph invariants, and required indexes. \\
  \today & Ch.~\ref{ch:semops} & Semantic operations chapter written:
  operation framework with dual command/delta layering, primitive and
  compound operation distinction, transactions as atomic user units,
  precondition classification (invariant vs. advisory) with strict
  authoring and lenient replay modes, the normative re-anchoring rule
  table, inverse-based undo with single-replica scope, semantic merge
  semantics with five merge families, and representative operations
  illustrating each major category. The full operation catalog is
  deferred to a separate conformance specification. \\
  \today & Ch.~\ref{ch:layout-ir} & Layout IR chapter written: pipeline
  of four stages (LogicalLayoutIR, ConstrainedLayoutIR,
  ResolvedLayoutIR, RenderIR-at-interface), shared provenance discipline,
  spatial primitives in staff spaces with f32 precision, polymorphic
  time axis trait with metric/proportional/aleatoric implementations,
  composite logical objects flattened to glyph-level constrained
  objects, spring-based horizontal and vertical layout, explicit
  engraving decision recording, hard and soft override mechanism with
  preservation rules, fine-grained dependency tracking with invalidation
  rules, glyph catalog interface, and uniform region containers across
  kinds. \\
  \today & Ch.~\ref{ch:format} & File format chapter written: bundle
  container structure with header, manifest, chunk store, and blob
  store; custom schema-defined binary format with self-describing
  extension envelope; content-addressed immutable chunks with manifest
  as the only mutable index; multi-level chunking (score-level,
  staff-region, locality-grouped cross-cutting, attachment chunks);
  operation log as canonical source of truth with dotted version
  vectors and hybrid logical clocks; snapshots as derivative caches;
  schema versioning with lazy migration and forward compatibility;
  format profiles (Full, ReadOnly, Lite); per-chunk zstd compression;
  atomic write protocol; deterministic s-expression text projection;
  streaming read guarantees. DRM and encryption explicitly excluded. \\
  \today & Ch.~\ref{ch:solver} & Constraint-solver interface chapter
  written: interface-and-contract specification rather than algorithm;
  the solver trait with from-scratch and incremental modes; Cassowary-
  style Required and Preferred constraint strengths; the constraint
  family catalog (spring, collision, alignment, containment, break,
  aesthetic, extension); the normative quality metric set (optical
  spacing, collisions, page fill, casting-off, beam angle, slur
  curvature, vertical balance, symbol density uniformity); Pareto-
  bounded quality target with configurable slack and tie-breaking
  weights for ambiguous cases; within-implementation byte-determinism
  with cross-implementation Pareto equivalence; incremental solving
  contract with propagation bounds; structured error reporting with
  partial solutions on budget exhaustion; reference suite for
  conformance; non-normative recommendation of Cassowary as a
  starting algorithm. \\
  \today & Chs.~\ref{ch:perf},~\ref{ch:extension} &
  Performance Requirements chapter expanded from skeleton: reference
  hardware profile (2026 desktop and tablet baselines), frame budgets
  for interactive edits at p99, cold-open targets, memory budgets
  separating core data from evictable caches, audio thread discipline
  requirements crossing the core/audio boundary, constraint-solver
  performance targets, file format performance targets, measurement
  methodology, regression testing recommendations. Extension Points
  chapter written from skeleton: consolidated catalog of every
  extension point in the spec organized by chapter of origin
  (notation grammar, tuning system, glyph, score graph, layout IR,
  constraint solver, file format extensions), the extension registry
  contract requiring versioning and identifier stability, plugin
  runtime considerations, conformance interactions. \\
  \today & Appendices & Glossary expanded with 50+ defined terms
  cross-referenced to their originating chapters. Bibliography
  populated with engraving treatises, music theory references,
  tuning and microtonality sources, constraint-based layout
  literature, CRDT references, and prior-art notation format
  citations. \\
  \today & All & Platform named \textit{Epiphany}. Typography refreshed
  to TeX Gyre Pagella (body), Heros (sans), and Cursor (mono).
  Color palette refined to deep teal, antique gold, warm parchment
  neutrals, and deep crimson accents. Title page redesigned with
  ornamental rules and an oldstyle classical feel. Status section,
  introduction, and PDF metadata updated to reflect the platform's
  name. \\
  \today & Pass 1 revision (Chs.~\ref{ch:pitch},~\ref{ch:time},~\ref{ch:tuning},~\ref{ch:graph},~\ref{ch:semops}) &
  Type-seam corrections in response to external review.
  Introduced \texttt{IdentifiedPitch} wrapping every embedded pitch
  with a stable \texttt{PitchId}; updated \texttt{PitchedEvent},
  \texttt{TrajectoryEndpoint}, and \texttt{TrajectoryShape::Stepwise}
  accordingly. Introduced \texttt{EventDuration} union covering
  musical, wall-clock, and indeterminate cases, with concrete-only
  \texttt{DurationBounds} avoiding recursive indeterminacy.
  Aleatoric regions now declare an explicit
  \texttt{AleatoricAnchoringDiscipline}. \texttt{PitchSpelling}
  carries a stack of accidentals (\texttt{Vec<AccidentalId>}) with
  normative additive semantics and duplicate-detection invariants,
  replacing the prior \texttt{Option<AccidentalId>}.
  \texttt{PitchSpacePosition} gains a \texttt{JiVector} variant for
  first-class just-intonation lattice positions. The default CMN
  pitch space is now correctly described as
  \texttt{cmn-12} (diatonic-over-chromatic), not pure chromatic.
  Chapter~\ref{ch:pitch}'s reference-pitch text rewritten to defer
  to Chapter~\ref{ch:tuning}, resolving the prior ownership
  contradiction. \texttt{MusicalPosition} and \texttt{MusicalDuration}
  now derive \texttt{Clone}, not \texttt{Copy}, with a permitted
  inline-or-pointer-tagged implementation noted as conforming.
  \texttt{Tie::pitch\_pairing} now uses \texttt{Vec<(PitchId, PitchId)>}
  instead of positional indices. The duplicate \texttt{AnalysisLayer}
  definition in Chapter~\ref{ch:pitch} was removed in favor of the
  authoritative definition in Chapter~\ref{ch:graph}. Graph
  invariants expanded to cover identified pitches, tombstoning,
  time-model agreement, and \texttt{PitchId}-pairing tie targets.
  Representative operations (\texttt{InsertEvent},
  \texttt{DeleteEvent}, \texttt{RespellPitch}, \texttt{Transpose})
  updated for atomic ID minting, pitch tombstoning,
  voice-promotion on concurrent insertion, and \texttt{PitchId}
  preservation through transposition. Status downgraded to
  ``structural working draft.'' \\
  \today & Pass 2 revision (Chs.~\ref{ch:time},~\ref{ch:graph},~\ref{ch:semops}) &
  Graph and time invariants strengthened in response to external
  review. \texttt{TimeAnchor} offsets generalized via the
  \texttt{AnchorOffset} union (Musical, WallClock, Zero), with
  normative agreement between offset variant and target region's
  time model. \texttt{TempoSegment} restructured to carry explicit
  \texttt{end: Option<TimeAnchor>}, \texttt{start\_tempo},
  \texttt{end\_tempo: Option<Tempo>}, and \texttt{shape}; tempo
  segments are required to be monotonic in musical time so that
  wall-clock-to-musical conversion is single-valued. Polymeter
  admitted natively: \texttt{StaffBasedContent} now carries
  \texttt{Vec<StaffInstance>} (region-local), plus
  \texttt{default\_metric\_grid: Option<MetricGrid>} and
  \texttt{barline\_alignment\_groups: Vec<BarlineAlignmentGroup>}.
  Measures and metric grids are per-staff-instance, not per-region.
  Staff identity split into the global \texttt{Staff} (declared once
  at score level) and the region-local \texttt{StaffInstance};
  \texttt{Score} top-level structure expanded with
  \texttt{staves: Vec<Staff>} and
  \texttt{staff\_groups: Vec<StaffGroup>}.
  \texttt{DeleteEvent} preconditions strengthened: target events
  belonging to tuplets require explicit
  \texttt{TupletCompensation} (\texttt{NotInTuplet},
  \texttt{ReplaceWithRest}, \texttt{RewriteTuplets},
  \texttt{CascadeDeleteTuplets}). Tie validation reframed from
  universal immediate-adjacency to class-specific profiles:
  \texttt{TieClass} (Standard, Editorial, CrossVoice,
  LaissezVibrer, Registered) admits editorial, cross-voice, and
  laissez-vibrer ties without relaxing standard-tie invariants.
  System-promoted voices formalized: \texttt{VoiceOrigin} enum
  (UserDeclared, Imported, SystemPromoted); promoted voices have
  deterministically-derived \texttt{VoiceId}s, are first-class
  graph objects, are visible to users, and are normalized only by
  explicit user operation. Graph invariants expanded from 14 to 18
  normative invariants reflecting staff-instance ontology,
  per-staff measures, anchor-offset agreement, system-promoted
  voice identity, and barline-alignment-group containment. New
  identifier types: \texttt{StaffInstanceId},
  \texttt{StaffGroupId}, \texttt{BarlineAlignmentGroupId},
  \texttt{TimeSignatureId}. \\
  \today & Pass 3 revision (Ch.~\ref{ch:semops},~\ref{ch:format}) &
  Concurrent semantics rewritten from pairwise operation merge to
  operation-set deterministic reduction. The replicated operation
  set is now framed as a true CRDT (a grow-only set of envelopes);
  the materialized score graph is its deterministic reduction, not
  itself a CRDT. \texttt{OperationId} (replica + counter) separated
  from \texttt{OperationStamp} (HLC + id); canonical reduction order
  is causal-then-HLC-then-lexicographic. Causal contexts expressed as
  dotted version vectors (\texttt{CausalContext}), not exhaustive
  predecessor lists; exhaustive lists are explicitly forbidden as
  canonical representation. \texttt{OperationEnvelope} introduced
  with id, author, stamp, causal context, optional transaction,
  payload. Four-phase operation lifecycle specified: Prepare,
  Commit, Reduce, Report. Operation effects formalized:
  \texttt{Applied}, \texttt{AppliedWithRepair}, \texttt{Conflicted},
  \texttt{TombstonedTarget}, \texttt{NoOp\{reason\}}; every replica
  reducing the same set records the same effects with the same
  reasons. Conflict records made first-class graph objects in a
  score-level \texttt{ConflictRegistry}: stable, addressable,
  user-visible, resolvable by subsequent
  \texttt{ResolveConflict} operations. Tombstones formalized with
  explicit retention rules (pending replication, undo, conflict
  resolution, attachment resolution). Re-anchoring rewritten as a
  total deterministic function: ``nearest'' is now a strict
  lexicographic minimum over (containment proximity, time distance,
  direction preference, object id), with no implementation
  discretion. Re-anchoring rule table updated for the new
  staff-instance ontology and the deterministic ``nearest''
  function. Transactions formalized as atomic replicated groups:
  primitive members \MUSTNOT{} materialize as independently
  visible intermediate states; failures conflict the whole
  transaction. Undo redefined as a forward compensating operation
  (\texttt{UndoTransaction}) with three policies
  (\texttt{StrictInverse}, \texttt{BestEffort}, \texttt{Cascade});
  bit-identical undo demoted to content equivalence. LWW reduction
  restricted to explicitly-marked \texttt{LwwAdvisory} scalar
  fields (metadata, view preferences, layout hints, system/page
  break preferences, optional cross-cutting display hints);
  structural musical fields \MUSTNOT{} default to LWW. Cache
  invalidation rule: caches are acceleration only, never canonical;
  invalidated on operation-set, causal-order, reduction-algorithm,
  or profile change. Per-operation merge families renamed to
  reduction rules and rewritten in the new model;
  \texttt{ChangeRegionTimeModel} now correctly conflicts rather
  than silently coercing incompatible contained events;
  \texttt{SetUserSystemBreak} explicitly marked LwwAdvisory;
  \texttt{RespellPitch} concurrent collisions preserve the losing
  spelling in a structural-field-collision conflict record. New
  identifier types: \texttt{ConflictId}, \texttt{TransactionId}.
  Operation log restructured to store envelopes rather than
  primitive operations with separate stored-sequence metadata. \\
  \today & Pass 4 revision (Ch.~\ref{ch:format}) &
  File-format restructured around Pass 3's operation-set
  canonical-document model. The earlier ``single mutable manifest''
  story replaced with a fixed 64-byte header plus two fixed-size
  256-byte superblock slots; the superblocks are the only mutable
  on-disk objects, and atomic commit is achieved by writing the
  inactive superblock slot with generation $+ 1$. Active superblock
  selected by highest valid generation. Generation gap greater than
  one detected as integrity anomaly with read-only recovery mode.
  Manifest reframed as a table of roots and declarations
  (operation\_roots, retained\_snapshots, blob\_roots,
  profile\_declarations, extension\_declarations, text\_projection
  root, integrity\_root); user-facing metadata explicitly removed
  from the manifest in favor of operation envelopes and reduced
  state. BLAKE3 fixed as the single content-hashing algorithm,
  replacing the prior ``BLAKE3 or SHA-256'' uncertainty. Hash
  preimage domain-separated and explicitly excludes
  \texttt{compression\_algorithm} to preserve deduplication across
  compression choices. Operation-envelope blocks specified as the
  canonical storage unit with a soft 1\,MiB uncompressed target
  size and a default 64\,MiB upper bound per profile. Snapshots
  separated into acceleration snapshots (caches, discardable) and
  canonical-base snapshots (declared in
  \texttt{retained\_snapshots}, required for pruning). Pruning
  protocol specified: canonical state preserved as (base snapshot
  $+$ post-frontier envelopes). Atomic write protocol specified
  step-by-step with explicit durable-flush points; the durable
  flush is platform-abstracted to cover \texttt{fsync},
  \texttt{FlushFileBuffers}, and equivalents. Crash recovery rules
  enumerated by failure point. Garbage collection specified as
  conservative, optional, deferred; never on the commit critical
  path. Edit barriers formalized as a structured first-class
  concept (\texttt{EditBarrier} with scope, affected object kinds,
  prohibited operation kinds, and condition), with an
  unsafe-edit escape hatch that tombstones extension data rather
  than silently corrupting it. Text projection reframed:
  semantic round-trip required (envelopes, declarations, reduced
  state); physical-layout fidelity (chunk offsets, compression
  choices, cache chunks, garbage regions, superblock generation)
  explicitly not required. Schema versioning, format profiles
  (Full, ReadOnly, Lite, Custom), and compression sections updated
  for the new architecture. Streaming-read cold-open procedure
  rewritten to use snapshot $+$ envelope-stream model. \\
  \today & Pass 5 revision (Ch.~\ref{ch:solver}) &
  Solver conformance reframed from universal Pareto optimality to
  reference-suite-based conformance with profile-specific
  thresholds. The earlier ``no better Pareto layout exists''
  obligation, which quantified over an undecidable space of
  possible layouts, has been replaced with checkable, falsifiable
  thresholds on a reference suite. Solver obligations split into
  four classes: validity conformance (hard constraints, provenance,
  unit system, invariants), determinism conformance, diagnostic
  conformance (\texttt{SolveStatus}, \texttt{SolveReport}, warnings,
  unsatisfied-constraint reporting, metric vectors), and quality
  conformance (reference-suite thresholds). Three solver tiers
  introduced (\texttt{Minimal}, \texttt{Standard}, \texttt{Advanced})
  with declared obligations per tier; tiers are orthogonal to
  file-format profiles. Deterministic budgets formalized
  (\texttt{max\_iterations}, \texttt{max\_nodes},
  \texttt{max\_constraint\_evaluations}); wall-clock time demoted
  to advisory only. Wall-clock \MUSTNOT{} change the canonical
  layout. Quality metrics normalized: \texttt{NormalizedMetric} is
  a finite \texttt{f64} in $[0.0, 1.0]$ with $0.0$ best and $1.0$
  worst; orientation and range fixed across all metrics including
  extension metrics. Hard constraints explicitly separated from
  quality metrics: a solver \MUSTNOT{} trade off a hard-constraint
  violation for a better quality score. Within-implementation
  determinism (byte-equal output for identical input + version)
  separated from cross-implementation conformance (reference-suite
  thresholds, not byte equality). \texttt{SolveReport} unified the
  prior \texttt{Result<SolveOutput, SolverError>} return: every
  invocation yields a report with a \texttt{SolveStatus}
  (\texttt{Solved}, \texttt{SolvedWithWarnings},
  \texttt{PartialBudgetExhausted}, \texttt{Unsatisfiable},
  \texttt{InternalError}). \texttt{PartialBudgetExhausted}
  preserves hard-constraint validity; budget exhaustion that
  cannot satisfy hard constraints returns \texttt{Unsatisfiable}.
  Six invalidation scopes specified (\texttt{ObjectLocal},
  \texttt{MeasureLocal}, \texttt{SystemLocal}, \texttt{PageLocal},
  \texttt{RegionLocal}, \texttt{WholeScore}); incremental solving
  contract reframed as observational equivalence to scoped full
  solving (solvers \MAY{} widen scopes conservatively but
  \MAY NOT{} narrow them). Reference algorithm demoted from
  ``recommended starting point'' to explanatory and diagnostic
  only; conforming implementations \MAY{} choose any algorithm.
  Pass 4 rollback-language corrected: rollback is achieved by
  retained manifests in the chunk store, not by retaining multiple
  on-disk superblock generations (the fixed layout has only two
  slots). \\
  \today & Pass 6 revision (App.~\ref{app:determinism}) &
  New Determinism Contract appendix consolidates reproducibility
  obligations scattered across Chapters~\ref{ch:pitch}--\ref{ch:solver}
  into a single normative reference. Five layers of determinism
  distinguished: canonical score determinism (byte-equal
  cross-implementation), canonical serialization determinism
  (byte-equal chunk hashes and text projection), layout
  determinism (byte-equal within implementation version),
  conformance determinism (reference-suite-based, not byte-equal
  cross-implementation), and non-canonical cache determinism
  (caches may vary, MUST NOT affect canonical state). Canonical
  floating-point values restricted to finite IEEE 754 binary64;
  NaN and infinity prohibited from canonical chunks;
  \texttt{-0.0} canonicalized to \texttt{+0.0} before storage and
  hashing. Little-endian IEEE 754 binary64 serialization fixed.
  IEEE 754 round-to-nearest, ties-to-even mandated for canonical
  numerical algorithms; \texttt{-ffast-math} and equivalents
  prohibited; reassociation, distribution, and implicit
  fused-multiply-add forbidden where they could change canonical
  output. Quantized layout coordinates introduced
  (\texttt{QuantizedCoord} at $1/1024$ staff space) so canonical
  \texttt{ResolvedLayoutIR} positions are stable across
  implementations whose internal computations agree to within
  half-quantum. Tolerance classes formalized
  (\texttt{AcousticCents}, \texttt{LayoutCoordinate},
  \texttt{QualityMetric}, \texttt{TempoIntegration},
  \texttt{SolverResidual}) with governance category
  (\texttt{Equality}, \texttt{Validation}, \texttt{Diagnostic});
  ad-hoc epsilons prohibited; tolerances MUST NOT apply to
  identity. Canonical iteration order specified for every
  enumerable collection (identifiers lexicographic on canonical
  byte form; operations causal-then-HLC-then-id; chunk references
  by kind, hash, offset; conflicts by id; extensions by id then
  version; strings byte-lexicographic on NFC UTF-8). Text
  normalization fixed at UTF-8 NFC; locale prohibited from
  affecting canonical parsing, serialization, sorting, decimal
  formatting, or projection. Randomness MUST NOT enter canonical
  state unless the seed is an explicit canonical input.
  Parallelism permitted but MUST produce results identical to a
  strict single-threaded baseline through deterministic ordered
  reductions. Compression confirmed as non-canonical (identical
  uncompressed content has identical hash regardless of
  compression choice); recompression permitted without altering
  canonical state. Open algorithm hooks (spelling pre-pass,
  decomposition, tempo integration, root-finding) required to be
  either normatively specified, profile-declared by versioned
  identifier, or explicitly non-canonical/advisory; vendor
  freedom over canonical algorithms prohibited. Conformance
  statement enumerates the required disclosures implementations
  must publish (libm, parallelism strategy, NFC implementation,
  declared algorithms for open-question areas). Status section
  updated to reference Appendix~\ref{app:determinism}. \\
  \today & Pass 7 revision (editorial, all chapters) &
  Editorial consistency pass: no architectural changes, only
  cleanup. Resolved the contradiction between
  Appendix~\ref{app:determinism}'s ``byte-equal canonical state
  cross-implementation'' and the existence of open canonical
  algorithms by adding an explicit conditional on the disposition
  of Section~\ref{sec:det:open}. Stale section-number cross-
  references replaced with symbolic labels (Section~2.4 $\to$
  \texttt{sec:pitch:spelling}; Section~2.4.4 $\to$
  \texttt{sec:pitch:precedence}; Section~5.9.3 $\to$
  \texttt{sec:graph:timebinding}; Section~6.6 $\to$
  \texttt{sec:semops:conflicts}; Section~6.8 $\to$
  \texttt{sec:semops:undo}). Terminology normalized:
  ``merge family'' $\to$ ``reduction rule'' where the operation
  framework was meant; ``Pareto-bounded quality model'' $\to$
  reference-suite/normalized-metric language in extension points;
  glossary entries for Operation log, Operation set, Operation
  envelope, Canonical reduction, Pareto frontier (design target),
  Staff, Staff instance, Conflict record, CRDT, Snapshot, and
  Transaction rewritten or added. New Appendix~\ref{app:deferred}
  (Intentionally Deferred Types and Specifications) catalogs
  companion specifications (Binary Format, Text Projection,
  Profile Conformance, Quality Metric Catalog, Reference Suite,
  Performance Reference Suite, Operation Catalog, Reference
  Algorithm), externally provided components (SMuFL, font metric
  provider, audio engine, plugin runtime, collaboration
  transport, user interface), open canonical algorithms, and
  extension registry catalogs; distinguishes deliberately external
  content from accidentally missing content. Status section
  rewritten to reflect structural stabilization: substantive
  review passes complete, architecture not under active redesign,
  remaining blockers are companion specifications and open
  canonical algorithms; the section enumerates what conformance
  can be established now and what awaits companion delivery.
  Revision history entries for Passes 1--7 read as the document's
  architectural-stabilization log, consolidating the design
  decisions taken in response to external review. \\
  \today & Pass 8 revision (post-review correctness pass) &
  Determinism correctness pass addressing review of the
  Pass 7 output. \texttt{ConflictId} reframed as content-derived
  via BLAKE3 truncation with domain tag \texttt{"MUSCCONF"},
  fixing a determinism leak where each replica would have minted
  its own conflict identifiers from local replica-counter
  sequences; conflict records produced during reduction now have
  identifiers that agree across replicas by construction.
  \texttt{RepairRecord} defined (was referenced by
  \texttt{OperationEffect::AppliedWithRepair} without
  definition), with a \texttt{RepairKind} enumeration covering
  \texttt{Reanchored}, \texttt{SpannerTruncated},
  \texttt{Orphaned}, \texttt{CascadeDeleted},
  \texttt{AttachmentTombstoned}, \texttt{VoicePromoted},
  \texttt{TupletCompensated}, and \texttt{Registered}.
  \texttt{TypedObjectId} defined as a tagged union over every
  identifier kind in the score graph, with canonical-bytes
  serialization for hashing/ordering/equality.
  \texttt{OperationKind} defined as a tagged enumeration of
  primitive operation variants, plus a \texttt{Registered}
  variant for extension-defined primitives. System-derived
  identifier namespace formalized:
  \texttt{ReplicaId::SYSTEM\_DERIVED} reserved at
  \texttt{0xffff\_ffff\_ffff\_ffff}, with the 64-bit counter
  derived deterministically via BLAKE3 truncation; domain tags
  for system identifiers (\texttt{"MUSCSVCE"} for
  system-promoted voices, \texttt{"MUSCSPCH"} for system-derived
  pitches) enumerated; user-authored replicas \MUSTNOT{} use
  the reserved value. Envelope-acceptance rules added: a
  well-formed envelope requires non-system replica id, finite
  physical time, well-formed DVV, deserializable payload;
  per-replica monotonicity enforced; clock-skew warnings allowed
  but reordering for plausibility prohibited; authentication
  deferred to collaboration transport. Canonical-base snapshot
  semantics repaired: manifest split into single
  \texttt{canonical\_base: Option<SnapshotRef>} plus
  \texttt{acceleration\_snapshots: Vec<SnapshotRef>}; coverage
  test specified as DVV-frontier membership, explicitly
  forbidding HLC-stamp comparison for canonicality determination;
  pruning replaces the canonical base atomically. Header CRC
  moved to the last 4 bytes of the 64-byte fixed header,
  resolving a circular-definition bug where the CRC field
  appeared inside its own coverage range. Manifest chunks
  \MUST{} be stored uncompressed in this format version,
  resolving the bootstrap chicken-and-egg problem (a reader
  needs to know how to decode the manifest before it has any
  manifest information). \texttt{ChunkId} defined as a newtype
  around \texttt{ContentHash}, clarifying the prior comment that
  the hash and id were redundant. Three-layer identity model for
  documents: \texttt{file\_uuid} (physical bundle, changes on
  Save As), \texttt{document\_id} (logical work, persists across
  Save As by default), optional \texttt{lineage\_id} (shared
  ancestor for genealogy/version-control purposes); Save As
  semantics specified normatively. \texttt{RetentionPolicy}
  defined as a first-class concept with
  \texttt{retain\_previous\_manifests},
  \texttt{retain\_duration}, \texttt{retain\_named\_checkpoints};
  rollback achieved via retained manifests in the chunk store,
  not multiple on-disk superblock generations.
  \texttt{BarrierCondition} defined (\texttt{Always},
  \texttt{ObjectExists}, \texttt{ObjectHasExtensionData},
  \texttt{All}, \texttt{Any}, \texttt{Not}, \texttt{Registered}),
  making the edit-barrier model implementable for unknown
  extensions. \texttt{GlyphCatalogIdentity} defined in
  Chapter~\ref{ch:layout-ir}, required for any byte-equal-output
  layout conformance claim; without it, layout output is
  implementation-specific. \texttt{Box<dyn TimeAxis>} replaced
  with \texttt{TimeAxisModel} enum (Metric, Proportional,
  Aleatoric, Registered), giving the layout IR a canonical
  representation for serialization, hashing, and conformance
  comparison while still permitting internal trait-object
  dispatch. Performance chapter restructured around an explicit
  core/product/Performance-Reference-Suite obligation split:
  core obligations bound the core's own work (reduction, layout,
  file writes, snapshot acquisition), while end-to-end UI/audio
  targets are product-layer obligations specified in companion
  documents. Editorial sweeps: remaining ``merge semantics''
  prose in Chapter~\ref{ch:intro}, the compound transpose
  primitive, the \texttt{VoiceOrigin::SystemPromoted} comment,
  and the extension-point design principles updated to
  ``reduction semantics'' or ``reduction rules''; tempo curve,
  decomposition, and spelling pre-pass open questions updated
  to reference Appendix~\ref{app:determinism}
  Section~\ref{sec:det:open} dispositions. Glossary additions:
  Edit barrier, Glyph catalog identity, Operation kind, Repair
  record, Retention policy, Typed object identifier, Document
  identifier (with cross-references to FileUuid and LineageId).
  Reading Guide canonical-vs-non-canonical table updated to
  reflect the canonical\_base / acceleration\_snapshots split. \\
  \today & Pass 9 revision (correctness and consistency pass) &
  Operation-envelope duplicate handling formalized into three
  outcomes: \texttt{IdempotentDuplicate} (same id, byte-identical
  envelope: silently dropped), \texttt{FirstAcceptance} (new id:
  accepted after well-formedness check), and
  \texttt{EnvelopeEquivocation} (same id, different canonical
  bytes: divergent envelope rejected, integrity diagnostic
  surfaced, bundle non-conforming if equivocation appears
  on-disk). Well-formedness invariant
  \texttt{envelope.stamp.id == envelope.id} made normative. HLC
  monotonicity reframed as a property of the set of accepted
  envelopes, not arrival order: out-of-order delivery is fine;
  on detection of inconsistency, both envelopes are retained but
  the offending replica's stream from the violating counter
  onward is quarantined and the document is flagged with a
  replica-equivocation anomaly. The free-form
  \texttt{NoOpReason::PreconditionFailedUnderReduction { reason:
  String }} replaced with a typed \texttt{PreconditionFailureReason}
  enum, closing a determinism leak (free-form strings in
  canonical state could otherwise differ across implementations).
  Transaction descriptor ordering integrated into canonical
  reduction via causal dependency: every primitive member of a
  transaction \MUST{} causally depend on its
  \texttt{DeclareTransaction} envelope; members lacking such a
  dependency reduce to \texttt{NoOp\{TransactionConflict\}} with
  a synthesized conflict record. The fixed-header
  \texttt{file\_uuid} requirement reworded to match the Pass 8
  identity model: fixed for the lifetime of a physical bundle,
  preserved by filesystem byte-copy, freshly minted on Save As.
  Manifest \texttt{integrity\_root} made \texttt{Option<ChunkRef>}
  (it is an accelerator); new ``Canonical and Non-Canonical
  Manifest Roots'' subsection partitions the manifest's fields
  by their role in canonical state, with the rule that failed
  verification of non-canonical chunks is rebuildable rather
  than corruption. Cold-open procedure rewritten to reflect the
  canonical-base / operation-envelope-block model (the prior
  text retained pre-Pass-3 region-chunks language). Streaming-
  read locatability list updated. System-derived counter
  collisions: mandatory collision check during reduction;
  collision becomes hard corruption via new
  \texttt{ConflictKind::SystemIdentifierCollision} record naming
  both colliding canonical-input sets. \texttt{GlyphCatalogIdentity}
  \texttt{metrics\_hash} scope broadened: hash covers glyphs
  referenced by the \texttt{ConstrainedLayoutIR} (solver inputs),
  not just glyphs in the final \texttt{ResolvedLayoutIR}.
  \texttt{CommitState} validation added to superblock selection:
  non-\texttt{Committed} superblock is not valid for ordinary
  selection. Schema-version fallback refined into three cases by
  chunk role: canonical chunks with unsupported major schema
  require read-only preservation mode (or refusal); non-canonical
  chunks may be opaquely preserved; unknown canonical operation
  payloads require extension support or open read-only without
  materialization. New \texttt{OperationKindTag} (discriminator-
  only enum) introduced for edit barriers, conformance reports,
  and diagnostic surfaces; \texttt{EditBarrier} switched to
  \texttt{prohibited\_operation\_kinds: Vec<OperationKindTag>}.
  ``Absolute musical positions \MUSTNOT{} appear in stored
  data'' rule narrowed: stored \emph{references to external
  time points} use \texttt{TimeAnchor}; events' \emph{own
  positions} within their owning voice and region may be stored
  as \texttt{EventPosition}. New Appendix~\ref{app:deferred}
  Section~\ref{sec:deferred:support-types} catalogs the semantic
  identity rules for \texttt{FileUuid}, \texttt{DocumentId},
  \texttt{LineageId}, \texttt{ManifestId}, \texttt{SnapshotId},
  \texttt{BlobId}, \texttt{AuthorId},
  \texttt{ReductionAlgorithmVersion}, \texttt{SolverProfile},
  \texttt{ProfileId}, \texttt{ContentHash}, \texttt{ChunkId},
  \texttt{OperationId}, \texttt{TransactionId},
  \texttt{ConflictId}, and the typed identifier family.
  Editorial: \texttt{MUSTNOT} and \texttt{SHOULDNOT} macros
  switched to \texttt{nobreakspace} so \texttt{pdftotext}
  extracts ``MUST NOT'' as two tokens; Manifest, Operation log,
  and Snapshot glossary entries updated to match the
  \texttt{canonical\_base} / \texttt{acceleration\_snapshots}
  manifest field names. \\
  \today & Pass 10 revision (order-independence and integrity-
  anomaly cleanup) &
  \texttt{OperationId} equivocation reworked from a first-arrival
  rule (which made canonical state arrival-order-dependent across
  replicas observing the same envelopes in different orders) to an
  order-independent \texttt{OperationSlot} model. The slot is
  either \texttt{Single(OperationEnvelope)} or
  \texttt{Equivocated{} \{ operation\_id, candidates:
  BTreeSet<EnvelopeHash> \}}; observing a second distinct
  canonical envelope under an \texttt{OperationId} transitions
  the slot to \texttt{Equivocated} regardless of arrival order.
  Equivocated slots produce no canonical reduction effect; both
  candidates are retained for diagnostic recovery. Operations
  causally depending on an equivocated id are held pending
  resolution. Resolution is by explicit policy: transport-level
  credential revocation, local user reconciliation, or a
  profile-declared deterministic selection function. No
  arrival-order default applies. HLC monotonicity-anomaly
  handling sharpened: a new \texttt{AnomalousReplicaSegment}
  type captures the offending replica, the first violating
  counter, the reason, and the excluded envelope ids. Operations
  from the violating counter onward are retained but
  \emph{excluded from canonical reduction}; canonical state is
  derived only from envelopes that satisfy the monotonicity
  invariant. The earlier ``quarantined'' wording is now backed
  by an explicit structural exclusion rule. System-derived ID
  collisions moved out of \texttt{ConflictKind} and into a new
  \texttt{IntegrityAnomaly} type with kinds
  \texttt{SystemIdentifierCollision},
  \texttt{OperationSlotEquivocated}, and
  \texttt{ReplicaStreamQuarantined}. The distinction is that
  conflict records are ordinary canonical-state facts (a
  representable musical conflict between two valid edits) while
  integrity anomalies indicate the structural assumptions
  underlying canonical state have failed and require external
  recovery; an integrity anomaly takes the document out of
  ordinary canonical operation. Stale ``refined in a later
  revision'' parenthetical removed from the \texttt{DeleteEvent}
  reduction rule; the re-anchoring total order is part of the
  current architecture and needs no further refinement. Typo
  fixed in the superblock-selection requirement
  (\texttt{MAY{}-be-valid} replaced with the prose ``invalid for
  ordinary selection but may be inspected in diagnostic recovery
  mode''). Glossary chapter references for ``Dotted version
  vector'' and ``Hybrid logical clock'' updated to ``Chapters
  6, 8'' with explicit ``defined operationally in
  Chapter~\ref{ch:semops}; serialized in
  Chapter~\ref{ch:format}'' wording, reflecting that the
  Pass 3 architectural shift made causal context part of
  operation semantics rather than only file storage. \\
  \today & Pass 11 revision (spec ratification) &
  Converted the v0 implementation's provisional, golden-locked byte
  choices into ratified normative spec text; the architecture is
  unchanged. \emph{Adopt-and-pin (bytes):} pinned the
  \texttt{TypedObjectId} 16-bit big-endian discriminant table
  (\texttt{0}--\texttt{27}, \texttt{Registered}~=~27; added the
  \texttt{Tuplet}/\texttt{RepeatStructure}/\texttt{LyricLine}/%
  \texttt{ChordSymbol}/\texttt{View} variants the implementation
  carried); the promoted-voice (\texttt{MUSCSVCE}), system-pitch
  (\texttt{MUSCSPCH}, tuning always part of identity), and
  integrity-anomaly (\texttt{MUSCSANM}, promoted to a reserved
  built-in tag) derivations; the \texttt{ChunkKind} /
  \texttt{ProfileId} / \texttt{CompressionAlgorithm} discriminants;
  the \texttt{ManifestId} preimage (with \texttt{manifest\_id}
  excluded); and the \texttt{RationalTime}/scalar primitive layouts
  plus the shared codec convention baseline the Binary Format
  companion inherits. \emph{Decide-and-pin:} \texttt{TempoShape::Linear}
  interpolates \emph{speed} (whole notes/second), not bpm; a
  \texttt{StructuralFieldCollision} tags the \emph{winner}
  \texttt{Conflicted}; lifted the $>$2-way voice-promotion
  generalization (partial-overlap, lowest-id-survivor) to normative;
  pinned the \texttt{TransactionCategory} and \texttt{ObjectKind} core
  vocabularies; added \texttt{ResolutionAction::Dismiss} so the
  \texttt{Dismissed} resolution state is reachable by an authored
  operation; and specified the (non-canonical) \texttt{LayoutObjectId}
  derivation with a \texttt{MUSCLOID} tag as the Track~A target (then
  deferred to Track~A; subsequently wired by the reference
  implementation in P12-I2). \emph{Fixes:} blob hashing
  is the bare \texttt{"MUSCBLOB" || payload} (deleted the
  contradictory ``identically to chunks'' phrasing); added the
  equal-generation superblock rule (\texttt{DivergentSameGeneration}
  on load-bearing-field divergence); defined \texttt{ProfileConstraints}
  and placed the required \texttt{RetentionPolicy} in it with
  first-declared multi-profile precedence; made the DVV zero-based
  counter floor normative; reconciled the graph-invariant count to 19
  (was 18 in QUICKSTART) and named the three construction-time MUSTs
  (time-signature beat-group sum, ordering-DAG acyclicity, and
  non-degenerate \texttt{TupletRatio}, now enforced at construction).
  \emph{Consolidation:} added the Canonical Byte-Layout Reference
  (Appendix~\ref{app:bytes}) gathering every ratified discriminant,
  derivation preimage, and primitive encoding into one place for the
  Binary Format companion to import, with a domain-tag registry and
  golden-lock anchors; and closed two follow-through gaps where the
  discriminant \emph{bytes} of \texttt{ObjectKind} (\texttt{Voice}~=~0,
  \texttt{Pitch}~=~1, \texttt{Registered}~=~2; feeds the
  \texttt{IntegrityAnomalyId} preimage) and \texttt{ResolutionAction}
  (\texttt{AcceptLoser}~=~0~$\dots$~\texttt{Registered}~=~5; feeds the
  operation content hash) were golden-locked in code but unpinned in
  spec text.
  \\
  \today & Pass 12 tranche 1 (audit spec-alignment) &
  Aligned spec text with dispositions the project had already made
  and the reference implementation exercises; no byte-layout or
  architecture change. The re-anchoring rule table's Slur and
  Spanner rows now state the Operation Catalog's surviving-endpoint
  semantics (re-anchor while $\geq 1$ endpoint survives,
  cascade-delete only when none does; proximity-aware re-targeting
  deferred). \texttt{TupletCompensation::RewriteTuplets} adopts the
  catalog's id-only v1 payload (\texttt{TupletRewrite} removed):
  graph-aware reduction refuses the variant rather than fabricate
  rewritten values, and a value-carrying payload revision reopens
  it. Operation-envelope blocks are pure envelope vectors (a
  content-addressed chunk cannot embed its own id); the per-block
  summary metadata
  (\texttt{dvv\_summary}/\texttt{min\_stamp}/\texttt{max\_stamp})
  relocated to the manifest's \texttt{operation\_block\_summaries}
  map (opaque ops-supplied bytes, non-canonical, ChunkId-keyed),
  and the order-blocks-by-\texttt{min\_stamp} \SHOULD{} was dropped
  because the canonical manifest encoding sorts chunk references by
  encoded form. The spelling and decomposition pre-passes are now
  specified as producing \emph{canonical derived annotations} ---
  deterministic functions of (materialized graph, profile,
  versioned algorithm id), recomputed on materialization and never
  stored --- replacing the stored-\texttt{Inferred}-attachment
  model; the incremental re-run requirement is demoted from
  \MUST{} to \MAY{} with caching and incrementality required to be
  unobservable. Spelling-precedence ties now break by canonical
  attachment order (attachments carry no creation timestamp). The
  \texttt{ConflictRegistry} moved off the \texttt{Score} root into
  canonical materialized state, matching Chapter~6.
  \texttt{SolverTier} gains the non-conformance \texttt{Stub}
  variant. Companion: Operation Catalog 0.2.0 $\rightarrow$ 0.3.0
  (ModifyEvent materializes metric placement changes behind a
  voice-occupancy placement precondition). Dispositions logged in
  \texttt{spec/PASS12\_RATIFICATION\_LOG.md}; the Pass 12 batch
  itself remains open.
  \\
  \today & Pass 12 tranche 1 addendum (Push-3 enablers) &
  Two enabler ratifications for the machinery-wiring work.
  \texttt{ResolveEquivocation}, previously a dangling prose
  reference in the equivocation section, is now a defined
  meta-operation: \texttt{OperationPayload} gains the variant here,
  and the Operation Catalog (0.3.0 $\rightarrow$ 0.4.0) pins its
  payload schema (\texttt{\{ target: OperationId, chosen:
  EnvelopeHash \}}), the order-independent
  earliest-resolve-governs promotion, idempotence, differing-resolve
  meta-conflict, and precondition no-ops; the profile-declared
  selection-policy path remains open (P12-K5). Engraving break
  overrides became position-addressable:
  \texttt{OverrideKind::SystemBreak}/\texttt{PageBreak} carry the
  break's \texttt{TimeAnchor} (the \texttt{ScoreGraph} target names
  the owning region; projected break overrides carry
  \texttt{Internal} origin pending P11-C8 authorship surfacing),
  closing the representational gap that blocked the logical-stage
  override projection §Engraving Overrides requires.
  \\
  \today & Phase-3 first tranche (casting-off + catalog) &
  Companion movements, no core-spec byte or architecture change. The
  Operation Catalog (0.4.0 $\rightarrow$ 0.5.0) gains the Phase-3
  schema-fill tranche --- \texttt{CreateStaff} (set-union global-staff
  mint, plus a staff-liveness precondition on
  \texttt{CreateStaffInstance}), \texttt{SetTimeSignature} and
  \texttt{SetTempoSegment} (LWW structural overwrites beneath the
  whole-grid \texttt{SetMetricGrid}), \texttt{SetStaffLayout} (LWW
  advisory) --- and a rewritten \sectionsc{UndoTransaction}:
  \texttt{StrictInverse}/\texttt{BestEffort} now perform
  \emph{value restoration} for overwrite primitives via
  canonical-order write chains (delete resurrection, Transpose
  inversion, and Cascade closure stay deferred, P11-C8 narrowed).
  The Binary Format companion (0.1.0 $\rightarrow$ 0.2.0) appends
  the corresponding wire discriminants. This document's
  \texttt{OperationKind} listing gains the four kinds. The create
  score/canvas slots remain deliberately unavailable pending an
  addressable root model (P12-K8).
  \\
  \today & Schema major 1 (data-model expansion) &
  Defines the two referenced-but-undefined types
  \texttt{CanvasLayoutDefaults} (with \texttt{CanvasSize} /
  \texttt{CanvasMargins}, staff-space page geometry, A4/8\,mm default;
  closes P12-I7) and \texttt{PitchRange} (advisory pitch compass, used by
  \texttt{Instrument.range} and \texttt{IndeterminacyHints}); adds
  \texttt{Region.permits\_spanning\_slurs} (default false, P12-K7). Records
  that schema major~1 is the first data-model expansion major: the
  canonical-base \texttt{MaterializedState} embeds none of these values and
  stays major~0, byte-identical, but \texttt{Region.permits\_spanning\_slurs}
  reaches the \emph{canonical} \texttt{CreateRegion} operation payload, so a
  major-0 reader opens a bundle carrying v1 \texttt{CreateRegion} ops
  read-only. The wire form, the resolved-layout length-prefix unification
  (the manifest-embedded barrier blobs stay \texttt{u64}), and the
  byte-for-byte v0${\to}$v1 migration are ratified in the Binary Format
  companion (0.2.0 $\rightarrow$ 0.3.0).
  \\
  \today & Schema major 1 Phase F (engrave ratify-only) &
  Ratifies three engrave/layout-ir dispositions the implementation already made
  (Pass 12 batch I8/I9/I10, now struck), with no byte-layout or architecture
  change: the \emph{break-constraint satisfaction} predicate
  (Requirement~\ref{req:layoutir:break-satisfaction} --- a \texttt{SystemBreakAt}
  / \texttt{PageBreakAt} is satisfied iff the final \texttt{ResolvedLayoutIR}
  starts a system/page at that slot, a region-first slot trivially); \emph{break-override
  attribution} threaded through a \texttt{ConstrainedLayoutIR.break\_origins}
  sidecar, declining to widen the normalized constraint record
  (Requirement~\ref{req:layoutir:break-origin-attribution}); and the
  \emph{system-continuation} synthesis kind
  \texttt{Registered(SYSTEM\_CONTINUATION\_SYNTHESIS)} with an
  \texttt{(original,\,ordinal)} instance key
  (Requirement~\ref{req:layoutir:continuation-synthesis}). Flags one open item:
  which region governs a cross-region slur's spanning permission (implemented as
  conservative AND, P12-K12).
  \\
  \today & Pass 12 G-ratification (the batch pass) &
  Retires the open Pass-12 batch (28 rows; dispositions in
  \texttt{PASS12\_RATIFICATION\_LOG.md}). Algorithm dispositions: the spelling
  pre-pass and notational decomposition move to profile-declared-by-versioned-
  identifier with ratified v1 defaults
  (Requirements~\ref{req:pitch:spelling-algorithm}
  and~\ref{req:time:decomposition-algorithm}; closes P12-H1/H3/H4/H5/C5);
  decomposition precedence pinned fixed (P12-H6); authored annotations for
  inference-ineligible targets surface
  (Requirement~\ref{req:pitch:authored-uninferred}, P12-H7). Identity:
  system-derived intrinsic content immutable under reduction (P12-K3). Genesis
  ratified outside the operation set (P12-K8). Cross-region slur permission =
  AND (P12-K12). Re-anchoring: cue cascade-on-any-source, Range truncate
  defined, annotation orphaning sanctioned, \texttt{ReanchorReason::
  SameCanvasNearer} appended (P12-C1/C2/C3/C4). Barriers: target-free and
  opaque-operation matching (Requirement~\ref{req:format:barrier-matching},
  P12-E4); unsafe-edit tombstone semantics
  (Requirement~\ref{req:format:unsafe-tombstone}, P12-E5; encoding deferred to
  the Binary Format companion). Solver: kind-determined strength
  (Requirement~\ref{req:solver:kind-strength}, P12-I4), sub-conformant report
  shape (Requirement~\ref{req:solver:subconformant-report}, P12-I5), the
  Minimal-tier constraint floor
  (Requirement~\ref{req:layoutir:constraint-floor}, P12-I6). The stale
  \texttt{OperationKindTag} listing gains the eleven appended tags. Companion
  movements: Operation Catalog 0.5.0~$\rightarrow$~0.6.0, Binary Format
  0.3.0~$\rightarrow$~0.4.0.
  \\
  \today & Schema major 2 (data-model expansion, Phase A) &
  Defines the leaf types Chapter~5's ratified shapes referenced but never
  pinned: \texttt{SlurKind}, \texttt{CurveDirection},
  \texttt{CurvatureOverride}, the shared \texttt{SpanStyle} record
  (consolidating the formerly-sketched identical
  \texttt{SlurStyle}/\texttt{TieStyle}/\texttt{SpannerStyle} triplet),
  \texttt{SubBeam}, \texttt{BeamGeometryOverride}, the
  \texttt{SpannerKind} payload types (\texttt{HairpinDirection},
  \texttt{OctaveOffset}, \texttt{PedalKind}, \texttt{TextLineDefinition},
  \texttt{BracketKind}) --- \texttt{SpannerKind}
  gains a leading \texttt{Generic} variant as the honest migration default
  --- \texttt{Volta}, \texttt{MetadataValue}, \texttt{Timestamp},
  \texttt{SoundConfiguration}, \texttt{TranspositionInterval} (structural,
  advisory until the Chapter~4 tuning catalog per the P12-K2 discipline),
  \texttt{UnpitchedMember}, \texttt{SpaceUnit}, \texttt{LineStyle}, and
  \texttt{StaffBracketKind}. Ratifies the clef/key content model into
  Chapter~5 (\texttt{Clef}, \texttt{ClefShape}, \texttt{KeySignature},
  \texttt{ClefChange}, \texttt{KeySignatureChange} --- the visible-slice
  value types) and retires the \texttt{ClefId} identifier sketch:
  \texttt{Staff.default\_clef} and \texttt{Instrument.default\_clef} embed
  \texttt{Clef} values, and \texttt{Instrument.default\_staff\_lines} is
  reconciled from \texttt{u8} to \texttt{StaffLineConfiguration}. Metadata
  timestamps are pinned strictly authored
  (Requirement~\ref{req:graph:metadata-timestamps}: nothing writes them
  implicitly). Chapter~8 gains the schema-major-2 paragraph (canonical base
  stays major~0; the filled values reach the snapshot and the seven embedding
  operation payloads, \texttt{InsertRegion} transitively). Wire forms and the
  total default-filling migration: Binary Format companion
  0.4.0~$\rightarrow$~0.5.0.
  \\
  \today & Schema major 2 (repeat authoring, Phase D) &
  The dedicated \texttt{CreateRepeatStructure}/%
  \texttt{DeleteRepeatStructure} pair enters the operation set:
  Chapter~6's \texttt{OperationKind} enum and Chapter~8's
  \texttt{OperationKindTag} listing gain the two variants (repeats are not
  \texttt{CrossCuttingValue} kinds on the wire --- the reconciliation
  note follows the \texttt{CrossCuttingStructure} listing), and the
  re-anchoring rule table gains the Repeat-structure~/ Anchor row
  (re-anchor to the nearest surviving anchor across \emph{every}
  event-referencing anchor site --- \texttt{start}/\texttt{end}, jump
  targets, volta spans --- cascade-delete only when none survives).
  Six-part schemas: Operation Catalog \sectionsc{Repeat Structures},
  0.6.0~$\rightarrow$~0.7.0. Wire discriminants 28/29 and the
  per-payload stamping (create born at v2; delete a minor kind append
  at major~0): Binary Format companion 0.5.0~$\rightarrow$~0.6.0.
  \\
  \today & Schema major 2 (rendering, Phase F) &
  The rendering consumers of the schema-major-2 data model are
  ratified into Chapter~7. \texttt{ResolvedLayoutIR} gains the
  non-glyph primitive vocabulary it has always carried in
  implementation --- line \texttt{Stroke}s (staff lines, stems,
  barlines, volta brackets) and cubic-B\'ezier \texttt{Curve}s
  (slurs) --- as \ref{req:layoutir:resolved-primitives}
  (provenance-traced, non-canonical). Two Minimal-tier rendering
  floors join the constraint floor: \ref{req:layoutir:repeat-render}
  (E1: repeat barlines at resolved boundaries, volta brackets with
  ending numbers, honest non-drawing of unresolvable boundaries,
  jump-kind marks and cross-region repeats deferred) and
  \ref{req:layoutir:slur-curve} (E2: a slur renders as a cubic
  B\'ezier honoring \texttt{CurvatureOverride}, honest non-drawing of
  unresolvable or cross-staff spans, a non-\texttt{Solid} line style
  surfaced not silently rendered solid). The curvature-computing
  algorithm and the \texttt{RenderIR} encoding stay forward-referenced
  out (\ref{sec:layoutir:forward}, \ref{sec:layoutir:render}). Layout
  geometry is non-canonical, so no wire form and no companion version
  moves --- an engrave-side ratify-as-implemented, mirroring the
  schema-major-1 Phase~F precedent.
  \\
  \today & Layout IR (slur rendering completion, Push 3) &
  The three slur refinements \ref{req:layoutir:slur-curve} deferred are
  now landed and its requirement extended: an authored non-\texttt{Solid}
  \texttt{SpanStyle} line renders faithfully (dashed / dotted --- the
  \texttt{LineStyle} rides the \texttt{Curve}, whose listing gains the
  field), a slur spanning a system break splits into per-system
  sub-curves (de~Casteljau; first segment keeps the slur's provenance,
  the rest synthesized continuations), and \texttt{slur\_shape\_penalty}
  moves off its pinned $0.0$ to a real measurement (Quality Metric
  Catalog \sectionsc{slur\_shape\_penalty} rationale refreshed; formula
  unchanged, no version move). Layout geometry stays non-canonical --- no
  wire form, no companion version moves.
  \\
  \midrule
  \today & Layout IR (primitive band ownership, Push 3) &
  Chapter 7 gains \ref{req:layoutir:primitive-band-ownership}: every
  primitive the projection presents to the solver ---
  \texttt{GlyphObject}, \texttt{Stroke}, \texttt{Curve} --- declares the
  \texttt{VerticalBandId} that owns it, naming a band that exists, and a
  vertical solver takes ownership from that declaration rather than
  inferring it from geometry. Ratifies-as-implemented the contract a
  vertical solver needs and an inference cannot supply: a slur's
  endpoints are lifted clear of its own staff, so its owner is not
  recoverable from where it is drawn. The \texttt{ConstrainedLayoutIR}
  listing gains \texttt{strokes} / \texttt{curves} (present since staff
  lines, never listed) and the \texttt{Stroke} / \texttt{Curve} listings
  gain \texttt{vertical\_band}; a band now exists for every staff of a
  region and for its margin, members or not.
  \ref{req:layoutir:resolved-band-ownership} carries the declaration
  through the resolved stage on strokes and curves (a
  \texttt{ResolvedGlyph} keeps none --- its ownership is already baked
  into its position). Band ownership is non-canonical attribution
  metadata: it enters no content hash and no canonical encoding, so no
  wire form and no companion version moves.
  \\
  \bottomrule
\end{longtable}

\end{document}
